\documentclass[11pt,a4paper]{article}
\usepackage[margin=25mm]{geometry}
\usepackage{fontspec}
\tracinglostchars=3
\setmainfont{Latin Modern Roman}
\setsansfont{Latin Modern Sans}
\setmonofont{DejaVu Sans Mono}[Scale=0.78]
\usepackage{amsmath,amssymb,amsthm,mathtools}
\usepackage{longtable,booktabs,array,calc,enumitem}
\usepackage{xcolor}
\usepackage[unicode,colorlinks=true,linkcolor=blue!45!black,urlcolor=blue!55!black]{hyperref}
\usepackage{xurl}
\usepackage{newunicodechar}
\newunicodechar{ℓ}{\ensuremath{\ell}}
\newunicodechar{π}{\ensuremath{\pi}}
\newunicodechar{χ}{\ensuremath{\chi}}
\newunicodechar{φ}{\ensuremath{\varphi}}
\newunicodechar{ζ}{\ensuremath{\zeta}}
\newunicodechar{′}{\ensuremath{\prime}}
\newunicodechar{α}{\ensuremath{\alpha}}
\newunicodechar{β}{\ensuremath{\beta}}
\newunicodechar{γ}{\ensuremath{\gamma}}
\newunicodechar{δ}{\ensuremath{\delta}}
\newunicodechar{ε}{\ensuremath{\varepsilon}}
\newunicodechar{θ}{\ensuremath{\theta}}
\newunicodechar{λ}{\ensuremath{\lambda}}
\newunicodechar{μ}{\ensuremath{\mu}}
\newunicodechar{ν}{\ensuremath{\nu}}
\newunicodechar{ξ}{\ensuremath{\xi}}
\newunicodechar{ρ}{\ensuremath{\rho}}
\newunicodechar{σ}{\ensuremath{\sigma}}
\newunicodechar{τ}{\ensuremath{\tau}}
\newunicodechar{ψ}{\ensuremath{\psi}}
\newunicodechar{ω}{\ensuremath{\omega}}
\newunicodechar{Γ}{\ensuremath{\Gamma}}
\newunicodechar{Δ}{\ensuremath{\Delta}}
\newunicodechar{Θ}{\ensuremath{\Theta}}
\newunicodechar{Λ}{\ensuremath{\Lambda}}
\newunicodechar{Ξ}{\ensuremath{\Xi}}
\newunicodechar{Π}{\ensuremath{\Pi}}
\newunicodechar{Σ}{\ensuremath{\Sigma}}
\newunicodechar{Φ}{\ensuremath{\Phi}}
\newunicodechar{Ψ}{\ensuremath{\Psi}}
\newunicodechar{Ω}{\ensuremath{\Omega}}
\newunicodechar{₀}{\ensuremath{_0}}
\newunicodechar{₁}{\ensuremath{_1}}
\newunicodechar{₂}{\ensuremath{_2}}
\newunicodechar{∞}{\ensuremath{\infty}}
\newunicodechar{≤}{\ensuremath{\leq}}
\newunicodechar{≥}{\ensuremath{\geq}}
\newunicodechar{∈}{\ensuremath{\in}}
\newunicodechar{⊗}{\ensuremath{\otimes}}
\newunicodechar{≅}{\ensuremath{\cong}}
\newunicodechar{≃}{\ensuremath{\simeq}}
\newunicodechar{∼}{\ensuremath{\sim}}
\newunicodechar{≠}{\ensuremath{\ne}}
\newunicodechar{→}{\ensuremath{\to}}
\newunicodechar{↔}{\ensuremath{\leftrightarrow}}
\newunicodechar{⊂}{\ensuremath{\subset}}
\newunicodechar{∩}{\ensuremath{\cap}}
\newunicodechar{∪}{\ensuremath{\cup}}
\usepackage{microtype}
\setlength{\emergencystretch}{3em}
\setlength{\parskip}{0.35em}
\setlength{\parindent}{0pt}
\setcounter{tocdepth}{2}
\providecommand{\tightlist}{\setlength{\itemsep}{0pt}\setlength{\parskip}{0pt}}
\newtheorem{theorem}{Theorem}[section]
\newtheorem{proposition}[theorem]{Proposition}
\newtheorem{lemma}[theorem]{Lemma}
\newtheorem{corollary}[theorem]{Corollary}
\theoremstyle{definition}\newtheorem{definition}[theorem]{Definition}
\theoremstyle{remark}\newtheorem{remark}[theorem]{Remark}
\title{ALPOGE\_BRIDGES\\The Riemann hypothesis and Dirichlet GRH\\in Grothendieck's programme}
\author{Research synthesis and proof checks in the present run}
\date{21 November 2026\quad\textnormal{(Europe/Berlin)}}
\begin{document}
\maketitle
\begin{abstract}
We adopt the supplied Alp\"oge RH and Dirichlet-GRH declarations as
theorems and determine their precise analytic and motivic consequences.
The main work concerns questions and unfinished statements of Grothendieck
and Deligne, using original archive sheets and exact literature locators.
Full proofs are supplied for the deductions and counterexamples established
in this run. Existing results are cited with their operative hypotheses.
The arithmetic consequences include shifted critical lines for abelian
Artin--Tate objects, effective finite Frobenius recognition, the corresponding
Borger realizations, explicit classical consequences, and trace positivity.
Separate cycle-theoretic, Hodge-theoretic and categorical arguments settle
several archive statements without using RH. Counterexamples locate the
information lost by passing to semisimple Euler factors or numerical traces.
Unresolved questions retain explicit gates and attempted approaches.
This working manuscript records ongoing work; its inclusion of a gate
is not a claim that the corresponding problem is solved.
\end{abstract}

\section*{Reading the claims and the check record}
The mathematical input is Alp\"oge (2026), \emph{Remark}, together with
the five named declarations in \nolinkurl{ALPOGE_Main.lean} and their
targets in \nolinkurl{ALPOGE_Statements.lean}. The supplied Mathlib revision
is recorded in the source-scope section. The declarations have been
inspected; the imported proof tree is not in the supplied packet, and
no fresh Lean kernel compilation is claimed. This does not weaken or
replace the input adopted by the task. In this manuscript, ``under the
adopted theorem'' describes the provenance of a consequence; it is not
an additional unproved RH hypothesis within the task's adopted theory.

An identifier denotes a finding established in this run, not a claim
of first publication. We do not claim an exhaustive novelty search.
The check table records actual full-proof readings, primary-source
statement checks, computations, and carried background separately.
Prior intake check reports are never relabelled as checks performed
here. The word ``refuted'' always names the precise statement or
proposed inference for which a counterexample is given.

The strongest arithmetic results concern primitive Dirichlet factors
and integral Tate twists. General Hecke characters over other fields,
nonabelian Artin representations, ramified mixed extensions, integral
lattices, period torsors and monodromy require the additional data
specified in their sections. A proved failure of one reconstruction
rule is not a logical-independence theorem about RH.

\section*{Ranking by the question settled}
\begin{enumerate}
\item \textbf{Original archive statements.} The cote 14 implication and
its proposed converse are addressed over an algebraically closed field
for the precisely defined generated Albanese quotient. All smooth
quasiprojective parameters are treated in every characteristic for the
cohomological implication; their integral relation quotient is zero in
characteristic zero and is proved to be primary torsion at the
characteristic in positive characteristic. The finite-field rational-point reading
has a separate counterexample. Cote 15's abelian--Artin motivic sector is
now treated using the cited 2026 all-CM Hodge theorem and its exact Tate
and standard-conjecture consequences; arbitrary motives and rational
fibre functors retain their distinct gates. Cote 18's real-Frobenius,
weight-shift and strong-period formulations, and cote 67's capping and
involution claims, receive explicit proofs or counterexamples.
The canonical capping map has an exact real-analytic image but fails
the natural holomorphic formulation. Cote 123's dimension, rational
singularity, nilpotent uniqueness and differential-rank questions have
exact classical answers and checked short deductions. Its strict
nonsimply-laced subregular rank prediction fails: the rank is always
one less than the semisimple rank in every type. The group quotient has a complete
central-character formula; type A has a closed minimal-polynomial
formula. The apposition quadruples are classified completely in type A.
Cote 66's named affine torsion-coefficient comparison cases are
affirmative; its universal degree-zero test is exactly henselianity.
\item \textbf{Deligne and the geometric gate.} The actual problem
Weil II 1.8.15 is answered by the precise admissible-Hodge-theory
theorems. Singular-fibre weights and nearby-cycle weights are computed
separately. The independently checked finite exterior-power criterion reduces the weight--monodromy gate to specified
inertia-invariant trace bounds. The universal stable curve has an explicit
geometric perverse decomposition, and its actual exit functor has the
proved right-Kan direct-image description. The linear theorem of cote
134-8 is established for all countable scalar rings, more generally
under a finite scalar pure-injective bound, or under either of two
finite-generation conditions, with the specified exact adjunction,
scalar compatibility and normalization retained. The unrestricted
large-ring case keeps its explicit counit gate.
\item \textbf{New effective arithmetic access.} AQ-FINITE gives a
finite recognition bound for semisimple abelian Artin--Tate objects;
BR-AT-01 and CC-F01/F01b give the exact shifted analytic factors.
The effective Siegel deduction, abelian-field Chebotarev and prime-search
bounds discharge their original conditional hypotheses.
\item \textbf{Criteria and failed routes.} Trace-class positivity,
moment inequalities, Li and approximation criteria are distinguished
from independent geometric constructions. Explicit mixed-extension,
monodromy, cross-ratio, divided-power and spectral counterexamples
identify the missing information. The 2026 rational Deligne--Drinfeld
theorem is an external result, separate from profinite GT surjectivity.
\end{enumerate}
\tableofcontents
\clearpage
\part{Checked findings and exactly scoped known results}

% BEGIN archive_questions.tex
% Agent contribution: /root/archive_questions.
% AQ identifiers denote findings of this run, not claims of priority.
\section{Archive questions: exact statements and mathematical scope}

\subsection{AQ-FINITE: finite Frobenius recognition of abelian Artin--Tate motives}
\paragraph{Statement.}
Adopt the Dirichlet GRH theorem supplied in \nolinkurl{ALPOGE_Main.lean},
\nolinkurl{LFunction_eq_zero_re}. Let $N\geq1$, put
\[
 Q=\operatorname{lcm}(N,4),\qquad B_Q=(\varphi(Q)\log Q)^2,
\]
and let $M,M'$ be finite direct sums of Tate twists of abelian Artin
motives over $\mathbb Q$, whose Artin representations factor through
$G_N=(\mathbb Z/N\mathbb Z)^\times$. Work initially with a coefficient
number field $E\subset\mathbb C$ splitting $G_N$. Write
\[
 M=\bigoplus_{\chi,n}(\chi\otimes E(n))^{m_{\chi,n}},
 \qquad
 M'=\bigoplus_{\chi,n}(\chi\otimes E(n))^{m'_{\chi,n}}.
\]
Then $M\simeq M'$ if and only if
\[
 \det(1-TF_p\mid M)=\det(1-TF_p\mid M')
 \quad\text{for every prime }p\leq B_Q,\ p\nmid Q.                 \tag{AQ.1}
\]
Here $E(n)$ has Frobenius eigenvalue $p^{-n}$. These finitely many
polynomials recover every multiplicity $m_{\chi,n}$. The same
isomorphism criterion holds with rational coefficients.
The characters are labelled by their Euler-factor values
$\chi(p)$; changing the Artin-reciprocity convention inverts
all character labels and leaves the criterion unchanged.
One may use a single auxiliary $\ell>B_Q$ for all the
displayed Frobenius factors, avoiding the prime $\ell$
among the tested primes.

\paragraph{Known input and exact locator.}
Lamzouri, Li, and Soundararajan, \emph{Conditional bounds for the least
quadratic non-residue and related problems}, arXiv:1309.3595v3,
Corollary 1.2, p.~3, proves under Dirichlet GRH that every reduced residue
class $a\bmod Q$, $Q>3$, contains a prime
$p_a\leq(\varphi(Q)\log Q)^2$. This result is cited, not reproved.

\begin{proof}
The forward direction is functoriality. Inflate all characters to $G_Q$.
For each $a\in G_Q$ choose such a prime $p_a$.
Every eigenvalue at $p_a$ is $\chi(a)p_a^{-n}$, of complex absolute
value $p_a^{-n}$. Since $p_a>1$, different integers $n$ give different
absolute values. Equality in (AQ.1) therefore gives equality,
for every $a$ and $n$, of the sums of eigenvalues of absolute
value $p_a^{-n}$. Multiplying this sum by $p_a^n$ gives
\[
 t_n(a)=\sum_\chi m_{\chi,n}\chi(a).
\]
Finite character orthogonality yields the reconstruction
\[
 m_{\chi,n}=\frac1{|G_Q|}\sum_{a\in G_Q}
             t_n(a)\chi(a)^{-1}.                                \tag{AQ.2}
\]
Indeed, translating $\sum_a\theta(a)$ by an element $b$ with
$\theta(b)\ne1$ makes that sum zero for every nontrivial character
$\theta$; it is $|G_Q|$ for the trivial character. Thus (AQ.2)
gives the same multiplicities for $M'$.
For rational coefficients, the semisimple commutative algebra
$\mathbb Q[G_Q]$ is a product of cyclotomic fields. The
multiplicity of each simple rational module is recovered from
the common multiplicity of its characters over a splitting field.
The rational representations, weight by weight, are therefore
isomorphic as well.
\end{proof}

This is a finite stopping criterion for recognition within the
abelian Artin--Tate part of the programme. It needs every residue
class, or equivalent full character information: values on an
arbitrary generating set do not determine a general character.
It does not reconstruct extension classes, integral lattices,
period comparisons, or nonabelian fundamental groups.

\subsection{AQ-14-FORWARD: the cohomological implication at cote 14, scan 38}
\paragraph{Archive question.}
The sheet labels \emph{Th\'eor\`eme 3.1}, marks its proof
\emph{d\'em. pas faite}, and asks whether $u_Z=0$ implies,
modulo torsion, the vanishing of
\[
 a({}^tZ):H^{2(d-i)+1}(X,\mathbb Q_\ell(d-i))
                     \longrightarrow H^1(A,\mathbb Q_\ell).
                                                               \tag{AQ.3}
\]
Here $d=\dim X$ and $A$ is an abelian parameter variety.
Its margin asks \emph{A-t-on une r\'eciproque}.
The lower display $H^{(n-i)+1}$ is inconsistent with the
degree in the preceding paragraph; (AQ.3) uses the latter.

\paragraph{Precise modern formulation.}
For a smooth projective geometrically connected $X$ over an
algebraically closed field $k$, let $N^i(X)\subset CH^i(X)_{\rm alg}$
be generated by
\[
 \Gamma_*z,\quad
 z\in\ker\bigl(CH_0(T)^0\longrightarrow\operatorname{Alb}_T(k)\bigr),
 \quad\Gamma\in CH^i(T\times X),                                \tag{AQ.4}
\]
where $T$ is smooth projective geometrically connected.
Set $P^i(X)=CH^i(X)_{\rm alg}/N^i(X)$, the generated Albanese
quotient. For $Z\in CH^i(A\times X)$
the map on degree-zero cycles descends to
$u_Z:A(k)\to P^i(X)$, since two lifts of the same point of
$\operatorname{Alb}_A=A$ differ by an Albanese-kernel class.
If $u_Z=0$ integrally, then (AQ.3) is zero for every
$\ell\ne\operatorname{char}k$.
The sole initially supplied sheet 14/038 does not define
$A^i_{\rm alb}(X)$, but the official full cote now resolves
the definition. In its physical PDF page 24 (footer and
manuscript page 23), Definition 1.1 declares a cycle
Albanese-equivalent to zero when it has the form
$Z(v)$ for a smooth connected projective parameter $T$,
a degree-zero zero-cycle $v$ with $S_T(v)=0$, and a
correspondence $Z$ for which the intersection is defined.
The word \emph{projectif} is bracketed in the sheet,
between \emph{un sch\'ema} and \emph{lisse connexe}.
We first retain that projective-parameter formulation.
AQ-14-QUASIPROJECTIVE below proves equality with the
quasiprojective-parameter relation in characteristic zero,
with the nonproper Albanese and zero-cycle conventions explicit.
Here $S_T$ is the Albanese summation map.
Proposition 1.2(i), on the same page, says these cycles
form a subgroup containing rational equivalences and
defines the quotient of the Chow group by its image.
Proposition 1.2(v), physical PDF page 25 (footer 24),
says Albanese equivalence implies algebraic equivalence.
Finally Corollary 2.3, physical PDF page 34 (footer 33),
defines $P^i(X)$ itself as the algebraically trivial
subgroup of $A^i_{\rm alb}(X)$. Thus this is the
original quotient in the projective formulation,
not a quotient defined by regular maps to abelian
varieties. The Chow moving lemma makes the finitely
many intersections in (AQ.4) defined; since rational
equivalences are already killed, the cycle-level and
Chow formulations agree. Theorem 2.1, physical PDF
page 31 (footer 30), defines $u_Z$ through the
Albanese summation map. Its image lies in $P^i(X)$,
so vanishing in this subgroup and in the original
total quotient are equivalent. The supplied
\nolinkurl{scan038} uses physical PDF numbering:
it is footer/manuscript page 37, after the cover.
Definition 1.1 itself says only that $k$ is a field.
The preceding handwritten sheets inspected in the full
cote concern other motive conventions and preparatory
cycle calculations; no explicit algebraic-closure
convention was found there. Accordingly the
algebraically closed hypothesis in our theorem is
retained as an explicit mathematical scope condition,
rather than attributed to an unseen original convention.

\paragraph{Primary inputs.}
Bloch (1979), \emph{Torsion algebraic cycles and a theorem of Roitman},
Propositions 3.5 (correspondences), 3.8 (specialization), and
Theorem 4.2 (prime-to-characteristic Roitman), provide the inputs.
They are also explicitly recorded in Achter, Casalaina-Martin,
and Vial, \emph{On the image of the second $\ell$-adic Bloch map},
arXiv:2007.07180v2, Propositions A.23, A.24, A.26, and A.29.
These appendix statements apply in every codimension.

\paragraph{Torsion lemma.}
For Bloch's map
\[
 \lambda_X^i:CH^i(X)[\ell^\infty]\longrightarrow
                    H^{2i-1}(X,\mathbb Q_\ell/\mathbb Z_\ell(i)),
\]
one has
\[
 N^i(X)\cap CH^i(X)[\ell^\infty]\subseteq\ker\lambda_X^i.          \tag{AQ.5}
\]
\begin{proof}
Choose $x$ in the intersection, a finite expression
$x=\sum_j\Gamma_{j*}z_j$ of the form (AQ.4), and $r$ with
$\ell^r x=0$. Spread out this entire expression, its
rational-equivalence certificates, its varieties and
correspondences, their Albanese maps, and the equalities
$\operatorname{alb}(z_j)=0$, over an integral finitely generated
ring over the prime ring. Shrink so that the varieties are
smooth projective, the chosen generic Albanese varieties extend
to abelian schemes $B_j$, and the chosen maps $T_j\to B_j$
extend. It is unnecessary to assert that these are universal
relative Albanese maps. Here the base can be chosen regular:
the finitely generated field of definition has a separating
transcendence basis over its prime field, and a model becomes
smooth over $\mathbb F_p$, or over $\mathbb Z[1/M]$ in
characteristic zero, after a further shrinking. All the
finitely many conditions just listed survive that shrinking.
In characteristic zero invert $\ell$ as well.
Choose a closed point with finite residue field of
characteristic $p\ne\ell$. A regular system of parameters
in its local ring gives a flag of prime ideals with regular
successive quotients. Localizing at consecutive primes
gives discrete valuation rings. Specialization along this
flag, after choosing algebraic closures of the residue
fields, gives geometric special fibres over
$\overline{\mathbb F}_p$, denoted by subscripts $0$.

Every degree-zero zero-cycle on a smooth projective connected
variety over $\overline{\mathbb F}_p$ is torsion: join its
finite support by finitely many projective curves and lift
the differences of points to their normalizations. Every
degree-zero divisor class on those curves is a point of a
Jacobian defined over a finite extension of a finite field,
and has finite order.

The specialized map $T_{j,0}\to B_{j,0}$ factors through
$h_j:\operatorname{Alb}(T_{j,0})\to B_{j,0}$. This map is an
isomorphism on $\ell$-power torsion. Indeed Kummer theory gives
$H^1(T,\mathbb Z_\ell(1))=T_\ell\operatorname{Pic}(T)
=T_\ell\operatorname{Pic}^0(T)_{\rm red}$: the N\'eron--Severi
group is finitely generated and has zero Tate module.
Duality of $\operatorname{Pic}^0_{\rm red}$ and Albanese
therefore gives
$H^1(T,\mathbb Z_\ell)=\operatorname{Hom}(T_\ell
\operatorname{Alb}_T,\mathbb Z_\ell)$, compatibly with maps.
This is also the codimension-one calculation in
Achter--Casalaina-Martin--Vial, Propositions A.25 and A.28.
The generic Albanese map thus induces an isomorphism on
$H^1(-,\mathbb Z_\ell)$, and smooth proper base change
preserves this isomorphism. The latter is obtained
by taking inverse limits of the finite-coefficient base-change
isomorphisms for each variety; no assertion that the Albanese
map is an isomorphism on every finite-coefficient $H^1$ is
needed. Equivalently $T_\ell h_j$ is an isomorphism.
It follows that $h_j$ is an isogeny with no $\ell$-primary
kernel, and hence an isomorphism on $\ell$-power torsion.
Since the specialization of the equality
$\operatorname{alb}(z_j)=0$ says that $z_{j,0}$ maps to zero
in $B_{j,0}$, the $\ell$-primary part of its actual Albanese
class is zero. Prime-to-$p$ Roitman implies that the
$\ell$-primary part of $z_{j,0}$ itself is zero. Thus every
$z_{j,0}$, and hence $x_0$, has order prime to $\ell$.
The equality $\ell^r x_0=0$ also specializes, so $x_0=0$.

Bloch's compatibility with specialization sends
$\lambda_X^i(x)$ to $\lambda_{X_0}^i(x_0)=0$. The
cohomological specialization map is an isomorphism by
smooth proper base change, first for finite coefficients
and then for their direct limit $\mathbb Q_\ell/\mathbb Z_\ell(i)$.
Consequently $\lambda_X^i(x)=0$. Only finitely many data are
spread for each $x$; no simultaneous model for all of $N^i(X)$
has been assumed.
\end{proof}

\begin{proof}[Proof of AQ-14-FORWARD]
Put $g=\dim A$. Roitman identifies
$CH_0(A)[\ell^\infty]$ with $A[\ell^\infty](k)$.
Use the torsion zero-cycle lift of a torsion point under this
isomorphism; it is unnecessary to identify it with the
literal class $[a]-[0]$. For every such torsion cycle $z$,
the premise gives $Z_*z\in N^i(X)\cap CH^i(X)[\ell^\infty]$.
Thus (AQ.5) gives $\lambda_X^i(Z_*z)=0$.
Correspondence compatibility, followed by taking Tate modules
and tensoring with $\mathbb Q_\ell$, makes
\[
 Z_*:H^{2g-1}(A,\mathbb Q_\ell(g))
             \longrightarrow H^{2i-1}(X,\mathbb Q_\ell(i))
\]
zero: the degree-zero Bloch map identifies the source
with $V_\ell A$. Under the perfect Poincar\'e pairings its
transpose is (AQ.3), so (AQ.3) is zero.
\end{proof}

This proves the stated generated-Albanese-quotient formulation
in every characteristic, independently of RH.

\paragraph{AQ-14-RATIONAL-PREMISE.}
The integral premise cannot be replaced by $u_Z\otimes\mathbb Q=0$.
Take an elliptic curve $E/\overline{\mathbb F}_p$ and
$X=A=E$, $Z=\Delta_E-E\times\{0\}$.
Then $P^1(E)=\operatorname{Pic}^0(E)$ and
$u_Z(a)=[a]-[0]$. The group $E(\overline{\mathbb F}_p)$ is
torsion, so the rationalized premise holds. The normalized
diagonal acts as the identity on the nonzero two-dimensional
$H^1(E,\mathbb Q_\ell)$, refuting the weakened assertion.
The phrase \emph{modulo torsion} in the original sheet
qualifies its cohomological conclusion.
For the asserted $P^1$ identification, a divisor
correspondence on $T\times E$ defines a morphism
$T\to\operatorname{Pic}^0(E)$ after choosing a base point,
and hence factors through $\operatorname{Alb}_T$.
Every generator of $N^1(E)$ is therefore zero.

\paragraph{AQ-14-FIELD-SCOPE.}
The algebraically closed base is also substantive.
If the premise is tested only on $k$-rational points
over an arbitrary field, it can fail to detect the
cohomological map. Take $k=\mathbb F_5$ and
$E:y^2=x^3+1$. The affine fibres over
$x=0,1,2,3,4$ have respectively $2,0,2,0,1$
points, so $\#E(\mathbb F_5)=6$, including infinity.
For $Z=6(\Delta_E-E\times\{0\})$, the map
$u_Z$ on $E(\mathbb F_5)$ is zero by Lagrange's
theorem, while the induced rational $\ell$-adic
map is $6\operatorname{id}\ne0$ for every $\ell\ne5$.
The original context's use of divisibility of
$A_Y(k)$ in physical PDF page 34 (footer 33)
is consistent with the algebraically closed
convention used here. This counterexample prevents
an unintended arbitrary-field reading.

\subsection{AQ-14-QUASIPROJECTIVE: removing the parameter qualification in characteristic zero}
Let $k$ be algebraically closed and let $X/k$ be smooth projective.
For a smooth connected parameter variety $T$, write
$Z_0(T)^0$ for the group of finite formal sums
$v=\sum_jn_j[t_j]$ of $k$-points with $\sum_jn_j=0$.
In this paragraph the Albanese of a nonproper $T$ means its
\emph{abelian Albanese}: the universal pointed morphism
$a_T:T\to\operatorname{Alb}_{\mathrm{ab}}(T)$ to an abelian
variety. Put $a_T(v)=\sum_jn_ja_T(t_j)$; this does not depend
on the base point used for the pointed morphism.

For $\Gamma\in\mathrm{CH}^i(X\times T)$ define
\[
 \Gamma(v)=\sum_j n_j\,\iota_{t_j}^{!}\Gamma
                         \in\mathrm{CH}^i(X),
\]
where $\iota_t:X\times\{t\}\hookrightarrow X\times T$
is the regular immersion of the fibre. Thus this is a
sum of fibre restrictions; no nonproper pushforward
$X\times T\to X$ is being used. Let $N^i_{\mathrm{qp}}(X)$
be generated by all such $\Gamma(v)$ with $T$ smooth connected
quasiprojective and $a_T(v)=0$. Let $N^i_{\mathrm{proj}}(X)$
use smooth connected projective $T$ instead.

If every permitted $T$ has a smooth projective compactification,
then
\[
 N^i_{\mathrm{qp}}(X)=N^i_{\mathrm{proj}}(X).                 \tag{AQ.14.1}
\]
In particular this equality holds in characteristic zero.
It also holds if the convention for a nonproper parameter is
the universal semiabelian Albanese, with the kernel imposed
on $Z_0(T)^0$. This last convention must be formulated on
zero-cycles, not automatically on the ordinary Chow group
of the nonproper variety.

\begin{proof}
The inclusion from right to left is immediate, since a
smooth projective parameter is quasiprojective. For the
other inclusion take a generator $\Gamma(v)$ as above.
Choose a smooth projective compactification
$j:T\hookrightarrow\overline T$. Represent $\Gamma$ by
a finite integral linear combination of integral
codimension-$i$ subvarieties of $X\times T$, and take
their closures in $X\times\overline T$. This gives a class
$\overline\Gamma\in\mathrm{CH}^i(X\times\overline T)$
whose restriction is $\Gamma$.

The morphism $a_{\overline T}\circ j:T\to
\operatorname{Alb}(\overline T)$, normalized to send the
chosen base point to zero, factors by universality through
$a_T$. Therefore
$a_{\overline T}(j_*v)=0$. Here $j_*v$ denotes the same
finite formal zero-cycle regarded on $\overline T$;
it does not assert the existence of a pushforward
$j_*:\mathrm{CH}_0(T)\to\mathrm{CH}_0(\overline T)$.
For each point in the support of $v$, refined fibre
restriction commutes with restriction to the open subset
$X\times T$; equivalently, both restrictions are computed
in the same neighbourhood of $X\times\{t_j\}$. Hence
\[
 \overline\Gamma(j_*v)=\Gamma(v).
\]
The left side is a generator allowed in
$N^i_{\mathrm{proj}}(X)$, proving the reverse inclusion.
The same proof works for the semiabelian convention,
because the target $\operatorname{Alb}(\overline T)$ is
also semiabelian, and because the semiabelian Albanese
of a smooth projective parameter is its ordinary
abelian Albanese.

In characteristic zero, embed the quasiprojective $T$
in a projective variety and resolve that projective
closure by a projective resolution that is an
isomorphism over its smooth locus, in particular over
$T$. The resulting $\overline T$ is smooth projective.
This is the strong compactification consequence of
Hironaka's resolution theorem, stated in precisely
this form in Kahn and Sujatha, \emph{Birational geometry
and localisation of categories}, May 22, 2008 author
version, Theorem 1.6.3(a), p.~11.
\end{proof}

Thus AQ-14-FORWARD applies in characteristic zero to the
quotient generated using all smooth quasiprojective
parameters, and AQ-14-CONVERSE over $\mathbb C$ remains
a counterexample for that larger parameter class as
well. The algebraically closed base-field qualification
remains. In positive characteristic this comparison
applies to parameters admitting a smooth projective
compactification; no existence of such compactifications
in all dimensions is asserted here. The next theorem proves the
rational cohomological forward implication without needing that
positive-characteristic equality of relations. AQ-14-ALTERATION then proves that the discrepancy between those relations is $p$-primary torsion.

\subsection{AQ-14-NONPROPER: the rational cohomological implication for all smooth quasiprojective parameters}
Retain the algebraically closed base $k$ and smooth projective $X$.
Let $N^i_{\mathrm{qp}}(X)$ be the relation defined in
AQ-14-QUASIPROJECTIVE, using the \emph{abelian} Albanese of a
smooth connected quasiprojective parameter and free degree-zero
zero-cycles. Put
$P^i_{\mathrm{qp}}(X)=\mathrm{CH}^i(X)_{\mathrm{alg}}/
N^i_{\mathrm{qp}}(X)$. For every prime
$\ell\ne\operatorname{char}k$, one has
\[
 N^i_{\mathrm{qp}}(X)\cap\mathrm{CH}^i(X)[\ell^\infty]
       \subseteq\ker\lambda_X^i.                            \tag{AQ.14.2}
\]
Consequently the implication of original 14/038 holds for this
quasiprojective-parameter quotient in every characteristic:
if $A$ is abelian, $Z\in\mathrm{CH}^i(A\times X)$, and
$u_Z:A(k)\to P^i_{\mathrm{qp}}(X)$ is zero integrally, then
\[
 a({}^tZ):H^{2(\dim X-i)+1}(X,\mathbb Q_\ell(\dim X-i))
                       \longrightarrow H^1(A,\mathbb Q_\ell)
\]
is zero. The same implication holds for the smaller relation
using the semiabelian Albanese kernel. No equality
$N^i_{\mathrm{qp}}=N^i_{\mathrm{proj}}$ in arbitrary positive
characteristic is asserted.

\paragraph{Inputs and coefficient conventions.}
Write $H_j^S(T,R)$ for Suslin homology: it is the homology of
the complex generated in degree $n$ by integral closed
subvarieties of $T\times\Delta^n$ finite and surjective over
$\Delta^n$, tensored with $R$. Its degree-zero cycles are
finite formal sums of closed points. Write $G_T$ for the
pointed Serre semiabelian Albanese and
\[
 0\longrightarrow S_T\longrightarrow G_T\longrightarrow B_T
       \longrightarrow0
\]
for its torus and abelian quotient. Universality shows that
$B_T$ is the abelian Albanese of $T$: a homomorphism from a
torus to an abelian variety is zero. The notation $[\ell^\infty]$
means the subgroup annihilated by some power of $\ell$.

Geisser, \emph{Rojtman's theorem for normal schemes},
arXiv:1402.1831v2, Theorem 1.1, p.~2, gives, for normal $T$,
\[
 H_0^S(T,\mathbb Z)[\ell^\infty]
       \xrightarrow{\sim}G_T(k)[\ell^\infty],\qquad
 H_1^S(T,\mathbb Z)\otimes\mathbb Q_\ell/\mathbb Z_\ell=0.
                                                               \tag{AQ.14.3}
\]
The degree-zero convention on the left is automatic for
torsion. This prime-to-characteristic theorem does not
require a smooth compactification. The finite-coefficient
perfect pairing used below is Geisser--Schmidt (2016),
Theorem 1.1, journal p.~92; their Theorem 7.1, p.~118, identifies it
with the Suslin--Voevodsky comparison. We invoke its
prime-to-characteristic case, where tame and ordinary
abelian $\ell$-primary covers agree. Kahn (2014),
\emph{Algebraic tori as Nisnevich sheaves with transfers},
Lemma 4.1(b), journal p.~707 (July 4, 2013 author version,
Lemma 3(b), p.~8), states
$\operatorname{Alb}_{\mathrm{sa}}(G)=G$ for a semiabelian $G$.

The integral generic-base-change input is Sun,
arXiv:1511.01393v1, Definition 2.2, p.~3, Lemma 2.3(iii),
pp.~4--5, and Theorem 2.5(i), p.~7, with
$\Lambda=\mathbb Z_\ell$. In Lemma 2.3(iii)'s proof, one
dense open works for \emph{all} finite coefficient levels,
by the finite set of Jordan--H\"older constituents. The
constant $\mathbb Z_\ell$ system has the single simple
constituent $\mathbb F_\ell$. Constructibility and preservation
of stratifiability are part of the stated setup and
Remark 2.5.1(ii). We use this only over a finite-type
$\mathbb F_p$-base, $p\ne\ell$, which satisfies the finite
cohomological-dimension condition (LO).

\begin{proof}
Characteristic zero follows already from
AQ-14-QUASIPROJECTIVE and AQ-14-FORWARD. Thus assume
$\operatorname{char}k=p>0$ and $\ell\ne p$.
We first give the three facts needed for specialization.

\emph{1. A correspondence with proper target descends to
Suslin zero-homology.}
For $\Gamma\in\mathrm{CH}^i(X\times T)$, the fibre-evaluation
map $Z_0(T)\to\mathrm{CH}^i(X)$ kills every Suslin boundary.
Indeed, a generator of a Suslin one-chain is an integral
$W\subset\mathbb A^1\times T$ finite and surjective over
$\mathbb A^1$. Its normalization $C$ is a smooth curve,
finite over $W$ and over $\mathbb A^1$. Let $\overline C$
be its smooth projective completion. Pull $\Gamma$ back
along the morphism $X\times C\to X\times T$; pullback is
well defined between smooth varieties by the graph/refined
Gysin construction. Extend a representing cycle by closure
to a class on $X\times\overline C$.

The function $t/(t-1)$ on $\overline C$ has divisor equal
to the fibre over zero minus the fibre over one, including
multiplicities, and no contribution over infinity. The
normalization projection formula identifies their images
in $Z_0(T)$ with the two faces of $W$. Applying the extended
correspondence to this principal divisor gives zero in
$\mathrm{CH}^i(X)$. Its evaluations at points of $C$ agree
with the original fibre evaluations, since refined pullback
is local around those fibres. This proves the claim and
produces a homomorphism
\[
 \Gamma_*:H_0^S(T,\mathbb Z)\longrightarrow\mathrm{CH}^i(X).
                                                               \tag{AQ.14.4}
\]
The construction is compatible with extension of the base
field and with automorphisms of that field. No pushforward
along the nonproper projection $X\times T\to X$ was used.

\emph{2. The integral $H^1$ of the semiabelian Albanese.}
The homology coefficient sequence and (AQ.14.3) give a
natural isomorphism
\[
 H_1^S(T,\mathbb Q_\ell/\mathbb Z_\ell)
             \simeq G_T(k)[\ell^\infty].                    \tag{AQ.14.5}
\]
For clarity, the coefficient sequence has left term
$H_1^S(T,\mathbb Z)\otimes\mathbb Q_\ell/\mathbb Z_\ell$
and right term
$\operatorname{Tor}_1^{\mathbb Z}(H_0^S(T,\mathbb Z),
\mathbb Q_\ell/\mathbb Z_\ell)=H_0^S(T,\mathbb Z)[\ell^\infty]$;
the left term is zero by (AQ.14.3).

The prime-to-$p$ finite pairings of Geisser--Schmidt are
natural and compatible with the coefficient homomorphisms.
Dualizing the direct system
$H_1^S(T,\mathbb Z/\ell^n)$, whose transition is induced by
$\mathbb Z/\ell^n\hookrightarrow\mathbb Z/\ell^{n+1}$,
identifies its dual with the inverse system
$H^1_{\mathrm{et}}(T,\mathbb Z/\ell^n)$ under reduction.
The pairings are viewed in $\mathbb Q_\ell/\mathbb Z_\ell$
via $1\bmod\ell^n\mapsto\ell^{-n}\bmod\mathbb Z_\ell$;
this makes the transition compatibility explicit. Thus
\[
 H^1_{\mathrm{et}}(T,\mathbb Z_\ell)
  \simeq\operatorname{Hom}\bigl(
       G_T(k)[\ell^\infty],\mathbb Q_\ell/\mathbb Z_\ell\bigr)
  \simeq\operatorname{Hom}_{\mathbb Z_\ell}
       (T_\ell G_T,\mathbb Z_\ell).                         \tag{AQ.14.6}
\]
The left side is the inverse limit of the finite groups;
there is no $\varprojlim^1 H^0$ obstruction, because the
constant $H^0$ transition maps are surjective. All
isomorphisms are functorial. Since $G_T$ is its own Serre
Albanese, the Albanese morphism $a_T:T\to G_T$ induces an
isomorphism on $H^1(-,\mathbb Z_\ell)$. This conclusion
is deliberately integral-adic, not an assertion about
$H^1(-,\mathbb Z/\ell^n)$ for each $n$.

\emph{3. The finite-field vanishing.}
Suppose for the moment that $k=\overline{\mathbb F}_p$.
Every class in $H_0^S(T,\mathbb Z)^0$ is torsion.
To see this, join the finitely many points in a representative
by finitely many integral curves in the connected smooth
quasiprojective variety $T$, and pass to their normalizations.
Such curves can be chosen through any prescribed finite
set by hyperplane sections of sufficiently high degree;
only the resulting curves and their normalizations are
used. Differences of points lift to degree-zero classes
on these smooth open curves. For a smooth curve $C$,
$H_0^S(C,\mathbb Z)^0$ is the group of points of its
Rosenlicht generalized Jacobian, equivalently the curve
case of Geisser, Theorem 1.3, p.~3. Each such Jacobian and
each of its points are defined over some finite extension
of a finite field, and its rational points over that
finite field form a finite group. Each lifted class,
and hence each original degree-zero class, is torsion.

Let $v\in H_0^S(T,\mathbb Z)^0$ have zero image in $B_T(k)$,
and let $v_\ell$ denote its $\ell$-primary component.
Under (AQ.14.3), the class $v_\ell$ corresponds to an
element of $S_T(k)[\ell^\infty]$. Consider the homomorphism
\[
 S_T(k)[\ell^\infty]\longrightarrow
 H_0^S(T,\mathbb Z)[\ell^\infty]
 \xrightarrow{\Gamma_*}\mathrm{CH}^i(X)[\ell^\infty]
 \xrightarrow{\lambda_X^i}
 H^{2i-1}(X,\mathbb Q_\ell/\mathbb Z_\ell(i)).                \tag{AQ.14.7}
\]
All varieties, maps and $\Gamma$ descend to a common finite
field $\mathbb F_q$, after enlarging it if needed. Thus
(AQ.14.7) is Frobenius-equivariant. On
$V_\ell S_T=T_\ell S_T\otimes\mathbb Q_\ell$ geometric
Frobenius has eigenvalues $q^{-1}$ times roots of unity:
a torus splits over a finite extension of $\mathbb F_q$.
On
$H^{2i-1}(X,\mathbb Q_\ell(i))$ its eigenvalues have
complex absolute value $q^{-1/2}$, by Deligne (1974),
Theorem (1.6), printed p.~276, after the indicated Tate
twist. These spectra are disjoint. A linear map
commuting with Frobenius between these two spaces is zero:
apply relatively prime characteristic polynomials and
B\'ezout's identity, or use generalized eigenspaces
after extending scalars. Semisimplicity is not
required for this step.

The rational Tate module of the last group in (AQ.14.7)
is exactly $H^{2i-1}(X,\mathbb Q_\ell(i))$. More explicitly,
that group has maximal divisible subgroup
$H^{2i-1}(X,\mathbb Z_\ell(i))_{\mathrm{free}}
 \otimes\mathbb Q_\ell/\mathbb Z_\ell$ and finite quotient
contained in $H^{2i}(X,\mathbb Z_\ell(i))[\ell^\infty]$.
The domain of (AQ.14.7) is divisible. Its image is therefore
in the maximal divisible subgroup, and a homomorphism
between these finite-corank divisible $\ell$-groups is
zero if its map of rational Tate modules is zero.
Consequently (AQ.14.7) is zero, and
$\lambda_X^i(\Gamma_*v_\ell)=0$.

It follows that over $\overline{\mathbb F}_p$ the conclusion
(AQ.14.2) holds: in a finite sum of generators each
$H_0^S$ class is torsion; its prime-to-$\ell$ component
maps to prime-to-$\ell$ torsion, whereas the $\ell$-primary
component has zero Bloch image by the preceding argument.
Taking the $\ell$-primary component of an expression for
$x\in N^i_{\mathrm{qp}}(X)\cap\mathrm{CH}^i(X)[\ell^\infty]$
gives the claimed vanishing.

\emph{4. Spreading to the finite-field situation.}
Now take arbitrary algebraically closed $k$ of characteristic
$p>0$ and a fixed
$x\in N^i_{\mathrm{qp}}(X)\cap\mathrm{CH}^i(X)[\ell^\infty]$.
Choose a finite expression
$x=\sum_j\Gamma_{j*}v_j$, with $v_j\in Z_0(T_j)^0$
mapping to zero in $B_{T_j}(k)$, and choose $r$ with
$\ell^r x=0$. Spread all varieties, finite zero-cycles,
correspondences, rational-equivalence witnesses for these
equalities, and the pointed maps
$T_j\to G_{T_j}\to B_{T_j}$ over an integral finite-type
$\mathbb F_p$-scheme $S$. After shrinking, $S$ is regular,
$X/S$ is smooth projective, the $T_j/S$ are smooth with
geometrically connected fibres, and the $G_j/S$ are
semiabelian schemes in exact sequences
$0\to S_j\to G_j\to B_j\to0$ with tori $S_j$ and
abelian schemes $B_j$. The sections representing the
zero-cycles remain inside $T_j$, and the relations
$B_j(v_j)=0$ persist. The spread maps need not be
relative universal Albanese morphisms.

Apply Sun's integral generic base change simultaneously to
the finitely many $T_j\to S$ and $G_j\to S$. Their
$R^1(-)_*\mathbb Z_\ell$ are constructible adic sheaves,
and become lisse on a further dense open. By (AQ.14.6)
the map induced by $T_j\to G_j$ is an isomorphism at
a geometric generic point. Indeed its base extension
to $k$ is the chosen Albanese map, and prime-to-$p$
\'etale cohomology is unchanged by extension of
algebraically closed fields. Its kernel and cokernel
vanish on a further dense open, so this is an
isomorphism of integral adic sheaves there. Generic
base change now identifies its stalk at \emph{every}
geometric point of that open with the actual map
on $H^1(-,\mathbb Z_\ell)$ of the fibre. This uses one
open for the full integral system, not separate opens
for infinitely many coefficient levels.

Choose a closed point $s$ in this open; its residue
field is finite. Over $\overline{k(s)}$, the map
$T_{j,s}\to G_{j,s}$ factors by universality as
\[
 T_{j,s}\longrightarrow G_{T_{j,s}}
                  \xrightarrow{h_j}G_{j,s}.
\]
By (AQ.14.6), $T_\ell h_j$ is an isomorphism. A
homomorphism of semiabelian varieties whose Tate module
is an isomorphism is an isogeny with no $\ell$-primary
kernel: a positive-dimensional connected kernel or
quotient has a nonzero Tate module, and a nonzero finite
$\ell$-primary kernel makes the Tate-module image a
proper sublattice. The induced homomorphism of abelian
quotients
\[
 b_j:B_{T_{j,s}}\longrightarrow B_{j,s}
\]
is also an isogeny with no $\ell$-primary kernel.
Indeed an isogeny of semiabelian varieties maps the
torus onto the torus; the kernel of its map of abelian
quotients is a quotient of its finite kernel. The first
assertion follows since the resulting quotient of the
target by the image torus is a quotient of the proper
source abelian variety, and hence has zero toric part.

The equality $B_{j,s}(v_{j,s})=0$ and injectivity of
$b_j$ on $\ell$-primary torsion imply that the
$\ell$-primary component of the actual abelian Albanese
class of $v_{j,s}$ is zero. Although the whole class
need not vanish, this is precisely what step 3 needs:
the $\ell$-primary Suslin class of $v_{j,s}$ corresponds
to the actual torus $S_{T_{j,s}}$. Thus the argument
in step 3 gives $\lambda_{X_s}^i(x_s)=0$.

Finally specialize from a geometric generic point to
$s$ along a regular flag through $s$, as in the complete
flag construction in AQ-14-FORWARD. The finite cycle
and rational-equivalence data make $x_s$ the
specialization of $x$, preserve $\ell^r x_s=0$, and
identify each specialized fibre evaluation with
$\Gamma_{j,s*}v_{j,s}$. Refined Gysin functoriality
provides this compatibility locally at the sections.
Bloch's specialization theorem sends
$\lambda_X^i(x)$ to $\lambda_{X_s}^i(x_s)$, and smooth
proper base change for $X/S$ makes the cohomological
specialization an isomorphism. Cohomological invariance
under algebraically closed extension also identifies
the original $k$-fibre with the geometric generic
fibre for this calculation. Therefore
$\lambda_X^i(x)=0$, proving (AQ.14.2).

For the final cohomological implication, lift each
$\ell$-power torsion point of $A$ to its torsion
zero-cycle class by prime-to-$p$ Roitman. The integral
premise puts its image under $Z$ in the left side of
(AQ.14.2). Bloch correspondence compatibility, followed
by Tate modules and tensoring with $\mathbb Q_\ell$,
therefore kills
$Z_*:H^{2\dim A-1}(A,\mathbb Q_\ell(\dim A))
\to H^{2i-1}(X,\mathbb Q_\ell(i))$.
Poincar\'e duality gives exactly the displayed transpose.
This is the same final argument as AQ-14-FORWARD, now
with the broader relation. For the semiabelian-kernel
convention, its relation is contained in the
abelian-kernel relation, so the torsion vanishing and
the implication follow as well.
\end{proof}

\subsection{AQ-14-ALTERATION: the two parameter relations agree away from the characteristic}
\paragraph{Statement.}
Let $k$ be algebraically closed of characteristic $p>0$,
and let $X/k$ be smooth projective.  Define
$N^i_{\rm proj}(X)\subset N^i_{\rm qp}(X)\subset CH^i(X)$
using the projective and smooth quasiprojective parameters,
respectively, and the ordinary abelian Albanese, exactly as
in AQ-14-QUASIPROJECTIVE.  Then
\[
 N^i_{\rm qp}(X)/N^i_{\rm proj}(X)
 \quad\hbox{is a $p$-primary torsion group.}               \tag{AQ.ALT.1}
\]
In particular the relations agree after tensoring with
$\mathbb Z[1/p]$, hence also with $\mathbb Q$.
The natural map between the two cycle quotients has
$p$-primary kernel and induces an isomorphism on
$\ell$-primary torsion for every $\ell\ne p$.
The same assertions hold if the nonproper relation uses
zero cycles annihilated by the semiabelian Albanese.
Together with AQ-14-QUASIPROJECTIVE this gives equality
with rational coefficients in every characteristic.
No bounded exponent for the possible $p$-primary kernel
is asserted.

\paragraph{Exact external inputs.}
The separability refinement is needed here: Illusie and
Temkin (2014), Expos\'e X, Theorem 3.5(iii), printed p.~202,
gives a \emph{separable} projective prime-to-$\ell$ alteration
with regular source over a perfect field.  Their
Theorem 2.1 alone is not the separability statement;
the distinction is explicit on p.~167 and in \S3.5.4,
p.~207.  The other input is the extension of a rational
map from a smooth variety to an abelian variety to a
morphism on the whole source: Milne, \emph{Abelian
Varieties}, version 2.0 (March 16, 2008), Chapter I,
Theorem 3.2, printed p.~17.  We use these as cited
classical inputs.  The moving and trace arguments
needed for their application are proved below.

\begin{proof}
\emph{1. Moving zero cycles inside a dense open using Suslin relations.}
For a smooth connected quasiprojective $T$ and a nonempty
open $U\subset T$, every class in $H^S_0(T,\mathbb Z)$
is represented by a zero cycle supported in $U$.
The replacement can be chosen with the same degree.
We prove this directly on point generators.
There is nothing to prove for a point already in $U$,
or if $\dim T=0$.
Fix $x\in T\setminus U$.  An affine neighborhood of $x$
embeds into affine space.  Successive general hyperplanes
through $x$ cut its dimension down to one and cut
$T\setminus U$ down to dimension zero in that neighborhood.
At each step one avoids finitely many positive-dimensional
components; this is possible because $k$ is infinite and
the linear system has only the point $x$ as its base locus.
An irreducible one-dimensional component through $x$
therefore meets $U$.  Its closure in $T$ is an integral
curve $C_0$ through $x$, still meeting $U$.
Let $C\to C_0$ be its normalization and choose
$x'\in C$ above $x$.  The normalization is finite,
and $C$ is a smooth curve because $k$ is perfect.
Let $\overline C$ be its smooth projective completion,
$D=\overline C\setminus C$, and
$F=C\setminus(C\times_T U)$.
Both $D$ and $F$ are finite sets.  Choose
$y\in C\times_T U$.

For sufficiently large $n$, Riemann--Roch gives a rational
function $f$ on $\overline C$, with poles only at $y$,
whose first-order jet at $x'$ is a uniformizer and whose
value at every point of $D\cup(F\setminus\{x'\})$
is $1$.  Explicitly the required evaluation map is
surjective because the first cohomology of
\[
 \mathcal O_{\overline C}
 \left(ny-2x'-\sum_{z\in D\cup(F\setminus\{x'\})}z\right)
\]
vanishes once its degree exceeds $2g(\overline C)-2$.
Thus $f$ has a simple zero at $x'$ and no zero or pole
at the other points of $D\cup F$.
Its divisor is of the form
\[
 \operatorname{div}(f)=[x']+z-a[y],
 \qquad \operatorname{supp}(z)\subset C\times_T U,
\]
where $z$ is effective and $a>0$.
This divisor is a Suslin boundary on $C$:
the nonconstant morphism
$g=f/(f-1):\overline C\to\mathbb P^1$ is finite,
and $g^{-1}(\mathbb A^1)\subset C$ since $f=1$ on $D$.
Its graph gives a finite correspondence from
$\mathbb A^1$ to $C$; its fibres at $0$ and $1$
differ by $\operatorname{div}(f)$, up to the harmless
choice of boundary sign.  Push this correspondence to
$T$.  It shows that $[x]$ has a same-degree representative
supported in $U$.  Linearity proves the moving assertion.

Both maps used in our relation respect these Suslin
boundaries.  For the correspondence action
$\Gamma_*:Z_0(T)\to CH^i(X)$, this is the proper-target
Suslin descent proved in Step 1 of AQ-14-NONPROPER:
pull to a normalized curve, close on its projective
completion, and evaluate the principal divisor there.
For the ordinary Albanese map it follows directly as well:
the morphism $C\to\operatorname{Alb}_{\rm ab}(T)$
extends to $\overline C$ by properness, and any principal
divisor has zero sum in that abelian variety.
Consequently, if $v$ has degree zero and Albanese sum zero,
we may replace it by a zero cycle $v_U$ supported in $U$
with
\[
 \deg v_U=0,\qquad\operatorname{alb}_T(v_U)=0,
 \qquad \Gamma_*(v_U)=\Gamma_*(v).                \tag{AQ.ALT.2}
\]

\emph{2. The prime-to-$\ell$ transfer.}
Fix a prime $\ell\ne p$ and a generator
$\Gamma_*(v)$ of $N^i_{\rm qp}(X)$.
Choose a projective integral compactification
$T\subset\overline T$.
The cited Theorem 3.5(iii), applied over the perfect field
$k$, gives a regular projective source and a separable
projective prime-to-$\ell$ alteration of $\overline T$.
Choose a connected component that dominates
$\overline T$; write it as
\[
 f:S\longrightarrow\overline T,\qquad
 d=[k(S):k(T)],\qquad (d,\ell)=1.
\]
It is smooth over $k$, projective and integral.
There exists a nonempty open $U\subset T$ over which
$f^{-1}(U)\to U$ is finite \'etale of degree $d$.
Use Step 1 to obtain $v_U$ satisfying (AQ.ALT.2).

Fix a point of $S$ to normalize
$a_S:S\to A_S=\operatorname{Alb}(S)$.
The fibre sum
\[
 \nu:U\longrightarrow A_S,\qquad
 u\longmapsto\sum_{s\in f^{-1}(u)}a_S(s)
\]
is a morphism: on an \'etale cover of $U$ the finite
\'etale cover splits into $d$ sections, their sum is
regular, and these sums descend because they are
invariant under permutations of the sections.
It is a rational map from smooth $T$ to an abelian
variety, so the cited extension theorem gives a morphism
$\overline\nu:T\to A_S$.
After subtracting its value at a base point, universality
of $\operatorname{Alb}_{\rm ab}(T)$ makes this morphism
factor through a homomorphism
$\operatorname{Alb}_{\rm ab}(T)\to A_S$.
Write $v_U=\sum_u n_u[u]$.  Since it has degree and Albanese sum zero,
\[
 w=f_U^*v_U\in Z_0(S)^0,
 \qquad \operatorname{alb}_S(w)
        =\sum_u n_u\nu(u)=0.                       \tag{AQ.ALT.3}
\]

Restrict a cycle representing $\Gamma$ to $X\times U$,
pull it back along the finite \'etale map
$X\times f^{-1}(U)\to X\times U$, and take its
cycle-theoretic closure on $X\times S$.
Call the resulting codimension-$i$ cycle $\Gamma'$.
Fibre Gysin evaluation commutes with this \'etale
pullback and with restriction to the open set.
Each of the $d$ points above $u\in U(k)$ therefore
has the same fibre class $\Gamma_u\in CH^i(X)$.
It follows that
\[
 \Gamma'_*(w)=d\,\Gamma_*(v_U)
             =d\,\Gamma_*(v)
             \in N^i_{\rm proj}(X),               \tag{AQ.ALT.4}
\]
where membership uses the projective parameter $S$
and (AQ.ALT.3).

\emph{3. The quotient and its torsion.}
Every element $c$ of $N^i_{\rm qp}/N^i_{\rm proj}$
is a finite sum of such generators.  For each
$\ell\ne p$, take the product of their degrees in
(AQ.ALT.4); it is prime to $\ell$ and annihilates $c$.
Taking any one $\ell\ne p$ first shows that $c$ is
torsion, of some finite order $n$.  For every
$\ell\ne p$ a prime-to-$\ell$ annihilator forces
$\ell\nmid n$.  Thus $n$ is a power of $p$, proving
(AQ.ALT.1).

Put $Q_{\rm proj}=CH^i(X)/N^i_{\rm proj}$ and
$Q_{\rm qp}=CH^i(X)/N^i_{\rm qp}$, and let
$K=N^i_{\rm qp}/N^i_{\rm proj}$.
The exact sequence
$0\to K\to Q_{\rm proj}\to Q_{\rm qp}\to0$
has $p$-primary torsion kernel.  Multiplication by
$\ell$ is an automorphism of $K$ when $\ell\ne p$:
this holds on every finite cyclic subgroup and hence
on their union.  Therefore the induced map on
$\ell$-primary torsion is injective.  It is surjective
as well: if $y\in Q_{\rm qp}$ is killed by
$\ell^a$, lift it to $x\in Q_{\rm proj}$;
then $\ell^a x\in K$, and choose $z\in K$ with
$\ell^a z=\ell^a x$.  The lift $x-z$ is killed by
$\ell^a$.  This proves the assertion about torsion.
Localization is flat and kills $K$, proving the
claimed equalities after inverting $p$.

Finally, a zero cycle in the semiabelian Albanese
kernel is in the abelian Albanese kernel.  Its relation
is therefore between $N^i_{\rm proj}$ and the relation
just treated.  The same $p$-primary conclusion follows
for that intermediate subgroup and quotient.
\end{proof}

\paragraph{Consequence for the original cohomological question.}
This also gives a second route to the forward implication
of 14/038.  If an $\ell$-primary torsion cycle
$x\in CH^i(X)$ belongs to $N^i_{\rm qp}$, (AQ.ALT.4),
applied to a finite expression for $x$, supplies
$d$ prime to $\ell$ with $dx\in N^i_{\rm proj}$.
AQ-14-FORWARD gives $d\lambda_X^i(x)=0$.
Since $\lambda_X^i(x)$ is $\ell$-primary and
multiplication by $d$ is invertible on that group,
$\lambda_X^i(x)=0$.
The original integral-premise and rational-transpose
argument then applies.  This deduction uses the
projective torsion lemma and the alteration theorem;
the earlier direct Suslin-weight proof is retained
as a separately checked argument.

\subsection{AQ-14-CONVERSE: the generated-quotient converse fails}
For the generated quotient just defined, over $\mathbb C$
the converse fails, even rationally in its target.
The following is an application of Walker (2007), Theorem 5.7(4),
p.~930, and Theorems 6.1 and 6.2, pp.~933--934.
Take a smooth split six-dimensional quadric $W$,
a sufficiently general high-degree smooth hypersurface section
$Y\subset W$, and the two-cycle $\gamma$ supplied by Theorem 6.1
with $r=2$: its singular homology class is zero, whereas
$s\gamma\ne0$ in $L_1H_4(Y,\mathbb Q)$.
The two families of maximal $\mathbb P^3$'s on $W$
give independent middle-dimensional algebraic
homology classes, supplying the rank-at-least-two
hypothesis of Theorem 6.1.
Here $\dim Y=5$. For any complex elliptic curve $E$ put
$X=Y\times E$, so $\dim X=6$, and
\[
 Z=\gamma\times(\Delta_E-E\times\{0\})\in CH^4(E\times X).
\]
Then (AQ.3) is zero but $u_Z$ is nonzero in $P^4(X)$,
even after tensoring with $\mathbb Q$.
\begin{proof}
The cycle class of $Z$ is zero because that of $\gamma$ is zero,
so its cohomological correspondence vanishes.
The needed factorization through the generated Albanese
quotient follows explicitly from Walker's Theorem 5.7(2),(4),
p.~930, and Example 5.5, p.~928. If $d=\dim X$,
$r=d-i$, and $\Gamma\in CH^i(T\times X)$, those statements give
\[
 \Phi^{\rm mor}_{X,r}(\Gamma_*z)
    =\Gamma_{J,*}\Phi^{\rm mor}_{T,0}(z)
    =\Gamma_{J,*}\operatorname{alb}_T(z).
\]
Theorem 5.7(4) applies to every Chow correspondence, not just
an equidimensional representative. The displayed expression
vanishes for each generator (AQ.4). Thus
$\Phi^{\rm mor}_{X,r}$ descends to $P^i(X)$.
The product construction in the proof of Walker's Theorem 6.2
gives a nonzero cylinder
\[
 \gamma\otimes H_1(E,\mathbb Q)\longrightarrow L_2H_5(X,\mathbb Q).
\]
Pairing its image with a dual class on $E$ and pushing to $Y$
gives a nonzero multiple of $s\gamma$, verifying nonvanishing.
This is a morphism of the (ind-)Hodge structures defining
the morphic Jacobian. The source has type $(-1,0),(0,-1)$
after twisting. A nonzero rational morphism from this source
has nonzero tangent map on the associated Jacobians: if the
relevant complex Hodge summand were killed, its conjugate
would also be killed, so the whole morphism would vanish.
More explicitly, strictness of the Hodge filtration gives
a nonzero complex linear map after quotient by $F^0$;
if it were zero, the whole real image would lie in
$F^0\cap\overline{F^0}$, which is zero in weight $-1$.
The underlying Lawson group is countable, as used in
Walker's proof of Theorem 6.2, p.~934. Quotienting the
nonzero vector-space image by its intersection with
this countable period group still gives uncountably
many image points. Torsion in a complex vector space
modulo a countable subgroup is countable, because
its inverse image is the countable union of the
groups $\tfrac1n\Lambda$. Hence
$a\mapsto\gamma\times([a]-[0])$ has nontorsion image
in the morphic Jacobian.
Since it factors through $P^4(X)$, the converse fails
even rationally.
\end{proof}
No unverified effective degree, such as degree 14, is
needed: the primary theorem's sufficiently-large-degree
construction supplies the counterexample. This is a
translation of a known phenomenon to the original
projective quotient of 14/038, independent of RH.
Walker's morphic Jacobian
need not be a finite-dimensional abelian variety.
Consequently this argument does not refute the separate
universality question for the classical Abel--Jacobi
map among finite-dimensional regular abelian targets;
Walker discusses that distinction in \S7.

\subsection{AQ-WEIL-EXIST: the strong conjecture at 15/100}
The source says that its preceding necessary conditions
are sufficient, following the boxed condition
$|\sigma\alpha|=q^{i/2}$. In the unambiguous modern form,
every nonzero algebraic integer with these absolute values
occurs on $H^i$ of a smooth projective scheme over
$\mathbb F_q$. This is true for every $i\geq0$.

The unconditional existence argument is the surjectivity
part of Milne (1994), Proposition 2.6, p.~15 of the author's
manuscript. The whole proposition classifies simple motives
under the Tate conjecture, explicitly assumed at the start
of \S2, p.~12; its full classification must retain that
assumption. Honda (1968), \emph{Isogeny classes of abelian
varieties over finite fields}, Main Theorem, p.~83, is
the unconditional existence input.
\begin{proof}[Short proof of the specified $H^i$ form]
For $i>0$ choose $\beta$ with $\beta^i=\alpha$.
It is integral over the algebraic integers, hence an
algebraic integer, and all its conjugates have absolute
value $q^{1/2}$. Honda's theorem gives an abelian variety
$A/\mathbb F_q$ with eigenvalue $\beta$ on $H^1$.
The K\"unneth summand $H^1(A)^{\otimes i}$ of $H^i(A^i)$
then contains $\beta^i=\alpha$.
For $i=0$, Kronecker's theorem makes $\alpha$ a root of
unity of some order $m$. The permutation action on
$H^0(\operatorname{Spec}\mathbb F_{q^m}
\times_{\mathbb F_q}\overline{\mathbb F}_q,\mathbb Q_\ell)$
contains all $m$th roots of unity. This scheme is smooth
projective. If geometric connectedness were imposed,
the weight-zero assertion would instead allow only $1$.
\end{proof}
The weak assertion of 15/106 follows in this Weil-number
reading, even without a further extension of the finite
field. The claimed identification of the entire motivic
Galois group is separate: realization's full faithfulness
and algebraicity of invariant correspondences are not
proved merely by the existence of each eigenvalue.
Milne's Proposition 2.6 and Remark 2.7 retain Tate precisely
at that point.
More explicitly, his Conjecture 1.14 is the equality
between the dimension of numerical codimension-$r$
cycle classes and the order of the pole at $s=r$.
Proposition 1.15 derives algebraicity of all Frobenius
invariants and semisimplicity; Corollary 1.16 derives
full faithfulness of the $\ell$-adic realization.
These are the additional inputs used for the
injectivity part of Proposition 2.6 and for the
generation assertion of Remark 2.7. The eigenvalue
existence proof above supplies neither of them.
The integral condition also forces $i\geq0$:
for $i<0$ the nonzero integral norm of $\alpha$
would have absolute value less than one.

\paragraph{AQ-TATE-2026: the abelian sector is cleared in current literature.}
The hypothesis required for abelian varieties has now
been supplied in the primary literature. Van Hoften
and Sempliner (2026), Hexagon:2610.00259v2,
Theorem I, p.~1, prove the rational Hodge conjecture
for every complex CM abelian variety. Applying
Milne (1999), \emph{Lefschetz motives and the Tate
conjecture}, Theorem 7.1, journal p.~72, with the definition (0.1)
and finite-model descent in the Introduction, gives
the full Tate conjecture for every abelian variety
over every finite field. Milne's November 9, 2026
author note states this explicitly as Theorem 3.1,
p.~6. His 2002 paper's final author version,
Theorem 3.3(a),(b), p.~8, also gives numerical
equals $\ell$-adic homological equivalence and the
Hodge standard conjecture for abelian varieties
over algebraically closed fields; Remark 3.5,
p.~9, treats good Weil cohomology theories.
These are cited results, not independently reproved
in this module. The older long-version numbering
\emph{Theorem 4.6} used in the 2026 paper must not
be confused with Theorem 3.3 in the specified final
author PDF.

For comparison, Broe (2026), arXiv:2608.28651v1, Theorem 1.1,
printed p.~1, proves for every abelian fourfold
$A/\mathbb F_q$ and $\ell\ne\operatorname{char}\mathbb F_q$ that
\[
 CH^2(A)\otimes\mathbb Q_\ell\longrightarrow
 H^4_{\rm et}(A_{\overline{\mathbb F}_q},\mathbb Q_\ell(2))
 ^{\operatorname{Gal}(\overline{\mathbb F}_q/\mathbb F_q)}
\]
is surjective. With the divisor theorem and hard
Lefschetz this gives Tate in every codimension on $A$.
Corollary 1.2, p.~2, and Corollary 8.2, p.~20,
give agreement of numerical and $\ell$-adic homological
equivalence for abelian fourfolds over arbitrary fields
where $\ell$ is invertible. Corollary 8.1(i), p.~20,
also gives rational equals numerical equivalence with
rational coefficients over a finite field.
These are cited preprint results whose exact statements
were inspected here; their complete proofs are not
claimed independently checked. The fourfold result
alone does not address all powers of a fourfold;
the October--November all-abelian result does.
It therefore clears the realization issue on the
abelian--Artin category below. It does not prove
Tate for arbitrary smooth projective varieties,
or that every numerical motive belongs to that category.
Milne's Remark 7.6(b), journal p.~76, explicitly
records the extension to homological motives of
abelian type.

\paragraph{AQ-ABELIAN-RECOGNITION.}
Fix $k=\mathbb F_q$, $\ell\ne\operatorname{char}k$,
and let $\mathcal A_q$ be the rational numerical
category obtained from the motives of abelian
varieties over $k$, finite \'etale $k$-schemes, and
Tate motives by finite sums, tensor products,
duals, and direct summands. Equivalently, its
objects are direct summands of finite sums
\[
 h^i(A)(n)\otimes h^0(S),\qquad
 A/k\text{ abelian},\quad S/k\text{ finite \'etale},
 \quad n\in\mathbb Z.
\]
We use geometric Frobenius $F$ and
$\mathbb Q_\ell(1)$ with eigenvalue $q^{-1}$.
The summand above has weight $w=i-2n$ and
parity $i\bmod2=w\bmod2$.
Let $R_\ell(M)$ be its ordinary direct-sum
realization, with the weight decomposition remembered.
Then
\[
 \operatorname{Hom}_{\mathcal A_q}(M,N)\otimes_{\mathbb Q}\mathbb Q_\ell
 \xrightarrow{\ \sim\ }
 \operatorname{Hom}_{\mathbb Q_\ell[F]}
       (R_\ell(M),R_\ell(N)).                         \tag{AQ.15}
\]
Moreover, for $M,N\in\mathcal A_q$,
\[
 M\simeq N\quad\Longleftrightarrow\quad
 P_M(t)=P_N(t),\qquad
 P_M(t)=\det(1-tF\mid R_\ell(M)).                    \tag{AQ.16}
\]
The same equivalence holds with the weight-signed
Euler factor
\[
 Z_M(t)=\prod_w
 \det(1-tF\mid R_\ell(M)_w)^{(-1)^{w+1}}            \tag{AQ.17}
\]
in place of $P_M$. These statements are in the
rational numerical category, and concern actual
objects with their full multiplicities.

\begin{proof}
First we explain the scalar extension in (AQ.15).
For an abelian variety $X/k$, numerical and
$\ell$-adic homological equivalence agree by the
cited all-abelian theorem. Let $A^r_{\rm num}(X)_{\mathbb Q}$
denote the space of rational cycles modulo numerical
equivalence. The numerical intersection pairing
with $A^{\dim X-r}_{\rm num}(X)_{\mathbb Q}$ is nondegenerate.
For any finite linearly independent collection
$\alpha_1,\ldots,\alpha_s$ of numerical classes,
one can choose complementary classes
$\beta_1,\ldots,\beta_s$ making
$(\deg(\alpha_i\beta_j))_{i,j}$ invertible over
$\mathbb Q$. Their cohomology classes are then
linearly independent over $\mathbb Q_\ell$, by
pairing against the $\beta_j$. This also bounds
$s$ by the finite Betti number, proving finiteness
and injectivity after extension to $\mathbb Q_\ell$.
The full Tate theorem supplies surjectivity onto
the Frobenius-invariant subspace. Hence
\[
 A^r_{\rm num}(X)_{\mathbb Q}\otimes\mathbb Q_\ell
 \simeq H^{2r}(X_{\bar k},\mathbb Q_\ell(r))^{F=1}.
                                                               \tag{AQ.18}
\]

The cohomological projectors of an abelian variety
are algebraic here without a further conjecture.
For any integer $a\geq2$, multiplication by $a$
acts on $H^j(A)$ as $a^j$. The polynomials
\[
 e_i=\prod_{\substack{0\leq j\leq2\dim A\\j\ne i}}
       \frac{[a]^*-a^j}{a^i-a^j}
\]
give its projectors in numerical correspondences:
their required identities hold in cohomology,
hence numerically by the established equivalence
on $A\times A$. Tensor products are realized
using $A\times B$, and Poincar\'e duality expresses
duals as complementary cohomological degrees
with Tate twists. This proves the stated description
of $\mathcal A_q$.

Consider first $M=h^i(A)(n)$ and $N=h^j(B)(m)$,
and set $a=\dim A$. A correspondence between them
has codimension $a+m-n$ on $A\times B$. Its
relevant K\"unneth component, by Poincar\'e duality,
is
\[
 H^{2a-i}(A)\otimes H^j(B)(a+m-n)
 \simeq\operatorname{Hom}\bigl(H^i(A)(n),H^j(B)(m)\bigr).
\]
It lies in the required cohomological degree
precisely when $j-i=2(m-n)$, equivalently when
the weights agree. In that case (AQ.18) for
$A\times B$, followed by the algebraic projectors,
gives exactly (AQ.15). For unequal weights every
Frobenius-equivariant map is zero: the eigenvalues
have complex absolute values $q^{(i-2n)/2}$ and
$q^{(j-2m)/2}$, so their spectra are disjoint.
The numerical correspondence group is also zero,
since its cohomological component has the wrong
degree and homological equals numerical equivalence.

For finite \'etale factors $S,T$, the schemes
$A\times B\times S\times T$ are disjoint unions
of abelian varieties viewed over finite extensions
of $k$. A Frobenius-invariant class on a component
orbit is determined by a class invariant under the
appropriate power of Frobenius on one component.
The all-abelian Tate theorem over that finite
extension supplies the cycle. Intersection pairings
are multiplied by the nonzero extension degree,
which is invertible over $\mathbb Q$. Thus the
same injection and surjection apply to these
factors. Passing to finite sums and applying
source and target idempotents proves (AQ.15)
for every object of $\mathcal A_q$.

Frobenius acts semisimply on all these realizations.
For $H^1(A)$ this follows from Tate (1966),
Theorem 2(a), printed p.~140: $\mathbb Q[\pi_A]$
is the centre of the semisimple algebra
$\operatorname{End}_k(A)\otimes\mathbb Q$.
It is a product of number fields, so its Frobenius
element has square-free minimal polynomial;
dualizing or inverting preserves this property.
The identity $H^i(A)=\bigwedge^i H^1(A)$ shows
semisimplicity in every degree. Finite \'etale
Frobenius has finite order, and tensor products,
Tate twists and invariant direct summands of
semisimple single-operator representations remain
semisimple in characteristic zero.

Equality in (AQ.16) therefore gives an isomorphism
$R_\ell(M)\simeq R_\ell(N)$ as
$\mathbb Q_\ell[F]$-modules: each irreducible
polynomial in a semisimple operator's characteristic
polynomial records the multiplicity of its simple
$\mathbb Q_\ell[F]$-module. By (AQ.15) this is
an invertible element of
$\operatorname{Hom}_{\mathcal A_q}(M,N)\otimes\mathbb Q_\ell$.
Choose a rational basis $f_1,\ldots,f_s$ of this
Hom space. The polynomial
\[
 D(x_1,\ldots,x_s)=
 \det R_\ell\!\left(\sum_\nu x_\nu f_\nu\right)
\]
over $\mathbb Q_\ell$ is not zero. Since
$\mathbb Q$ is infinite, some rational tuple
has $D\ne0$; this elementary fact follows by
induction on the number of variables, using the
finite number of roots of a nonzero polynomial
in one variable. It gives a rational morphism
$f:M\to N$ with invertible realization.
By full faithfulness its realization's inverse
comes from a morphism in
$\operatorname{Hom}_{\mathcal A_q}(N,M)\otimes\mathbb Q_\ell$.
The equations $gf=1_M$ and $fg=1_N$ form a
finite linear system over $\mathbb Q$ in the
coefficients of $g$. A rational linear system
solvable after extension to $\mathbb Q_\ell$
is solvable over $\mathbb Q$, by row reduction.
Thus $f$ has a rational inverse. The converse
in (AQ.16) is immediate.

Finally a reciprocal root contributed by weight
$w$ has complex absolute value $q^{-w/2}$.
Weights of different parity cannot contribute
the same root. Hence the numerator and denominator
of (AQ.17) are coprime, each with constant term
one. Equality of those rational functions
recovers their numerator and denominator separately,
and their product is $P_M(t)$. This reduces
(AQ.17) to (AQ.16) and completes the proof.
\end{proof}
No essential-surjectivity claim onto all Frobenius
representations is needed here. In particular,
this recognition theorem does not identify an
arbitrary smooth projective motive with an object
of $\mathcal A_q$, does not retain integral lattices,
and does not classify mixed extensions. It is a
genuine strengthened answer to the abelian sector
of cote 15, supplied by the new cycle-algebraicity
theorem rather than by analytic zero location.

\paragraph{AQ-RATIONAL-FIBRE: an explicit neutralization obstruction.}
Even this now-established abelian category need not
admit a fibre functor with values in rational vector
spaces. This is a concrete restriction on any reading
of the archive's motivic Galois group as a neutral
group over $\mathbb Q$. It is not a refutation of a
group defined after scalar extension or of a motivic
gerbe formulation.
\begin{proof}
Consider $E:y^2=x^3-x$ over $\mathbb F_3$, then
base-change it to $k=\mathbb F_9$. It is smooth:
the derivative of $x^3-x$ is $-1$ in characteristic
three. For each of $x=0,1,2$, the right side is
zero, so there is one affine point; with infinity,
$\#E(\mathbb F_3)=4$. Thus its $3$-Frobenius
polynomial on $H^1$ is $T^2+3$. The square of
Frobenius is $-3$ on $H^1$, and injectivity of
the Tate-module action of endomorphisms implies
that the $9$-Frobenius endomorphism is $[-3]$.
Tate (1966), Theorem 2(d), p.~140, identifies
\[
 D=\operatorname{End}_{\mathbb F_9}(E)\otimes\mathbb Q
   =D_{3,\infty},
\]
the quaternion division algebra ramified at $3$
and infinity. Alternatively, the main theorem
on p.~134 gives $\dim_{\mathbb Q}D=4$, since
Frobenius is scalar on a two-dimensional Tate
module; every nonzero endomorphism of an elliptic
curve is an isogeny, so this algebra is a division
algebra.

Put $M=h^1(E)$ in $\mathcal A_9$. Pullback
endomorphisms give an injective algebra map
$D^{\rm op}\to\operatorname{End}(M)$.
By (AQ.15), the target has rational dimension
four, since its scalar extension is
$M_2(\mathbb Q_\ell)$. This map is therefore
an isomorphism, so $\operatorname{End}(M)$ is
a quaternion division algebra.

Use the usual parity-corrected symmetry on
numerical motives, for which the ordinary
$\ell$-adic realization is a symmetric tensor
functor. The categorical trace of $1_M$ is
then $2$: its image is the ordinary trace of
the identity on $H^1(E,\mathbb Q_\ell)$,
and the scalar map $\mathbb Q\to\mathbb Q_\ell$
is injective. A rational fibre functor would
preserve this trace, hence send $M$ to a
two-dimensional rational vector space. Its
faithfulness would embed the four-dimensional
division algebra $\operatorname{End}(M)$ into
$M_2(\mathbb Q)$. Equal dimensions force this
embedding to be an isomorphism. That is
impossible because $M_2(\mathbb Q)$ has
nonzero noninvertible elements. This proves
the asserted obstruction.

It also exhibits an object restriction invisible
in a list of eigenvalues alone. Here
$P_M(t)=(1+3t)^2$. There cannot be an object
$L\in\mathcal A_9$ with $P_L(t)=1+3t$:
otherwise (AQ.16) would give $L\oplus L\simeq M$,
but the former has a nontrivial idempotent
endomorphism and the latter's division algebra
has none. The eigenvalue $-3$ occurs, as
AQ-WEIL-EXIST predicts, but its multiplicity
cannot be discarded in a rational motivic
classification.
\end{proof}

\paragraph{AQ-WEIL-NOTATION: the recovered $Q'$ and the norm-one subgroup.}
The official cote 15 PDF has now been inspected beyond the
supplied excerpts.  Physical page 94, printed footer 93,
opens with the explicit definition
\emph{``$Q'$: cl\^oture tot[alement] r\'eelle de $\mathbb Q$''};
the bracket expands the handwritten abbreviation.
Thus $Q'=\mathbb Q^{\rm tr}$, the field of totally real
algebraic numbers, and the $i$ in $Q'(i)$ means
$\sqrt{-1}$, distinct from the integral weight index.
Physical page 98/footer 97 discusses a totally imaginary
quadratic extension of a totally real field and writes
$\lambda\overline\lambda=q^w$, followed by the
normalization $\lambda=q^{w/2}u$ with
$u\overline u=1$.
Physical page 100/footer 99 contains the boxed correction
requiring every conjugate to have absolute value $q^{w/2}$,
as already preserved in the supplied transcription.
These recovered norm and purity conditions identify the
weight-zero interpretation used in AQ-WEIL-LIMIT.

More precisely, let
\[
 \mathbb Q^{\rm CM}=\bigcup_{L\ {\mathrm{CM}}}L
 \quad\hbox{inside the fixed }\overline{\mathbb Q}\subset\mathbb C.
\]
Then
\[
 \mathbb Q^{\rm CM}=\mathbb Q^{\rm tr}(\sqrt{-1}).
 \tag{AQ.WN.1}
\]
For a prime power $q=p^n$, the group of algebraic
$p$-units with all conjugate absolute values $q^{w/2}$,
for some integer $w$, is exactly
\[
 \left\{\lambda\in\mathbb Q^{\rm tr}(\sqrt{-1})^*:
   \lambda\text{ is a $p$-unit},\quad
   \lambda\overline\lambda=q^w\text{ for some }w\in\mathbb Z
 \right\}.                                                   \tag{AQ.WN.2}
\]
Its weight-zero subgroup is $W_0(p)$, independent of $n$.
Thus the historical group $M(q)$ and its weight-zero part
agree with $W(q)$ and $W_0(p)$ when the explicit
norm/purity conditions on pages 98 and 100 are retained.

\begin{proof}
For a CM field $L$, let $L^+$ be its totally real
subfield.  Choose $b\in L$ with
$\delta=b-\overline b\ne0$.  Then
$\delta^2\in L^+$, and $a=-\delta^2$ is totally positive:
every embedding sends $\delta$ to a nonzero purely
imaginary number.  The algebraic number $\sqrt a$
is totally real, because all its conjugates are real
square roots of positive conjugates of $a$.
Hence
$L=L^+(\delta)\subset\mathbb Q^{\rm tr}(\sqrt{-1})$.
Conversely any element of the latter field belongs to
$F(\sqrt{-1})$ for a finite totally real field $F$,
and $F(\sqrt{-1})$ is a CM field.  This proves (AQ.WN.1),
including the assertion that the displayed union is a field.

Let $\lambda$ have the absolute values in (AQ.WN.2),
and let $L$ be the Galois closure of $\mathbb Q(\lambda)$.
Write $c$ for complex conjugation on $L$.
For every $\sigma\in\operatorname{Gal}(L/\mathbb Q)$,
\[
 c(\sigma\lambda)=q^w/(\sigma\lambda).
\]
For every $\tau$ in the same Galois group,
$c\tau$ and $\tau c$ consequently agree on all the
conjugates $\sigma\lambda$, which generate $L$.
Thus $c$ is central.
If $c=1$, the Galois field $L$ is totally real.
Otherwise $L^+=L^{\langle c\rangle}$ is totally real,
$[L:L^+]=2$, and $L$ has no real embedding, so $L$ is CM.
In either case (AQ.WN.1) places $\lambda$ in
$\mathbb Q^{\rm tr}(\sqrt{-1})$; the norm identity follows
from its specified absolute value at the fixed embedding.

Conversely this field is Galois over $\mathbb Q$, and
complex conjugation fixes $\mathbb Q^{\rm tr}$ and
negates $\sqrt{-1}$.  It therefore commutes with every
embedding.  Applying any embedding to
$\lambda\overline\lambda=q^w$ gives
$|\sigma\lambda|^2=q^w$, as required.
The $p$-unit condition is the same on both sides,
so (AQ.WN.2) follows.  Setting $w=0$ gives the claimed
weight-zero subgroup, independently of $q=p^n$.
\end{proof}

Finite-prime valuations alone must not be substituted
for the norm/purity condition in this identification.
For example $\varepsilon=2+\sqrt3$ and its inverse
$2-\sqrt3$ are algebraic integers, so $\varepsilon$
is a unit at every finite place.  In particular every
finite-prime sum $v_{\mathfrak p}(\varepsilon)+
v_{\overline{\mathfrak p}}(\varepsilon)$ is zero,
for every chosen $p$.  But
$\varepsilon\overline\varepsilon=\varepsilon^2\ne1$
and its two real conjugates have unequal absolute values.
No positive power is a weight-zero Weil number.
This checks the limitation of a \emph{valuation-only}
paraphrase of the margin on 15/102; it does not assert
that the manuscript's terminology \emph{unit\'e absolue}
itself discards archimedean conditions.
The corrected group interpretation uses the explicit
norm condition actually present on the preceding sheet.

\paragraph{AQ-WEIL-LIMIT: the erased question at 15/104.}
The text \emph{Trouve-t-on $P^0/U$ tout entier?}
and its subsequent root-extraction test are inside a
deletion in the supplied transcription. In the standard
Weil-number interpretation its algebraic assertion is true.
Let $W_0(p)$ be the multiplicative group of algebraic
$p$-units whose complex conjugates all have absolute
value one. Then the directed system with transition
maps $u\mapsto u^m$ has limit
\[
 \varinjlim W_0(p)=W_0(p)\otimes_{\mathbb Z}\mathbb Q
                 \simeq W_0(p)/\mu_\infty .
\]
\begin{proof}
For any abelian group $A$, the system indexed by
positive integers, with multiplication by $m$ from
index $n$ to index $mn$, has colimit
$A\otimes\mathbb Q$: the compatible map sends
the class of $a$ at index $n$ to $a/n$.
Every root of an element of $W_0(p)$ is again
a $p$-unit with all conjugate absolute values one.
Thus $W_0(p)$ is divisible, and its torsion subgroup
is exactly $\mu_\infty$. Its quotient by torsion
is uniquely divisible, giving the last isomorphism.
\end{proof}
The recovered original definition $Q'=\mathbb Q^{\rm tr}$
and the preceding explicit norm-one condition identify
this weight-zero interpretation by AQ-WEIL-NOTATION.
The literal deleted calculation therefore has the stated
affirmative answer with its norm/purity conditions retained.
Its identification with the motivic group of arbitrary
motives still requires the separate Tate, generation,
fibre and base-field transition conventions described above;
the abelian realization sector has been cleared by
AQ-TATE-2026 and CC-AG01--03.

\paragraph{Coefficient-extended group: the completed abelian-sector answer.}
The companion proof CC-AG01--03, in
\texttt{research/consequences\_criteria\_abelian\_group.tex},
now strengthens the rational recognition result to a full
category and group description. Fix
$E=\overline{\mathbb Q}_\ell$ and an embedding
$\overline{\mathbb Q}\hookrightarrow E$. Let $W(q)$ be the
group of algebraic $p$-units whose conjugate absolute values
are $q^{w/2}$ for one integer $w$. With the parity-corrected
symmetry and idempotent completion after extending Hom spaces,
the result is
\[
 \operatorname{Kar}(\mathcal A_q\otimes_{\mathbb Q}E)
  \simeq \operatorname{Gr}^{\mathrm{fin}}_{W(q)}(E),\qquad
 \operatorname{Aut}^{\otimes}(\omega_E)=D_E(W(q)).
\]
The target means finite-dimensional vector spaces graded by
$W(q)$, with multiplication of labels as tensor product;
$D_E(M)=\operatorname{Spec}E[M]$ for an abelian group $M$.
Choosing $\sqrt q$ gives
$D_E(W(q))\simeq D_E(W_0(p))\times\mathbb G_{m,E}$,
with Tate character $(1,-2)$. The full proof uses the
Hom-space theorem proved above, Honda occurrence, and
lifting of the individual eigenspace projectors over $E$.
It also proves the distinguished Frobenius has Zariski-dense
powers. The finite-field group is not a protorus, since its
character group contains $\mu_\infty$. This answers the
abelian-sector group question with explicit coefficient and
fibre conventions, while preserving the rational quaternion
obstruction and distinguishing passage to
$\overline{\mathbb F}_p$ from extension of coefficients.
It does not assert that all numerical motives are of
abelian type or identify the fixed-finite-field category
with the separate category of Frobenius germs over
$\overline{\mathbb F}_p$.  The previously missing field
notation is now identified by AQ-WEIL-NOTATION.

\subsection{AQ-TWIST: the weight-index correction at 18/009}
The sheet's conventions give
\[
 W_j(V(n))=W_{j+2n}(V),\quad F^j(V(n))=F^{j+n}(V),\quad
 \operatorname{gr}^W_j(V(n))
 =(\operatorname{gr}^W_{j+2n}V)(n).                             \tag{AQ.6}
\]
The last equality follows by taking the quotient of the
first two consecutive weight steps, with the shifted Hodge
filtration. The transcribed unshifted graded formula fails
for $V=\mathbb Q(0)$, $n=-1$, $j=2$: the left side is
$\mathbb Q(-1)$ and the unshifted right side is zero.

\subsection{AQ-REAL: the even-Tate assertion at 18/038 is false}
The sheet claims that $f_{\infty,M}=1$ forces a sum of
tensor powers of $\mathbb Q(2)$.
Let $A$ be the sign Artin motive over $\mathbb R$ and
$L=A\otimes\mathbb Q(-1)$. Its real involution is
$(-1)(-1)=1$, but its weight $2$ is not the weight $-4r$
of any $\mathbb Q(2)^{\otimes r}$. This refutes the claim.
In the semisimple real Artin--Tate category the correct
objects are exactly finite sums of $L^{\otimes r}$,
$r\in\mathbb Z$.
\begin{proof}
The simple objects are $A^\epsilon\otimes\mathbb Q(n)$,
$\epsilon\in\{0,1\}$, with real involution
$(-1)^{\epsilon+n}$. It is trivial exactly when
$\epsilon\equiv n\pmod2$, in which case this object is
$L^{\otimes(-n)}$. Direct-sum decomposition proves both
directions.
\end{proof}
The classification is asserted here only for Artin--Tate
motives. Extending it from Hodge types to all pure motives
would require an additional algebraicity input.

\subsection{AQ-PERIOD-STRONG: a counterexample to 18/011}
The strong transcendence conjecture on 18/011 predicts
$\operatorname{trdeg}_k k(\alpha,\beta)=g+u$.
In the usual reading $T_{\rm Hdg}=\operatorname{gr}_F T_{\rm dR}$
and $Q$ parametrizes normalized splittings, this fails.
Take $E:y^2=x^3-x+1$ over $\mathbb Q$.
Its $j=-6912/23$ is not integral, so it is non-CM;
for $M=h^1(E)$ one has $G=\mathrm{GL}_2$, $g=4$, $u=1$.
Here we use the classical integrality of CM
$j$-invariants, recalled in Habegger (2014),
\S1, p.~1, and the elliptic motivic-group
identification in Kreutz--Shen--Vial,
Proposition 5.14.
Here $\alpha$ is the comparison from Betti to de Rham
cohomology, $T_{\rm Hdg}=\operatorname{gr}_F T_{\rm dR}$,
$\beta$ is the Hodge splitting normalized to induce the
identity on that associated graded, and
$\gamma=(\beta\alpha)^{-1}$, as in the source triangle.
The parabolic preserving its Hodge line has
one-dimensional unipotent radical, explaining $u=1$.
\begin{proof}
Choose rational Betti eigenvectors for the real involution
and a rational de Rham basis with Hodge line $\mathbb Ce_1$.
Write
\[
 P=\begin{pmatrix}a&b\\c&d\end{pmatrix},\quad
 \delta=ad-bc\ne0,\quad C=\operatorname{diag}(1,-1).
\]
The opposite Hodge line is $PCP^{-1}\mathbb Ce_1$,
generated by $(ad+bc,2cd)^t$. Both $c,d$ are nonzero:
otherwise a real-involution eigenline would be the
Hodge line, forcing it to equal its opposite.
Hence the normalized splitting is
\[
 \beta=\begin{pmatrix}1&-(ad+bc)/(2cd)\\0&1\end{pmatrix}.
\]
It is rational in $P$, so $k(P,\beta)=k(P)$ has
transcendence degree at most four, not $g+u=5$.
Moreover
\[
 \gamma=(\beta P)^{-1}
 =\begin{pmatrix}d/\delta&1/(2c)\\-c/\delta&1/(2d)\end{pmatrix},
 \qquad \gamma_{11}\gamma_{22}+\gamma_{12}\gamma_{21}=0.
\]
Thus its transcendence degree is at most three,
also contradicting the predicted value $g=4$.
\end{proof}
This refutation does not settle the ordinary non-CM
elliptic period conjecture or algebraicity of its
Hodge-splitting scalar.

\subsection{AQ-PERIOD-GENERATION: 18/013(b) is a known theorem}
Question (a) asks for period-field transcendence degree
$\dim G$, namely four without CM or two with CM.
Question (b) asks whether the group generated by a
translate of the period locus is all of $G$.
The answer to (b) for one elliptic curve is yes:
Kreutz, Shen, and Vial (2026), \emph{De Rham--Betti classes
with coefficients}, arXiv:2206.08618v3, Proposition 3.1
and Theorem 5.13(ii) identify the generated period torsor
with the motivic torsor; Proposition 5.14 identifies
the elliptic motivic group. The coefficient restriction
in Theorem 5.13(ii) is satisfied for a single elliptic
curve. This known result is cited, not reproved.
It is weaker than density of the period locus.

Question (a) is known for CM curves by Chudnovsky (1980),
\S2A, Theorem 1 and Corollary 2, printed p.~343
in the 1978 Helsinki ICM proceedings. In particular,
a nonzero CM period and $\pi$ are algebraically
independent. The non-CM case remains the full
period-conjecture gate. Question (c) asks whether the
Hodge-splitting scalar can be algebraic without CM.
In the normalization $y^2=4x^3-g_2x-g_3$,
put $\omega=dx/y$, $\xi=x\,dx/y$, take full periods
$\omega_1,\omega_2$ with $\operatorname{Im}
(\omega_2/\omega_1)>0$, and write the periods of
$\xi$ as $-\eta_1,-\eta_2$. The opposite Hodge line
is generated by $\xi+h\omega$. Its periods must be
proportional to $\overline{\omega_1},
\overline{\omega_2}$, which gives
\[
 h=\frac{\eta_1\overline{\omega_2}-\eta_2\overline{\omega_1}}
          {\omega_1\overline{\omega_2}-\omega_2\overline{\omega_1}}
   =\frac{\pi^2}{3\omega_1^2}E_2^*(\tau).
\]
Here $\tau=\omega_2/\omega_1$ and
$E_2^*(\tau)=E_2(\tau)-3/(\pi\operatorname{Im}\tau)$.
The second equality uses
$\eta_1\omega_2-\eta_2\omega_1=2\pi i$ and
$\eta_1\omega_1=\pi^2E_2(\tau)/3$.
The CM conclusions in (a) and (c) can be separated
from the remaining non-CM problem.
\begin{proof}[The CM period field and algebraicity of the splitting]
Take a nonreal CM endomorphism $u$, defined over
$\overline{\mathbb Q}$. In the algebraic de Rham
basis its action has the form
$u^*\omega=u\omega$ and
$u^*\xi=a\omega+\overline u\xi$, with
$a\in\overline{\mathbb Q}$. The opposite Hodge line
is the $\overline u$-eigenline, so
\[
 a+hu=\overline u h,\qquad
 h=\frac{a}{\overline u-u}\in\overline{\mathbb Q}.
\]
The period ratio $\tau$ is imaginary quadratic.
The preceding formula for $h$ and the Legendre
relation give
\[
 \omega_2=\tau\omega_1,\qquad
 \eta_1=h\omega_1+\frac{\pi}{\omega_1\operatorname{Im}\tau},
 \qquad \eta_2=\tau\eta_1-\frac{2\pi i}{\omega_1}.
\]
Thus the full period and quasi-period field is
$\overline{\mathbb Q}(\omega_1,\pi)$; the reverse
inclusion follows from the Legendre relation.
Chudnovsky's cited Corollary 2 makes its
transcendence degree exactly two.
\end{proof}
The remaining non-CM assertion is the nonholomorphic weight-two
Eisenstein-transcendence question attributed to Katz
in Waldschmidt (2004), p.~277, immediately after
Conjecture 3.22. That passage states the same remaining
condition for a non-CM lattice with algebraic invariants.
Its non-CM transcendence is not established by the
supplied proofs. AQ-FINITE cannot supply that missing
algebraic-independence assertion.

\paragraph{AQ-PERIOD-2026: current-source and normalization check.}
The current primary literature does not identify the
three different weight-two functions in these questions.
Kreutz--Shen--Vial, arXiv:2206.08618v3, Introduction
and \S4.2, retain the full period conjecture separately
from their generated-torsor theorem. Mohajer (2026),
arXiv:2611.25399v1, \S1.1, equation (1.3), pp.~2--3,
still formulates the non-CM complex four-period
transcendence-degree assertion as conjectural.
His Theorems 4.2 and 4.4 concern comparison periods
in $B_{\rm dR}$ and their normalized images in
$\mathbb C_p$, respectively. His Proposition 2.1
concerns zeros of the holomorphic $E_2$, whereas
the scalar above involves the almost holomorphic
$E_2^*=E_2-3/(\pi\operatorname{Im}\tau)$.

Sebbar and Besrour (2026), arXiv:2611.14502v1,
\S8.4, Remark 8.11, Proposition 8.12 and
Question 8.13, pp.~32--33, consider the ordinary
$p$-adic Katz form $E^{\rm Katz}_2$. They leave
the existence of a non-CM zero with algebraic
$j$-invariant over $\mathbb Q$ as a question;
their nonvanishing proposition assumes the crystalline
period conjecture of Ancona--Fr\u a\c til\u a.
Algebraicity over a finite extension of $\mathbb Q_p$
does not answer algebraicity over $\mathbb Q$.
Thus these November 2026 results do not supply the
missing complex algebraic-independence assertion in
18/013(a), or the non-CM algebraicity classification of
the complex splitting scalar in 18/013(c). This is a
source-scope conclusion, not a proof of logical
nonimplication from the adopted analytic theorems.

\paragraph{AQ-CM-PERIOD-2026: the new rank-one comparison results.}
Kahn and Otsubo (2026), arXiv:2611.13723v1,
Corollary 3.21, p.~24, show that equal Hecke characters
give equal complex period classes for rank-one abelian-type
$K$-coefficient motives over a number field $k\supset K^{\rm gal}$.
Here the class lies in
$(\mathbb C\otimes_{\mathbb Q}K)^*/(k\otimes_{\mathbb Q}K)^*$
(\S2.4, pp.~6--7).
Remark 3.22, p.~25, relates the converse, up to a finite-order
character, to Rohrlich's conjecture.
For rank-one algebraic motives, equality of crystalline period
classes at almost all places implies equality of Hecke characters;
the converse holds under the preceding abelian-type and base-field
hypotheses (Corollary 3.35, pp.~30--31).
Corollary 3.25, p.~26, proves Mumford--Tate when the rank-one
$K$-motive becomes algebraic after finite extension.
For an abelian CM field $K$ of conductor $m$ and
$k\supset\mathbb Q(\mu_m)$,
Theorems 4.36 and 4.39, pp.~46--47,
refine Gross--Deligne, and Theorem 4.44, p.~49,
gives the corresponding $p$-adic beta-period formula.
These are positive comparison results with specified hypotheses.

For a non-CM elliptic curve, $h^1(E)$ has rational rank two
and endomorphism algebra $\mathbb Q$.  A rank-one structure
over a coefficient field would require a quadratic field to
embed in that endomorphism algebra, which is impossible.
The determinant is a Tate object and records only the
determinant period.  Thus these rank-one statements do not
give the four-period algebraic independence or the
complex $E_2^*$ assertion in 18/013.

\paragraph{AQ-KO-RANK: a correction to an ancillary dimension claim.}
With Definition 2.5, p.~8, of the same paper, the dimension
sentence in Lemma 4.6, p.~34, requires one additional rank:
for $K=\mathbb Q(\mu_m)$ with $m>2$,
\[
 \operatorname{rank}_{\mathbb Z}\Xi(K)
       =\frac{\varphi(m)}2+1,
 \tag{AQ.KO.1}
\]
where the paper prints $\varphi(m)/2$.
\begin{proof}
Put $r=\varphi(m)/2$ and choose one embedding
$\sigma_j$ from each of the $r$ complex-conjugate pairs.
Since $K$ itself is a CM field, the additional CM-subfield
condition in that definition is automatic.
A function $\phi:\Sigma_K\to\mathbb Z$ belongs to
$\Xi(K)$ precisely when all the sums
$\phi(\sigma_j)+\phi(\overline\sigma_j)$ equal
one common integer $w$.
Consequently
\[
 \Xi(K)\longrightarrow\mathbb Z^{r+1},\qquad
 \phi\longmapsto(w,\phi(\sigma_1),\ldots,\phi(\sigma_r))
\]
is an isomorphism of abelian groups: the inverse sends
$(w,a_1,\ldots,a_r)$ to the function with values
$a_j$ and $w-a_j$ on the $j$th pair.
This proves (AQ.KO.1).  The map
$\phi\mapsto2\phi-w$ used in the cited rank argument
has the full group of constant functions as its kernel,
rather than being injective.  Already $m=3$ gives
$\Xi(\mathbb Q(\mu_3))=\mathbb Z^2$.
\end{proof}
This calculation concerns that rank sentence; the
period-comparison results above are cited with their
own hypotheses, and their complete external proofs
are not certified by this source-scope check.

\subsection{AQ-COCYCLE: the image-density statement at 18/059}
Let $P$ be an affine torsor under an algebraic group $G$
over $F\subset\mathbb C$, and $p\in P(\mathbb C)$.
The closure of the image of its automorphism cocycle
is the translated $F$-locus of $p$. Thus the image is
Zariski dense in $G_{\mathbb C}$ exactly when $p$ is
generic in $P$ over $F$.
\begin{proof}
Embed $P$ in affine space over $F$, and let $I$ be the
ideal over $\mathbb C$ of the
$\operatorname{Aut}(\mathbb C/F)$-orbit of $p$.
Each finite-dimensional degree-bounded part of $I$
is stable under the coefficientwise semilinear action.
Its unique reduced-row-echelon basis, relative to an
ordered monomial basis, must therefore have entries
fixed by every automorphism. The fixed field of
$\operatorname{Aut}(\mathbb C/F)$ is $F$: any nontrivial
embedding of $F(a)$ over $F$, for $a\notin F$, extends
first across a transcendence basis and then an algebraic
closure to an automorphism of $\mathbb C$.
Consequently every degree-bounded part has a basis
over $F$, and $I=(I\cap F[\mathbf x])\mathbb C[\mathbf x]$.
An $F$-polynomial vanishes on the orbit exactly when
it vanishes at $p$. Thus $I$ is precisely the ideal of
$Z_F(p)_{\mathbb C}$, proving the density assertion.
The torsor difference map $D_p:P_{\mathbb C}\to G_{\mathbb C}$
is an algebraic isomorphism and the cocycle is
$D_p(\sigma p)$, up to inversion according to convention.
Its closure is consequently $D_p(Z_F(p)_{\mathbb C})$.
\end{proof}
The generated-subgroup version is false for general
torsors. In $G=\mathbb G_m^2$ take $F=\mathbb Q$ and
$p=(t,t+1)$ with $t$ transcendental. Its locus is a curve.
The translated cocycle curve is
$u\mapsto(u/t,(u+1)/(t+1))$.
A character $x^a y^b$ identically one there would have
orders zero at $u=0$ and $u=-1$, giving $a=b=0$.
Thus this proper curve generates the whole torus.
This refutes an abstract torsor inference, without
claiming an arithmetic-motivic counterexample to the
period conjecture.

\subsection{AQ-CM: the hexagonal elliptic curve and a genuine GRH consequence}
The hexagonal real-elliptic classification recorded in SA22,
\S2(A.7), concerns the $j=0$ period lattice and its real
forms. It is not the Grothendieck--Teichm\"uller hexagon.
For example, the latter is the identity
\[
 f(z,x)z^m f(y,z)y^m f(x,y)x^m=1,\qquad
 z=(xy)^{-1},\quad m=(\lambda-1)/2,
\]
in a profinite free group, as in Boavida de Brito, Horel,
and Robertson (2019), Definition 7.2(II). The two uses of
\emph{hexagonal} name different mathematical data.

\paragraph{Hodge obstruction.}
For $E:y^2=x^3+1$, the motive $h^1(E)$ is not Artin--Tate,
even after adjoining its CM field $\mathbb Q(\sqrt{-3})$.
\begin{proof}
An Artin motive has Hodge type $(0,0)$, and $\mathbb Q(n)$
has type $(-n,-n)$ in our convention.
Tensor products, sums, subobjects, and quotients of such
Hodge structures have only diagonal types. But
$H^1(E)$ has nonzero types $(1,0)$ and $(0,1)$.
Base extension to an algebraic number field does not
change these complex Hodge types. This proves the claim.
\end{proof}
Thus the fact that the CM field is abelian does not turn
its weight-one Hecke character into a finite-order
Dirichlet character with an integral Tate twist.
No conclusion about the zeros of $L(h^1(E),s)$ is obtained
merely by substituting that field in the supplied
Dirichlet theorem.

\paragraph{Effective distribution of trace-zero primes.}
Let $a_p=p+1-\#E(\mathbb F_p)$, for $p>3$, and put
$N_0(x)=\#\{3<p\leq x:a_p=0\}$. Then
\[
 a_p=0\ \Longleftrightarrow\ p\equiv2\pmod3,\qquad
 N_0(x)=\tfrac12\operatorname{Li}(x)+O(\sqrt x\log x)            \tag{AQ.7}
\]
under the adopted Dirichlet GRH. The implied constant is
effective.
\begin{proof}
With the quadratic character extended by $\chi_p(0)=0$,
\[
 a_p=-\sum_{x\in\mathbb F_p}\chi_p(x^3+1).
\]
If $p\equiv2\bmod3$, cubing permutes $\mathbb F_p$,
so the sum is $\sum_u\chi_p(u+1)=0$.
If $p\equiv1\bmod3$, set $m=(p-1)/2$.
Euler's criterion and the binomial theorem show
\[
 a_p\equiv-\sum_{j=0}^m {m\choose j}
                 \sum_{x\in\mathbb F_p}x^{3j}
       ={m\choose (p-1)/3}\pmod p.
\]
The $j=0$ sum is $p=0$ modulo $p$.
Among the positive exponents $3j\leq3(p-1)/2$, the
only multiple of $p-1$ is $p-1$, and its power sum is
$-1$ modulo $p$. The displayed binomial coefficient
is nonzero modulo $p$ because all its factorial arguments
are less than $p$. This proves the equivalence.

Now $\zeta_{\mathbb Q(\sqrt{-3})}(s)=
\zeta(s)L(s,\chi_{-3})$, so the adopted theorem gives
GRH for this quadratic Dedekind zeta function.
Lagarias--Odlyzko (1977), Theorem 1.1, equation (1.3),
applied to its nontrivial Frobenius class, gives
\[
 \left|\pi(x;3,2)-\tfrac12\operatorname{Li}(x)\right|
 \leq c_1\left(\tfrac12\sqrt x\log(3x^2)+\log3\right),
 \qquad x>2.
\]
Here $c_1$ is their effective absolute constant.
For $x\geq3$, $N_0(x)=\pi(x;3,2)-1$ because $p=2$
is a bad prime of $E$. This proves (AQ.7).
\end{proof}
This gives an effective arithmetic statement about the
hexagonal curve using the quadratic splitting character.
It does not replace that curve's weight-one Hecke
$L$-function by a Dirichlet $L$-function.

\subsection{AQ-67-RANGE: the global isomorphism in 67/078 is false}
The literal question is \emph{Ces fl\`eches sont elles
des iso --- et si oui, comment expliciter un iso en
sens inverse ?} In the stable nodal $(0,4)$ case,
normalize both three-pointed rational components with
the node at $0$ and the other two marked points at
$1,\infty$. In the tangent-tensor normalization of
the capping construction, a tensor $q$ in the image
has absolute value equal to the product of the two
node-cap conformal radii. Consequently
\[
 |q|\leq16,\qquad q=17\ \hbox{has no preimage}.                  \tag{AQ.8}
\]
If the target is written with dual tangent tensors,
the equivalent statement uses the inverse tensor.
This normalization is fixed by the original sheets:
67/075 calls the propagation datum a nonzero basis of
$T_{x'}\otimes T_{x''}$; 67/078 writes its inverse
as a dual pairing and asks for mutually polar disks
in the two tangent lines. In coordinates, polarity
is $|q^{-1}|r'r''=1$, hence $r'r''=|q|$.
The target in 67/076--077 imposes no bound on this
nonzero tensor. We use the map without additional
marked-point disks in 67/077.

\begin{proof}
For completeness, the Koebe bound needed here follows
from the area argument. If
$F(z)=z+b_0+\sum_{n\geq1}b_nz^{-n}$ is univalent for
$|z|>1$, Green's formula gives the area inside
$F(|z|=r)$ as
$\pi(r^2-\sum_{n\geq1}n|b_n|^2r^{-2n})\geq0$.
Letting $r\downarrow1$ proves
$\sum_{n\geq1}n|b_n|^2\leq1$.
For a univalent $f(z)=z+a_2z^2+\cdots$ in the unit disk,
$g(z)=\sqrt{f(z^2)}=z+(a_2/2)z^3+\cdots$ is an odd
univalent function there. The branch exists because
$f(z^2)/z^2$ never vanishes in the disk. Equality of
two $g$-values forces equality of the squared
arguments under $f$, and oddness then proves
injectivity. Its reciprocal
$1/g(1/z)$ is univalent outside the disk, so the area
inequality gives $|a_2|\leq2$.
If $w\ne0$ is omitted by $f$, the normalized univalent
function $f/(1-f/w)$ has second coefficient $a_2+1/w$.
Thus $|1/w|\leq|a_2|+|a_2+1/w|\leq4$.
This is the Koebe quarter theorem (Duren, 1983,
Theorem 2.3, p.~31).

A node cap is the image of a univalent map $h$ from
the unit disk with $h(0)=0$ and conformal radius
$r=|h'(0)|$. It omits $1$ and $\infty$.
Apply the preceding conclusion to $h/h'(0)$ and
the omitted value $1/h'(0)$ to get $r\leq4$.
Both caps therefore have radii at most four.
The tensor is the product of the two derivative
vectors, up to the unit-modulus boundary identification,
so its norm is at most sixteen. The nonzero tensor
$17$ belongs to the stated unrestricted target,
establishing failure of surjectivity.
\end{proof}
The intake SA22 reports the sharper exact image
$0<|q|<16$, tangent invertibility, and a $\pi_1$-isomorphism.
Its full prior proofs are not in the supplied packet.
The next result independently supplies complete proofs
for the labelled geodesic-cut $(0,4)$ sector, while
AQ-67-NONHOLO separates the real and complex tangent claims.
The elementary obstruction (AQ.8) already answers
67/078 negatively and is independent of RH.

\subsection{AQ-67-EXACT: the canonical radius and the exact labelled $(0,4)$ image}
\paragraph{Scope and normalization.}
Consider the geodesic-cut interpretation of the map without marked-point
disks in cote 67/077. Fix a partition of four labelled punctures into two
pairs and remember the separating geodesic. Its length is $\ell>0$.
Cutting produces two hyperbolic pairs of pants of boundary lengths
$(\ell,0,0)$, where zero denotes a cusp. Cap the geodesic boundary by the
unit disk using its normalized hyperbolic arclength, as prescribed in
67/074. On each resulting three-pointed sphere use the unique coordinate
with cap centre $0$ and the two labelled cusps $1,\infty$.
Let $r(\ell)>0$ be the conformal radius of the cap at zero in that
coordinate. A cap coordinate $w$ has $|w|=1$ at the seam and
$z=r(\ell)w+O(w^2)$ after a rotation of $w$.
The tensor coordinate $q$ of 67/075 then has modulus $r(\ell)^2$.

Put
\[
 x=\frac{\ell}{4\pi},\qquad
 \phi(x)=\operatorname{Im}\bigl(\log\Gamma(1+ix)
                       -\log\Gamma(\tfrac12+ix)\bigr),
\]
where the analytic logarithm of $\Gamma$ on the right half-plane is
real on the positive real axis. Then
\[
 \boxed{\quad r(\ell)=4\exp\!\left(\frac{\phi(x)-\pi/4}{x}\right).\quad}
                                                               \tag{AQ.67.1}
\]
The function $r:(0,\infty)\longrightarrow(0,4)$ is a real-analytic
diffeomorphism. With the orientation of the normalized twist angle
$\theta\in\mathbb R/2\pi\mathbb Z$ chosen so that its increase multiplies
the tangent tensor by $e^{i\theta}$, the capping map is
\[
 (\ell,\theta)\longmapsto r(\ell)^2e^{i\theta}.
                                                               \tag{AQ.67.2}
\]
It is a real-analytic diffeomorphism onto $0<|q|<16$.
Its real tangent map is invertible everywhere, and its composite with
the inclusion $\{0<|q|<16\}\hookrightarrow\mathbb C^\times$ induces an
isomorphism of fundamental groups. These are statements about this
labelled geodesic-cut sector. They do not assert complex holomorphicity
or algebraic \emph{\'etaleness} of the map in 67/077--082, nor do they
settle the full multi-node question.

\paragraph{Geometric inputs and source locators.}
The hyperbolic pair of pants with specified nonnegative boundary
lengths is unique, with zero lengths interpreted as cusps. Its
Fenchel--Nielsen length and twist parametrization is recalled in
Mondello (2011), \S2.2, p.~640. The canonical attachment here is the
infinite flat-cylinder grafting of Mondello, \S5.4,
Theorem 5.4, p.~660; a punctured flat cylinder is conformally the
punctured disk with the arclength parametrization used above.
For an explicit independent source for the projective-connection
calculation below, see F\i rat (2021), equations (3.11), (3.13),
(3.14a), (3.23), and (3.30)--(3.34), arXiv:2102.03936v2,
pp.~14--17 and 24--25. We specialize the two other ends to cusps
and determine the branch directly below; we do not infer the branch
from that paper's integer-valued numerical prescription (3.27).

\begin{proof}
Write $\lambda=\ell/(2\pi)$, so $x=\lambda/2$.
The developing map $D$ of the hyperbolic metric is a local isometry
to a half-plane. Near a geodesic of length $2\pi\lambda$, diagonalize
its hyperbolic monodromy, with fixed points $0,\infty$.
In the arclength cap coordinate its analytic continuation is, up
to a nonzero constant, $w^{i\lambda}$. Indeed, increasing
$\arg w$ by $2\pi$ multiplies this function by $e^{-2\pi\lambda}$,
and along $|w|=1$ its hyperbolic arclength is
$\lambda\,d\arg w$. The same expression supplies the analytic
continuation across the seam. Its Schwarzian derivative therefore
has a double pole at the cap centre with leading coefficient
$(1+\lambda^2)/2$. At a cusp its leading coefficient is $1/2$,
as follows from the local developing coordinate $\log(z-z_0)$.
There are no other poles. On the sphere with centres and cusps
$0,1,\infty$, the resulting Schwarzian is consequently
\[
 T(z)=\frac{1+\lambda^2}{2z^2}
       +\frac1{2(z-1)^2}
       -\frac{1+\lambda^2}{2z(z-1)}.                           \tag{AQ.67.3}
\]
To see uniqueness, any two such projective connections differ by
a quadratic differential with at most simple poles at three points;
this space is $H^0(\mathbb P^1,\mathcal O(-1))=0$.

Set $a=(1+i\lambda)/2$. Direct substitution in
$y''+\tfrac12T(z)y=0$ gives the two solutions
\[
 y_\pm(z)=z^{(1\pm i\lambda)/2}(1-z)^{1/2}
     {}_2F_1\!\left(\frac{1\pm i\lambda}{2},
                    \frac{1\pm i\lambda}{2};
                    1\pm i\lambda;z\right).
\]
The ratio, normalized by its leading term at zero, is
\[
 D(z)=\frac{y_+(z)}{y_-(z)}
     =z^{i\lambda}
       \frac{{}_2F_1(a,a;2a;z)}
            {{}_2F_1(\bar a,\bar a;2\bar a;z)}.                \tag{AQ.67.4}
\]
All developing maps with the prescribed Schwarzian differ by a
M\"obius transformation. The ratio in (AQ.67.4) has the same
diagonal monodromy at zero; hence it is a valid developing map
after changing the target half-plane by that transformation.
Its boundary still contains the two fixed points $0,\infty$;
thus it is a straight half-plane bounded by a line through zero,
rather than a disk with a different boundary circle.

We now fix the phase without an unspecified multiple of $\pi$.
On the real interval $0<z<1$, $y_- =\overline{y_+}$.
Their Wronskian, computed at zero, is
$y_+'y_- -y_+y_-'=i\lambda$. In particular neither solution
vanishes there, and the continuously lifted argument $\Theta(z)$
of $D(z)$, normalized by
$\Theta(z)-\lambda\log z\longrightarrow0$ as $z\downarrow0$,
satisfies
\[
 \Theta'(z)=\frac{\lambda}{|y_+(z)|^2}>0.
                                                               \tag{AQ.67.5}
\]
The zero-balanced hypergeometric asymptotic
\[
 {}_2F_1(a,a;2a;z)
   \sim\frac{\Gamma(2a)}{\Gamma(a)^2}[-\log(1-z)]
       \qquad (z\uparrow1)
\]
is DLMF (15.4.21). Thus
\[
 \lim_{z\uparrow1}\Theta(z)=\vartheta(\lambda),\qquad
 \vartheta(\lambda)=2\operatorname{Im}
       \left(\log\Gamma(1+i\lambda)
                    -2\log\Gamma(\tfrac12+i\lambda/2)\right).
                                                               \tag{AQ.67.6}
\]
This is an equality of lifted arguments. For clarity, the
asymptotic is locally uniform in $\lambda\geq0$. The hypergeometric
factor on $0<z<1$ has no zeros when $\lambda>0$ by the Wronskian,
and at $\lambda=0$ it is positive. Its continuous logarithm,
zero at $z=0$, therefore has limiting phase equal to the indicated
continuous gamma phase, rather than that phase plus a nonzero
multiple of $2\pi$. Equivalently that possible integer is locally
constant in $\lambda$ and is zero at $\lambda=0$.

The pair of pants has its reflection fixing the perpendicular
seams. After capping, this antiholomorphic involution fixes
$0,1,\infty$, so it is complex conjugation. The perpendicular
seam from the cusp at $1$ to the geodesic boundary is the terminal
segment of the real interval $(0,1)$. Its developing image lies
on the unit circle, by (AQ.67.4). The cusp endpoint has phase
$\vartheta$. By (AQ.67.5), approaching that cusp from the core
uses phases less than $\vartheta$. Hence the developing half-plane
of the core is the half-plane with argument interval
$(\vartheta-\pi,\vartheta)$. The axis of the boundary monodromy,
whose fixed points are $0,\infty$, is its ray of argument
$\vartheta-\pi/2$. The seam meets that axis perpendicularly.
Starting from the cusp, this is the first intersection of the
seam with the axis; the further formal intersections obtained
after crossing the boundary line of that half-plane are not in
the hyperbolic core. This selects the lifted phase
$\vartheta-\pi/2$, with no integer ambiguity.

Choose $w=1$ at this perpendicular foot. On a collar of the
geodesic, and thence by analytic continuation into its cap,
\[
 w=\exp\left(\frac{\log D-i(\vartheta-\pi/2)}{i\lambda}\right).
\]
Here the logarithm has the lifted phase just specified.
Its boundary modulus is one and its boundary argument has the
normalized arclength, so this is precisely the prescribed cap
coordinate. At zero, (AQ.67.4) gives
\[
 w=\exp\left(-\frac{\vartheta-\pi/2}{\lambda}\right)z+O(z^2),
 \qquad
 r(\ell)=\exp\left(\frac{\vartheta-\pi/2}{\lambda}\right).
                                                               \tag{AQ.67.7}
\]
The duplication formula, DLMF (5.5.5), yields
$\vartheta=2\lambda\log2+2\phi(\lambda/2)$, proving
(AQ.67.1).

It remains to establish the exact range and the nonvanishing
derivative, rather than only an upper bound. The digamma integral
DLMF (5.9.16) gives, for $x\geq0$,
\[
 \phi'(x)=\operatorname{Re}\bigl(\psi(1+ix)
                            -\psi(\tfrac12+ix)\bigr)
           =\int_0^\infty\frac{\cos(xt)}{e^{t/2}+1}\,dt.
                                                               \tag{AQ.67.8}
\]
Let $f(t)=(e^{t/2}+1)^{-1}$. Then
\[
 f''(t)=\frac{e^{t/2}(e^{t/2}-1)}{4(e^{t/2}+1)^3}>0
               \quad(t>0).
\]
The function and its derivatives decay exponentially. Two
integrations by parts, with their endpoint terms included, give
\[
 \phi'(x)=\frac1{x^2}\int_0^\infty f''(t)(1-\cos xt)\,dt>0
                  \quad(x>0).                                \tag{AQ.67.9}
\]
Stirling's expansion for the same analytic logarithms, DLMF
(5.11.1), or the corresponding gamma-ratio expansion (5.11.13),
gives $\phi(x)\longrightarrow\pi/4$ as $x\to\infty$.
Since $\phi(0)=0$, it follows that
$0<\phi(x)<\pi/4$ for $x>0$. Therefore
\[
 \frac{d}{dx}\log r(4\pi x)
       =\frac{x\phi'(x)+\pi/4-\phi(x)}{x^2}>0.                 \tag{AQ.67.10}
\]
At zero, $\phi'(0)=2\log2$, so
\[
 r(\ell)=16\exp\left(-\frac{\pi^2}{\ell}+O(\ell^2)\right)
              \quad(\ell\downarrow0),
 \qquad r(\ell)\longrightarrow4\quad(\ell\to\infty).
\]
Real analyticity and (AQ.67.10) prove that $r$ is the asserted
diffeomorphism.

Finally Fenchel--Nielsen gluing of the two labelled pairs of
pants has exactly one length $\ell>0$ and one twist modulo
$\ell$, equivalently an angle $\theta$ modulo $2\pi$.
The two derivative vectors each have norm $r(\ell)$, and
the twist contributes their tensor's unit phase. This proves
(AQ.67.2). Its inverse is
$\ell=r^{-1}(\sqrt{|q|})$, $\theta=\arg q$ with the latter
viewed as a circle-valued coordinate. In logarithmic polar
coordinates the real Jacobian is diagonal, with entries
$2r'(\ell)/r(\ell)>0$ and $1$. Both the punctured disk and
$\mathbb C^\times$ deformation retract to the same circle
$|q|=1$; the inclusion therefore induces an isomorphism on
$\pi_1$, taking one full twist to the positive generator.
\end{proof}

\subsection{AQ-67-NONHOLO: the natural complex-analytic formulation fails}
Give the labelled geodesic-cut $(0,4)$ space its natural complex
structure, namely the covering
\[
 \mathcal C=\mathcal T(S_{0,4})/\langle D_\alpha\rangle
\]
of the moduli space of four labelled punctures, where $\alpha$ is
the chosen separating curve and $D_\alpha$ is its Dehn twist.
Equivalently, this is the component of the cover that remembers
the isotopy class of the cut. The two sides and all punctures
remain labelled. Fenchel--Nielsen coordinates upstairs are
$(\ell,\tau)\in\mathbb R_{>0}\times\mathbb R$, with
$D_\alpha(\ell,\tau)=(\ell,\tau+\ell)$; downstairs the
normalized angle is $\theta=2\pi\tau/\ell$ modulo $2\pi$.
The canonical capping map of AQ-67-EXACT is not holomorphic
from this complex manifold to $\mathbb C^\times$.
Thus its everywhere invertible real tangent map does not make
it a complex-analytic or $\mathbb C$-algebraic \emph{\'etale}
map. This refutes that stronger interpretation of 67/077--082
for the natural Teichm\"uller complex structure already in
the labelled $(0,4)$ sector.

\begin{proof}
Suppose it were holomorphic. By AQ-67-EXACT it is a bijection
with nonzero real Jacobian onto the punctured disk
$\mathcal D=\{0<|q|<16\}$. The holomorphic inverse function
theorem then makes it a biholomorphism with $\mathcal D$.
Conjugating the rotations $q\mapsto e^{is}q$ by it gives
holomorphic automorphisms of $\mathcal C$. Their lifts to
$\mathcal T(S_{0,4})$, continuously chosen to equal the identity
at $s=0$, are the explicit global diffeomorphisms
\[
 H_s(\ell,\tau)=
       \left(\ell,\tau+\frac{s\ell}{2\pi}\right),
                      \qquad s\in\mathbb R.                  \tag{AQ.67.11}
\]
This follows directly from the tensor's unit twist phase;
any choice of zero-angle section changes $\theta$ by a function
of $\ell$ and leaves (AQ.67.11) unchanged. Since the quotient
projection is a holomorphic covering, every $H_s$ is
holomorphic as well.

Let $\omega_{\mathrm{WP}}$ and $g_{\mathrm{WP}}$ be the
Weil--Petersson K\"ahler form and metric. Wolpert (1983),
Lemma 4.1, p.~229, says that the arclength twist vector field
$t_\alpha=\partial/\partial\tau$ has Hamiltonian potential
$-\ell$, in his normalization:
\[
 \iota_{t_\alpha}\omega_{\mathrm{WP}}=-d\ell.
\]
The vector field of (AQ.67.11) is
$V=(\ell/2\pi)t_\alpha$, so
\[
 \iota_V\omega_{\mathrm{WP}}
   =-\frac{\ell}{2\pi}d\ell
   =-d\!\left(\frac{\ell^2}{4\pi}\right).
\]
Cartan's formula and $d\omega_{\mathrm{WP}}=0$ imply
$\mathcal L_V\omega_{\mathrm{WP}}=0$; hence every global
$H_s$ preserves $\omega_{\mathrm{WP}}$. A different fixed
normalization of the Weil--Petersson form only multiplies
these identities by a nonzero constant and gives the same
conclusion. A holomorphic map preserving the K\"ahler form
preserves the metric, since that metric is recovered from
the form and the complex structure. Thus every $H_s$ would
be a Weil--Petersson isometry of $\mathcal T(S_{0,4})$.

Brock and Margalit, \emph{Weil--Petersson isometries via the
pants complex}, Main Theorem, p.~2 of arXiv:math/0412499v1,
prove that every such isometry is induced by an element of
the extended mapping class group; their theorem explicitly
includes $S_{0,4}$. That group is countable. On the other
hand (AQ.67.11) gives distinct maps for every two distinct
real numbers $s$, because $\ell>0$ and the upstairs twist
coordinate is real, not taken modulo $\ell$. This is an
uncountable family of isometries, a contradiction.
\end{proof}

\subsection{AQ-67-INVOLUTION: the proposition holds; its corollary fails}
The proposition of 67/091 gives an equivalence between
two-regular modules with involution and triples
$(P,Q,m)$, where $P,Q$ are two-regular and
$m\subset P/2P\oplus Q/2Q$ meets each summand trivially.
This proposition is correct, with morphisms of triples
preserving $m$.
\begin{proof}
For an involution on $M$, let $P=M_+$ and $Q=M_-$.
Two-regularity gives $P\cap Q=0$ and
$2M\subset P\oplus Q\subset M$.
Let $m=2M/(2P\oplus2Q)$. If an element of $2M$
has zero $Q$-coordinate, say $2x\in P$, then
$2(\sigma x-x)=0$, so $x\in P$ and $2x\in2P$.
This proves $m\cap(P/2P)=0$, and the other summand
is identical.

Conversely take the inverse image
$L\subset P\oplus Q$ of $m$, with involution
$(p,q)\mapsto(p,-q)$. It is stable because modulo
two this involution is the identity. It is
two-regular, and the intersection conditions give
$L_+=2P\oplus0$, $L_-=0\oplus2Q$.
Multiplication by two canonically identifies these
eigenspaces with $P,Q$ and recovers the original
$m$ from $2L/(4P\oplus4Q)$.
Starting with $M$, the reconstructed $L$ is $2M$,
canonically isomorphic to $M$ by multiplication
by two. Maps restrict to the eigenspaces, and
maps of triples preserve the indicated inverse
images. These constructions are inverse on maps
as well as objects, proving the equivalence.
\end{proof}

The corollary claiming invertible rank-one eigenspaces
for a nontrivial involution on a rank-two projective
module over a commutative two-regular ring is false
even when the ring is an integral domain. Set
\[
 R=\mathbb Z[x,y,z]/(x^2+yz-1),\quad M=R^2,\qquad
 \sigma=\begin{pmatrix}x&y\\z&-x\end{pmatrix}.                   \tag{AQ.9}
\]
\begin{proof}
The polynomial $f=x^2+yz-1$, viewed in
$\mathbb Z[x,y][z]$, is primitive and linear over
the fraction field; hence it is irreducible in
the unique factorization domain $\mathbb Z[x,y,z]$.
Thus $R$ is a domain, in particular two-regular.
Clearly $\sigma^2=1$ and $\sigma\ne\pm1$.

In $S=\mathbb Z[x,y,z]$ set $A=\sigma-I$ and
$B=\sigma+I$. Then $AB=BA=fI$.
If $\bar A\bar v=0$ over $R$, lift $v$ to $S^2$
and write $Av=fw$. Multiplication by $B$ gives
$fv=fBw$; since $f$ is a non-zero-divisor,
$v=Bw$. Therefore $\ker\bar A=\operatorname{im}\bar B$.
Interchanging $A,B$ proves
$\ker\bar B=\operatorname{im}\bar A$, and hence
\[
 P=\ker(\sigma-I)\simeq\operatorname{coker}(\sigma-I).
\]
At the maximal ideal
$\mathfrak m=(2,x-1,y,z)$ every entry of $A$
vanishes, so
$P/\mathfrak mP\simeq\mathbb F_2^2$.
Generically $A$ has rank one, giving $P$ rank one.
An invertible module would have a one-dimensional
fibre at $\mathfrak m$, a contradiction.
The same argument applies to $Q$.
\end{proof}
Connectedness of $\operatorname{Spec}R$ alone cannot
repair the corollary. This algebraic counterexample
does not involve RH, and cannot be excluded by it.

\subsection{AQ-154-COUNT: what can actually be counted over finite fields}
The supplied 154/120 says
\[
 \chi(\mathcal X)=s_0-s_1+s_2=1-g,\qquad
 2s_1=4s_0,\qquad \chi(\mathcal X)=s_2-s_0.
\]
Its $\mathcal X_2$ is a cellular surface of relative
positions of pseudolines. Neither that sheet nor SA11
supplies it with a scheme over a finite field.
The original Euler-characteristic calculation is
therefore not an assertion about
$\#\mathcal X_2(\mathbb F_q)$.

There is a genuine arithmetic consequence for a
different, explicitly specified object.
Let $L_1,\dots,L_n$ be distinct rational projective
lines in $\mathbb P^2_{\mathbb Q}$, let $U$ be their
complement, and let $S$ contain the finitely many
primes where lines coincide or their incidence
lattice changes. If $v$ runs through their
intersection points and $m_v$ lines meet at $v$, set
\[
 b_2=\sum_v(m_v-1)-n+1.
\]
For every $p\notin S$ and $r\geq1$,
\[
 \#U(\mathbb F_{p^r})=p^{2r}-(n-1)p^r+b_2,\qquad
 Z(U_{\mathbb F_p},t)
  =\frac{(1-pt)^{n-1}}{(1-p^2t)(1-t)^{b_2}}.                  \tag{AQ.10}
\]
Consequently its incomplete arithmetic zeta function is
\[
 \zeta^S(U,s)=
 \frac{\zeta^S(s-2)\zeta^S(s)^{b_2}}
      {\zeta^S(s-1)^{n-1}},\qquad
 \zeta^S(s)=\prod_{p\notin S}(1-p^{-s})^{-1}.                 \tag{AQ.11}
\]
\begin{proof}
Every intersection point is rational because the
lines are rational. Over $\mathbb F_{p^r}$ the sum
of the $n$ line counts is $n(p^r+1)$.
It overcounts an intersection $v$ exactly $m_v-1$
times. Subtracting the corrected union count from
$p^{2r}+p^r+1$ proves the first formula.
Insert it in
$Z(U_{\mathbb F_p},t)=
\exp(\sum_{r\geq1}\#U(\mathbb F_{p^r})t^r/r)$
and use $-\log(1-z)=\sum_{r\geq1}z^r/r$ to get
the second. Multiplying its Euler factors gives
(AQ.11), initially for $\operatorname{Re}s>3$
and then by meromorphic continuation.
\end{proof}
Under the adopted RH, the nontrivial-zero contribution
to the divisor of (AQ.11) is on
$\operatorname{Re}s=\tfrac12,\tfrac32,\tfrac52$:
the middle line gives poles when $n>1$, while the
other two give zeros when their exponents are
positive. Trivial zeros, the zeta pole, and the
explicit omitted Euler factors remain separately
specified by (AQ.11). These shifted divisors do not
give a formula for the pseudoline counts or the
topology of $\mathcal X_2$.

\subsection{AQ-SA08-LIFT: a concrete divided-power lifting obstruction}
The supplied SA08 intake concerns the uniqueness problem
of cote 136/117--121 and a proposed lifting input.
The original cote-136 sheets are absent from this packet.
The auxiliary lifting claim can nevertheless be checked
directly in its primary source: Caradot--Lin (2026),
arXiv:2511.22614v4, \S2.1, pp.~5--6, and
Proposition 3.10, pp.~19--20.
There \emph{free complex} means free in each degree,
and surjectivity of the PD ideals is degreewise.
Under these hypotheses the asserted lift need not exist,
over any nonzero commutative ring $k$.
\begin{proof}
Use cochain degrees and the square-zero PD dg algebras
\[
 R'=k\oplus(ku\xrightarrow{\,d\,}kv),\qquad
 |u|=-2,\quad |v|=-1,\quad du=v,
 \qquad R=k\oplus k\epsilon,\quad|\epsilon|=-2,\quad d\epsilon=0.
\]
Their augmentation ideals have zero multiplication,
and $\gamma_n=0$ for $n\geq2$ defines divided powers.
All divided-power identities hold: products of two
ideal elements vanish, and the differential identity
$d\gamma_n(u)=\gamma_{n-1}(u)\,du$ is zero for $n\geq2$,
including $u\,du=uv=0$ when $n=2$.
Define $\pi(u)=\epsilon$, $\pi(v)=0$, and $\pi|_k=1$.
It is a PD dg map surjective on the ideals.
Take the termwise free complex $X=k\epsilon$,
concentrated in degree $-2$, and the identity chain
map $X\to k\epsilon\subset R$.
A PD dg algebra lift from the free PD dg algebra
on $X$ would restrict on its generator to a chain
lift $\epsilon\mapsto cu$.
Compatibility with $\pi$ forces $c=1$, whereas
compatibility with the differential forces $cv=0$,
hence $c=0$ in $k$. This contradicts $k\ne0$.
\end{proof}
The notation $X$ here denotes just $k$ placed in
degree $-2$; no resolution or projectivity in the
category of complexes is being assumed.
The failed implication is precisely from a graded
lift to a chain lift. Theorem 3.15 of the same
version, pp.~23--24, assumes a quasi-isomorphism.
Our $\pi$ is not one:
$H^*(R')=k$ but $H^*(R)=k\oplus k\epsilon$.
Consequently this counterexample does not refute
that later theorem.

The unresolved SA08 statements are the actual
integral common-source assertion A-CS, terminality
of the all-arrow localization A-TL, and application
of the SA09 comparison structures A-SA09.
The intake states them but does not supply their
complete definitions and 15,000-line proof context.
Neither a weight grading on a divided-power algebra
nor the adopted analytic zero theorem supplies a
chain lift in the displayed counterexample.
That is a specific false lifting route; it is not
a disproof of the remaining existence or
localization conjectures.



% END archive_questions.tex

% BEGIN consequences_criteria_Tate2026.tex
\section{Current primary inputs on CM Hodge, abelian Tate, and the
standard
conjectures}\label{current-primary-inputs-on-cm-hodge-abelian-tate-and-the-standard-conjectures}

Research module \nolinkurl{consequences_criteria}, checked 20--21 November
2026 UTC (literature cutoff: 20 November), GPT-6 Astra Pro; effort not
exposed. This module updates the archive analysis using identified 2026
primary manuscripts. Their major theorems are \textbf{cited literature
inputs}, not theorems independently reproved in this run. The adopted
Alpöge RH and Dirichlet GRH hypotheses play no role in the deductions
below.

\subsection{CC-TATE01. The current all-CM input and its exact
provenance}\label{cc-tate01.-the-current-all-cm-input-and-its-exact-provenance}

\textbf{Cited 2026 preprint theorem.} For every complex abelian variety
of CM type, every rational Hodge class in every codimension is
algebraic. The precise locator is van Hoften and Sempliner, \emph{The
Hodge Conjecture for CM abelian varieties}, \textbf{Theorem I, p.~1},
Hexagon \textbf{2610.00259v2}. This includes powers and products of CM
abelian varieties, because those varieties are themselves of CM type.

The version history dates v1 to 15 October and v2 to \textbf{18 October
2026}. The current record's header displays the original 15 October
date; the downloaded v2 PDF creation and modification timestamps both
read 18 October 2026, 10:11:24 UTC. Thus the edition used here is
specified by the immutable v2 identifier, not merely by the header date.
The named authors are Pol van Hoften and Jack Sempliner. Their §1.3,
p.~4, explains that they give a variant proof inspired by the OpenAI
manuscript and do not certify every detail of that earlier manuscript.
Our cited input is their own Theorem I. Neither a journal publication
nor a complete independent verification of their proof is claimed here.

\textbf{Dependency inspection.} The reduction is Theorem 2.2.1, p.~6;
the algebraic surface criterion is Theorem 2.3.4, pp.~7--8; the required
nonzero cup product is Theorem 4.0.2, p.~12. We read the
surface-criterion proof completely, the cited Albanese statement Theorem
3.1.11 and Corollary 3.1.13, pp.~11--12, and the final argument §4.2.8,
pp.~19--22, including the diagonal-period Lemma 4.2.7, pp.~17--18, and
the archimedean nonvanishing Lemma 4.2.9, pp.~20--21. Those analytic
lemmas retain their cited theta-correspondence inputs; these inputs were
not independently reconstructed. The reduction's reference to Hazama is
recorded as a dependency, including the authors' p.~6 footnote about the
formulation of Hazama's Lemmas 7.11--7.12. This is a dependency read,
not a claim to have checked every underlying paper.

Milne's \textbf{9 November 2026} note, \emph{On the Hodge, Tate, and
standard conjectures for abelian varieties}, is linked from his own
mathematics homepage and available there as
\texttt{math/xnotes/HTS.pdf}. Its Hexagon deposit is
\textbf{2611.00244v1}, submitted 9 November and displayed as published
10 November. Milne states the CM result as Theorem 2.1, p.~4, and
explicitly treats the formerly conditional abelian consequences as
unconditional in §3, p.~6. We read the ten-page note. It is useful
corroboration and context; the numbered primary implications below are
also checked in the original papers.

\subsection{CC-TATE02. Exact abelian consequences of the cited CM
theorem}\label{cc-tate02.-exact-abelian-consequences-of-the-cited-cm-theorem}

Let \(A/\mathbb F_q\) be an abelian variety of arbitrary dimension, let
\(\ell\ne\operatorname{char}\mathbb F_q\), and let
\(N^r(A)_{\mathbb Q}\) denote codimension-\(r\) cycles with rational
coefficients modulo numerical equivalence. The resulting finite-field
statement is

\[
N^r(A)_{\mathbb Q}\otimes_{\mathbb Q}\mathbb Q_\ell
\xrightarrow{\ \sim\ }
H^{2r}_{\mathrm{et}}(A_{\overline{\mathbb F}_q},\mathbb Q_\ell(r))^{G_{\mathbb F_q}}
\qquad\text{for every }r.
\]

In particular, the order of the pole of \(Z(A,t)\) at \(t=q^{-r}\) is
\(\dim_{\mathbb Q}N^r(A)_{\mathbb Q}\). These assertions concern
rational coefficients.

\textbf{Exact original implication.} Milne (1999), \emph{Lefschetz
motives and the Tate conjecture}, \textbf{Theorem 7.1, journal p.~72},
assumes the Hodge conjecture for all complex CM abelian varieties and
proves his Tate statement (0.1) for all abelian varieties over
\(\overline{\mathbb F}_p\). Statement (0.1), \textbf{journal p.~46},
includes both numerical versus homological equivalence and spanning of
the Tate subspace. The following paragraph explicitly gives the
corresponding assertion for each finite-field model and the zeta pole
order. The theorem's proof is organized through Propositions \textbf{7.3
and 7.4, pp.~73--75}. These statements and the complete proofs of these
two propositions were read. The auxiliary theory in §§1--6 is used as
cited literature, not independently reproved. The journal PDF, whose
pagination is 45--76, was inspected; the author's separate erratum
contains an inaccurate bibliographic range 47--81, which is not
propagated here.

\textbf{Finite-field descent check.} The passage from geometric Tate
classes to classes over a given finite field can also be made explicit.
Write an invariant class as a finite \(\mathbb Q_\ell\)-linear
combination of classes of cycles over one common finite extension
\(\mathbb F_{q^m}\). Push each cycle down to \(A\) and divide by \(m\).
On geometric cohomology this is the averaging operator

\[
\frac1m\sum_{j=0}^{m-1}F^j.
\]

It fixes the given \(F\)-invariant class. Thus no enlargement of the
stated ground field is required in the displayed surjectivity assertion.
The numerical-kernel assertion comes from Milne's full Tate statement.
This argument is a formal verification of the field descent, not a new
proof of the Tate input.

\textbf{Standard conjectures.} For every algebraically closed field
\(k\) and every prime \(\ell\ne\operatorname{char}k\), numerical and
\(\ell\)-adic homological equivalence coincide on every abelian variety
over \(k\). For a polarized abelian variety of dimension \(g\), the
primitive rational algebraic classes in codimension \(r\le g/2\) carry
the positive-definite pairing

\[
(a,b)\longmapsto(-1)^r\deg\bigl(a\cdot b\cdot L^{g-2r}\bigr).
\]

The exact inspected locator is Milne (2002), \textbf{Theorem 3.3(a),(b),
pp.~8--9 of the final author manuscript dated 19 September 2001},
\nolinkurl{2002aP.pdf}. Its hypothesis is precisely the CM Hodge theorem in
CC-TATE01. We read the theorem and its full proof, including the
specialization argument to arbitrary algebraically closed \(k\).
\textbf{Corollary 3.4, p.~9}, gives the corresponding polarization of
the category of abelian numerical motives.

For cycles defined over an arbitrary field \(k\), these comparisons
follow by base change to \(\bar k\). Indeed, a complementary cycle over
\(\bar k\) is defined over a finite extension \(k'/k\); the degree
pairing with a cycle defined over \(k\) can be computed using its
pushforward from \(k'\), up to the nonzero extension degree. Thus
numerical triviality over \(k\) is equivalent to numerical triviality
after algebraic closure, with rational coefficients. The \(\ell\)-adic
homological condition is geometric by definition. Positivity on the
rational classes over \(k\) is the restriction of the geometric
positive-definite form. This gives the stated abelian comparison and
positivity over arbitrary fields, for \(\ell\) invertible in the field.

Two source precisions matter. First, van Hoften--Sempliner's §1.1.2
cites \textbf{Theorem 4.6}, the earlier long-version numbering; it
should not be silently represented as the number in the final
manuscript. Second, Milne's abbreviated 2026 display in §3.6 omits the
factor \((-1)^r\); the correct signed form is explicit in the 2002
paper's introduction, p.~2, and §3, p.~6. We use that signed form. The
final manuscript's \textbf{Remark 3.5, p.~9}, concerns an arbitrary
\emph{good} Weil cohomology, where ``good'' means that homological and
numerical equivalence agree on all varieties in the chosen
product-closed class. It does not remove that condition for an
unspecified Weil cohomology theory.

Milne's 2026 note states the finite-field conclusion as \textbf{Theorem
3.1, p.~6}, and the Hodge-standard conclusion as \textbf{Theorem 3.5,
p.~8}. The chain is an application of the cited 2026 theorem to old
conditional theorems. It does not use RH, Dirichlet GRH, a number-field
generalized RH, or an unproved automorphy assertion.

\subsection{CC-TATE03. Rational specialization: a separate 2026
input}\label{cc-tate03.-rational-specialization-a-separate-2026-input}

The Hexagon record \textbf{2610.00226v1}, contributed by \textbf{Keerthi
Madapusi}, is dated \textbf{13 October 2026}. Its actual PDF title is
\emph{Milne's rationality conjecture}, with the byline \textbf{``ChatGPT
and Claude''}; §1.1, p.~3, identifies Madapusi's questions and feedback.
Milne's November note cites this source under Madapusi's name and a
longer descriptive title. The bibliography below records the PDF
attribution and the contributor explicitly.

\textbf{Theorem 1.1, p.~1}, concerns an abelian variety
\(A/\overline{\mathbb Q}\) with good reduction
\(A_0/\overline{\mathbb F}_p\) at a fixed place, a rational Hodge class
\(\gamma\) of codimension \(r\), and a rational Lefschetz class
\(\delta\) of complementary codimension on \(A_0\). The traces

\[
\operatorname{Tr}_v(\gamma_{0,v}\delta_v)
\]

are the images of one rational number for all \(v=\ell\ne p\) and for
the crystalline realization \(v=p\), with the specified coefficient
extension in the latter. Theorem 5.1, pp.~11--12, supplies the
specialized Weil-class dichotomy; §5.1 deduces the main statement. The
lifting input is recorded precisely in \textbf{Proposition A.1,
pp.~13--14}, which invokes Kisin--Zhou, including a quasi-split
replacement and marked tensor transport. These are inspected source
hypotheses and dependency locators; this run does not claim an
independent proof of their lifting arguments.

Combining CC-TATE01 with this manuscript's \textbf{Corollary 6.2,
p.~13}, gives a single cycle \(z_\gamma\in CH^r(A_0)_{\mathbb Q}\)
representing every specialized realization of \(\gamma\). The
corollary's proof was read: CM transport gives
\(\gamma_0=f_0^*\kappa_0\); algebraicity of the CM class \(\kappa\)
permits specialization of a representing cycle; for an actual
homomorphism \(h=Nf_0\), the cycle is \(N^{-2r}h^*Z_0\). This is
algebraicity on the \textbf{special fibre}. It does not assert that
\(\gamma\) is algebraic on a general non-CM characteristic-zero \(A\).

\subsection{CC-TATE04. What this resolves in cote 15, and what it leaves
to
prove}\label{cc-tate04.-what-this-resolves-in-cote-15-and-what-it-leaves-to-prove}

The supplied transcription \nolinkurl{ARCHIVE_15_scan100.tex} asks for
sufficient arithmetic conditions for a Frobenius eigenvalue. Scans 102,
104, and 106 then seek a motivic Galois description and its geometric
version. These are distinct mathematical assertions.

\begin{longtable}[]{@{}
  >{\raggedright\arraybackslash}p{(\columnwidth - 4\tabcolsep) * \real{0.3333}}
  >{\raggedright\arraybackslash}p{(\columnwidth - 4\tabcolsep) * \real{0.3333}}
  >{\raggedright\arraybackslash}p{(\columnwidth - 4\tabcolsep) * \real{0.3333}}@{}}
\toprule\noalign{}
\begin{minipage}[b]{\linewidth}\raggedright
Archive issue
\end{minipage} & \begin{minipage}[b]{\linewidth}\raggedright
Consequence supported by the inspected chain
\end{minipage} & \begin{minipage}[b]{\linewidth}\raggedright
Remaining condition
\end{minipage} \\
\midrule\noalign{}
\endhead
\bottomrule\noalign{}
\endlastfoot
Existence of an individual Weil-integer eigenvalue & The Honda--Tate
construction already supplies the occurrence result analyzed separately
in the AQ module. & An occurrence statement does not identify all
correspondences on another variety with that eigenvalue. \\
Tate classes on abelian varieties and their products over finite fields
& The finite-field rational cycle isomorphism in CC-TATE02 holds in
every dimension and codimension. & The 2026 CM theorem is a cited
primary input, not an independently reproved theorem of this report. \\
A smooth projective variety with a motive of abelian type & Milne
(1999), \textbf{Remark 7.6(b), journal p.~76}, explicitly extends Tate
to a variety whose motive \textbf{modulo homological equivalence} lies
in the Tannakian subcategory generated by abelian varieties; products of
curves are an example. & One must establish this actual motivic
membership for the particular variety. Matching Frobenius characteristic
polynomials alone does not establish it. \\
The geometric motivic group in the abelian category & Milne's
\textbf{Proposition 7.3, pp.~73--74}, gives \(M=P\) for the category of
abelian numerical motives over \(\overline{\mathbb F}_p\), with \(P\)
his Weil-number protorus. & This is a statement about the abelian
category and its Tannakian group, not a neutralization of every motivic
gerbe over \(\mathbb Q\). \\
Primitive positivity and the numerical/homological comparison &
CC-TATE02 discharges these conditions for abelian varieties and the
indicated realizations. & No inspected theorem here proves these
statements for an arbitrary smooth projective variety in positive
characteristic. \\
Integral, mixed, or global analytic strengthening & No such
strengthening follows from the rational pure statements just listed. &
Integral comparison, extension data, and number-field analytic
continuation/zeros require their own hypotheses and constructions. \\
\end{longtable}

\textbf{The group-theoretic application in full.} Let \(M\) be the
fundamental group used in §7 of Milne (1999) for abelian numerical
motives over \(\overline{\mathbb F}_p\). The pre-existing construction
gives \(P\subseteq M\subseteq L\). The CM Hodge theorem identifies the
CM Hodge correspondences with algebraic correspondences, so
specialization gives the additional inclusion \(M\subseteq S\) inside
the common Lefschetz group \(T\). Milne's intersection theorem gives
\(S\cap L=P\). Therefore

\[
P\subseteq M\subseteq S\cap L=P,
\]

and hence \(M=P\). This is exactly the known-result application in
Proposition 7.3; it neither assumes nor proves that all smooth
projective motives belong to this category. Milne defines \(P\) by the
Galois module of germs of Weil numbers in §4, \textbf{author-manuscript
p.~14}. These germs use the power transition maps as the finite field
grows. A convention for homological versus cohomological weights must be
retained when comparing his notation with the archive.

\textbf{Gerbe precision.} Milne's separate \textbf{1999b
addendum/erratum, p.~1}, retracts an optimistic introductory claim that
an intersection of Tannakian groups alone automatically constructs the
desired quotient category. Its stated issue is the nonneutral setting.
It does not retract Theorem 7.1; its §7, pp.~3--4, even supplies
Deligne's alternative proof route. Thus the present all-CM theorem
discharges the hypotheses of the inspected abelian Tate theorem, while
an asserted canonical rational fibre functor or the archive's literal
displayed product decomposition must still be justified on its own
terms.

The AQ module makes both sides of this distinction concrete:
\textbf{AQ-ABELIAN-RECOGNITION} proves that full Frobenius determinants
recognize actual rational numerical objects in its explicitly defined
abelian--Artin category, whereas \textbf{AQ-RATIONAL-FIBRE} gives an
elliptic quaternionic obstruction to a rational fibre functor and to
discarding eigenvalue multiplicities. Both complete proofs were
independently read in this run and passed on their stated inputs.

The safe global status statement is therefore precise: \textbf{the
inspected 2026 literature resolves the abelian Tate/standard inputs and
the stated abelian-type Tate input; the unrestricted archive
identification is not established by those inputs.} No assertion that
every broader version is currently open is needed for this conclusion.

\subsection{CC-TATE05. Fourfold confirmation and revision
control}\label{cc-tate05.-fourfold-confirmation-and-revision-control}

Broe's \emph{The Tate conjecture for abelian fourfolds over finite
fields}, \textbf{arXiv:2608.28651v1, 19 August 2026}, gives
\textbf{Theorem 1.1, p.~1}, for every abelian fourfold over a finite
field and \textbf{Corollary 1.2, p.~2}, for standard D on abelian
fourfolds over arbitrary fields, with \(\ell\) invertible. Corollary
8.2, p.~20, repeats the latter with its proof. The complete argument in
§§7--8 was read, with earlier reduction and lifting results retained as
cited inputs. Its \textbf{Remark 1.3, p.~2}, distinguishes positivity on
the numerical quotient from Grothendieck's original homological
formulation. This verifies and strengthens the previously uninspected
Broe entry in SA22; the later all-CM chain makes the abelian dimension
restriction unnecessary for the Tate/standard conclusions used here.

Ningyi Li's separate \textbf{arXiv:2609.27916v2} is explicitly
\textbf{withdrawn on 8 October 2026}. Its abstract must not be used as a
current theorem citation. The withdrawal does not concern the distinct
fivefold manuscript arXiv:2609.06265v2. That manuscript was examined
only selectively and is not required as an input anywhere in this
module.

\subsection{Bibliography and stable primary
URLs}\label{bibliography-and-stable-primary-urls}

Broe, M. (2026, August 19). \emph{The Tate conjecture for abelian
fourfolds over finite fields} (arXiv:2608.28651v1) {[}Preprint{]}.
\url{https://arxiv.org/pdf/2608.28651v1}

ChatGPT, \& Claude. (2026, October 13). \emph{Milne's rationality
conjecture} (Hexagon:2610.00226v1) {[}Preprint contributed by Keerthi
Madapusi; authorship and supervision explained in §1.1{]}.
\url{https://hexagonmath.org/pdf/2610.00226v1}

Li, N. (2026, October 8). \emph{The Tate conjecture for powers of
abelian fourfolds and the Hodge conjecture for powers of CM fourfolds}
(arXiv:2609.27916v2) {[}Withdrawn preprint record{]}.
\url{https://arxiv.org/abs/2609.27916}

Milne, J. S. (1999). Lefschetz motives and the Tate conjecture.
\emph{Compositio Mathematica, 117}(1), 45--76.
\url{https://doi.org/10.1023/A:1000776613765}

\begin{itemize}
\tightlist
\item
  Inspected journal PDF:
  \url{https://www.cambridge.org/core/services/aop-cambridge-core/content/view/962D833FAAC3219F55A58B193323B52A/S0010437X99000317a.pdf/lefschetz-motives-and-the-tate-conjecture.pdf}
\item
  Author manuscript: \url{https://jmilne.org/math/articles/1999bP.pdf}
\item
  Author addendum/erratum:
  \url{https://jmilne.org/math/articles/add/1999b.pdf}
\end{itemize}

Milne, J. S. (2002). Polarizations and Grothendieck's standard
conjectures. \emph{Annals of Mathematics, 155}(2), 599--610.
\url{https://doi.org/10.2307/3062126}

\begin{itemize}
\tightlist
\item
  Inspected final author manuscript, dated September 19, 2001:
  \url{https://www.jmilne.org/math/articles/2002aP.pdf}
\item
  Publisher metadata: \url{https://annals.math.princeton.edu/2002/155-2/p07}
\item
  Author's separate explanatory addendum:
  \url{https://www.jmilne.org/math/articles/2002a.html}
\end{itemize}

Milne, J. S. (2026, November 9). \emph{On the Hodge, Tate, and standard
conjectures for abelian varieties} (Hexagon:2611.00244v1) {[}Expository
preprint{]}. \url{https://www.jmilne.org/math/xnotes/HTS.pdf}

\begin{itemize}
\tightlist
\item
  Versioned deposit: \url{https://hexagonmath.org/2611.00244v1}
\end{itemize}

van Hoften, P., \& Sempliner, J. (2026, October 18). \emph{The Hodge
conjecture for CM abelian varieties} (Hexagon:2610.00259v2)
{[}Preprint{]}. \url{https://hexagonmath.org/pdf/2610.00259v2}

The source snapshots, extracted text, and selected rendered theorem
pages are preserved in \texttt{research/cc\_sources/}. Exact
web-reference identifiers and inspection methods are in the accompanying
citation and check-event ledgers.

% END consequences_criteria_Tate2026.tex

% BEGIN consequences_criteria_abelian_group.tex
\section{An explicit motivic group for the abelian sector after
extending
coefficients}\label{an-explicit-motivic-group-for-the-abelian-sector-after-extending-coefficients}

\textbf{Scope.} This is a complete deduction from the all-abelian Tate
and numerical--homological equivalence inputs recorded in CC-TATE01--02
and the fully proved realization theorem AQ-ABELIAN-RECOGNITION. It
concerns the abelian-generated category over a \textbf{fixed finite
field}, after extending coefficients. It neither classifies arbitrary
numerical motives nor identifies the uncertain historical notation in
cote 15 with a modern group by fiat. Literature cutoff: 20 November
2026; checks: 21 November 2026 UTC. No claim of historical novelty is
made.

\subsection{CC-AG01 --- The abelian category with algebraically closed
coefficients}\label{cc-ag01-the-abelian-category-with-algebraically-closed-coefficients}

Let \(q=p^a\), let \(\ell\ne p\), set \(E=\overline{\mathbf Q}_\ell\),
and fix an embedding \(\iota:\overline{\mathbf Q}\hookrightarrow E\).
Let \(\mathcal A_q\) be the rational numerical category generated by
abelian varieties over \(\mathbf F_q\), finite étale schemes, and Tate
motives under finite sums, tensor products, duals, and direct summands,
exactly as in AQ-ABELIAN-RECOGNITION. Use the parity-corrected symmetry,
so its ordinary direct-sum realization is symmetric monoidal. Define

\[
 \mathcal A_{q,E}
 =\operatorname{Kar}(\mathcal A_q\otimes_{\mathbf Q}E),
\]

where the tensor product first extends every Hom space and
\(\operatorname{Kar}\) then adds images of idempotents. Thus coefficient
extension adds summands; it does not add motives of arbitrary varieties
as generators.

Define the multiplicative abelian group

\[
 W(q)=\left\{\alpha\in\overline{\mathbf Q}^{\,\times}\ \middle|\
 \begin{array}{l}
 \alpha\text{ is a unit at every finite prime away from }p,\\
 \exists w\in\mathbf Z\ \forall\sigma:\ |\sigma(\alpha)|=q^{w/2}
 \end{array}\right\}.
\]

The integer \(w=w_q(\alpha)\) is unique. Let
\(\operatorname{Gr}^{\mathrm{fin}}_{W(q)}(E)\) be the category of
\(W(q)\)-graded \(E\)-vector spaces with finite total dimension. Its
tensor grading is multiplication in \(W(q)\), and its symmetry is the
ordinary interchange of vector factors.

\textbf{Theorem.} Geometric Frobenius gives a symmetric monoidal
equivalence

\[
 \mathcal A_{q,E}\ \xrightarrow{\ \sim\ }\
 \operatorname{Gr}^{\mathrm{fin}}_{W(q)}(E),\qquad
 M\longmapsto
 \bigoplus_{\alpha\in W(q)}
 R_E(M)_{F=\iota(\alpha)}.                         \tag{CC-AG.1}
\]

In particular, every \(\alpha\in W(q)\) labels an invertible rank-one
object \(L_\alpha\), unique up to isomorphism, and every object is a
finite direct sum of these, with arbitrary nonnegative multiplicities.
Tensor product satisfies
\(L_\alpha\otimes L_\beta\simeq L_{\alpha\beta}\). The weight is
\(w_q(\alpha)\), and \(\mathbf Q(1)\) corresponds to the label
\(q^{-1}\). Also

\[
 K_0(\mathcal A_{q,E})\simeq\mathbf Z[W(q)]
\]

as rings, with the classes of actual objects corresponding to the
finitely supported nonnegative sums of group elements.

\textbf{Complete proof.} AQ-ABELIAN-RECOGNITION proves, for objects
\(M,N\) of the rational category,

\[
 \operatorname{Hom}_{\mathcal A_q}(M,N)
          \otimes_{\mathbf Q}\mathbf Q_\ell
 \simeq
 \operatorname{Hom}_{\mathbf Q_\ell[F]}
       (R_\ell M,R_\ell N),                       \tag{CC-AG.2}
\]

and proves semisimplicity of Frobenius. Extend (CC-AG.2) to \(E\). The
commutant on the right is the kernel of the finite linear system
\(TF_M=F_NT\); field extension commutes with this kernel. The left side
becomes \(\operatorname{Hom}_{\mathcal A_q}(M,N)\otimes_{\mathbf Q}E\).
Thus realization remains fully faithful on the category with extended
Hom spaces. For idempotents \(e\in\operatorname{End}(M)\otimes E\) and
\(f\in\operatorname{End}(N)\otimes E\), its Hom space is

\[
 f\bigl(\operatorname{Hom}_{\mathcal A_q}(M,N)\otimes E\bigr)e.
\]

The identical formula holds for maps between their realized images.
Hence realization is fully faithful on the entire idempotent completion.
Frobenius remains semisimple after extension and on every summand.

Its eigenvalues lie in \(W(q)\). Indeed, eigenvalues \(\pi\) on \(H^1\)
of an abelian variety are integral Weil numbers of weight one. Choose an
identification of \(\overline{\mathbf Q}\) with the algebraic numbers in
\(\mathbf C\) when using complex conjugation. Then
\(\pi\overline{\pi}=q\), so integrality of both factors forces every
non-\(p\) finite valuation to vanish. They are therefore \(p\)-units.
The eigenvalues on exterior powers are products of such numbers; Tate
twists multiply them by integral powers of \(q\); finite étale factors
supply roots of unity. These operations, duals and summands preserve the
indicated group of numbers and their integral weights. This proves that
the displayed functor lands in the target category. Since a map between
semisimple Frobenius modules commutes with Frobenius exactly when it
preserves its eigenspaces, the functor is fully faithful into that
category.

It remains to prove that every label occurs and can be split off with
multiplicity one. Let \(\alpha\in W(q)\), of weight \(w\). Choose
\(n\ge0\) large enough that \(\beta=q^n\alpha\) is integral and
\(m=w+2n>0\). Such \(n\) exists: outside \(p\) all valuations are
already zero, and multiplication by \(q^n\) clears the finitely many
negative valuations above \(p\) in a number field containing \(\alpha\).

Choose \(\gamma\in\overline{\mathbf Q}\) with \(\gamma^m=\beta\). It is
integral by transitivity of integrality, and every conjugate has
absolute value \(q^{1/2}\). Honda's Main Theorem (1968, p.~83) therefore
supplies an abelian variety \(B/\mathbf F_q\) for which \(\gamma\) is a
Frobenius eigenvalue on \(H^1\). This is exactly the occurrence argument
already proved in AQ-WEIL-EXIST. The object

\[
 M_\alpha=h^1(B)^{\otimes m}(n)\in\mathcal A_q
\]

has eigenvalue \(\gamma^m q^{-n}=\alpha\), with weight \(m-2n=w\).

Its realized \(\alpha\)-eigenspace is nonzero and Frobenius is scalar
there. Choose a line in that eigenspace and a vector-space complement to
it, including all other eigenspaces. Projection onto the line commutes
with Frobenius. Full faithfulness over \(E\) lifts this rank-one
projector to an idempotent in
\(\operatorname{End}_{\mathcal A_q}(M_\alpha)\otimes E\). The idempotent
completion supplies its image \(L_\alpha\). This is why occurrence alone
would not have sufficed in the rational category: the projector is now
allowed to have extended coefficients.

Every finite-dimensional \(W(q)\)-graded vector space is a finite sum of
these one-dimensional realizations. This proves essential surjectivity.
The tensor-product eigenspace decomposition is compatible with
multiplication of labels, and the corrected symmetry realizes as
ordinary symmetry, so the equivalence is symmetric monoidal. Full
faithfulness and realization give
\(L_\alpha\otimes L_\beta\simeq L_{\alpha\beta}\), in particular
\(L_{\alpha^{-1}}\) is the inverse. The claims about weights and Tate
twists follow from the eigenvalue convention. Finally the target is
semisimple with these simple objects; its Grothendieck ring is freely
generated additively by their classes, with the multiplication just
computed. This proves the ring assertion and the whole theorem.
\(\square\)

\textbf{Comparison with the primary literature.} Milne's Definition 2.5,
author-manuscript pp.~14--15, uses the condition that some \(q^n\alpha\)
is integral. Under the common archimedean absolute-value condition, this
is equivalent to the \(p\)-unit condition above: complex conjugation
gives \(\alpha\overline\alpha=q^w\); away from \(p\), integrality of
both conjugate factors and this identity force both valuations to be
zero. The converse is the denominator-clearing argument in the proof.
Milne's Proposition 2.6, p.~15, contains the occurrence argument, while
his Proposition 3.7, p.~36, gives the broader coefficient-extended
classification under the full Tate conjecture imposed in that part of
the paper. Here (CC-AG.2) is already proved for \(\mathcal A_q\) from
the cited all-abelian theorem; no full Tate conjecture for arbitrary
varieties is used.

\subsection{CC-AG02 --- The split group, its weight factor, and its
Frobenius
point}\label{cc-ag02-the-split-group-its-weight-factor-and-its-frobenius-point}

Write \(\omega_E\) for the underlying-vector-space functor in (CC-AG.1).
For any abelian group \(M\), let

\[
 D_E(M)=\operatorname{Spec}E[M],
\]

with group algebra Hopf structure
\(\Delta([\alpha])=[\alpha]\otimes[\alpha]\).

\textbf{Theorem.} There is a canonical identification, relative to the
chosen realization and \(\iota\),

\[
 \operatorname{Aut}^{\otimes}(\omega_E)
       \simeq D_E(W(q)).                            \tag{CC-AG.3}
\]

Let

\[
 W_0(p)=\left\{\alpha\in\overline{\mathbf Q}^{\,\times}\ \middle|\
 \begin{array}{l}
 \alpha\text{ is a }p\text{-unit},\\
 |\sigma(\alpha)|=1\text{ for every }\sigma
 \end{array}\right\}.
\]

Choosing a square root \(u=\sqrt q\) gives

\[
 W(q)\simeq W_0(p)\times\mathbf Z,\qquad
 \alpha\longmapsto(\alpha u^{-w_q(\alpha)},w_q(\alpha)),
\]

and consequently

\[
 \operatorname{Aut}^{\otimes}(\omega_E)
       \simeq D_E(W_0(p))\times\mathbf G_{m,E}.       \tag{CC-AG.4}
\]

The weight cocharacter is the inclusion of the displayed \(\mathbf G_m\)
factor. The Tate object has character \((1,-2)\). Frobenius is the
distinguished point \(f\) with character values
\(\chi_\alpha(f)=\iota(\alpha)\), and its nonnegative powers are Zariski
dense in this affine group.

\textbf{Complete proof.} Let \(R\) be any commutative \(E\)-algebra. A
tensor automorphism of \(\omega_E\otimes_E R\) acts on each \(L_\alpha\)
by a scalar \(g_\alpha\in R^\times\). Naturality with respect to every
matrix endomorphism of \(L_\alpha^{\oplus r}\), and to the inclusions
and projections of summands, forces its action on every object to be
this scalar on its entire \(\alpha\)-graded summand. Compatibility with
tensor product and the unit says

\[
 g_{\alpha\beta}=g_\alpha g_\beta,\qquad g_1=1.
\]

Conversely any group homomorphism \(W(q)\to R^\times\) defines such a
natural tensor automorphism by those scalars. This correspondence is
functorial in \(R\), and the group algebra represents precisely this
functor. It proves (CC-AG.3).

The function \(w_q:W(q)\to\mathbf Z\) is a group homomorphism with
kernel \(W_0(p)\). The number \(u\) is a \(p\)-unit of weight one, so
\(n\mapsto u^n\) splits this homomorphism. Explicitly, the displayed map
and \((a,n)\mapsto au^n\) are inverse group homomorphisms. Taking group
algebras gives (CC-AG.4). A scalar \(z\) in the weight factor acts on
label \(\alpha\) by \(z^{w_q(\alpha)}\), which proves the cocharacter
statement. Since \(q^{-1}=u^{-2}\), the Tate character is \((1,-2)\).

The map \(\alpha\mapsto\iota(\alpha)\) is multiplicative, so it defines
\(f\in D_E(W(q))(E)\). To prove density, let a regular function

\[
 h=\sum_{j=1}^s c_j[\alpha_j]\in E[W(q)]
\]

with distinct \(\alpha_j\) vanish at every \(f^r\), \(r\ge0\). For
\(r=0,\ldots,s-1\), this gives

\[
 \sum_{j=1}^s c_j\,\iota(\alpha_j)^r=0.
\]

The Vandermonde determinant is nonzero because \(\iota\) is injective
and the labels are distinct. Thus every \(c_j\) is zero. No nonzero
regular function vanishes on these points; their Zariski closure is the
whole affine group. \(\square\)

\subsection{CC-AG03 --- Fixed-field torsion, coefficients, and the
limits of the
description}\label{cc-ag03-fixed-field-torsion-coefficients-and-the-limits-of-the-description}

The character group \(W(q)\) contains all roots of unity \(\mu_\infty\),
precisely its torsion subgroup. Thus \(D_E(W(q))\) is a diagonalizable
affine group of multiplicative type, generally of infinite type; it is
\textbf{not a protorus}. For example, the inclusion
\(\mu_n\subset W(q)\) gives a quotient
\(D_E(W(q))\twoheadrightarrow D_E(\mu_n)\), a nontrivial finite étale
group when \(n>1\). This also shows why the weight-zero factor
\(D_E(W_0(p))\) must not be confused with an identity component: it
retains the same root-of-unity characters. These are the finite-field
Artin contributions.

\textbf{Justification of the quotient assertion.} For a subgroup
\(H\subset M\), choose representatives for the cosets of \(H\) in \(M\).
The group algebra \(E[M]\) is then a free \(E[H]\)-module, with those
representatives as a basis. Its rank is nonzero, so the induced map
\(D_E(M)\to D_E(H)\) is faithfully flat. For \(H=\mu_n\), the abstract
character group is cyclic of order \(n\), whose group algebra is
\(E[t]/(t^n-1)\). In characteristic zero it is finite étale and has
\(n\) geometric points. A protorus, being an inverse limit of tori, is
connected and cannot have this nontrivial finite étale quotient.
\(\square\)

This description is compatible with AQ-RATIONAL-FIBRE. Over
\(\mathbf F_9\), the quaternionic elliptic-curve motive of
AQ-RATIONAL-FIBRE has eigenvalue \(-3\) with multiplicity two, and there
is no rational numerical object having only one copy. After coefficient
extension, (CC-AG.1) supplies the corresponding rank-one summand. Thus
(CC-AG.3) does not assert a neutral \(\mathbf Q\)-valued fibre functor
or remove rational multiplicity restrictions.

There is also a separate base-field operation: passing from a fixed
\(\mathbf F_q\) to \(\overline{\mathbf F}_p\) uses power transitions on
Frobenius characters and their germs. Milne treats this distinction in
\textbf{Proposition 3.2}, p.~34. We have held the finite base field
fixed throughout. The category of arbitrary numerical motives, integral
lattices, mixed extensions, and the archive's uncertain \(Q'\), \(M_0\),
or quotient conventions remain outside (CC-AG.1)--(CC-AG.4).

\subsection{Bibliography and exact
inputs}\label{bibliography-and-exact-inputs}

\begin{itemize}
\tightlist
\item
  Honda, T. (1968). Isogeny classes of abelian varieties over finite
  fields. \emph{Journal of the Mathematical Society of Japan, 20}(1--2),
  83--95. \url{https://doi.org/10.2969/jmsj/02010083}. \textbf{Main Theorem,
  p.~83}; the definition immediately preceding it identifies the
  relevant integral numbers by their conjugate absolute values. Primary
  PDF theorem text inspected through
  \url{https://people.math.rochester.edu/faculty/doug/otherpapers/Honda-Ab.pdf}.
\item
  Milne, J. S. (1994). Motives over finite fields. In U. Jannsen, S. L.
  Kleiman, \& J.-P. Serre (Eds.), \emph{Motives} (Proceedings of
  Symposia in Pure Mathematics, Vol. 55, Part 1, pp.~401--459). American
  Mathematical Society. \url{https://jmilne.org/math/articles/1994aP.pdf}.
  Author-manuscript locators: parity-corrected symmetry, pp.~5--6;
  scalar extension (1.3.2), p.~7; Definition 2.5 and Proposition 2.6,
  pp.~14--15; Proposition 3.2, p.~34; Proposition 3.7, p.~36. The global
  classification in the latter source retains its stated Tate
  hypothesis.
\item
  Internal proved inputs: \textbf{AQ-WEIL-EXIST} and
  \textbf{AQ-ABELIAN-RECOGNITION}, including (AQ.15), in
  research/archive\_questions.tex; the all-abelian Tate and
  numerical--homological input and its exact primary bibliography are
  recorded in \textbf{CC-TATE01--02},
  research/consequences\_criteria\_Tate2026.tex. The major 2026 all-CM
  theorem remains an inspected cited input, not a theorem whose full
  foundational proof has been independently reconstructed in this run.
\end{itemize}

% END consequences_criteria_abelian_group.tex

% BEGIN consequences_criteria_cross_cycles.tex
\section{A finite cross-correspondence criterion for a specified
variety}\label{a-finite-cross-correspondence-criterion-for-a-specified-variety}

\textbf{Scope and input.} Fix \(k=\mathbf F_q\), a prime
\(\ell\ne\operatorname{char}k\), and \(E=\overline{\mathbf Q}_\ell\),
with the algebraic-number embedding used in CC-AG01. This section uses
the proved abelian-sector equivalence CC-AG01, and therefore the
inspected cited all-abelian Tate and numerical--homological inputs
underlying AQ.15. It gives an exact remaining criterion for a specified
smooth projective variety; it does not assert that its cycle condition
always holds. Literature cutoff: 20 November 2026; proof checks: 21
November 2026 UTC.

\subsection{The homological category used
here}\label{the-homological-category-used-here}

Let \(\mathcal H_{\ell,E}(k)\) be the category of pure motives formed
from \textbf{\(E\)-linear algebraic correspondences modulo the kernel of
their \(E\)-valued \(\ell\)-adic cycle class}, followed by idempotent
completion and inversion of the Tate object. Use the usual correction of
commutativity by cohomological parity. Its realization \(R\) on the
ordinary direct sum of cohomology is faithful. Indeed, Poincaré duality
and the Künneth formula identify the class of a correspondence with its
graded linear action; being zero under the action is exactly the
equivalence relation imposed.

This coefficient convention matters. For arbitrary varieties we do not
identify this category with a naive scalar extension of rational
homological Hom spaces without proving scalar injectivity. Nor do we
begin by realizing arbitrary numerical motives. On the abelian--Artin
subcategory, the numerical--homological equality and scalar-injectivity
proof in AQ.15 do identify it with \(\mathcal A_{q,E}\); hence that
subcategory embeds fully faithfully in \(\mathcal H_{\ell,E}(k)\).

All cohomological projectors needed here are algebraic in the
homological category. The characteristic polynomials of Frobenius on
distinct cohomological degrees are relatively prime by their different
weights. The Chinese remainder theorem produces polynomials in the
Frobenius correspondence acting as identity in one degree and zero in
the others. Their idempotent identities hold after realization and thus
in this category. Frobenius semisimplicity is not required for this
construction: the full characteristic polynomials already annihilate the
relevant generalized eigenspaces. This is the finite-field construction
discussed by Milne (1994), author-manuscript pp.~5--6.

\subsection{CC-XC01 --- An abelian target realizing the
semisimplification}\label{cc-xc01-an-abelian-target-realizing-the-semisimplification}

Let \(X/k\) be nonempty, smooth and projective of pure dimension \(d\),
and put

\[
 H=h_{\ell,E}(X),\qquad V=R(H)=\bigoplus_i H^i_{\mathrm{et}}(X_{\bar k},E).
\]

There exist finitely many \(Y_\nu/k\), each an abelian variety or a
finite étale scheme, and an idempotent \(e\) on

\[
 T=\bigoplus_{\nu=1}^s h(Y_\nu),\qquad M=(T,e)\in\mathcal A_{q,E},
\]

such that \(W=R(M)\) is isomorphic, as a Frobenius module, to
\(V^{\mathrm{ss}}\). In particular \(\dim_E W=\dim_E V\). No
semisimplicity assumption on \(V\) is made.

\textbf{Proof.} Deligne's Weil theorem (1974, Theorem 1.6, p.~276) gives
integral algebraic Frobenius eigenvalues of weight \(i\) on \(H^i(X)\).
Their products with their complex conjugates are \(q^i\), so they are
\(p\)-units. Thus every eigenvalue belongs to \(W(q)\), independently of
possible Jordan blocks.

For an eigenvalue \(\alpha\) in a positive degree \(i\), choose
\(\gamma\) with \(\gamma^i=\alpha\). As in AQ-WEIL-EXIST, \(\gamma\) is
an integral weight-one Weil number. Honda's Main Theorem gives an
abelian variety \(B_\alpha/k\) with eigenvalue \(\gamma\) on \(H^1\).
Then \(H^i(B_\alpha^i)\) contains the Künneth summand
\(H^1(B_\alpha)^{\otimes i}\) and hence the eigenvalue \(\alpha\). The
eigenvalues in degree zero are roots of unity, since Frobenius permutes
the finitely many geometric connected components of \(X\). A root of
unity of order \(r\) occurs in
\(H^0(\operatorname{Spec}\mathbf F_{q^r})\).

Use enough copies of these varieties to accommodate every algebraic
multiplicity in the characteristic polynomial of \(V\). This gives a
finite \(T\) whose semisimple realization contains the desired
eigenspaces with at least those multiplicities. Project onto the desired
subspaces. This projector commutes with Frobenius; AQ.15 after
coefficient extension lifts it to \(e\). Its image has exactly
\(V^{\mathrm{ss}}\) as realization. No Tate twist was necessary in this
construction, because all degrees of \(X\) are nonnegative and its
eigenvalues are integral. \(\square\)

\subsection{CC-XC02 --- A determinant condition on one projected
cross-cycle
map}\label{cc-xc02-a-determinant-condition-on-one-projected-cross-cycle-map}

Fix \(T,e,M,W\) as in CC-XC01. A codimension-\(d\) cycle \(z_\nu\) on
\(X\times Y_\nu\) acts in degree zero from \(h(X)\) to \(h(Y_\nu)\).
Define the projected cross-cycle map

\[
 \begin{split}
 \kappa:\ 
 \bigoplus_{\nu=1}^s
 \bigl(\mathrm{CH}^d(X\times Y_\nu)_{\mathbf Q}
              \otimes_{\mathbf Q}E\bigr)
 &\longrightarrow \operatorname{Hom}_{E[F]}(V,W),\\
 (z_\nu)_\nu&\longmapsto R(e)\circ R((z_\nu)_\nu).
 \end{split}                                      \tag{CC-XC.1}
\]

Write \(I=\operatorname{im}\kappa\). It is a finite-dimensional
\(E\)-linear subspace of \(\operatorname{Hom}_E(V,W)\). Choose bases of
the equal-dimensional spaces \(V,W\); the condition
\(\det|_I\not\equiv0\) is independent of those choices.

\textbf{Theorem.} The following conditions are equivalent:

\begin{enumerate}
\def\labelenumi{\arabic{enumi}.}
\tightlist
\item
  \(H\) is isomorphic in \(\mathcal H_{\ell,E}(k)\) to an object of
  \(\mathcal A_{q,E}\).
\item
  \(H\simeq M\) in that category.
\item
  \(I\) contains an invertible linear map.
\item
  The determinant polynomial is not identically zero on \(I\).
\end{enumerate}

When they hold, Frobenius on \(V\) is automatically semisimple. Thus the
determinant condition incorporates both semisimplicity and the needed
cross-correspondence algebraicity.

\textbf{Complete proof.} By the definition of the homological category
and the idempotent completion,

\[
 I=R\bigl(\operatorname{Hom}_{\mathcal H_{\ell,E}}(H,M)\bigr).
                                                        \tag{CC-XC.2}
\]

The codimension in (CC-XC.1) is \(d=\dim X\), independently of
\(\dim Y_\nu\): the action

\[
 a\longmapsto(p_{Y_\nu})_*\bigl(p_X^*a\smile\operatorname{cl}(z_\nu)\bigr)
\]

adds \(2d\) and then subtracts \(2d\) from cohomological degree, with
cancelling Tate twists. The cycle is defined over \(k\), so its action
commutes with Frobenius. Different cohomological degrees have disjoint
Frobenius spectra, so a Frobenius-equivariant map between these
realizations preserves the degrees.

Suppose (3). By (CC-XC.2), choose \(f:H\to M\) with invertible
realization \(\phi:V\to W\). Its Frobenius equivariance shows that \(V\)
is semisimple, since \(W\) is. We must still prove that \(f\) has an
algebraic inverse; faithfulness alone would not establish this.

Poincaré duality gives an isomorphism

\[
 \delta_H:H\xrightarrow{\sim}H^\vee(-d).
\]

Transfer its realization across \(\phi\) to the Frobenius-equivariant
isomorphism

\[
 b=(\phi^\vee(-d))^{-1}\,R(\delta_H)\,\phi^{-1}:
       W\longrightarrow W^\vee(-d).               \tag{CC-XC.3}
\]

Both \(M\) and \(M^\vee(-d)\) belong to \(\mathcal A_{q,E}\). Its full
realization theorem, AQ.15, therefore lifts \(b\) to an algebraic
morphism \(\delta_M:M\to M^\vee(-d)\). This step requires fullness only
inside the already established abelian sector.

Now form the algebraic morphism

\[
 g=\delta_H^{-1}\circ f^\vee(-d)\circ\delta_M:
          M\longrightarrow H.                    \tag{CC-XC.4}
\]

Using (CC-XC.3), its realization is

\[
 R(g)=R(\delta_H)^{-1}\,
       \phi^\vee(-d)\,
       (\phi^\vee(-d))^{-1}\,
       R(\delta_H)\,\phi^{-1}
       =\phi^{-1}.
\]

Faithfulness of the homological realization now gives \(gf=1_H\) and
\(fg=1_M\). Thus (3) implies (2). This argument needs neither a separate
surjectivity assumption for reverse cross-correspondences nor
numerical--homological equivalence for \(X\).

Condition (2) implies (1), and an isomorphism in (2) has invertible
realization, so it implies (3). If (1) holds, say \(H\simeq M'\) with
\(M'\in\mathcal A_{q,E}\), then \(V\) is semisimple. Hence
\(R(M')\simeq V\simeq V^{\mathrm{ss}}\simeq W\). Full faithfulness
within \(\mathcal A_{q,E}\) lifts this Frobenius-module isomorphism and
its inverse to \(M'\simeq M\); therefore (1) implies (2).

Finally (3) implies (4). Conversely, a nonzero polynomial on a
finite-dimensional vector space over the infinite field \(E\) is nonzero
at some point: after choosing coordinates this follows by induction on
the number of variables, using the finitely many roots of a nonzero
one-variable polynomial. Apply this to \(\det|_I\) to deduce (3) from
(4). This proves all equivalences and the semisimplicity conclusion.
\(\square\)

\textbf{Finite witnesses.} The image \(I\) is spanned by the images of
algebraic cycles and is finite-dimensional. Consequently it has a finite
spanning collection of such images. Equivalently, the condition is
witnessed by finitely many cycles \(z_1,\ldots,z_t\) for which

\[
 \det\!\left(\sum_{j=1}^t x_j\kappa(z_j)\right)
\]

is not the zero polynomial. This is an existence criterion; it supplies
no effective method or complexity bound for finding the cycles.

\subsection{CC-XC03 --- The semisimple version as a finite projected
Tate
condition}\label{cc-xc03-the-semisimple-version-as-a-finite-projected-tate-condition}

\textbf{Corollary.} If Frobenius on \(V\) is semisimple, then the
conditions in CC-XC02 are also equivalent to the surjectivity of
\(\kappa\).

\textbf{Proof.} Under semisimplicity, \(V\simeq W\) as Frobenius
modules. Surjectivity then puts such an isomorphism in \(I\), so CC-XC02
applies. Conversely, if \(H\simeq M\), transport any
Frobenius-equivariant map \(V\to W\) across that isomorphism to an
endomorphism of \(W\). Full faithfulness on \(M\) supplies its algebraic
endomorphism. Composing gives the required morphism \(H\to M\), and
(CC-XC.2) places its realization in \(I\). Thus \(\kappa\) is onto.
\(\square\)

One sufficient hypothesis for this projected surjectivity is the usual
Tate surjectivity in \textbf{codimension \(d\)} on each of the finitely
many \(X\times Y_\nu\):

\[
 \mathrm{CH}^d(X\times Y_\nu)_{\mathbf Q}
       \otimes_{\mathbf Q}\mathbf Q_\ell
 \twoheadrightarrow
 H^{2d}((X\times Y_\nu)_{\bar k},\mathbf Q_\ell(d))^{F=1}.
                                                        \tag{CC-XC.5}
\]

Indeed, Poincaré duality and Künneth identify this invariant space with
the Frobenius-equivariant maps from \(H^*(X)\) to \(H^*(Y_\nu)\). Extend
scalars to \(E\), sum, and project by \(e\). Every map to \(W=eR(T)\) is
obtained by projecting its composition with the inclusion
\(W\hookrightarrow R(T)\), so (CC-XC.5) implies the surjectivity of
(CC-XC.1). This remains a condition on a specified finite family and
specified codimension. It requires no Tate assertion for all varieties
or all codimensions. The projected condition can be weaker than full
surjectivity in (CC-XC.5).

Without semisimplicity, surjectivity onto
\(\operatorname{Hom}_{E[F]}(V,W)\) alone would not guarantee an
invertible map: all such maps could have nontrivial kernel. The
determinant criterion is the exact formulation retaining that issue.

\subsection{CC-XC04 --- Tate and numerical--homological equality after
the
criterion}\label{cc-xc04-tate-and-numericalhomological-equality-after-the-criterion}

\textbf{Corollary.} Suppose the equivalent conditions of CC-XC02 hold
for \(X\). Let \(k'/k\) be any finite extension, let \(n\ge0\), let
\(A/k'\) be an abelian variety and \(S/k'\) finite étale, and put

\[
 Z=X_{k'}^{\,n}\times A\times S.
\]

For the fixed prime \(\ell\), Frobenius on
\(H^*(Z_{\bar k},\mathbf Q_\ell)\) is semisimple; the rational Tate
cycle map is surjective in every codimension; and numerical equivalence
and \(\ell\)-adic homological equivalence agree on rational algebraic
cycles of \(Z\).

\textbf{Complete proof.} If \(Z\) is empty, all assertions are vacuous,
so assume it is nonempty. Base change of algebraic correspondences
preserves the isomorphism \(h_{\ell,E}(X)\simeq M\) proved in CC-XC02.
Its base change remains an abelian--Artin motive over \(k'\). Tensoring,
and then multiplying by \(h(A)\) and \(h(S)\), gives an isomorphism of
\(h_{\ell,E}(Z)\) with an object of \(\mathcal A_{q',E}\), where
\(q'=\#k'\). AQ.15 applies over every finite field. It proves
semisimplicity after extension to \(E\), and semisimplicity over
\(\mathbf Q_\ell\) follows because a characteristic-zero operator is
semisimple if and only if its extension to an algebraic closure is.

For codimension \(r\), fullness on the abelian sector identifies
algebraic morphisms \(\mathbf 1\to h_{\ell,E}(Z)(r)\) with the
Frobenius-invariant vectors of its realization. The weight decomposition
shows that only cohomological degree \(2r\) contributes. Consequently

\[
 \operatorname{im}\!\left(
 \mathrm{CH}^r(Z)_{\mathbf Q}\otimes_{\mathbf Q}E
 \longrightarrow H^{2r}(Z_{\bar k},E(r))
 \right)
 =H^{2r}(Z_{\bar k},E(r))^{F=1}.                 \tag{CC-XC.6}
\]

Here is the descent to the usual \(\mathbf Q_\ell\)-valued Tate
statement. Let \(C_\ell\) denote the image of
\(\mathrm{CH}^r(Z)_{\mathbf Q}\otimes_{\mathbf Q}\mathbf Q_\ell\) in
\(U=H^{2r}(Z_{\bar k},\mathbf Q_\ell(r))^{F=1}\). Flat scalar extension
commutes both with that image and with the kernel defining \(F=1\)
invariants. Equation (CC-XC.6) therefore says
\(C_\ell\otimes_{\mathbf Q_\ell}E=U\otimes_{\mathbf Q_\ell}E\). Faithful
flatness of a field extension gives \(U/C_\ell=0\). This argument uses
the image subspace; it does not presume injectivity after scalar
extension of rational homological classes.

Write \(D=\dim Z\). The Poincaré pairing

\[
 H^{2r}(Z_{\bar k},\mathbf Q_\ell(r))
 \times
 H^{2D-2r}(Z_{\bar k},\mathbf Q_\ell(D-r))
 \longrightarrow\mathbf Q_\ell
\]

is perfect and Frobenius invariant. Its restriction to the two invariant
subspaces is perfect: after diagonalizing semisimple Frobenius over
\(E\), eigenvalues \(\lambda,\mu\) can pair nontrivially only when
\(\lambda\mu=1\); the \(1\)-eigenspace is therefore paired perfectly
with the other \(1\)-eigenspace. Nondegeneracy descends to
\(\mathbf Q_\ell\).

If a rational cycle \(z\) is numerically trivial but
\(a=\operatorname{cl}_\ell(z)\ne0\), this invariant perfect pairing
provides an invariant complementary class \(b\) with
\(\langle a,b\rangle\ne0\). The already proved rational Tate
surjectivity expresses \(b\) as a finite \(\mathbf Q_\ell\)-linear
combination of classes of rational complementary cycles. Numerical
triviality makes every corresponding intersection degree zero,
contradicting the chosen pairing. Thus numerical triviality implies
homological triviality. The converse follows directly from compatibility
of the cycle class with intersection degrees. This proves the assertion
on rational cycles without assuming numerical realization at the start.
\(\square\)

The conclusion is for the specified prime \(\ell\) and this class of
products and base extensions. It does not assert the determinant
condition for arbitrary \(X\), or a comparison with another prime
without an additional argument.

\subsection{Relation to Milne and the remaining
question}\label{relation-to-milne-and-the-remaining-question}

Milne (1994), \textbf{Proposition 1.15 and Corollary 1.16},
author-manuscript pp.~11--12, derive semisimplicity,
numerical--homological equality and full realization from the full Tate
conjecture. \textbf{Proposition 2.6 and Remark 2.7}, p.~15, then give
the global abelian/Artin-generation conclusion under that hypothesis.
Those inspected statements do not themselves state the particular
one-sided finite-target criterion above. Here we prove that criterion
using only the already settled abelian-sector realization and the
specified cycle condition for \(X\). No historical priority claim is
made.

For a particular \(X\), the unresolved requirement has therefore been
made concrete: construct finitely many cross-correspondences whose
projected realization has nonzero determinant. Pure numerical Frobenius
eigenvalues alone do not establish this. The adopted RH and Dirichlet
GRH do not provide these cycles or prove the determinant condition.

\subsection{Sources and locators}\label{sources-and-locators}

\begin{itemize}
\tightlist
\item
  Deligne, P. (1974). La conjecture de Weil: I. \emph{Publications
  Mathématiques de l'IHÉS, 43}, 273--307. \textbf{Theorem 1.6, p.~276}.
  \url{https://www.numdam.org/item/PMIHES\_1974\_\_43\_\_273\_0.pdf}.
\item
  Honda, T. (1968). Isogeny classes of abelian varieties over finite
  fields. \emph{Journal of the Mathematical Society of Japan, 20}(1--2),
  83--95. \url{https://doi.org/10.2969/jmsj/02010083}. Main Theorem, p.~83;
  primary PDF:
  \url{https://people.math.rochester.edu/faculty/doug/otherpapers/Honda-Ab.pdf}.
\item
  Milne, J. S. (1994). Motives over finite fields. In U. Jannsen, S. L.
  Kleiman, \& J.-P. Serre (Eds.), \emph{Motives} (Proceedings of
  Symposia in Pure Mathematics, Vol. 55, Part 1, pp.~401--459). American
  Mathematical Society. \url{https://jmilne.org/math/articles/1994aP.pdf}.
  Author-manuscript pp.~5--6 for finite-field projectors, corrected
  symmetry and duality; pp.~11--12 for Proposition 1.15 and Corollary
  1.16; p.~15 for Proposition 2.6 and Remark 2.7.
\item
  Proved internal inputs: \textbf{CC-AG01} in
  research/consequences\_criteria\_abelian\_group.tex;
  \textbf{AQ.15/AQ-ABELIAN-RECOGNITION} and \textbf{AQ-WEIL-EXIST} in
  research/archive\_questions.tex. Their exact all-abelian primary
  inputs and limits are recorded in CC-TATE01--02 and the corresponding
  source ledger.
\end{itemize}

% END consequences_criteria_cross_cycles.tex

% BEGIN root_springer_archive.tex
\section{Cote 123: the tangent-rank prediction, Springer dimensions, and positivity}
\label{sec:cote123}

\subsection{The original letter and a corrected locator}
The original is Grothendieck's letter to Kostant dated 22 October 1969,
cote 123, physical PDF pages 22--23, manuscript and archive footers
21--22, in \url{https://grothendieck.umontpellier.fr/123.pdf}.
The official PDF has a cover. Thus the intake's locator ``123/023''
for the dimension question is a physical-PDF locator, not footer 23:
footer 23, physical page 24, is the nearly blank verso.
The entire two-page letter was inspected visually in this run.

On its second page Grothendieck asks how the tangent rank of the
adjoint quotient \(\chi:\mathfrak g\to\mathfrak h/W\), and of
the analogous group quotient, varies. For an orbit of dimension
\(\dim G-r-2\), he predicts rank \(r-1\) when all roots have
the same length, and strictly smaller rank otherwise. He also asks
for the dimension of the variety of Borel subgroups containing an
element, or whose Lie algebra contains a given Lie-algebra element.
Writing \(2N=\dim G-r\), an orbit of dimension \(2N-2i\)
is predicted to have such a variety of dimension \(i\).
The initial setting of the letter is a complex simple adjoint group.
These are mathematical questions in the letter itself, distinct
from the later intake's finite-field positivity interpretation.

\paragraph{Known answer to the dimension question.}
In the complex setting the asserted dimension is indeed
\[
 \dim\mathcal B_x=\frac{\dim Z_G(x)-r}{2}.
 \tag{123.1}
\]
Steinberg (1976), \textbf{Theorem 4.6, p.~217}, proves it for
unipotent group elements; \textbf{Theorem 6.5, p.~223}, proves
it for arbitrary group elements. His sufficient-characteristic
condition is automatic over \(\mathbf C\); it must not be
silently discarded when applying that paper in positive
characteristic. The Lie-algebra reduction is Springer's (1984)
\textbf{Lemma 1, p.~273}: split off the semisimple part and
pass to its connected reductive centralizer. In characteristic
zero a nilpotent element and its exponential have the same
centralizer and the same containing Borel subgroups. Applying
Steinberg's unipotent formula in that centralizer gives (123.1)
for the Lie-algebra question as well. This is an application
of known results, not a new proof of their geometry.

\subsection{BR-12301. The tangent rank in type A}
Over \(\mathbf C\), for \(X\in\mathfrak{gl}_n\), the tangent
rank of its adjoint quotient is the degree of its minimal
polynomial. For \(X\in\mathfrak{sl}_n\), it is that degree
minus one. In particular, a nilpotent matrix with largest
Jordan block of size \(b\) has ranks \(b\) and \(b-1\),
respectively. The same \(\mathfrak{sl}_n\) formula applies
to the Lie algebra of the adjoint group \(\mathrm{PGL}_n\).

\begin{proof}
The coefficients of the characteristic polynomial generate
the conjugation-invariant ring on \(\mathfrak{gl}_n\).
For clarity, one can check this assertion without choosing
generators from an abstract restriction theorem. Matrices
with distinct eigenvalues form a dense open subset; each is
conjugate to a diagonal matrix. The restriction of an
invariant polynomial is symmetric in the diagonal entries.
The elementary symmetric polynomials generate these symmetric
polynomials and lift to the characteristic coefficients.
Density proves that restriction is injective, which proves
the assertion. On trace-zero matrices the same argument uses
the hyperplane \(\sum_i\lambda_i=0\), and drops the first
characteristic coefficient. A polynomial on that hyperplane
invariant under the finite permutation group can be extended
invariantly to the whole diagonal space by averaging any
extension; hence its invariant ring is precisely the
symmetric-polynomial ring modulo the first coefficient.

In characteristic zero Newton identities replace the
characteristic coefficients, by a polynomial change of
coordinates with everywhere invertible triangular Jacobian,
with \(p_j(X)=\operatorname{Tr}(X^j)\), \(1\leq j\leq n\).
The differential is
\[
 dp_j|_X(Y)=j\operatorname{Tr}(X^{j-1}Y).
\]
The trace pairing on all matrices is nondegenerate, so their
span has dimension
\(\dim\langle 1,X,\ldots,X^{n-1}\rangle=\deg m_X\),
where \(m_X\) is the minimal polynomial. The last equality
is its defining minimality, together with division of
polynomials by \(m_X\).

The annihilator of \(\mathfrak{sl}_n\) for the trace pairing
is the scalar line: off-diagonal matrix units first force
a matrix in the annihilator to be diagonal, and differences
of diagonal units force all its entries to agree. The
quotient on \(\mathfrak{sl}_n\) uses \(p_2,\ldots,p_n\).
Their differentials therefore have rank equal to the
dimension of the image of \(\langle X,\ldots,X^{n-1}\rangle\)
modulo the scalar line. Together with that line these
powers span \(\mathbf C[X]\), so this rank is
\(\deg m_X-1\). For a nilpotent matrix \(m_X=t^b\).
Passing from \(\mathrm{SL}_n\) to its adjoint quotient
does not change the Lie algebra or its conjugation orbits.
This proves all assertions.
\end{proof}

\subsection{BR-12302. The strict nonsimply-laced rank prediction is false}
The strict inequality in the letter is refuted by the
complex simple adjoint group \(G=\mathrm{PSp}_4\), of type
\(C_2\). Put
\[
 J=\begin{pmatrix}0&I_2\\-I_2&0\end{pmatrix},\qquad
 N=\begin{pmatrix}0&I_2\\0&0\end{pmatrix}
       \in\mathfrak{sp}_4.
\]
Its orbit has dimension \(6=\dim G-r-2\), but
\(\operatorname{rank}(d\chi_N)=1=r-1\), despite the two
root lengths. The unipotent element represented by \(U=I+N\)
gives the same counterexample for the group quotient.
This refutes the stated strict inequality;
it does not assert a complete rank classification in every type.

\begin{proof}
Every element of \(\mathfrak{sp}_4\) has the form
\[
 Y=\begin{pmatrix}A&B\\C&-A^t\end{pmatrix},
 \qquad B=B^t,\quad C=C^t.
\]
Thus its dimension is \(4+3+3=10\) and its rank is two.
Direct multiplication gives
\[
 [Y,N]=\begin{pmatrix}-C&A+A^t\\0&C\end{pmatrix}.
\]
The Lie centralizer has \(C=0\), arbitrary symmetric
\(B\), and skew-symmetric \(A\), so dimension \(3+1=4\).
In characteristic zero an algebraic group stabilizer is
smooth; equivalently the conjugation orbit map has its
usual Lie-algebra tangent quotient. Therefore the orbit
has dimension \(10-4=6\). Passing through the finite
central isogeny to \(\mathrm{PSp}_4\) does not change it.

For completeness its adjoint quotient is explicitly
\(\chi(Y)=(c_2(Y),c_4(Y))\), where
\[
 \det(tI-Y)=t^4+c_2(Y)t^2+c_4(Y).
\]
Indeed, a regular semisimple symplectic matrix in the Lie
algebra has eigenvalues \(\lambda_1,\lambda_2,-\lambda_1,
-\lambda_2\). The symplectic form pairs opposite eigenspaces,
so an eigenbasis can be normalized to a symplectic basis.
Such matrices form a dense open subset. The normalizer
acts on the two \(\lambda_i\) by independent sign changes
and permutation. Its polynomial invariants are the
symmetric polynomials in \(\lambda_1^2,\lambda_2^2\),
generated by their sum and product. These lift to
\(-c_2\) and \(c_4\). Restriction and density prove
that these two coefficients generate the full invariant
ring, just as in the preceding proof.

Newton's identity gives \(c_2(Y)=-\tfrac12\operatorname{Tr}(Y^2)\).
Consequently \(dc_2|_N(Y)=-\operatorname{Tr}(NY)=-\operatorname{Tr}C\),
which is nonzero, for example when \(C=\operatorname{diag}(1,0)\).
Since \(N\) has rank two, all its three-by-three minors
vanish; the differential of \(c_4=\det\), whose entries
are those cofactors, is zero at \(N\). The quotient
differential thus has rank exactly one, as claimed.

Here is the group version, so that the letter's use of
unipotent orbits causes no ambiguity. Since \(N^2=0\) and
\(N^tJ=-JN\), direct expansion gives \(U^tJU=J\).
The centralizer of \(U=I+N\) is the centralizer of \(N\),
so its orbit also has dimension six. On a maximal torus
of \(\mathrm{Sp}_4\) the eigenvalues are
\(z_1,z_2,z_1^{-1},z_2^{-1}\). Put \(w_j=z_j+z_j^{-1}\).
Inversion-invariant Laurent polynomials in one variable
are polynomials in \(z+z^{-1}\), by the recurrence for
\(z^m+z^{-m}\). Independent inversions and permutation
therefore give invariant generators
\(t=w_1+w_2=\operatorname{Tr}g\) and
\(e=w_1w_2=c_2(g)-2\). Restriction to the dense regular
semisimple locus shows these generate the group invariant
ring, exactly as above. The central element \(-I\) sends
\((t,e)\) to \((-t,e)\); hence the adjoint group quotient
has coordinates \((t^2,e)\). At \(U\), \(t=4\), so
\((t,e)\mapsto(t^2,e)\) has invertible differential.

A tangent vector at \(U\) is \(UY\) with
\(Y\in\mathfrak{sp}_4\). Using
\(c_2(g)=((\operatorname{Tr}g)^2-\operatorname{Tr}(g^2))/2\),
\(U^2=I+2N\), and \(\operatorname{Tr}Y=0\), we obtain
\[
 dt_U(UY)=\operatorname{Tr}(NY),\qquad
 de_U(UY)=4\operatorname{Tr}(NY)-\operatorname{Tr}(U^2Y)
          =2\operatorname{Tr}(NY).
\]
The first functional is nonzero by the same choice of
\(C\), so the group quotient also has rank one, including
after descent to \(\mathrm{PSp}_4\).
\end{proof}

\subsection{BR-12303. The exact positive point-count consequence}
Let \(X_x\) be the Springer fibre of a nilpotent Lie-algebra
element \(x\) over \(\mathbf F_q\), with \(G\) connected reductive and residue
characteristic good for \(G\). Springer's (1984)
\textbf{Theorem 1, pp.~272--273}, says that \(H^j(X_x,\mathbf Q_\ell)\)
is pure of weight \(j\). His \textbf{Corollary 1 to Theorem 2,
p.~280}, gives odd vanishing, and \textbf{Lemma 4 on that page}
says that the eigenvalues on \(H^{2i}\) are \(q^i\) times
roots of unity. Here \(\ell\ne\operatorname{char}\mathbf F_q\).
These are the relevant finite-field inputs; purity alone
would not give the root-of-unity assertion.

There exists \(m\geq1\) such that, for all \(r\geq1\),
\[
 \#X_x(\mathbf F_{q^{mr}})
       =\sum_i b_{2i}q^{mri},\qquad
 Z(X_x/\mathbf F_{q^m},t)
       =\prod_i(1-q^{mi}t)^{-b_{2i}},
 \tag{123.2}
\]
where \(b_{2i}=\dim H^{2i}(X_x,\mathbf Q_\ell)\).
If all the root-of-unity factors are already one, one
can take \(m=1\). No Weyl-group trace is inserted in (123.2).

\begin{proof}
There are finitely many eigenvalues. Choose \(m\) divisible
by the orders of all their root-of-unity factors. The
eigenvalues of \(F^{mr}\) on \(H^{2i}\) are then all
\(q^{mri}\); this statement about traces needs no
semisimplicity assumption. The proper trace formula,
together with odd vanishing, proves the count. Substituting
it in \(Z(t)=\exp(\sum_{r\geq1}\#X_x(\mathbf F_{q^{mr}})t^r/r)\)
and using \(-\log(1-z)=\sum_{r\geq1}z^r/r\) proves the
displayed product as a formal power series. Every exponent
\(b_{2i}\) is nonnegative. This proves the stated consequence.
\end{proof}

\subsection{BR-12304. A Green function with a negative coefficient}
The intake SA20, \S4, D.018, goes from positive ordinary
counts to nonnegative coefficients of Green polynomials.
That implication is false. Use the Springer action and the
normalization of Springer (1984), \S5, p.~279:
\[
 Q_w(x)=\sum_i(-1)^i\operatorname{Tr}
       (Fw\mid H^i(X_x,\mathbf Q_\ell)).
 \tag{123.3}
\]
For \(G=\mathrm{SL}_2\), \(x=0\), odd \(q\), and the
nonidentity element \(s\in W=S_2\),
\[
 X_0=\mathbf P^1,\qquad Q_1(0)=1+q,
       \qquad Q_s(0)=1-q.
 \tag{123.4}
\]
The last polynomial has a negative coefficient. The
trace calculation holds for every odd \(q\). To invoke
the Green-function identification exactly as stated on
Springer's p.~279, take odd \(q\) sufficiently large;
the same polynomial is then a Green polynomial in that
source's terminology.

\begin{proof}
The regular semisimple orbit with eigenvalues \(1,-1\)
is the variety
\(U=(\mathbf P^1\times\mathbf P^1)\setminus\Delta\)
of ordered distinct eigenlines. The nonidentity Weyl element
interchanges the two lines. For \(x=0\), the simultaneous
resolution identifies its cohomology with that of \(X_0\);
this is precisely \textbf{Proposition 3 and the definition
of the Weyl action in \S5, p.~279}, of Springer (1984).
Thus it suffices to compute the swap action on \(H^*(U)\).

Projection \(U\to\mathbf P^1\) is a Zariski locally trivial
affine-line bundle: over a fixed line, complements to that
line form an affine line, and choosing a local complement
trivializes this torsor. Homotopy invariance gives cohomology
only in degrees zero and two, both of dimension one.
Let \(h_1,h_2\) be the two divisor classes on the product.
The divisor \(\Delta\) has class \(h_1+h_2\), as is seen
by intersecting it with each ruling. The localization
sequence for this smooth divisor yields
\[
 H^2(U,\mathbf Q_\ell)
   =\bigl(\mathbf Q_\ell h_1\oplus\mathbf Q_\ell h_2\bigr)
       /\mathbf Q_\ell(h_1+h_2).
\]
Here the adjacent \(H^1(\Delta)(-1)\) is zero and the
Gysin image is exactly the displayed divisor class.
The swap sends \(h_1\) to \(h_2=-h_1\), so it acts
by minus one on this one-dimensional group and by one
on \(H^0\). Both divisor classes are rational over
\(\mathbf F_q\), hence geometric Frobenius acts by
\(q\) on \(H^2\), and by one on \(H^0\).
Equation (123.3) now gives (123.4).
\end{proof}

The dimension formula and the ordinary positive count remain
valid. The failure occurs at the added Weyl trace, whose
character values need not be nonnegative. This is a concrete
boundary for any proposed passage from purity to positivity.
No RH or Dirichlet-GRH input is used in this section.

\subsection{Primary bibliography and inspection scope}
Grothendieck, A. (1969, October 22). \emph{Letter to B. Kostant}.
Montpellier archives, cote 123, physical pages 22--23,
footers 21--22. \url{https://grothendieck.umontpellier.fr/123.pdf}.
The letter also raises questions about opposed tori and rational
singularities; these are separate from the four statements proved
or cited above and are retained in the question register.

Springer, T. A. (1984). A purity result for fixed point varieties
in flag manifolds. \emph{Journal of the Faculty of Science,
University of Tokyo, Section IA, Mathematics, 31}(2), 271--282.
\url{https://doi.org/10.15083/00039529}.
Primary scan: \url{https://repository.dl.itc.u-tokyo.ac.jp/record/39538/files/jfs310202.pdf}.
Theorem 1, pp.~272--273; Lemma 1, p.~273; Proposition 3 and
the action convention, p.~279; Corollary 1 and Lemma 4, p.~280.
The stated results and the indicated action construction were read;
the full purity proof and its preceding geometric sources are
cited inputs, not independently reproved here.

Steinberg, R. (1976). On the desingularization of the unipotent
variety. \emph{Inventiones Mathematicae, 36}, 209--224.
\url{https://doi.org/10.1007/BF01390010}.
Theorem 4.6 and its characteristic convention, p.~217;
Remark 6.3 and Theorem 6.5, p.~223. Primary digitization:
\url{https://gdz.sub.uni-goettingen.de/download/pdf/PPN356556735_0036/LOG_0018.pdf}.
The complete short proofs of the two numbered statements were
read with their earlier classification and dimension results
retained as cited inputs.

% END root_springer_archive.tex

% BEGIN source_cote123_rational.tex
\section{Cote 123: the rational singularities of adjoint-quotient fibres}
\label{sec:SM-cote123-rational}

The letter to Kostant dated 22 October 1969, physical pages 22--23
of cote 123 (handwritten pages 21--22), is explicitly over
$\mathbb C$.  It asks whether the fibres of
$\mathfrak g\to\mathfrak h/W$, and of $G\to T/W$ when $G$ is
simply connected, have rational singularities.  Its stated test is
$R^if_*\mathcal O=0$ for $i>0$ on a desingularization.
This is a known positive result, independent of any adopted
analytic-number-theory input.

\paragraph{Finding F-SM-23 (known affirmative answer, with the reduction written out).}
Let $\mathfrak g$ be a complex semisimple Lie algebra.
Every scheme-theoretic fibre of its adjoint quotient
$\mathfrak g\to\mathfrak h/W$ is reduced and normal and has
rational singularities.  Let $G$ be a connected simply connected
complex semisimple algebraic group, with maximal torus $T$ and Weyl
group $W$.  The same assertions hold for every scheme-theoretic
fibre of $G\to T/W$.
Thus for each such fibre $F$ and every proper resolution
$f:Y\to F$,
\[
 f_*\mathcal O_Y=\mathcal O_F,\qquad
 R^if_*\mathcal O_Y=0\quad(i>0).
\]
These assertions use classical results; the following argument is a
complete deduction from the named inputs, not a new proof of those
inputs.

\begin{proof}
\emph{The Lie-algebra statement.}
Kraft and Procesi (1982, \S0.3, p.~540) state reducedness and
normality for every fibre.  Their immediately following \S0.4 on
the same page explicitly states that these fibres have rational
singularities, citing Hesselink.  The scope of \S0.2 is a reductive
complex group; these introductory assertions are not restricted
to the classical orbit-closure classification treated later in that
paper.  This supplies the first statement directly.

\emph{The group fibre as a scheme.}
Put $r=\operatorname{rank}G$.  Steinberg's fundamental-character
coordinates identify $T/W$ with $\mathbb A^r$ (1965,
Theorem~6.1(b), p.~62, and \S6.16, p.~65).
For a fibre $F$, Theorem~6.11, p.~64, gives an irreducible
underlying variety of codimension $r$, with a dense regular class
whose complement has codimension at least two.
The fibre ideal is generated by $r$ functions in the smooth
variety $G$.  Its height is $r$, so these functions form a regular
sequence locally.  Thus $F$ is Cohen--Macaulay.  Theorem~8.1,
p.~70, says that their differentials are independent on the
regular class.  Therefore $F$ is generically reduced; the absence
of embedded associated primes for a Cohen--Macaulay scheme implies
that $F$ is reduced.  Its complement to the smooth regular locus
has codimension at least two.  Serre's $R_1$ and $S_2$ criterion
then gives normality.  Normality is a deduction here; it is not
misattributed to the literal statement of Theorem~6.11.

\emph{Reduction to a unipotent variety.}
Let $s$ represent the unique closed semisimple class in $F$ and
put $L=Z_G(s)$.  This is a connected reductive group: reductivity
and connectedness are the semisimple-centralizer facts recorded
in Steinberg (1965, \S2.7 and Remark~2.10, p.~53).  The latter
remark states the connectedness result with proof deferred; we
use the established result rather than claim to reproduce that
proof.

Since $F$ is the closure of its regular class, the algebraic
semisimple-part map of Kraft--Procesi (\S0.2, pp.~539--540) gives
a $G$-equivariant morphism
\[
 j:F\longrightarrow G\cdot s=G/L.
\]
Its fibre over $s$ is $s\mathcal U(L)$, where $\mathcal U(L)$ is
the unipotent variety.  Indeed, Jordan decomposition gives that
description for every element with semisimple part $s$; conversely
every $su$, $u\in\mathcal U(L)$, has the same invariant-character
values as $s$ and belongs to $F$.  The standard homogeneous-space
descent, applied to $G\to G/L$, therefore gives
\[
 F \simeq G\times^L \mathcal U(L),\qquad
 [g,u]\longmapsto gsu g^{-1}.
 \tag{SM23.1}
\]
Here the right-hand fibre has its reduced structure.  This also
agrees with the scheme fibre: after the smooth faithfully flat
base change $G\to G/L$, the left side is
$G\times j^{-1}(s)$ and is reduced, since $F$ is reduced.
Faithfully flat descent of reducedness over $\mathbb C$ makes
$j^{-1}(s)$ reduced.  The preceding identification of reduced
closed subvarieties is consequently scheme-theoretic.

Hesselink (1976, Theorem~A and \S2, p.~249) supplies the
equivariant Springer resolution
\[
 \tau:Y_L\longrightarrow\mathcal U(L),\qquad
 \tau_*\mathcal O_{Y_L}=\mathcal O_{\mathcal U(L)},\qquad
 R^i\tau_*\mathcal O_{Y_L}=0\quad(i>0).
 \tag{SM23.2}
\]
The theorem is stated for a split reductive group in
characteristic zero, and hence applies to $L$.
The construction is the Springer resolution of the nilpotent
cone; the equivariant logarithm identifies that cone with
$\mathcal U(L)$ in characteristic zero.  Its equivariance is
essential for the next step.

Form the associated morphism
\[
 f:G\times^L Y_L\longrightarrow G\times^L\mathcal U(L)\simeq F.
\]
The smooth $L$-torsor $G\to G/L$ has sections \'etale locally.
For such a cover $V\to G/L$, the base change of $f$ is
$\operatorname{id}_V\times\tau$.
It is proper and birational, with smooth source, so these
properties descend and $f$ is a resolution.
Flat base change for $\tau$, followed by (SM23.2), gives
\[
 R^i(\operatorname{id}_V\times\tau)_*
       \mathcal O_{V\times Y_L}
 =
 \begin{cases}
 \mathcal O_{V\times\mathcal U(L)},&i=0,\\
 0,&i>0.
 \end{cases}
\]
The corresponding assertions for $Rf_*\mathcal O$ descend along
the resulting faithfully flat \'etale cover of $F$.
This proves rational singularities of $F$.

Finally, the independence of the chosen resolution in
characteristic zero is part of the standard rational-singularity
criterion: the definition with \emph{some} resolution is equivalent
to the definition with \emph{every} resolution.  A directly checked
primary research reference stating this equivalence is
Glazer (2019, \S2A, p.~488), which cites
Kempf--Knudsen--Mumford--Saint-Donat, I.3, pp.~50--51.
It yields the displayed conclusion for every resolution.
\end{proof}

\paragraph{Relation to the particular resolution in the letter.}
Steinberg (1976, Corollary~6.4, p.~223), explicitly attributed
there to Grothendieck, states that for the closure $F$ of a regular
conjugacy class and $t\in T\cap F$ the morphism
\[
 \{(x,gB):g^{-1}xg\in tU\}\longrightarrow F
\]
is a desingularization.  This applies under the present complex
simply connected hypothesis.  F-SM-23 and independence of the
resolution give the requested higher-direct-image vanishing for
this resolution as well.  Corollary~6.4 alone supplies a
resolution, not the higher-direct-image vanishing; the latter
comes from rational singularities as explained above.

\paragraph{Scope and status.}
This closes the rational-singularity question in the original
complex setting.  It is not presented as a new theorem of this
project, as a characteristic-positive analogue, or as a statement
about arbitrary individual orbit closures.  No general
non-simply-connected group assertion is needed here.

\subsection*{References for F-SM-23}
\begin{itemize}
\item Glazer, I. (2019). On rational singularities and counting points
of schemes over finite rings. \emph{Algebra \& Number Theory, 13}(2),
485--500. \url{https://doi.org/10.2140/ant.2019.13.485}.
Exact input: \S2A, p.~488.
\item Hesselink, W. H. (1976). Cohomology and the resolution of the
nilpotent variety. \emph{Mathematische Annalen, 223}(3), 249--252.
\url{https://doi.org/10.1007/BF01360956}.
Exact inputs: Theorem~A and \S2, p.~249.
\item Kraft, H., \& Procesi, C. (1982). On the geometry of conjugacy
classes in classical groups. \emph{Commentarii Mathematici
Helvetici, 57}, 539--602.
\url{https://doi.org/10.1007/BF02565876}.
Exact inputs: \S0.2, pp.~539--540; \S\S0.3--0.4, p.~540.
\item Steinberg, R. (1965). Regular elements of semi-simple algebraic
groups. \emph{Publications Math\'ematiques de l'IH\'ES, 25}, 49--80.
\url{https://doi.org/10.1007/BF02684397}.
Exact inputs: \S2.7 and Remark~2.10, p.~53;
Theorem~6.1(b), p.~62; Theorem~6.11, p.~64;
\S6.16, p.~65; Theorem~8.1, p.~70.
\item Steinberg, R. (1976). On the desingularization of the unipotent
variety. \emph{Inventiones Mathematicae, 36}, 209--224.
\url{https://doi.org/10.1007/BF01390010}.
Exact input: Corollary~6.4, p.~223.
\end{itemize}

% END source_cote123_rational.tex

% BEGIN consequences_criteria_apposition.tex
\section{Cote 123: apposition, opposition, and rigid quadruples in type
A}\label{cote-123-apposition-opposition-and-rigid-quadruples-in-type-a}

\textbf{Scope and status.} These results answer the relevant original
questions for the complex simple adjoint groups
\(G=\operatorname{PGL}_n(\mathbf C)\), \(n\geq2\). All subgroup
automorphisms are complex algebraic automorphisms. No RH or
Dirichlet-GRH input is used. The arithmetic assertion at the end
concerns subgroup schemes individually defined over a specified subfield
\(k\subset\mathbf C\). Literature cutoff: 20 November 2026; source
inspection and proofs: 21 November 2026 UTC. No historical novelty is
asserted.

Grothendieck's letter of 22 October 1969, cote 123, physical PDF page
22, printed footer 21, asks for the classification of quadruples
\((T,B,T',B')\) with the maximal tori in Kostant apposition, whether the
two Borels are opposite for some or all choices, whether a distinguished
conjugacy class can be invariant under outer automorphisms, and what
automorphism groups these quadruples have. The original page was
visually inspected. The distinctions between these questions are
preserved below.

\begin{longtable}[]{@{}
  >{\raggedright\arraybackslash}p{(\columnwidth - 2\tabcolsep) * \real{0.5000}}
  >{\raggedright\arraybackslash}p{(\columnwidth - 2\tabcolsep) * \real{0.5000}}@{}}
\toprule\noalign{}
\begin{minipage}[b]{\linewidth}\raggedright
Result
\end{minipage} & \begin{minipage}[b]{\linewidth}\raggedright
Precise conclusion in type \(A_{n-1}\)
\end{minipage} \\
\midrule\noalign{}
\endhead
\bottomrule\noalign{}
\endlastfoot
CC-AP01 & Every ordered apposition pair is conjugate to the diagonal and
Fourier pair. \\
CC-AP02 & Opposition is exactly a specified list of nonvanishing Fourier
minors. \\
CC-AP03 & Every pair admits opposite Borel choices. All choices are
opposite if and only if \(n\) is prime. For composite \(n\), an explicit
common unipotent subgroup disproves opposition. \\
CC-AP04 & The inner-conjugacy classes are the explicit free quotient
\(K_n\backslash(S_n\times S_n)\), where
\(K_n=(\mathbf Z/n)^2\rtimes(\mathbf Z/n)^\times\). Their number is
\((n!)^2/(n^2\varphi(n))\). \\
CC-AP05 & For every odd \(n\geq3\), no such class is invariant under the
nontrivial outer automorphism; all quadruples then have trivial full
automorphism group. For even \(n\geq4\), outer-invariant examples exist,
with automorphism group \(C_2\). An exact finite criterion covers every
class. \\
CC-AP06 & A quadruple whose four subgroup schemes are defined over
\(k\subset\mathbf C\) exists in split \(\operatorname{PGL}_{n,k}\) if
and only if \(k\) contains a primitive \(n\)-th root of unity. In
particular there is no such quadruple over \(\mathbf Q\) for
\(n\geq3\). \\
CC-AP07 & Adding a generating principal element of \(T\cap N_G(T')\)
gives the free quotient \((\mathbf Z/n)^2\backslash(S_n\times S_n)\),
with \((n!)^2/n^2\) classes. \\
\end{longtable}

\subsection{Inspected primary
convention}\label{inspected-primary-convention}

Use Kostant's convention of a Cartan \(\mathfrak t'\) in apposition to
\(\mathfrak t\) with respect to a principal element \(P\). The intrinsic
characterization needed here is

\[
 \mathfrak t=\mathfrak g^{\operatorname{Ad}P},
 \qquad
 \operatorname{Ad}P(\mathfrak t')=\mathfrak t',
 \qquad
 \operatorname{Ad}P|_{\mathfrak t'}\text{ is a Coxeter element}.
 \tag{CC-AP.1}
\]

A lift of a Weyl-group Coxeter element produces this configuration, and
apposition has this Coxeter action. These are Kostant (1959),
\textbf{Theorem 8.6 and the immediately following Corollary 8.6,
p.~1026}. Their statements were read in the OCR of the author's
collected-paper reprint. Kostant's own later account provides an
independently accessible primary check: \textbf{Theorem 0.1 and the
following apposition sentence, arXiv:1003.0046v1 p.~2; Theorems 1.2--1.3
p.~6; Proposition 1.8 and Theorem 1.9 p.~8}. His principal element is
\(\exp(2\pi i\rho^\vee/h)\); regularity alone is not its definition.

The original definition in 1959, §7.3, p.~1018, was located through
later references but was \textbf{not directly inspected} in the original
scan. No claim of such inspection is made. The proof below uses the
inspected equivalent characterization (CC-AP.1), together with the
cyclic-element construction explicitly identified as apposition by
Kostant himself in 2010. The short proof of his Theorem 1.9 was read;
the full 1959 structure theorem is a cited input.

When using (CC-AP.1) at group level, take the connected centralizer
determined by its Lie algebra. In an adjoint group the full centralizer
of a regular semisimple group element can be disconnected.

\subsection{CC-AP01 --- The diagonal and Fourier normal
form}\label{cc-ap01-the-diagonal-and-fourier-normal-form}

Put \(I=\mathbf Z/n\mathbf Z\), fix \(\zeta=e^{2\pi i/n}\), and use the
basis \(e_0,\ldots,e_{n-1}\) of \(V=\mathbf C^n\). Define

\[
 C e_j=e_{j+1},\qquad
 Qe_j=\zeta^j e_j,\qquad
 v_a=\sum_{j\in I}\zeta^{ja}e_j,
 \qquad F=(\zeta^{ja})_{j,a\in I}.
\]

Let \(T_0\) be the diagonal maximal torus of \(\operatorname{PGL}(V)\),
and set \(T_1=FT_0F^{-1}\).

\textbf{Proposition.} The ordered pair \((T_0,T_1)\) is in Kostant
apposition, and every ordered apposition pair in
\(\operatorname{PGL}_n(\mathbf C)\) is conjugate to it.

\textbf{Proof.} The vectors \(v_a\) form a basis because the usual
Vandermonde determinant of the distinct numbers \(\zeta^a\) is nonzero.
Direct calculation gives

\[
 Cv_a=\zeta^{-a}v_a,\qquad Qv_a=v_{a+1},
 \qquad QCQ^{-1}=\zeta C.
\]

Thus \(T_1\) is the connected centralizer of \(C\). The adjoint action
of \(Q\) fixes exactly the diagonal Cartan and acts on the \(n\)
eigenlines of \(T_1\) by an \(n\)-cycle. Such a cycle is a Coxeter
element of \(S_n\). This verifies apposition by (CC-AP.1). Equivalently,
\(C=\sum_{j=0}^{n-2}E_{j+1,j}+E_{0,n-1}\) is precisely the sum of the
simple-root vectors and the negative-highest-root vector for the reverse
standard positive system. All coefficients are nonzero, so Kostant's
cyclic-element construction applies. Projectively
\(Q=\exp(2\pi i\rho^\vee/n)\) for this positive system.

Conversely, take a pair \((T,T')\) and a principal \(P\) as in
(CC-AP.1). Conjugate \(T'\) to \(T_0\). The action of \(P\) on its
eigenlines is a Coxeter element, hence an \(n\)-cycle. After a
permutation of the eigenlines, a matrix representative of \(P\) has the
form

\[
 A e_j=c_j e_{j+1},\qquad c_j\ne0.
\]

Choose \(c\) with \(c^n=\prod_jc_j\) and choose nonzero diagonal entries
\(d_j\) recursively by \(d_{j+1}=c\,d_j/c_j\). The last equation agrees
with the first by the displayed product relation. Conjugation by
\(\operatorname{diag}(d_j)\) changes \(A\) to \(cC\), which represents
\(C\) in \(\operatorname{PGL}(V)\). Therefore the ordered pair becomes
\((T_1,T_0)\), since the fixed Lie algebra of \(\operatorname{Ad}C\) is
the Cartan centralizing \(C\).

Finally \(F^2=nR\), where \(Re_j=e_{-j}\). The permutation \(R\)
normalizes \(T_0\). It follows that \(F^{-1}T_0F=T_1\), while
\(F^{-1}T_1F=T_0\). Conjugating by \(F^{-1}\) therefore exchanges the
two tori and gives the asserted ordered normal form. \(\square\)

\subsection{CC-AP02 --- The exact opposition
test}\label{cc-ap02-the-exact-opposition-test}

Write permutations \(\pi,\sigma\in S_I\) as ordered lists. Every Borel
containing \(T_0\) stabilizes exactly one of the flags

\[
 {\cal E}^{\pi}_k=\langle e_{\pi(0)},\ldots,e_{\pi(k-1)}\rangle ,
\]

and every Borel containing \(T_1\) is obtained from

\[
 {\cal V}^{\sigma}_r=\langle v_{\sigma(0)},\ldots,v_{\sigma(r-1)}\rangle .
\]

Here \(0\leq k,r\leq n\). This is the full-flag description of Borels in
\(\operatorname{PGL}(V)\): the invariant subspaces for a maximal
diagonal torus are sums of its distinct one-dimensional weight spaces.

\textbf{Proposition.} The Borels
\(B_\pi=\operatorname{Stab}({\cal E}^\pi)\) and
\(B'_\sigma=\operatorname{Stab}({\cal V}^\sigma)\) are opposite if and
only if, for every \(1\leq k\leq n-1\),

\[
 \det\bigl(\zeta^{ja}\bigr)_{
 j\in\{\pi(k),\ldots,\pi(n-1)\},
 \ a\in\{\sigma(0),\ldots,\sigma(n-k-1)\}}\ne0.
 \tag{CC-AP.2}
\]

The order chosen on either displayed set affects only the sign.

\textbf{Proof.} First recall directly why opposite flags are
characterized by

\[
 {\cal E}_k\cap{\cal V}_{n-k}=0
 \quad(1\leq k\leq n-1).
 \tag{CC-AP.3}
\]

Two full flags admit a simultaneous adapted basis. One construction uses
the dimensions \(r_{ij}=\dim({\cal E}_i\cap{\cal V}_j)\): the
nonnegative integers

\[
 r_{ij}-r_{i-1,j}-r_{i,j-1}+r_{i-1,j-1}
\]

are the dimensions of the corresponding successive quotients of the
intersections. Their row and column sums are one, so the nonzero
positions form a permutation matrix. Choosing a vector at each of these
positions gives the simultaneous basis, by induction on the
intersections.

Order that basis according to the first flag, and let \(\eta(i)\) be the
position of its \(i\)-th vector in the second flag. If an increasing
pair \(i<j\), \(\eta(i)<\eta(j)\), occurs, the nontrivial root subgroup
sending the \(j\)-th basis vector to itself plus a multiple of the
\(i\)-th preserves both flags. Hence the common stabilizer contains a
nontrivial unipotent subgroup and is not a torus. If there is no
increasing pair, \(\eta\) is the reverse ordering. The common stabilizer
is then exactly the diagonal torus in this basis. The reverse ordering
is also equivalent to (CC-AP.3). This proves the equivalence with
opposition in the sense used in the letter.

For the flags in the statement, the quotient map
\(V\to V/{\cal E}^{\pi}_k\), restricted to \({\cal V}^{\sigma}_{n-k}\),
is represented by the square submatrix in (CC-AP.2). Its domain and
codomain both have dimension \(n-k\). Its kernel is the intersection in
(CC-AP.3). Thus (CC-AP.2) and (CC-AP.3) are equivalent. \(\square\)

\subsection{CC-AP03 --- ``Some'' always; ``all'' exactly in prime
dimension}\label{cc-ap03-some-always-all-exactly-in-prime-dimension}

\textbf{Theorem.} For \(G=\operatorname{PGL}_n(\mathbf C)\):

\begin{enumerate}
\def\labelenumi{\arabic{enumi}.}
\tightlist
\item
  Every apposition pair admits opposite containing Borels. More
  strongly, after the normal form of CC-AP01, the Fourier flag in the
  ordering \(0,1,\ldots,n-1\) is opposite to every coordinate flag, and
  the coordinate flag in that ordering is opposite to every Fourier
  flag.
\item
  Every quadruple with tori in apposition has opposite Borels if and
  only if \(n\) is prime.
\item
  If \(n\) is composite, the failure can be witnessed by a rank-one
  nilpotent matrix whose associated one-parameter unipotent group lies
  in both Borels.
\end{enumerate}

\textbf{Proof.} For the first assertion, fix the Fourier ordering
\(\sigma(a)=a\). For any coordinate ordering and any \(k\), the minor in
(CC-AP.2) has columns with exponents \(0,\ldots,n-k-1\) and distinct row
parameters \(\zeta^j\). It is a Vandermonde determinant and is nonzero.
Conversely, for the standard coordinate ordering its rows have
consecutive exponents \(k,\ldots,n-1\). Factoring \(\zeta^{ka}\) from
the column indexed by \(a\) again leaves a Vandermonde determinant in
the distinct column parameters \(\zeta^a\). This proves both
strengthened assertions.

If \(n=p\) is prime, every square Fourier minor is nonzero by the
Chebotarev minor theorem, in the exact primary formulation \textbf{Tao
(2005), Lemma 1.3, p.~123, proof pp.~124--125}; equivalently
arXiv:math/0308286v6, Lemma 1.3, pp.~3--4. Applying that theorem to
(CC-AP.2) proves opposition for every pair of containing Borels. CC-AP01
then covers every apposition pair. The cited minor theorem is a known
input; this application does not claim a new proof of it.

For the converse write \(n=ab\), where \(a,b>1\). The submatrix with
rows \(\{0,a\}\) and columns \(\{0,b\}\) is

\[
 \begin{pmatrix}1&1\\1&\zeta^{ab}\end{pmatrix}
 =\begin{pmatrix}1&1\\1&1\end{pmatrix}.
\]

Choose the coordinate ordering so that its last two indices are \(0,a\),
and choose the Fourier ordering so that its first two indices are
\(0,b\). Equation (CC-AP.2) fails at \(k=n-2\).

For a direct group-theoretic witness put

\[
 x=v_0-v_b,\qquad
 \lambda=e_0^*-e_a^*,\qquad
 N=x\otimes\lambda.
 \tag{CC-AP.4}
\]

The vector \(x\ne0\) belongs both to \({\cal E}_{n-2}^{\pi}\) and to
\({\cal V}_2^\sigma\), because its coordinates at \(0,a\) vanish. The
functional \(\lambda\) vanishes on both these subspaces: on the
coordinate subspace this is immediate, and on \(v_0,v_b\) it follows
from \(\zeta^{ab}=1\). Hence \(N\ne0\), \(N^2=0\), and \(I+tN\) is a
nontrivial one-parameter unipotent subgroup.

Every earlier step of the coordinate flag is contained in
\(\ker\lambda\), and every later step contains \(x\). The same holds for
the Fourier flag, with its dividing step at dimension two. Therefore
\(I+tN\) preserves both complete flags. Their common stabilizer contains
this nontrivial unipotent subgroup and cannot be a torus.

In the smallest example \(n=4\), take \(\zeta=i\), coordinate ordering
\((1,3,0,2)\), and Fourier ordering \((0,2,1,3)\). Rescaling (CC-AP.4)
gives the particularly simple witness

\[
 N=(e_1+e_3)(e_0^*-e_2^*).
\]

This is a counterexample inside the original complex simple adjoint
group of type \(A_3\), with apposition established before the flag
calculation. \(\square\)

\subsection{CC-AP04 --- Explicit classification and inner
rigidity}\label{cc-ap04-explicit-classification-and-inner-rigidity}

Let

\[
 K_n=N_G(T_0)\cap N_G(T_1),\qquad
 \Omega_n=S_I\times S_I.
\]

For \(a\in I^\times\) and \(b,d\in I\), define the projective class of
the matrix

\[
 g_{a,b,d}e_j=\zeta^{ad j}e_{aj+b}.
 \tag{CC-AP.5}
\]

\textbf{Theorem.} These are exactly the elements of \(K_n\). On the
coordinate and Fourier eigenline labels their action is

\[
 j\longmapsto aj+b,\qquad
 c\longmapsto a^{-1}c+d.
 \tag{CC-AP.6}
\]

Consequently

\[
 K_n\simeq(I\times I)\rtimes I^\times,\qquad
 a\cdot(b,d)=(ab,a^{-1}d).
\]

The inner-conjugacy classes of ordered quadruples are exactly

\[
 K_n\backslash\Omega_n,\qquad
 (a,b,d)\cdot(\pi,\sigma)
 =(a\pi+b,\ a^{-1}\sigma+d).
 \tag{CC-AP.7}
\]

This action is free. Thus every quadruple has trivial inner stabilizer,
and the number of classes is

\[
 \frac{(n!)^2}{n^2\varphi(n)}.
 \tag{CC-AP.8}
\]

\textbf{Proof.} An element normalizing \(T_0\) is monomial in the
coordinate basis, and an element normalizing \(T_1\) is monomial in the
Fourier basis. Thus a representative \(g\) of an element of \(K_n\)
satisfies, for permutations \(\pi_0,\sigma_0\) of \(I\) and nonzero
scalars \(u_j,w_c\),

\[
 ge_j=u_j e_{\pi_0(j)},\qquad
 gv_c=w_c v_{\sigma_0(c)}.
\]

Comparison of the coordinate indexed by \(\pi_0(j)\) yields

\[
 u_j\zeta^{jc}=w_c\zeta^{\pi_0(j)\sigma_0(c)}.
\]

Divide this equation by its instance at \(c=0\), and then use its
instance at \(j=0\) to remove the remaining scalar. The result is the
congruence

\[
 jc=(\pi_0(j)-\pi_0(0))(\sigma_0(c)-\sigma_0(0))
 \pmod n.
 \tag{CC-AP.9}
\]

At \(j=c=1\), the two factors on the right are units and are inverses.
Put \(a=\pi_0(1)-\pi_0(0)\), \(b=\pi_0(0)\), and \(d=\sigma_0(0)\).
Taking \(c=1\) and then \(j=1\) in (CC-AP.9) gives

\[
 \pi_0(j)=aj+b,\qquad \sigma_0(c)=a^{-1}c+d.
\]

The equation at \(c=0\) then gives \(u_j=u\,\zeta^{ad j}\), with a
scalar \(u\) independent of \(j\). This proves that the projective
matrix is (CC-AP.5).

Conversely (CC-AP.5) gives

\[
 g_{a,b,d}v_c
 =\zeta^{-b(a^{-1}c+d)}v_{a^{-1}c+d},
\]

so all the displayed matrices normalize both tori. Their action on the
two sets of eigenlines is faithful. Composition of the affine
permutations gives

\[
 (a,b,d)(a',b',d')
 =(aa',\,ab'+b,\,a^{-1}d'+d),
\]

which proves the claimed semidirect-product structure and its order
\(n^2\varphi(n)\).

By CC-AP01 every quadruple has a representative with torus pair
\((T_0,T_1)\), and precisely \(K_n\) remains available for conjugating
such representatives. This proves the orbit classification (CC-AP.7). If
a \(K_n\)-element preserves a full flag made from eigenlines, its
permutation of those eigenlines preserves every initial segment and is
therefore the identity. Preserving both flags makes both affine
permutations in (CC-AP.6) the identity. It follows that \(a=1,b=d=0\),
and the projective transformation is the identity. The action is free,
so counting proves (CC-AP.8). This also proves inner rigidity for every
original quadruple, not merely for an opposite one. \(\square\)

For example, there is one class in type \(A_1\), two in type \(A_2\),
and eighteen in type \(A_3\). Formula (CC-AP.7) is an explicit finite
classification in type A; it is not a classification in the other simple
types.

\subsection{CC-AP05 --- Outer action, distinguished classes, and full
automorphisms}\label{cc-ap05-outer-action-distinguished-classes-and-full-automorphisms}

Let \(\theta(g)=(g^{-1})^{\mathsf T}\). For \(n\geq3\) it represents the
nontrivial element of the outer automorphism group of \(G\).

\textbf{Theorem.} On the quotient (CC-AP.7), \(\theta\) acts by

\[
 (\pi,\sigma)\longmapsto
 \bigl(\pi^{\mathrm{rev}},-\sigma^{\mathrm{rev}}\bigr),
 \quad
 \pi^{\mathrm{rev}}(i)=\pi(n-1-i).
 \tag{CC-AP.10}
\]

A class is fixed if and only if some \(a\in I^\times\), \(b,d\in I\)
satisfy, for every \(0\leq i<n\),

\[
 a\pi(i)+b=\pi(n-1-i),\qquad
 a^{-1}\sigma(i)+d=-\sigma(n-1-i).
 \tag{CC-AP.11}
\]

For \(n\geq3\), the full automorphism group of an individual quadruple
is trivial if this condition fails and is \(C_2\) if it holds. If \(n\)
is odd, the condition always fails. If \(n\geq4\) is even, it holds for
explicit examples.

\textbf{Proof.} We first justify the outer-automorphism statement in the
precise algebraic category being used. Any algebraic automorphism of
\(G\), after an inner adjustment, preserves a chosen Borel and its
maximal torus. Its action on the roots then preserves the positive
system and induces an automorphism of the chain \(A_{n-1}\). There are
only identity and reversal, with only identity when \(n=2\). The
composition of \(\theta\) with conjugation by the reverse permutation
matrix preserves this Borel and torus and induces reversal.

If the induced root action is the identity, the automorphism is the
identity on the torus, since its character lattice for the adjoint group
is generated by the roots. On the Lie algebra root spaces it has the
form \(E_{ij}\mapsto c_{ij}E_{ij}\). Preservation of brackets gives
\(c_{ij}c_{ji}=1\), and, for distinct \(i,j,k\),
\(c_{ik}=c_{ij}c_{jk}\). Fixing one index shows \(c_{ij}=s_i/s_j\) for
suitable \(s_i\ne0\); the two-index case follows directly from
\(c_{ij}c_{ji}=1\). Thus conjugation by \(\operatorname{diag}(s_i)\)
gives the same differential. Automorphisms of the connected complex
group are determined by their differential, so the automorphism is
inner. Finally, for \(n\geq3\), reversal of the simple roots cannot
arise from an inner automorphism preserving this Borel and torus: its
conjugating element would belong to \(B\cap N_G(T)=T\), whose action on
the roots is the identity. This proves \(\operatorname{Out}(G)=C_2\)
with representative \(\theta\) when \(n\geq3\), and
\(\operatorname{Out}(G)=1\) when \(n=2\).

Both \(T_0\) and \(T_1\) are preserved by \(\theta\): inverse transpose
preserves diagonal matrices and sends the algebra of circulant matrices
to itself. On a full flag \(\theta\) sends the stabilizer to the
stabilizer of the flag of annihilators, in reverse order, using the
standard bilinear form to identify \(V^*\) with \(V\). The dual
coordinate basis is the same coordinate basis. For the Fourier basis one
has

\[
 v_a^{\mathsf T}v_b=n\,\mathbf 1_{a+b=0},
\]

so its dual basis has labels \(-a\). This proves (CC-AP.10). Combining
it with (CC-AP.7) proves the equivalence with (CC-AP.11).

By CC-AP04 the inner stabilizer is trivial. Consequently the
automorphism group of a quadruple injects into
\(\operatorname{Out}(G)\), and its image is exactly the stabilizer of
its inner-conjugacy class. In particular, if (CC-AP.11) holds, an inner
adjustment of \(\theta\) fixes the quadruple. Its square is inner and
fixes the quadruple, hence is the identity. This proves the asserted
automorphism-group classification.

Now suppose \(n\geq3\) is odd and (CC-AP.11) holds. The affine
permutation

\[
 R(j)=aj+b
\]

reverses an ordering of \(n\) elements. It is an involution with exactly
one fixed point. Therefore \(a^2=1\pmod n\), and the equation
\((a-1)j=-b\) has exactly one solution. The latter fact implies
\(\gcd(a-1,n)=1\). Multiplying the identity \((a-1)(a+1)=0\) by the
inverse of \(a-1\) gives \(a=-1\pmod n\).

The second equation of (CC-AP.11) says that

\[
 S(c)=-a^{-1}c-d
\]

also reverses an ordering of \(n\) elements. With \(a=-1\), this is the
translation \(S(c)=c-d\). Since it is an involution and \(n\) is odd,
\(2d=0\) implies \(d=0\). Then \(S\) is the identity, which cannot
reverse an ordering of at least three distinct elements. This
contradiction proves the assertion for every odd \(n\geq3\), including
composite \(n\).

For an even example write \(n=2m\), with \(m\geq2\), and take

\[
 \pi=(0,1,\ldots,m-1,\ 2m-1,2m-2,\ldots,m),\qquad
 \sigma=(0,1,\ldots,2m-1).
\]

Equation (CC-AP.11) holds with \(a=1,b=m,d=1\). The corresponding class
is outer-invariant, and its automorphism group is \(C_2\), as already
proved. \(\square\)

\textbf{Direct answer to the original distinguished-class question.} In
every type \(A_{n-1}\) with odd \(n\geq3\), there is no inner-conjugacy
class invariant under all outer automorphisms. Thus a distinguished
class satisfying Grothendieck's stated necessary invariance condition
cannot exist there. In type \(A_2\), the two classes of CC-AP04 are
exchanged. Every quadruple in these odd types has trivial full algebraic
automorphism group. In type \(A_1\) the sole class is invariant and its
automorphism group is trivial. For even \(n\geq4\) the theorem supplies
invariant classes and the exact automorphism test, but does not assert a
unique canonical choice among them.

For a completely explicit description in type \(A_2\), the two classes
are distinguished by
\(\operatorname{sgn}(\pi)\operatorname{sgn}(\sigma)\). Translations of
\(\mathbf Z/3\) are even permutations, while multiplication by \(-1\) is
odd on both sets of labels, so this product is \(K_3\)-invariant. The
count of two classes proves it is complete. Representatives are
\((\pi,\sigma)=((0,1,2),(0,1,2))\) and \(((0,1,2),(0,2,1))\). Reversal
changes the row sign, whereas negation followed by reversal leaves the
Fourier sign unchanged; thus \(\theta\) exchanges the two classes
directly. Both representatives have opposite Borels by CC-AP03.

\subsection{CC-AP06 --- The exact split-field
obstruction}\label{cc-ap06-the-exact-split-field-obstruction}

Let \(k\subset\mathbf C\) be a field, and retain \(n\geq2\). Here a
\(k\)-defined quadruple means that each of \(T,B,T',B'\) is a subgroup
scheme of the split group \(\operatorname{PGL}_{n,k}\) defined over
\(k\), and that after extending to \(\mathbf C\) it has the apposition
property above.

\textbf{Theorem.} Such a quadruple exists if and only if
\(\zeta_n\in k\) for a primitive \(n\)-th root of unity \(\zeta_n\).

\textbf{Proof.} A \(k\)-defined Borel in split \(\operatorname{PGL}_n\)
stabilizes a full flag over \(k\). A \(k\)-defined maximal torus inside
it is split. To see the latter assertion concretely, lift the torus and
Borel to \(\operatorname{GL}_n\) and choose a \(k\)-basis adapted to the
flag. The lifted torus has dimension \(n\) and is a torus after base
change. Projection of the upper triangular group onto its diagonal
quotient restricts to an injective map on this torus: its kernel would
be both a torus subgroup and a unipotent subgroup, and hence trivial in
characteristic zero. By equal dimension this map is an isomorphism onto
\((\mathbf G_m)^n\). Thus the torus is split, and its one-dimensional
weight spaces admit a basis over \(k\).

Apply this argument to both \(T\subset B\) and \(T'\subset B'\). Choose
a \(k\)-basis diagonalizing the first torus. There is an
\(A\in\operatorname{GL}_n(k)\) such that the second torus is
\(AT_0A^{-1}\). CC-AP01 gives a complex transformation preserving the
first torus and carrying the second to \(FT_0F^{-1}\). Transformations
preserving \(T_0\) are monomial, so this says exactly that \(A\) is
obtained from \(F\) by row and column permutations and nonzero row and
column scalings.

All entries of \(A\) are therefore nonzero. Its cross-ratios

\[
 \frac{A_{ic}A_{jd}}{A_{id}A_{jc}}\qquad(i\ne j,\ c\ne d)
 \tag{CC-AP.12}
\]

belong to \(k\). Row and column scalings leave these ratios unchanged,
and row and column permutations only relabel them. For the Fourier
matrix, taking rows \(0,1\) and columns \(0,1\) gives the ratio
\(\zeta\). Hence \(\zeta\in k\).

Conversely, if a primitive \(\zeta\) belongs to \(k\), the Fourier
matrix, both tori, and the coordinate and Fourier flags are all defined
over \(k\). They give the required quadruple by CC-AP01. \(\square\)

For \(n=3\) and \(k=\mathbf Q\), this is an independently proved
arithmetic nonsplitness obstruction: a full rational apposition
quadruple does not exist in split \(\operatorname{PGL}_3\), even though
the diagonal torus and the connected centralizer of the rational
cyclic-shift matrix give a rationally defined pair of tori. The second
torus in that pair is not split over \(\mathbf Q\).

The supplied SA20 intake states an arithmetic nonsplitness result
labelled D.016 for ``RigKos in \(A_2/\mathbf Q\)'' but explicitly says
that Part D was read only at statement level. The full definition and
proof of that object were not supplied here. CC-AP06 therefore
establishes the concrete subgroup-scheme assertion above and does not
identify it with an unseen moduli functor, a coarse moduli point, or
every possible twisted form.

\subsection{CC-AP07 --- Classification with a specified generating
principal
element}\label{cc-ap07-classification-with-a-specified-generating-principal-element}

The original letter also asks for the classification when a generating
principal element of the finite group \(T\cap N_G(T')\) is included in
the data.

\textbf{Corollary.} In type \(A_{n-1}\), this group is cyclic of order
\(n\). The inner-conjugacy classes of the ordered data
\((T,B,T',B',P)\), where \(P\) generates that group, are

\[
 K_n\backslash(\Omega_n\times I^\times),
 \quad
 (a,b,d)\cdot(\pi,\sigma,u)
 =(a\pi+b,\ a^{-1}\sigma+d,\ a^{-1}u).
 \tag{CC-AP.13}
\]

Equivalently, after normalizing \(P=Q\), they are

\[
 (I\times I)\backslash\Omega_n,\qquad
 (b,d)\cdot(\pi,\sigma)=(\pi+b,\sigma+d).
 \tag{CC-AP.14}
\]

Both actions are free and their common number of orbits is
\((n!)^2/n^2\).

\textbf{Proof.} In the normal form, an element of \(T_0\cap N_G(T_1)\)
belongs to \(K_n\) and has trivial permutation of the coordinate
eigenlines. By (CC-AP.6), this forces \(a=1,b=0\), so the element is
\(g_{1,0,d}=Q^d\). Conversely every such power normalizes \(T_1\). Hence

\[
 T_0\cap N_G(T_1)=\langle Q\rangle\simeq C_n.
\]

Its generators are precisely \(Q^u\), \(u\in I^\times\). They are
principal: each is conjugate to \(Q\) by permuting its \(n\) distinct
eigenvalues. Direct conjugation, or (CC-AP.5), gives in the projective
group

\[
 g_{a,b,d}Q^u g_{a,b,d}^{-1}=Q^{a^{-1}u}.
\]

This proves (CC-AP.13). Its action is free because its action on the
first two coordinates is already free by CC-AP04. Therefore its orbit
count is

\[
 \frac{(n!)^2\varphi(n)}{n^2\varphi(n)}
 =\frac{(n!)^2}{n^2}.
\]

Given any generator exponent \(u\), choose \(a=u\) to change it to
\(1\). The remaining transformations fixing \(u=1\) have \(a=1\) and are
exactly the two independent translations in (CC-AP.14). This proves the
alternative classification, its freeness, and its count. \(\square\)

These two counts agree with the letter's \(|W|^2/(hz\varphi(h))\) and
\(|W|^2/(hz)\): in type \(A_{n-1}\), \(W=S_n\), \(h=n\), and \(z=n\) is
the order of the center of the simply connected cover
\(\operatorname{SL}_n\).

\subsection{Primary bibliography and inspection
record}\label{primary-bibliography-and-inspection-record}

Grothendieck, A. (1969, October 22). \emph{Letter to B. Kostant}.
Montpellier archives, cote 123, physical PDF page 22, footer 21.
\url{https://grothendieck.umontpellier.fr/123.pdf}. The original page was
visually read, including the distinct classification, opposition,
distinguished-class, and automorphism questions.

Kostant, B. (1959). The principal three-dimensional subgroup and the
Betti numbers of a complex simple Lie group. \emph{American Journal of
Mathematics, 81}(4), 973--1032. \url{https://doi.org/10.2307/2372999}. Theorem
8.6 and the following Corollary 8.6, p.~1026, were read in the OCR of
the author's reprint in \emph{Collected Papers}, Vol. I, pp.~130--189
(2009), \url{https://doi.org/10.1007/b94535\_11}. The source scan itself was
not obtained; original §7.3, p.~1018, is a located but uninspected
reference. OCR inspected at
\url{https://dokumen.pub/collected-papers-volume-i-1955-1966-2009nbsped-9780387095820-0387095829.html}.

Kostant, B. (2010). Experimental evidence for the occurrence of \(E_8\)
in nature and the radii of the Gosset circles. \emph{Selecta
Mathematica, 16}, 419--438. \url{https://doi.org/10.1007/s00029-010-0023-1}.
Inspected author preprint: \url{https://arxiv.org/pdf/1003.0046v1} (1 March
2010). Exact locators in this version: Theorem 0.1 and apposition
sentence p.~2; Proposition 1.1 p.~5; Theorems 1.2--1.3 p.~6; Theorem 1.7
and principal-element discussion p.~7; Proposition 1.8 and Theorem 1.9,
including its complete short proof, p.~8.

Tao, T. (2005). An uncertainty principle for cyclic groups of prime
order. \emph{Mathematical Research Letters, 12}(1), 121--127.
\url{https://doi.org/10.4310/MRL.2005.v12.n1.a11}. Lemma 1.3 p.~123; proof
pp.~124--125; preliminary Lemma 1.2 p.~123. The journal PDF was directly
downloaded and the statements and complete indicated proofs read:
\url{https://www.intlpress.com/site/pub/files/\_fulltext/journals/mrl/2005/0012/0001/MRL-2005-0012-0001-a011.pdf}.
Also inspected: \url{https://arxiv.org/pdf/math/0308286v6} (22 July 2004),
Lemmas 1.2--1.3 pp.~3--4. The prime Fourier-minor statement is credited
there to Chebotarev; no original Chebotarev theorem number is invented.

\textbf{Proof provenance.} CC-AP01--07 give complete proofs of the
stated type-A deductions, with Kostant's apposition/Coxeter
characterization and the prime Fourier-minor theorem explicitly retained
as known inputs. The finite normalizer and outer-action calculations,
the composite unipotent witness, the field-of-definition obstruction,
and the classification with a specified generator are proved here. This
is neither a proof of RH nor a classification of apposition quadruples
in every simple type.

% END consequences_criteria_apposition.tex

% BEGIN root_effectivity.tex
\section{Spectral effectivity and the full Tate conjecture: the precise numerical-motive equivalence}

The SA06 intake, row D.4f and its successor-work instruction at
lines 501--503, asks whether effectivity detected by Frobenius is
equivalent to the Tate conjecture for all motives.  Its earlier proof
file is absent, but this question has the following precise answer in
the category of rational numerical motives.  This is a deduction from
classical results, not a new solution of the Tate conjecture.

Fix $k=\mathbb F_q$.  Let $\mathcal N(k)$ be the pseudo-abelian
category of pure numerical motives with rational coefficients and all
integer Tate twists.  Call an object \emph{effective} if it is a
direct summand of a finite direct sum of motives $h(X)$ with $X/k$
smooth projective, without a positive Tate twist.  The Frobenius
endomorphism $\pi$ is central.  On a simple object $S$, the algebra
$\mathbb Q[\pi_S]$ is a number field, so $\pi_S$ defines an algebraic
number up to conjugacy.  Call $M$ \emph{spectrally integral} if all
these numbers for its simple constituents are algebraic integers.
This definition does not presume that an $\ell$-adic realization
factors through numerical equivalence.

The structural inputs are Jannsen's semisimplicity theorem and
Frobenius centrality, in the exact formulations of Geisser (1998),
Theorem 2.1, p.~111, and Propositions 2.3--2.4, p.~112.
The decisive known result is Geisser, Theorem 1.1, p.~109,
equivalently Theorem 2.7, p.~114: the full Tate conjecture over $k$
is equivalent to the assertion that a simple numerical motive whose
Frobenius is $1$ is the unit.  We use this theorem as a cited input.
Here \emph{full Tate} means equality of the pole order at $s=r$
and the dimension of codimension-$r$ numerical cycle classes, for
all smooth projective $X/k$ and all $r$; it is not merely the
surjectivity assertion for one variety and one degree.

\subsection{BR-EFF01. The equivalence and its weight-zero test}
The following assertions are equivalent:
\begin{enumerate}
\item The full Tate conjecture holds over $k$.
\item Every spectrally integral object of $\mathcal N(k)$ is effective.
\item Every object $M$ of $\mathcal N(k)$ with $\pi_M=1_M$ is effective.
\item Every simple object $S$ with $\pi_S=1_S$ is effective.
\end{enumerate}
For every effective object, spectral integrality holds without
any Tate-conjecture assumption.  Thus item 2 can equivalently be
written ``effective if and only if spectrally integral.''

\paragraph{Complete deduction.}
For smooth projective $X$, its cohomological Frobenius characteristic
polynomials are monic integral polynomials, and their product annihilates Frobenius
in the category of homological motives.  They therefore annihilate
it numerically.  The minimal polynomial on a numerical summand
divides such a polynomial, so all its roots are algebraic integers.
This proves the unconditional necessity.

Assume item 4, and let $S$ be simple with $\pi_S=1$.  Since $S$
is effective, it is a summand of some finite sum of $h(X)$.
The algebraic K\"unneth projectors over $\mathbb F_q$ give
$h(X)=\bigoplus_{i\geq0}h^i(X)$; the precise cited input is
Katz--Messing (1974), Theorem 2(1), pp.~76--77.
For $i>0$, a polynomial annihilating $\pi_{h^i(X)}$ has all roots
of absolute value $q^{i/2}>1$.  In particular $\pi_{h^i(X)}-1$
is invertible: a B\'ezout identity between its annihilating
polynomial and $T-1$ constructs the inverse.  Centrality of
Frobenius now forces every map $S\to h^i(X)$ and $h^i(X)\to S$
to be zero.  Hence $S$ is a summand of a finite sum of $h^0(X)$.

These degree-zero objects are Artin motives: $h^0(X)$ is the
motive of the finite \'etale scheme of geometric connected
components.  For a geometrically connected component the
degree-zero projector is obtained from a degree-one rational
zero-cycle, obtained by dividing any closed point by its degree.
For a component with constant field $\mathbb F_{q^a}$ this gives
$h(\operatorname{Spec}\mathbb F_{q^a})$.  In its permutation
representation Frobenius is a cyclic permutation, and the
$1$-eigenspace is the constant line.  The rational projector is
$a^{-1}\sum_{j=0}^{a-1}\pi^j$.  Thus the part with Frobenius $1$
in any finite sum of such Artin objects is a sum of units.  The
simple object $S$ is consequently the unit.  Geisser's cited
theorem proves item 1.

Suppose now that item 1 holds.  The same cited theorem says that
simple numerical motives are determined by their Frobenius number.
Let $S$ be simple and spectrally integral, with Frobenius number
$\alpha$.  Its weight is an integer $m$, and all conjugates have
absolute value $q^{m/2}$.  Because $\alpha$ is a nonzero algebraic
integer, the norm shows $m\geq0$.

If $m=0$, Kronecker's theorem says that $\alpha$ is a root of
unity.  The rational Artin representation for its cyclotomic
polynomial is a summand of the regular representation of a finite
cyclic quotient of $\operatorname{Gal}(\overline k/k)$, hence is
an effective Artin motive.  Its simple constituent with Frobenius
$\alpha$ is isomorphic to $S$ by Geisser's characterization.

If $m>0$, choose an algebraic root $\beta$ with $\beta^m=\alpha$.
It is integral and each conjugate has absolute value $q^{1/2}$.
Honda--Tate supplies an abelian variety $B/k$ having $\beta$ among
its $H^1$ Frobenius eigenvalues.  The tensor power
$h^1(B)^{\otimes m}$ is an effective summand of $h(B^m)$, and its
realization has $\alpha$ among its eigenvalues.  Under full Tate,
numerical and homological equivalence agree and realization is
faithful.  Its decomposition into simple numerical summands
therefore contains a summand $S_\alpha$ with Frobenius number
$\alpha$.  Geisser's characterization gives $S\simeq S_\alpha$.
Thus $S$ is effective.  A finite direct sum of effective motives
is effective, proving item 2.  The implications $2\Rightarrow3
\Rightarrow4$ are immediate.  This proves all assertions.
\(\square\)

This proof never assumes a rational neutral fibre functor.
The quaternion obstruction in AQ-RATIONAL-FIBRE is compatible
with every step.  It also does not identify numerical and Chow
effectivity: a lift of a numerical projector to a Chow projector
requires additional information.

\subsection{BR-EFF02. The effective cone in the now established abelian sector}
Let $\mathcal A_q$ be the rational numerical category generated by
abelian, Artin, and Tate motives, with the same conventions as
CC-AG01.  Using the cited all-abelian Tate result in CC-TATE01--04,
an object of $\mathcal A_q$ is effective if and only if it is
spectrally integral.  In fact it then has an effective presentation
using only abelian varieties and finite \'etale schemes.  Over
$E=\overline{\mathbb Q}_\ell$, the equivalence of CC-AG01 identifies
this full subcategory with the finite $W(q)$-graded vector spaces
whose nonzero components are indexed by
\[
 W^+(q)=\{\alpha\in W(q):\alpha\text{ is an algebraic integer}\}.
 \tag{EFF.1}
\]
This is a tensor monoid of labels, not a group of labels.

\paragraph{Complete deduction.}
Necessity is the unconditional first paragraph of BR-EFF01.
For sufficiency, apply its construction to every simple constituent.
The required comparisons occur entirely inside $\mathcal A_q$.
Their validity follows from the all-abelian Tate theorem and its
full-faithfulness consequence AQ-ABELIAN-RECOGNITION: for simple
objects in this sector, an intersecting Frobenius spectrum gives a
nonzero morphism after scalar extension, hence a nonzero rational
morphism, and such a morphism is an isomorphism by semisimplicity.
In weight zero use the finite \'etale construction; in positive
weight use $h^1(B)^{\otimes m}$ as above.  These constructions
stay effective before any coefficient extension.

After extension to $E$, all rank-one $W(q)$-labels split by
CC-AG01.  Each label in $W^+(q)$ is effective by the same
construction, and each effective label is integral by necessity.
Finite sums and summands give (EFF.1).  Multiplication preserves
integrality, so $W^+(q)$ is a monoid.  For example $q\in W^+(q)$
whereas $q^{-1}\notin W^+(q)$, so it is not a group.
\(\square\)

\subsection{BR-EFF03. The slope test and the essential Artin correction}
For a Weil $q$-number $\alpha$ of weight $m$, spectral integrality
is equivalent to
\[
 \frac{v(\alpha)}{v(q)}\geq0
 \quad\text{for every finite place }v\mid p
 \text{ of a number field containing }\alpha.
 \tag{EFF.2}
\]
Consequently (EFF.2) is an exact effectivity test in the abelian
sector of BR-EFF02.  For all numerical motives the corresponding
universal assertion is equivalent to full Tate by BR-EFF01.

\paragraph{Complete proof.}
The definition of a Weil number gives $q^N\alpha$ integral for
some $N$.  At places away from $p$, this implies $v(\alpha)\geq0$.
Choose a normal field containing $\alpha$ and its conjugates.
For a complex-conjugation automorphism $c$ one has
$\alpha c(\alpha)=q^m$.  The same nonnegativity holds for
$c(\alpha)$ away from $p$, and their valuations sum to zero
there.  Both valuations are therefore zero.  Hence the only
possible denominators of $\alpha$ occur over $p$, and (EFF.2)
is precisely their absence.  This proves the equivalence.
\(\square\)

There is a small but necessary source correction in using an
effective realization theorem.  Milne--Ramachandran (2009),
Proposition 1.7 in the author final version of May 2, 2008,
p.~4, is printed for every effective semisimple Tate structure
with an abelian variety as its sole source.  Its unrestricted
weight-zero wording has the following counterexample.
The one-dimensional space with Frobenius $-1$ is effective,
semisimple, and of weight zero.  For any abelian variety $B/k$,
$H^0(B,\mathbb Q_\ell)$ has Frobenius $1$, whereas every
positive-degree eigenvalue has absolute value greater than one.
Thus $-1$ never occurs in $H^*(B,\mathbb Q_\ell)$.
The space does occur in
$H^0(\operatorname{Spec}\mathbb F_{q^2},\mathbb Q_\ell)$, as
the anti-invariant summand.  This proves the counterexample and
the required repair in that case.  The general weight-zero repair
is the finite cyclic Artin construction in BR-EFF01; the positive
weight construction there is complete.  The effective targets of
CC-XC01 already include these Artin summands.  No refutation of
the generalized Tate conjecture is inferred from this wording
correction.

\paragraph{Exact sources.}
\begin{itemize}
\item Geisser, T. (1998). Tate's conjecture, algebraic cycles and
rational $K$-theory in characteristic $p$. \emph{$K$-Theory,
13}(2), 109--122.
\url{https://www2.rikkyo.ac.jp/web/geisser/Tate_KtheoryJournal.pdf}.
Theorem 1.1, p.~109; Theorem 2.1 and Corollary 2.2, p.~111;
Propositions 2.3--2.4, p.~112; Theorem 2.7, p.~114.
\item Milne, J. S., \& Ramachandran, N. (2009). Motivic complexes
over finite fields and the ring of correspondences at the generic
point. \emph{Pure and Applied Mathematics Quarterly, 5}(4),
1219--1252. Operative author final version, May 2, 2008:
\url{https://www.jmilne.org/math/articles/2009a.pdf}.
Proposition 1.7, p.~4; Remark 1.12, p.~6; Lemma 2.1 and
Proposition 2.2, p.~8. The weight-zero correction above is explicit.
\item Milne, J. S. (2007, September 19). \emph{The Tate conjecture
over finite fields (AIM talk)}.
\url{https://www.jmilne.org/math/xnotes/TateAimv1.pdf}.
Section 1, pp.~1--4, particularly footnote 4, p.~4, on motivic
effectivity and generalized Tate. This contextual statement is
not substituted for the precise numerical-motive theorem proved here.
\end{itemize}

% END root_effectivity.tex

% BEGIN root_group_rank.tex
\section{The group quotient in cote 123: central stabilizers and an exact type-A rank formula}

The rank question in the letter, cote 123 physical page 23, footer 22,
mentions both the Lie-algebra quotient and $G\to T/W$ for the original
complex adjoint group.  The simply connected qualification occurs later,
for rational singularities.  These scopes matter: a regular element of
an adjoint group can have quotient differential of rank zero.  The
following complete deductions give a finite formula in every semisimple
type and a closed matrix formula in type A.

\subsection{BR-12305. The differential of a finite linear quotient}
Let a finite group $Z$ act linearly on a finite-dimensional complex
vector space $V$, let $\delta:V\to V/Z$ be the affine quotient, and
let $H=Z_x$ be the stabilizer of $x\in V$.  Identify $T_xV$ with $V$.
Then
\[
 \operatorname{im}(d\delta_x)^*=(V^*)^H.
 \tag{GR.1}
\]
Consequently, for any linear map $L:U\to V$, the rank of
$d\delta_x\circ L$ is the rank of the restrictions to $L(U)$ of
the $H$-invariant linear forms on $V$.

\paragraph{Complete proof.}
Every invariant polynomial has an $H$-invariant differential at $x$,
since $H$ fixes $x$.  This proves one inclusion in (GR.1).
Conversely, let $\lambda\in(V^*)^H$.  The maximal ideals of the
finitely many distinct points of $Zx$ are pairwise comaximal, as are
their squares.  The Chinese remainder theorem supplies a polynomial
$e$ with $e\equiv1\pmod{\mathfrak m_x^2}$ and
$e\equiv0\pmod{\mathfrak m_y^2}$ for every $y\in Zx$, $y\ne x$.
Put
\[
 f(v)=\frac1{|H|}\sum_{z\in Z}
       \lambda(z^{-1}v-x)e(z^{-1}v).
\]
Reindexing the sum shows that $f$ is $Z$-invariant.  At $v=x$,
the terms with $z\notin H$ have zero first jet.  A term with
$z\in H$ has differential $\lambda\circ z^{-1}=\lambda$;
its factor $\lambda(z^{-1}x-x)$ is zero, so differentiating $e$
adds nothing.  Therefore $df_x=\lambda$.  This proves equality.
The image of the dual of a composite is the restriction of this
image along $L$, so its dimension gives the final assertion.
\(\square\)

No smoothness of $V/Z$ at $\delta(x)$ was assumed.  Its Zariski
tangent space can have dimension larger than $\dim V$; (GR.1)
still computes the actual differential rank.

\subsection{BR-12306. A finite character formula for every central quotient}
Let $\widetilde G$ be a connected simply connected complex
semisimple group of rank $r$, let $Z$ be a subgroup of its center,
and put $G=\widetilde G/Z$.  For the fundamental representations
$\rho_i$ of $\widetilde G$, write
$\chi_i=\operatorname{Tr}\rho_i$ and let $\xi_i:Z\to\mathbb C^*$
be their central characters.  Steinberg (1965), Theorem 6.1(b),
p.~62, and the quotient discussion (6.9)--(6.10), pp.~63--64,
give the polynomial coordinates
\[
 \widetilde\chi=(\chi_1,\ldots,\chi_r):
 \widetilde G\longrightarrow\mathbb A^r
       =\widetilde T/W.
\]
For $g\in G$, choose a lift $\widetilde g$, put
$x_i=\chi_i(\widetilde g)$, and define
\[
 H_x=\{z\in Z:\xi_i(z)x_i=x_i\text{ for every }i\},
 \qquad J_x=\{i:\xi_i|_{H_x}=1\}.
 \tag{GR.2}
\]
If $\chi_G:G\to T/W$ is the group quotient, then
\[
 \operatorname{rank}d\chi_G|_g
 =\dim_{\mathbb C}\operatorname{span}
 \left\{Y\longmapsto
 \operatorname{Tr}\bigl(\rho_i(\widetilde g)d\rho_i(Y)\bigr)
 :i\in J_x\right\}\subseteq\mathfrak g^*.
 \tag{GR.3}
\]
This is an explicit finite formula: evaluate the fundamental
characters, determine the finite stabilizer (GR.2), and compute
the rank of the displayed $|J_x|$ linear forms on a Lie-algebra
basis.  It applies to arbitrary elements, including those with
nontrivial unipotent part.

\paragraph{Complete proof.}
The central isogeny $\widetilde G\to G$ is \'etale in characteristic
zero, and its differential is an isomorphism.  Its action on the
fundamental-character coordinates is diagonal with weights $\xi_i$,
because a central element acts as the scalar $\xi_i(z)$ in an
irreducible representation.  Taking invariants in either order gives
\[
 \mathbb C[G]^G
 =\bigl(\mathbb C[\widetilde G]^{\widetilde G}\bigr)^Z
 =\mathbb C[x_1,\ldots,x_r]^Z.
\]
Thus the lifted quotient map is the composite of
$\widetilde\chi$ with $\delta:\mathbb A^r\to\mathbb A^r/Z$.
The $H_x$-invariant linear forms are precisely the span of $dx_i$
with $i\in J_x$.  BR-12305 gives the rank as that of these
coordinate differentials composed with $d\widetilde\chi$.

Identify $T_{\widetilde g}\widetilde G$ with $\mathfrak g$ by
$Y\mapsto\widetilde gY$.  Differentiating
$\chi_i(\widetilde g\exp(tY))$ at zero gives exactly the trace
in (GR.3).  Replacing the lift by $z\widetilde g$ leaves $H_x$
and $J_x$ unchanged and multiplies each displayed form by the
nonzero scalar $\xi_i(z)$.  The rank is therefore independent
of the lift.  This completes the proof. \(\square\)

For a regular lift, Steinberg's Theorem 8.1, p.~70, makes all
$d\chi_i$ independent.  In that case (GR.3) simplifies to
$\operatorname{rank}d\chi_G|_g=|J_x|$.  In particular the
simply connected case has full rank $r$ at every regular element,
whereas the adjoint quotient can have a smaller rank.

\subsection{BR-12307. Closed formula for all type-A central quotients}
Let $c$ divide $n$, let $G=\operatorname{SL}_n/\mu_c$, and let
$A\in\operatorname{SL}_n(\mathbb C)$ lift $g\in G$.
Write its characteristic polynomial as
\[
 \det(tI-A)=\sum_{i=0}^{n}(-1)^i c_i(A)t^{n-i},
 \qquad c_0=c_n=1.
\]
Define
\[
 d=\gcd\bigl(c,\{i:1\leq i<n,\ c_i(A)\ne0\}\bigr),
 \tag{GR.4}
\]
with the convention that the gcd of the empty index set and $c$
is $c$.  If $m_{A^d}$ is the minimal polynomial of $A^d$, then
\[
 \boxed{\operatorname{rank}d\chi_G|_g
        =\deg m_{A^d}-1.}
 \tag{GR.5}
\]
The formula covers $c=1$ (the simply connected group) and $c=n$
(the adjoint group $\operatorname{PGL}_n$).

\paragraph{Complete proof.}
For $\operatorname{SL}_n$, the fundamental characters are
$c_i=\operatorname{Tr}(\bigwedge^i A)$, $1\leq i<n$.
The scalar $\zeta I\in\mu_c$ acts with central character
$\zeta^i$.  Hence (GR.2) gives $H_x=\mu_d$ and
\[
 J_x=\{d,2d,\ldots,n-d\}.
\]
All nonzero $c_i$ have index divisible by $d$, and $d\mid n$,
so $\det(tI-A)=f(t^d)$ for a monic polynomial $f$ of degree
$h=n/d$.  The Cayley--Hamilton identity gives $f(A^d)=0$.
In particular $\deg m_{A^d}\leq h$.

To compute the differentials without an eigenvalue-genericity
assumption, differentiate the identity
$\det(I+tA)=\sum_i c_i t^i$ in the tangent direction $AY$.
The determinant derivative and the formal geometric series give
\[
 dc_i|_A(AY)=\sum_{j=1}^{i}(-1)^{j-1}c_{i-j}(A)
                                  \operatorname{Tr}(A^jY).
 \tag{GR.6}
\]
For $i=md$, a nonzero coefficient on the right forces $d\mid j$.
Writing $B=A^d$, the forms $dc_{md}$, $1\leq m<h$, are
therefore related to
$Y\mapsto\operatorname{Tr}(B^jY)$, $1\leq j<h$, by a
triangular matrix whose diagonal entries are
$(-1)^{md-1}$, and are nonzero.

The annihilator of $\mathfrak{sl}_n$ in
$\mathfrak{gl}_n$ under the trace pairing is the scalar line.
The dimension of these forms restricted to $\mathfrak{sl}_n$
is consequently
\[
 \dim\operatorname{span}\{I,B,\ldots,B^{h-1}\}-1
 =\deg m_B-1.
\]
The last equality uses $\deg m_B\leq h$ and the defining
minimality of the minimal polynomial.  Equation (GR.3) now
proves (GR.5).  No Jordan block was discarded and no
semisimplicity of $A$ was used. \(\square\)

Equivalently, $\deg m_{A^d}$ is the sum, over distinct nonzero
eigenvalues $\eta$ of $A^d$, of the largest Jordan-block size
among eigenvalues $\lambda$ of $A$ with $\lambda^d=\eta$.
Indeed the polynomial $t^d$ has nonzero derivative at each
nonzero $\lambda$, so it preserves the sizes of the local
Jordan blocks, while it can merge their eigenvalue labels.
This gives a direct formula from the Jordan data of $A$.

\subsection{BR-12308. Regular adjoint elements with rank-zero quotient differential}
For every $n\geq2$, the group $\operatorname{PGL}_n(\mathbb C)$
has a regular semisimple element $g$ with
$\operatorname{rank}d\chi_G|_g=0$.
In the simply connected group $\operatorname{SL}_n$, the
corresponding regular lift has rank $n-1$.

\paragraph{Complete counterexample.}
Fix a primitive $n$th root of unity $\zeta$, set
$Q=\operatorname{diag}(1,\zeta,\ldots,\zeta^{n-1})$, and
take $g=[Q]$.  Its eigenvalues are distinct, so its connected
centralizer is a maximal torus of dimension $n-1$.  It is
therefore regular semisimple.  Choose a scalar $a\ne0$ with
$A=aQ\in\operatorname{SL}_n$.  Its characteristic polynomial
is $t^n-a^n$, so $c_i(A)=0$ for every $1\leq i<n$.
For $G=\operatorname{PGL}_n$ one has $c=n$, hence $d=n$ in
(GR.4), and $A^n=a^n I$ has minimal polynomial of degree one.
Equation (GR.5) gives rank zero.  For the simply connected
group one instead has $c=d=1$, and the distinct eigenvalues
give $\deg m_A=n$, hence rank $n-1$.
\(\square\)

This example concerns the finite central quotient of the
\emph{group} invariant space.  It does not contradict Kostant's
regular-point criterion for the Lie-algebra quotient or Steinberg's
simply connected group theorem.  The original rank questions
must therefore keep their two quotient maps separate.

\paragraph{Primary input.}
Steinberg, R. (1965). Regular elements of semi-simple algebraic
groups. \emph{Publications Math\'ematiques de l'IH\'ES, 25},
49--80. \url{https://doi.org/10.1007/BF02684397}.
Theorem 6.1(b), p.~62; (6.9)--(6.10), pp.~63--64;
Theorem 8.1, p.~70. The finite-quotient differential lemma and
the formulas (GR.3)--(GR.6) are completely proved above.

% END root_group_rank.tex

% BEGIN consequences_criteria_adjoint_rank.tex
\section{Cote 123: exact Lie-algebra quotient ranks and the subregular question}\label{cote-123-exact-lie-algebra-quotient-ranks-and-the-subregular-question}

\textbf{Scope.} Complex semisimple Lie algebras and complex simple adjoint algebraic groups. These are unconditional characteristic-zero results; neither RH nor Dirichlet GRH is used. The original question is Grothendieck's letter to Kostant of 22 October 1969, cote 123, physical PDF page 23, handwritten/printed footer 22, visually inspected. It asks about both the Lie-algebra and group quotient maps, proposes a strict rank drop below \(r-1\) for subregular elements in nonsimply-laced types, and suggests uniqueness of the subregular orbit might characterize ADE. The two proposed distinctions are false. The group quotient away from the unipotent locus must be treated separately, as in BR-12305--08.

\textbf{Literature and proof status.} Richardson (1987) is a primary announcement with proof outlines, as he expressly says on p.~228. Its stated results are used as known inputs, not represented as independently re-proved here. Its exceptional table has three demonstrable printing/omission errors and one explicitly unresolved entry. The errors are corrected below using its own theorems. The formerly unresolved entry is explicitly settled in Dinar (2020), Corollary 5.5; a separate short deduction from Johnson--Sommers' theorem is proved below. Johnson--Sommers was inspected in arXiv:1706.04820v2, 5 February 2018, and Dinar in arXiv:1905.01829v4, 7 April 2020; both journal publications are 2020. Literature cutoff: 20 November 2026. Actual inspection date: 21 November 2026 UTC. No novelty is claimed.

\subsection{CC-RK01 --- Reduction to nilpotent elements and precise classical locators}\label{cc-rk01-reduction-to-nilpotent-elements-and-precise-classical-locators}

Let \(\chi_{\mathfrak g}:\mathfrak g\to\mathfrak t/W\) be the adjoint quotient and write \(x=s+e\) for the Jordan decomposition. Set

\[
 \mathfrak l=\mathfrak g^s,
 \qquad \mathfrak l_{\mathrm{der}}=[\mathfrak l,\mathfrak l]
       =\bigoplus_j\mathfrak l_j,
 \qquad e=(e_j)_j.
\]

Richardson's \textbf{Proposition 1.2, p.~229}, gives

\[
 \operatorname{rank}(d\chi_{\mathfrak g})_x
   =\dim\mathfrak z(\mathfrak l)
       +\sum_j\operatorname{rank}(d\chi_{\mathfrak l_j})_{e_j}.
 \tag{CC-RK.1}
\]

His notation \(\mathfrak g(s)\) in that proposition is the fixed subspace of the group centralizer; for semisimple \(s\), the centralizer is a connected Levi subgroup and that subspace is \(\mathfrak z(\mathfrak l)\). The simple-factor sum is the product decomposition of the quotient. Thus the following nilpotent data, together with the exceptional table below, determine the rank at every element.

For a nilpotent matrix let \(\lambda=(\lambda_1\geq\cdots\geq\lambda_k)\) denote its Jordan partition.

\begin{longtable}{@{}p{.23\linewidth}p{.38\linewidth}p{.29\linewidth}@{}}
\toprule\noalign{}
\begin{minipage}[b]{\linewidth}\raggedright
Type
\end{minipage} & \begin{minipage}[b]{\linewidth}\raggedleft
Exact nilpotent rank
\end{minipage} & \begin{minipage}[b]{\linewidth}\raggedright
Primary locator
\end{minipage} \\
\midrule\noalign{}
\endhead
\bottomrule\noalign{}
\endlastfoot
\(A_r=\mathfrak{sl}_{r+1}\) & \(\lambda_1-1\) & Richardson, Theorem 4.2.1, p.~233 \\
\(B_r=\mathfrak{so}_{2r+1}\), \(C_r=\mathfrak{sp}_{2r}\) & \(\lfloor\lambda_1/2\rfloor\) & Richardson, Theorem 4.3.3, p.~233 \\
\(D_r=\mathfrak{so}_{2r}\) & \(\lfloor\lambda_1/2\rfloor+\mathbf1_{\{k=2,\ \lambda_1\text{ odd}\}}\) & Richardson, Lemma 4.4.1 and Theorem 4.4.2(a)--(c), p.~234 \\
\end{longtable}

The partition conditions are in \textbf{§4.1, p.~232}: even parts have even multiplicity in orthogonal types; odd parts have even multiplicity in symplectic type. In \(D_r\), a very even partition labels two orbits, and both have the displayed rank. The compact \(D_r\) formula is precisely a restatement of the three cases in Theorem 4.4.2: with two parts, either both are odd or the partition is \((r,r)\) with \(r\) even. This is a citation of a known computation, not a novelty claim for the all-element matrix formulas separately proved in the root module.

\subsection{CC-RK02 --- Completed exceptional nilpotent rank table}\label{cc-rk02-completed-exceptional-nilpotent-rank-table}

The following table groups the \textbf{Bala--Carter labels in Richardson's Carter convention} by the rank of the quotient differential. Every unlisted \textbf{nonzero} nilpotent orbit has rank \(1\); the zero orbit always has rank \(0\). This expressly excludes zero from the source's abbreviated ``all other'' rows.

\begin{longtable}{@{}p{.08\linewidth}p{.08\linewidth}p{.74\linewidth}@{}}
\toprule\noalign{}
\begin{minipage}[b]{\linewidth}\raggedright
Type
\end{minipage} & \begin{minipage}[b]{\linewidth}\raggedleft
Rank
\end{minipage} & \begin{minipage}[b]{\linewidth}\raggedright
Orbits
\end{minipage} \\
\midrule\noalign{}
\endhead
\bottomrule\noalign{}
\endlastfoot
\(E_6\) & 6 & \(E_6\) \\
& 5 & \(E_6(a_1)\) \\
& 4 & \(D_5\) \\
& 3 & \(E_6(a_3)\), \(D_5(a_1)\), \(A_5\) \\
& 2 & \(A_4+A_1\), \(D_4\), \(A_4\) \\
\(E_7\) & 7 & \(E_7\) \\
& 6 & \(E_7(a_1)\) \\
& 5 & \(E_7(a_2)\) \\
& 4 & \(E_7(a_3)\), \(E_6\), \(D_6\) \\
& 3 & \(E_7(a_4)\), \(D_6(a_1)\), \(E_6(a_1)\), \(D_5+A_1\), \(D_5\) \\
& 2 & \(A_6\), \(E_7(a_5)\), \(E_6(a_3)\), \(D_6(a_2)\), \(D_5(a_1)+A_1\), \(A_5+A_1\), \((A_5)'\), \(D_5(a_1)\), \(D_4+A_1\), \((A_5)''\), \(D_4\) \\
\(E_8\) & 8 & \(E_8\) \\
& 7 & \(E_8(a_1)\) \\
& 6 & \(E_8(a_2)\) \\
& 5 & \(E_8(a_3)\), \(E_7\) \\
& 4 & \(E_8(a_4)\), \(E_8(b_4)\), \(E_7(a_1)\) \\
& 3 & \(E_8(a_5)\), \(E_8(b_5)\), \(D_7\), \(E_7(a_2)\), \(E_6+A_1\), \(E_6\) \\
& 2 & \(E_8(a_6)\), \(D_7(a_1)\), \(E_8(b_6)\), \(E_7(a_3)\), \(E_6(a_1)+A_1\), \(A_7\), \(D_7(a_2)\), \(D_6\), \(D_5+A_2\), \(E_6(a_1)\), \(E_7(a_4)\), \(D_6(a_1)\), \(D_5+A_1\), \(D_5\) \\
\(F_4\) & 4 & \(F_4\) \\
& 3 & \(F_4(a_1)\) \\
& 2 & \(F_4(a_2)\), \(B_3\), \(C_3\) \\
\(G_2\) & 2 & \(G_2\) \\
\end{longtable}

\textbf{Derivation and source corrections.} Richardson's \textbf{Appendix, pp.~239--240}, lists the nonzero exponents at each orbit; \textbf{Lemma 2.1, Proposition 2.2 and Corollary 2.3, p.~230}, identify their number with the differential rank. The table above counts those lists, with the following explicit changes. Both appendix pages were visually inspected, so these are defects in the printed source, not OCR repairs.

\begin{enumerate}
\def\labelenumi{\arabic{enumi}.}
\item
  The regular \(E_8\) exponent list ends with printed \(27\); the correct last exponent is \(29\). The ordinary exponent list is \(1,7,11,13,17,19,23,29\), also directly inspected in Dinar's \textbf{Table 2, arXiv v4 p.~11}. The rank is \(8\) either way, but the exponent itself matters in CC-RK03.
\item
  The \(E_8(a_1)\) row includes a spurious \(9\), giving eight printed entries. Richardson's \textbf{Lemma 5.1, p.~234}, says the subregular exponents are precisely the first \(r-1\) ordinary exponents. Thus this row is \(1,7,11,13,17,19,23\), and its rank is \(7\).
\item
  The class \(E_6\subset E_8\) is omitted from the rank-\(3\) rows. A regular nilpotent in the standard \(E_6\) Levi has exactly the common ordinary exponents \(1,7,11\), by Richardson's \textbf{Propositions 3.1 and 6.1, pp.~231, 234}, or the multiplicity formulation \textbf{§8.1, p.~238}. Its rank is therefore \(3\). This is a theorem-based insertion, not a conjectural correction.
\item
  The \(E_8(a_2)\) entry has the source's explicit question mark at \(19\). The complete deduction below removes that question mark using Johnson--Sommers.
\end{enumerate}

No independent reconstruction of every exceptional-orbit computation is claimed. The table is the application of those known primary results with the four changes explicitly justified here.

\subsection{\texorpdfstring{CC-RK03 --- The formerly unresolved \(E_8(a_2)\) entry has rank six}{CC-RK03 --- The formerly unresolved E\_8(a\_2) entry has rank six}}\label{cc-rk03-the-formerly-unresolved-e_8a_2-entry-has-rank-six}

\textbf{Proposition.} For \(e\in\mathfrak e_8\) in the orbit \(E_8(a_2)\),

\[
 \operatorname{rank}(d\chi_{\mathfrak e_8})_e=6,
 \qquad
 \operatorname{Exp}(\mathfrak e_8,e)=\{1,7,11,13,17,19\}.
 \tag{CC-RK.2}
\]

\textbf{Direct known theorem.} Dinar, arXiv:1905.01829v4, \textbf{Corollary 5.5, p.~14, proof p.~15}, states that the quotient differential has rank \(s\) at the distinguished nilpotent orbit \(Z_r(a_{r-s})\) in the list fixed at the end of \textbf{§4, p.~9}. That list includes \(E_8(a_2)\), hence \(r=8,s=6\). The opening of \textbf{§5, p.~12}, explicitly identifies this as the case left undecided by Richardson. \textbf{Table 1, p.~10}, and \textbf{Table 2, p.~11}, give the relevant orbit and ordinary exponent data. The statement, its short proof, and the preceding local argument in Lemma 5.1, Proposition 5.3 and Lemma 5.4, pp.~12--14, were inspected. This is not a claim to have independently reconstructed the whole paper or its separate integrability theorem.

\textbf{Inputs for a separate complete deduction.} Johnson--Sommers, arXiv:1706.04820v2, \textbf{§1.3, equations (1.1)--(1.2), pp.~2--3; Definition 1.2 and Theorem 1.3(1)--(2), p.~3; Figure 1(b), p.~4; §5.9, p.~15}. In their notation this orbit is the first-family case \(s=2\), induced from a Levi with two orthogonal simple roots. Their parameter \(s\) counts the simple roots of the inducing Levi semisimple part and is different from Dinar's parameter \(s=6\). In \(E_8\), the short-root representation \(V_\phi\) is the adjoint representation. Their theorem identifies the lowest-degree adjoint copy in the ideal of the orbit closure inside the nilpotent cone at degree \(m_{8-2+1}=m_7=23\), and an additional copy at degree \(m_8=29\). The statement of \textbf{Corollary 3.2, p.~9}, also gives normality of this closure, but normality is not needed separately in the following proof.

\textbf{Complete proof of the deduction.} Let \(A=\mathbf C[\mathcal N]\), and let \(I\subset A\) be the homogeneous ideal defining \(\overline{G e}\). The adjoint isotypic multiplicity in \(A\), degree by degree, is one in each degree

\[
 1,7,11,13,17,19,23,29
\]

and zero in all other degrees. This is Kostant's result recalled by Johnson--Sommers in §1.3: these copies are the derivatives of the eight fundamental invariants. Their Theorem 1.3(1) says that \(I\) contains no adjoint copy of degree less than \(23\), and that it contains the copy in degree \(23\). Part (2) says it contains the copy in degree \(29\). Each homogeneous graded piece is finite dimensional and is a semisimple rational \(G\)-module over \(\mathbf C\). Applying \(\operatorname{Hom}_G(\mathfrak e_8,-)\) to

\[
 0\longrightarrow I_j\longrightarrow A_j
     \longrightarrow (A/I)_j\longrightarrow0
\]

is therefore exact. Consequently the adjoint multiplicity in \(A/I=\mathbf C[\overline{G e}]\) is one in the six degrees \(1,7,11,13,17,19\) and zero in the other degrees. Richardson's \textbf{Proposition 2.5 and Corollary 2.6, p.~231}, identify the total multiplicity with the quotient-differential rank and the graded multiplicities with the orbit exponents. This gives (CC-RK.2). \(\square\)

Equivalently, the degree-\(19\) derivative copy of a degree-\(20\) fundamental invariant cannot vanish on the entire orbit. If it vanished at \(e\), equivariance would make it vanish on \(Ge\), hence on its closure, contradicting the minimum degree \(23\) of an adjoint copy in \(I\). This explains directly why the missing exponent is present.

\textbf{Additional source precision.} The arXiv v2 line in §5.9 describing the \emph{subregular} \(E_8(a_1)\) ideal prints \(V_{17}\); Theorem 1.3 with \(s=1\) instead forces \(V_{29}\). Our argument uses the numbered theorem and the separate, consistent \(E_8(a_2)\) entry \(V_{23}\), not that erroneous subregular line. No unsupported claim about whether a later journal rendering corrects this typo is made.

\subsection{CC-RK04 --- Subregular uniqueness and rank do not characterize ADE}\label{cc-rk04-subregular-uniqueness-and-rank-do-not-characterize-ade}

\textbf{Known Lie-algebra statement.} For every complex simple Lie algebra \(\mathfrak g\), of dimension \(N\) and rank \(r\), there is exactly one nilpotent orbit of dimension \(N-r-2\). For \(x\) in that orbit,

\[
 \operatorname{rank}(d\chi_{\mathfrak g})_x=r-1.
 \tag{CC-RK.3}
\]

The uniqueness statement is Richardson's \textbf{§5 opening paragraph, p.~234}, and the rank statement is his \textbf{Lemma 5.1 on the same page}, for all simple types. Johnson--Sommers independently restate the uniqueness in \textbf{§1.2, p.~2}. Richardson attributes the rank theorem to Slodowy; the exact inspected numbered locator used here is Richardson's Lemma 5.1, not an invented Slodowy theorem number. In type \(A_1\), the subregular orbit is the zero orbit, and the same formula gives rank zero.

\textbf{Unipotent consequence, with complete proof.} The original adjoint group \(G\) also has exactly one unipotent orbit of dimension \(N-r-2\). Embed \(G\) faithfully in \(\operatorname{GL}(\mathfrak g)\) by its adjoint representation. For a nilpotent \(X\) in its Lie algebra, the finite matrix sum \(\exp(X)\) is unipotent and lies in \(G\), since it is the usual complex Lie-group exponential. Conversely, if \(u\in G\) is unipotent, put \(M=u-I\) and consider the polynomial curve

\[
 u^t=\sum_{j\geq0}\binom tj M^j.
\]

The sum is finite. At every nonnegative integer \(t\) its value is a power of \(u\), hence lies in the Zariski-closed subgroup \(G\). The inverse image of \(G\) under this curve is closed in \(\mathbf A^1\) and contains infinitely many points, so it is all of \(\mathbf A^1\). Differentiating at \(t=0\) gives \(\log u\in\operatorname{Lie}(G)\), and this matrix is nilpotent. The finite exponential and logarithm polynomials are inverse on the nilpotent and unipotent varieties, respectively, by their formal power-series identities modulo a sufficiently high power of a nilpotent matrix. Both commute with conjugation. They therefore give a \(G\)-equivariant isomorphism of these varieties, identifying their individual orbits and orbit dimensions. The known unique nilpotent orbit transfers to the claimed unique unipotent orbit. \(\square\)

Thus the uniqueness speculation is refuted in all simple types, as is the proposed strict inequality below \(r-1\) for the Lie-algebra quotient in nonsimply-laced types. The next proposition transfers the rank statement specifically at unipotent elements and therefore settles the original subregular rank assertion for both maps. The distinction between general Lie and group ranks remains essential away from this locus, as in BR-12305--08.

\subsection{CC-RK05 --- All nilpotent ranks transfer to the unipotent group quotient}\label{cc-rk05-all-nilpotent-ranks-transfer-to-the-unipotent-group-quotient}

Let \(G\) be any connected complex semisimple algebraic group, with maximal torus \(T\), Lie algebras \(\mathfrak g,\mathfrak t\), and Weyl group \(W\). The group need not be simply connected. Denote the two quotient maps by

\[
 \chi_{\mathfrak g}:\mathfrak g\longrightarrow\mathfrak t/W,
 \qquad
 \chi_G:G\longrightarrow T/W.
\]

\textbf{Proposition.} If \(e\in\mathfrak g\) is nilpotent and \(u=\exp(e)\), then

\[
 \operatorname{rank}(d\chi_G)_u
       =\operatorname{rank}(d\chi_{\mathfrak g})_e.
 \tag{CC-RK.4}
\]

Consequently the whole nilpotent table in CC-RK01--03 also gives the ranks at the corresponding unipotent elements in every central form of \(G\). In particular, the unique subregular unipotent orbit has quotient-differential rank \(r-1\) in \textbf{every} simple type. The ranks in \(B_r,C_r,F_4,G_2\) are \(r-1,r-1,3,1\), respectively, for both original maps.

\textbf{Primary structural inputs.} The Lie-algebra restriction isomorphism is recalled by Richardson in \textbf{§1, p.~229}. The group restriction isomorphism
\(\mathbf C[G]^G\simeq\mathbf C[T]^W\), without a simply-connected hypothesis, is Steinberg's \textbf{Corollary 6.4, p.~63}, proved from \textbf{Lemmas 6.2--6.3, p.~62}. These exact printed pages and statements were inspected. The following is a complete local analytic deduction; no global equivalence of the two ambient quotient maps is asserted.

\textbf{Proof.} The torus exponential has derivative the identity at \(0\). Choose a sufficiently small \(W\)-invariant open neighborhood \(V\subset\mathfrak t\) of \(0\) on which it is a biholomorphism onto its image \(V_T\subset T\), a \(W\)-invariant neighborhood of \(1\). Since the exponential is \(W\)-equivariant, it induces a biholomorphism

\[
 \psi:V/W\xrightarrow{\sim}V_T/W.
 \tag{CC-RK.5}
\]

Here the two finite-group quotients are open analytic neighborhoods in the analytifications of \(\mathfrak t/W\) and \(T/W\). Indeed a \(W\)-invariant open set is saturated for the finite orbit quotient, and its holomorphic functions on the quotient are the invariant holomorphic functions; applying this to the exponential and its equivariant inverse proves the assertion.

The open set \(U=\chi_{\mathfrak g}^{-1}(V/W)\) contains the nilpotent cone. A regular semisimple \(x\in U\) is conjugate to a \(t\in V\): its Cartan conjugacy class is a \(W\)-orbit, and its quotient value lies in \(V/W\). Equivariance of the exponential and the two restriction isomorphisms give

\[
 \chi_G(\exp x)=\psi(\chi_{\mathfrak g}(x))
 \tag{CC-RK.6}
\]

for all such \(x\). Regular semisimple elements form a Zariski-open dense subset of \(\mathfrak g\), hence are analytically dense in \(U\). Both sides of (CC-RK.6) are holomorphic on \(U\); after embedding the affine target \(T/W\) into an affine space, equality of each coordinate extends by continuity to all \(U\).

For completeness, the derivative of the exponential in a faithful matrix representation satisfies

\[
 (L_{u^{-1}})_*(d\exp)_e(Y)
   =\int_0^1 \exp(-t\,\operatorname{ad}e)Y\,dt
   =\left(\sum_{j\geq0}
       \frac{(-1)^j(\operatorname{ad}e)^j}{(j+1)!}\right)Y.
 \tag{CC-RK.7}
\]

The integral identity follows by differentiating the exponential power series term by term, or by integrating
\(\exp((1-t)e)Y\exp(te)\) before left translation. Since \(\operatorname{ad}e\) is nilpotent, the displayed sum is finite and has constant term \(1\). It is therefore an invertible endomorphism of \(\mathfrak g\), so \((d\exp)_e\) is invertible. Differentiating (CC-RK.6) at \(e\) gives

\[
 (d\chi_G)_u\circ(d\exp)_e
       =(d\psi)_0\circ(d\chi_{\mathfrak g})_e.
\]

Both \((d\exp)_e\) and \((d\psi)_0\) are isomorphisms. Taking ranks proves (CC-RK.4). Analytic and algebraic tangent differentials of these algebraic quotient maps agree, so this is the rank asked for in the original letter. \(\square\)

\subsection{Exact bibliography and available primary files}\label{cc-rk-bibliography}

Grothendieck, A. (1969, October 22). \emph{Letter to B. Kostant} {[}Manuscript, cote 123, physical PDF p.~23, footer 22{]}. Université de Montpellier. \url{https://grothendieck.umontpellier.fr/123.pdf}

Richardson, R. W. (1987). Derivatives of invariant polynomials on a semisimple Lie algebra. In M. Cowling, C. Meaney, \& W. Moran (Eds.), \emph{Miniconference on harmonic analysis and operator algebras} (Proceedings of the Centre for Mathematical Analysis, Australian National University, Vol. 15, pp.~228--241). Australian National University. \url{https://maths.anu.edu.au/files/CMAProcVol15-Richardson_1.pdf}

Johnson, B., \& Sommers, E. (2020). Equations for some nilpotent varieties. \emph{International Mathematics Research Notices, 2020}(14), 4433--4464. \url{https://doi.org/10.1093/imrn/rny132} . Inspected author version: arXiv:1706.04820v2, 5 February 2018, \url{https://arxiv.org/pdf/1706.04820v2} . All page locators in this module refer to that version.

Dinar, Y. (2020). On integrability of transverse Lie--Poisson structures at nilpotent elements. \emph{Journal of Geometry and Physics, 155}, Article 103690. \url{https://doi.org/10.1016/j.geomphys.2020.103690} . Inspected author version: arXiv:1905.01829v4, 7 April 2020, \url{https://arxiv.org/pdf/1905.01829v4} . All Dinar page locators in this module refer to that version.

Steinberg, R. (1965). Regular elements of semi-simple algebraic groups. \emph{Publications Mathématiques de l'IHÉS, 25}, 49--80. \url{https://doi.org/10.1007/BF02684397} . The general restriction isomorphism used here is Corollary 6.4, p.~63, with Lemmas 6.2--6.3, p.~62.

Local primary snapshots: \texttt{research/cc\_sources/Richardson1987.pdf}, its extracted text, and visually inspected table images \nolinkurl{Richardson1987_table-12.png}/\nolinkurl{Richardson1987_table-13.png}; \texttt{research/cc\_sources/JohnsonSommers2018.pdf} and its extracted text; \texttt{research/cc\_sources/Dinar2019v4.pdf} and its extracted text. The source and check ledgers accompanying this module record the exact retrieval references, inspection scopes, and actual proof-reading events.

% END consequences_criteria_adjoint_rank.tex

% BEGIN root_henselian.tex
\section{Cote 66, subject 21: the exact affine henselian comparison}

\subsection{Original statement and coefficient scope}
The official cote 66 PDF, physical page 4, archive footer 3, subject 21,
asks when
\[
 H^*(X,F)\longrightarrow H^*(X_0,F_0)
\]
is an isomorphism for a closed subscheme $X_0\subset X$, and asks
\emph{``Suffit-il que ce soit vrai en dim 0?''}
The margin identifies \emph{cohomologie \'etale}. The specified cases
are $X=\operatorname{Spec}A$, $X_0=\operatorname{Spec}(A/I)$ with
$A$ separated and complete for the $I$-adic topology, and \'etale
localization along a closed subscheme of an affine scheme.
The sheet does not explicitly write the quantifier on $F$ or restrict
it to torsion. We state those choices: the all-torsion-sheaf affine
question has the following exact answer, while a test on one fixed
sheaf does not suffice.

\subsection{BR-66-01. Degree zero for all torsion coefficients is an exact affine test}
Let $A$ be a commutative ring and $I\subset A$ any ideal, and set
$X=\operatorname{Spec}A$, $Z=\operatorname{Spec}(A/I)$. The following
conditions are equivalent:
\begin{enumerate}
\item $(A,I)$ is a henselian pair.
\item For every torsion abelian sheaf $F$ on $X_{\text{ét}}$,
      the restriction $H^0(X,F)\to H^0(Z,F|_Z)$ is an isomorphism.
\item For every such $F$ and every $q\geq0$,
      $H^q(X,F)\to H^q(Z,F|_Z)$ is an isomorphism.
\item For every finite morphism $f:Y\to X$, restriction gives a
      bijection from the clopen subsets of $Y$ to those of
      $Y\times_X Z$.
\end{enumerate}
In particular, the degree-zero condition can be tested just on the
sheaves $f_*\underline{\mathbf F}_2$ for all finite $f$.

\paragraph{Exact known inputs.}
Gabber's affine analogue of proper base change gives (1)$\Rightarrow$(3)
for arbitrary torsion sheaves, with no noetherian assumption and no
invertibility assumption on their torsion orders. The exact statement
used is The Stacks Project, Theorem 59.82.7, Tag \nolinkurl{09ZI}; the
original article is Gabber (1994), pp.~325--335. Its complete published
proof is not reproduced here. The equivalence (1)$\Leftrightarrow$(4)
is the idempotent criterion in The Stacks Project, Lemma 15.11.6(1),(3),
Tag \nolinkurl{09XI}. We prove the short implication from the stated
sheaf test to that criterion, to settle the quantifier precisely.

\paragraph{Complete proof of the additional implication and equivalence.}
The implication (3)$\Rightarrow$(2) is immediate. Assume (2), and
fix a finite morphism $f:Y\to X$. The sheaf
$F=f_*\underline{\mathbf F}_2$ is annihilated by $2$, so (2) applies.
Put $Y_Z=Y\times_X Z$ and $f_Z:Y_Z\to Z$. Proper base change in
degree zero for the finite morphism gives
\[
 (f_*\underline{\mathbf F}_2)|_Z
       \simeq (f_Z)_*\underline{\mathbf F}_2.
\]
An exact inspected locator for this known input is The Stacks
Project, Theorem 59.91.11, Tag \nolinkurl{095T}; only its degree-zero
finite-morphism case is used. By the definition of direct image,
condition (2) therefore gives
\[
 \Gamma(Y,\underline{\mathbf F}_2)
       \xrightarrow{\sim}\Gamma(Y_Z,\underline{\mathbf F}_2).
 \tag{HEN.1}
\]
A section of the constant two-element \'etale sheaf is the same as
a decomposition of the scheme into two open and closed subsets:
it is a section of the disjoint union of two copies of the scheme,
and the inverse images of its two components are the two subsets.
Thus (HEN.1) is precisely condition (4). Since $Y$ is affine,
these decompositions are also precisely the idempotents of its ring,
and restriction is reduction modulo $I$. The cited idempotent
criterion implies (1). The same proof works if (2) is assumed only
for this displayed family of finite direct images. Combining it with
Gabber's theorem proves all stated equivalences. \(\square\)

\subsection{BR-66-02. Both affine cases named on the sheet are affirmative}
If $A$ is $I$-adically separated and complete, the desired comparison
holds for every torsion abelian \'etale sheaf and every degree.
The same holds for the henselization $(A^h,IA^h)$ of any affine pair,
with closed subscheme $\operatorname{Spec}(A/I)$.
These are direct applications of known results: The Stacks Project,
Lemma 15.11.4, Tag \nolinkurl{0ALJ}, says an $I$-adically complete pair is
henselian; Lemma 15.12.1, Tag \nolinkurl{0A02}, constructs the henselization
as the filtered colimit of \'etale neighborhoods inducing an isomorphism
on the closed subscheme, proves it henselian, and identifies
$A^h/IA^h=A/I$. Apply BR-66-01. This identifies the affine \'etale
localization intended by the second case without asserting that an
arbitrary unrelated pair is henselian.

\subsection{BR-66-03. Testing only one coefficient sheaf in degree zero fails}
Even for smooth affine schemes over $\mathbf C$, comparison in degree
zero for one fixed finite constant sheaf does not imply comparison
in degree one for that same sheaf.

\paragraph{Complete counterexample.}
Take $X=\operatorname{Spec}\mathbf C[t,t^{-1}]$,
$Z=\{t=1\}=\operatorname{Spec}\mathbf C$, and
$F=\underline{\mathbf Z/2}$. Both schemes are connected, so
restriction on $H^0$ is the identity on $\mathbf Z/2$.
The morphism
\[
 \operatorname{Spec}\mathbf C[u,u^{-1}]\longrightarrow X,
 \qquad t=u^2,
\]
is a torsor for $\{1,-1\}\simeq\mathbf Z/2$. It is finite \'etale:
its ring is free of rank two over $\mathbf C[t,t^{-1}]$, and the
derivative $2u$ is a unit. The sign action is free and acts transitively
on every geometric fibre. It is not a trivial torsor, because a section
would give a unit $a(t)\in\mathbf C[t,t^{-1}]$ with $a(t)^2=t$.
Every Laurent-polynomial unit is $ct^m$: the least and greatest
exponents of a product are sums of the corresponding exponents,
so a product equal to $1$ forces the difference between the greatest
and least exponent of each factor to be zero. Then $a(t)^2=t$
would require $2m=1$, which is impossible.

The classification of abelian-sheaf torsors by $H^1$ (The Stacks
Project, Lemma 21.4.3, Tag \nolinkurl{03AJ}, with the section criterion
in Lemma 21.4.2, Tag \nolinkurl{03AI}) therefore gives
a nonzero class in $H^1(X,F)$. Its restriction to $Z$ is trivial:
the fibre consists of the two points $u=1,-1$ and has a section.
Consequently $H^1(X,F)\to H^1(Z,F|_Z)$ is not injective.
This proves the assertion. \(\square\)

The result answers the original affine torsion formulation, including
its degree-zero test with the universal coefficient quantifier.
It does not assert a general comparison theorem for arbitrary
quasicompact schemes or for nontorsion coefficients. The unseen
SA22 run's specific nonseparated counterexample and its residual
\texttt{U:henselian-comparison} are not identified with this elementary
fixed-sheaf counterexample.

\subsection*{Exact references and source inspection}
Grothendieck, A. (n.d.). \emph{List of subjects}, cote 66,
physical PDF p.~4, archive footer 3, subject 21.
\url{https://grothendieck.umontpellier.fr/66.pdf}.
The complete sheet and an enlarged subject-21 crop were visually read
in this run; ``66/004'' in the intake denotes the physical PDF page.

Gabber, O. (1994). Affine analog of the proper base change theorem.
\emph{Israel Journal of Mathematics, 87}, 325--335.
\url{https://doi.org/10.1007/BF02773001}.
The publisher abstract and bibliographic record were inspected;
the accessible complete theorem and proof used are the following
Stacks Project version, so no unretrieved theorem number from the
original article is asserted as newly checked.

The Stacks Project Authors. (2026). \emph{The Stacks Project}.
Theorem 59.82.7 (Gabber), \url{https://stacks.math.columbia.edu/tag/09ZI};
Lemma 15.11.6, \url{https://stacks.math.columbia.edu/tag/09XI};
Theorem 59.91.11, \url{https://stacks.math.columbia.edu/tag/095T};
Lemma 15.11.4, \url{https://stacks.math.columbia.edu/tag/0ALJ};
Lemma 15.12.1, \url{https://stacks.math.columbia.edu/tag/0A02}.
Lemmas 21.4.2--21.4.3 in \url{https://stacks.math.columbia.edu/tag/03AG}.
All displayed statements were directly inspected on 21 November 2026.

% END root_henselian.tex

% BEGIN root_deligne.tex
\section{Deligne's explicitly stated problem and the geometric weight boundary}
\label{sec:deligne-problem}

\subsection{D-WII-1.8.15: good degenerations of mixed Hodge structures}

There is an actual question next to the cited local-monodromy theorem.
Deligne (1980), \emph{La conjecture de Weil II}, 1.8.14--1.8.15,
printed pp.~178--179 (PDF pp.~42--43), asks to
``D\'egager une classe de `bonnes' variations de structures de Hodge mixtes''
on the punctured disc. In precise mathematical terms, for a good
unipotent variation $(V,W,F)$ with $N=\log T$, he asks for a finite
increasing filtration $M$ such that
\[
 NM_a\subseteq M_{a-2},\qquad
 N^b:\operatorname{gr}^{M}_{a+b}\operatorname{gr}^{W}_a V
       \xrightarrow{\sim}
       \operatorname{gr}^{M}_{a-b}\operatorname{gr}^{W}_a V
 \quad(b\geq0).
\]
He also asks for an asymptotic nilpotent orbit and for the filtration
$M$ to supply a mixed Hodge structure for the nilpotent orbit.
This is a faithful formal restatement of the displayed problem, not
an attribution of a question about the Riemann zeta function.

\paragraph{Answer: known geometric theorem, independent of the adopted RH.}
Admissible graded-polarizable variations provide the requested class.
For the precise definition use Kashiwara (1986), \S1.8,
conditions (1.8.1)--(1.8.4), pp.~993--994, and \S1.9: quasi-unipotence,
relative monodromy, and locally free extension of the Hodge and weight
graded bundles. Proposition 4.5.1 and Proposition 5.2.1,
pp.~1008--1010, identify the limiting data and the mixed Hodge
structure for $M$. His Theorem 4.4.1 and Theorem 4.5.2,
pp.~1007--1008, establish the compatibility of relative filtrations
and the normal-crossing extension criterion. These are theorems about
the defined class; the assertion is not that every variation satisfying
only Deligne's initial three axioms is automatically admissible.

For an explicit asymptotic estimate, Pearlstein (2000),
arXiv:math/0002030v3, Theorem 5.1(i)--(iii), p.~30,
gives a limiting filtration $F_\infty$ and the horizontal orbit
$F_{\rm nilp}(z)=\exp(zN)F_\infty$, with
\[
 d\bigl(F(z),F_{\rm nilp}(z)\bigr)
 \leq K(\operatorname{Im}z)^\beta
                  e^{-2\pi\operatorname{Im}z}
 \quad(\operatorname{Im}z\gg0)
\]
in the mixed Hodge metric. This supplies the asymptotic part in a
stronger quantitative form. The relative-monodromy mixed Hodge
structure and this estimate together answer the stated existence
problem. The question is settled by the cited Hodge-theoretic work;
no new proof or historical priority claim is made here.

\subsection{D-WII-1.9.5: the arithmetic base is a different gate}

At the end of 1.9.5, printed p.~181 (PDF p.~45), Deligne contrasts
normal-crossing divisors in a regular scheme of finite type over
$\mathbb Z$ with the preceding finite-field theory, already mentioning
$\operatorname{Spec}\mathbb Z[1/\ell]$. This observation must be
separated from the solved complex-analytic problem 1.8.15.
For the proper smooth local-field case a precise modern statement is
Scholze (2012), \emph{Perfectoid spaces}, Conjecture 9.3, p.~48
of arXiv:1111.4914:
for $H^i(X_{\overline K},\mathbb Q_\ell)$ the $j$th monodromy-graded
piece has Frobenius weight $i+j$. His Theorem 9.6, p.~49, proves
this in its stated set-theoretic complete-intersection situation
inside a projective smooth toric variety. We retain all of those
hypotheses; it is not a theorem for arbitrary proper smooth $X/K$.

BR-WM-01 identifies a concrete missing datum: a Frobenius
characteristic polynomial and the covariance equation do not determine
the monodromy filtration. Hence the adopted global zero-location
statements do not supply this datum through a reconstruction from
Euler polynomials. This is a proved failure of that proposed route,
not a metamathematical independence claim. The remaining gate for a
specified degeneration is the purity of each actual monodromy-graded
piece, with its actual $N$ and geometric realization.

\subsection{The universal stable curve: purity of a complex and weights of a fibre}

The RESIDUE intake, \S4(1), mixes two valid constructions with an
invalid identification. For a proper map from a smooth scheme over
a finite field, the constant complex is pure, and proper pushforward
preserves purity in the sense of mixed complexes. The exact source is
Beilinson, Bernstein, and Deligne (1982), Stabilit\'es 5.1.14 and its dual
form, Corollary 5.4.4, Theorem 5.4.5, and Corollary 5.4.6,
pp.~128--129 and 142. The splitting statements there hold after
base change to the algebraic closure: arithmetic Frobenius
semisimplicity is not silently added.

The universal stable-curve application uses Sun's stack decomposition
theorem, with its finite-diagonal and affine-stabilizer hypotheses,
as proved in CC-ST01 below. A splitting does not in general descend
merely from its existence on scheme charts. It does not say that every ordinary
cohomology group of a singular geometric fibre is pure of its
cohomological degree. BR-NODE-01 computes the actual group: its
weight-zero graph subspace and weight-one normalization part need
not have a single weight. Nearby smooth cohomology is a different
object. In the one-node genus-one smoothing its monodromy filtration
has graded pieces of weights zero and two, while the nodal special
fibre has only a one-dimensional weight-zero $H^1$.

For that last nearby-cycle statement the exact classical input is
SGA~7~II, Expos\'e XV, Th\'eor\`eme 3.4(ii)--(iii), printed
pp.~195--196: specialization, vanishing cycles, and the
Picard--Lefschetz inertia formula. The determinant
$\det H^1=\mathbb Q_\ell(-1)$ identifies the quotient twist once
the trivial specialization line has been identified. This explicit
case needs no general mixed-characteristic weight--monodromy theorem.
The entrance-path claims RESIDUE H.3, H.5, H.6 have their own
categorical proofs, unavailable in the supplied packet. Their
reported status is carried as an intake report, not promoted to a
newly checked theorem by this geometric calculation.

% END root_deligne.tex

% BEGIN root_core.tex
% Root mathematical contribution. Finding IDs identify deductions checked in this run.
\section{The exact abelian Artin--Tate consequence}
\label{sec:root-artin-tate}

We use the cohomological convention $\mathbb Q(1)$ of weight $-2$,
so that a geometric Frobenius at a good prime acts on its $\ell$-adic
realization by $p^{-1}$.  Coefficients may be extended to a number field
containing the values of the finite characters involved.  A rational
statement is recovered by grouping coefficient-conjugate characters.

\paragraph{Finding BR-AT-01 (the full shifted-zero consequence).}
Let $M$ be a finite direct sum of rank-one pure motives
\[
 M=\bigoplus_{\chi,n} (\chi\otimes\mathbb Q(n))^{\oplus m_{\chi,n}},
 \qquad m_{\chi,n}\in\mathbb Z_{\geq0},\quad n\in\mathbb Z,
\]
where $\chi$ is a finite-order character of
$\operatorname{Gal}(\overline{\mathbb Q}/\mathbb Q)$ with abelian image.
Use the primitive Dirichlet character associated with each $\chi$ for
its full Artin $L$-function.  Then
\[
 L(M,s)=\prod_{\chi,n}L(\chi,s+n)^{m_{\chi,n}}.
\]
Under the adopted Alp\"oge theorem, the nontrivial zeros contributed by
the summand of twist $n$ have real part $\tfrac12-n$.
These zero multiplicities add in the displayed product: none is cancelled
by a pole of another factor.  For a pure summand of weight $w=-2n$, this
is the expected line $\Re s=(w+1)/2$.

For a mixed object $V$ whose weight-graded pieces have this form, and a
finite set $S$ containing every ramified prime, the precise assertion is
\[
 L^S(V,s)=L^S(\operatorname{gr}^W V,s)
 =\prod_{\chi,n}L^S(\chi,s+n)^{m_{\chi,n}}.
\]
If Euler factors at primes in $S$ are deleted, the resulting extra zeros
lie on the lines $\Re s=-n$; they must not be called violations of the
nontrivial-zero assertion.  At ramified primes, no general multiplicativity
of inertia invariants for mixed extensions is being assumed.

\begin{proof}
At an unramified prime $p$, the eigenvalue for $\chi\otimes\mathbb Q(n)$
is $\chi(p)p^{-n}$, giving the Euler factor
$(1-\chi(p)p^{-s-n})^{-1}$.  The identity of the convergent Euler
products follows in a right half-plane and extends to their meromorphic
continuations.  This also proves the direct-sum formula.  On replacing
$z=s+n$, the strip $0<\Re z<1$ becomes
$-n<\Re s<1-n$; the adopted Dirichlet theorem places every zero there
on $\Re s=1/2-n$.  A primitive nonprincipal Dirichlet $L$-function is
entire; a principal primitive factor is $\zeta(s+n)$ and has its only
pole at $s=1-n$.  Such poles have integral real part, while the indicated
nontrivial zeros have half-integral real part.  Hence they cannot cancel.

For a filtration preserved by Frobenius, choose a basis respecting the
filtration.  Its matrix is block upper triangular; the determinant of
$1-p^{-s}\operatorname{Frob}_p$ is the product of the determinants of
the diagonal blocks.  At every $p\notin S$ this gives the mixed formula.
Deleting a nonzero local rank-one factor multiplies the full $L$-function
by $1-\alpha p^{-s-n}$ with $|\alpha|=1$.  Its zeros satisfy
$p^{-\Re s-n}=1$, hence $\Re s=-n$.  A ramified one-dimensional
character with zero inertia-invariant space contributes the factor $1$,
so there is nothing to delete in that case.  This proves all the assertions.
\end{proof}

The elementary Euler-factor argument is included because it specifies the
mixed, twisted, and incomplete forms not stated by the Lean declarations.
It does not assert that every general mixed-motivic $L$-function is
determined at every bad prime by its semisimplification.  In particular,
calling an object ``mixed Artin--Tate'' does not make its extension class
part of a Dirichlet character.

\paragraph{Archimedean convention.}
For $\chi(-1)=(-1)^a$, $a\in\{0,1\}$, the factor of
$\chi\otimes\mathbb Q(n)$ is $\Gamma_{\mathbb R}(s+n+a)$, where
$\Gamma_{\mathbb R}(s)=\pi^{-s/2}\Gamma(s/2)$.  Its real Frobenius
acts as $\chi(-1)(-1)^n$.  Both the Hodge type $(-n,-n)$ and that
sign are needed to recover the archimedean factor.  The sign alone does
not recover the weight: $\mathbb Q(0)$ and $\mathbb Q(2)$ both have
trivial real Frobenius.  The corrected real-Frobenius classification at
cote 18/038 is treated with the original text below.

\subsection{Mixed coefficients: a rigorous obstruction to reconstructing extensions}

Deligne--Goncharov (2005), \S1.1, equations (1.1.2), (1.1.4), and the
proof of Proposition 1.7, construct the Kummer extension $K(x)$ and
identify
\[
 \operatorname{Ext}^1_{\mathrm{MT}(\mathbb Q)}
   (\mathbb Q(0),\mathbb Q(1))=\mathbb Q^\times\otimes_{\mathbb Z}\mathbb Q.
\]
These exact locators are in the supplied file \nolinkurl{Sec1ccc.tex}; this
identification is a cited theorem, not a new result here.

\paragraph{Finding BR-MIX-01 (the Kummer extension survives identical
good Euler factors).}
There are nonisomorphic mixed Tate motives with the same weight-graded
object and the same Euler polynomials at every prime of simultaneous
good reduction.  Explicitly,
\[
 0\longrightarrow\mathbb Q(1)\longrightarrow K(2)
 \longrightarrow\mathbb Q(0)\longrightarrow0
\]
is nonsplit, whereas $\mathbb Q(0)\oplus\mathbb Q(1)$ is split.  For
$S$ containing $2$ and the realization prime $\ell$, both have
\[
 L^S(s)=\zeta^S(s)\zeta^S(s+1).
\]
Thus even complete knowledge of these good Euler factors, their
meromorphic continuation, and all their zeros does not reconstruct
the mixed extension.  This refutes that particular reconstruction
rule; it is not a logical independence theorem about RH.

\begin{proof}
The $2$-adic valuation gives a homomorphism
$v_2:\mathbb Q^\times\to\mathbb Z$ and hence a linear map
$\mathbb Q^\times\otimes\mathbb Q\to\mathbb Q$ taking $2\otimes1$
to $1$.  Consequently the class of $2$ is nonzero, and the cited
Ext-identification makes its extension nonsplit.  An isomorphism to
the split object would preserve the unique weight filtration and
give a splitting, a contradiction.  By Deligne--Goncharov,
Proposition 1.7 and its proof, this Kummer object is unramified away
from $2$ (and the $\ell$-adic realization is used away from $\ell$).
For $p\notin S$ its Frobenius matrix in a weight-adapted basis is
upper triangular with diagonal entries $p^{-1},1$.  Its Euler
polynomial is therefore $(1-T)(1-p^{-1}T)$, the polynomial of the
split object.  Taking their Euler products proves the identity.
\end{proof}

The same proof works for every $x\in\mathbb Q^\times$ with a nonzero
valuation somewhere, after deleting the primes dividing $x$.
The unipotent radical of the mixed Tate motivic Galois group, its
extension classes, and its period matrix are additional information.
This is a concrete reason why a theorem about the zero divisor of
these $L$-functions does not by itself compute the groups in step 17
of \emph{Remark}.  In particular the tentative footnote there gives
neither the degrees nor the primary parts required for a $K$-group
vanishing statement.  The assertion that \emph{all} positive-degree
$K$-groups of $\mathbb Z$ vanish is false already for
$K_1(\mathbb Z)=\{\pm1\}$; it is not an assertion made by that footnote.

\section{Cote 1: what the Hilbert-space trace theorem settles}
\label{sec:root-trace}

The supplied scans 1/079--080 contain the Fredholm determinant estimates
and the trace/eigenvalue formula; 1/111--112 contain the genus-zero
argument.  They do not mention $\xi$, $\Xi$, a Riemann-zero operator,
or RH.  Those are applications in the supplied SA23 intake, \S4.1,
and in this work.  An exact attribution therefore distinguishes the
original analytic theorem from its later arithmetic use.

Write
\[
 \xi(s)=\tfrac12s(s-1)\pi^{-s/2}\Gamma(s/2)\zeta(s),\qquad
 \Xi(t)=\xi(\tfrac12+it),\qquad F(w)=\frac{\Xi(\sqrt w)}{\Xi(0)}.
\]
Evenness of $\Xi$ makes $F$ single-valued and entire.  The classical
Hadamard product and zero-counting facts used here are Titchmarsh--
Heath-Brown (1986), \S2.12 (the $\xi$ product), Theorems 9.3--9.4,
and equation (9.4.3), p.~214 (the Riemann--von Mangoldt count).
We use their statements, not a
new proof of those analytic inputs.

\paragraph{Finding BR-TR-01 (a positive trace-class realization, and its limit).}
There is a diagonal trace-class operator $B$ with
\[
 F(w)=\det(I-wB).
\]
It can be constructed before RH by taking one zero $t_j$ from each
pair $\{t,-t\}$ of $\Xi$, counted with multiplicity, and putting
$B=\operatorname{diag}(t_j^{-2})$ on $\ell^2$.
For this normal realization, RH is equivalent to $B\geq0$.
Consequently the adopted RH gives such a positive realization.
Every trace-class operator $T$ with $\det(I-wT)=F(w)$ has the same
nonzero eigenvalues with algebraic multiplicity and satisfies
\[
 \operatorname{Tr}(T^k)=\sum_jt_j^{-2k}>0\quad(k\geq1)
\]
under RH.  It is false that every such $T$ is positive or self-adjoint.
Moreover, no trace-class operator has $\Xi(t)$ itself as an ordinary
Fredholm determinant in an affine variable $t$.

\begin{proof}
The nontrivial zeta zeros have counting function $O(T\log(T+2))$.
Partial summation, or a dyadic decomposition, gives
$\sum_j|t_j|^{-2}<\infty$; finitely many bounded zeros cause no issue
because $\Xi(0)\ne0$.  Thus the displayed diagonal operator is trace
class.  The entire even function $\Xi$ has order one, so $F$ has
order at most $1/2$.  Hadamard factorization in order less than one,
together with $F(0)=1$, gives the genus-zero product
$F(w)=\prod_j(1-wt_j^{-2})$, which is precisely its Fredholm determinant.

If RH holds, all $t_j$ are real and nonzero; thus every diagonal
entry of $B$ is positive.  Conversely $B\geq0$ implies
$t_j^{-2}>0$, so $t_j$ is real.  Indeed, for a nonzero complex number
$t$, $t^2$ positive real forces $t$ real; the purely imaginary
alternative has negative square.  Hence every zero of $\xi(s)$
has real part $1/2$.

For a general trace-class $T$, the Fredholm zero theorem identifies
its nonzero eigenvalues and their algebraic multiplicities with the
reciprocals of the zeros of its determinant.  The trace theorem then
gives the power-trace formula; the sum converges absolutely.  To see
why positivity of $T$ does not follow, start with $B$ under RH and
replace its first two diagonal entries $b_1,b_2$ by the block
$\left(\begin{smallmatrix}b_1&1\\0&b_2\end{smallmatrix}\right)$.
This is a finite-rank perturbation, remains trace class, has the
same determinant, and is not self-adjoint in the given Hilbert
inner product.  It therefore is not a positive operator there.

Finally suppose $\Xi(t)=c\det(I+(t-a)T)$ for trace-class $T$ and
$c\ne0$.  Its zeros would be $a-\lambda_j^{-1}$, with the nonzero
eigenvalues $\lambda_j$ of $T$ counted algebraically.  The trace-class
eigenvalue inequality forces $\sum_j|\lambda_j|<\infty$.
But the Riemann--von Mangoldt asymptotic gives divergent sum of the
reciprocals of the distances of the $\Xi$ zeros from any fixed $a$.
For example, in dyadic height intervals the count is of order
$2^k k$, whereas every such distance is at most $C2^k$ for large
$k$ (the imaginary parts of the $t$-zeros are bounded by $1/2$).
The sum over that interval is at least a positive constant times
$k$.  This contradicts the required summability.  A zero at $a$
would already contradict the nonzero determinant value at $t=a$.
\end{proof}

The construction uses the zero set; it does not solve the stronger
geometric Hilbert--P\'olya problem of producing the operator from
primes, correspondences, or a cohomology theory independently of its
prescribed spectrum.  The result settles the exact positive normal
realization and its trace inequalities, not that additional
construction problem.

\paragraph{Finding BR-TR-02 (a strictly positive moment hierarchy).}
Under RH put $b_j=t_j^{-2}>0$ and
$m_r=\sum_j b_j^{r+1}$ for $r\geq0$.  For every integer $d\geq0$,
both Hankel matrices $(m_{a+b})_{0\leq a,b\leq d}$ and
$(m_{a+b+1})_{0\leq a,b\leq d}$ are positive definite.
For $x\geq0$,
\[
 H(x)=-\frac{F'(-x)}{F(-x)}
      =\operatorname{Tr}\bigl(B(I+xB)^{-1}\bigr)
      =\sum_j\frac{b_j}{1+xb_j},
 \qquad (-1)^r H^{(r)}(x)>0\quad(r\geq0).
\]
These are quantitative inequalities in the cote 1 trace framework,
not assertions that its unspecified geometric operators are positive.

\begin{proof}
The positive measure $\mu=\sum_jb_j\delta_{b_j}$ is finite because
$B$ is trace class, and it has infinitely many distinct support points:
$\Xi$ has infinitely many zeros and finite multiplicity at each point.
For a nonzero polynomial $P(u)=\sum_{a=0}^d c_au^a$,
\[
 \sum_{a,b=0}^d\overline{c_a}c_bm_{a+b}
   =\int |P(u)|^2\,d\mu(u)>0.
\]
A nonzero polynomial cannot vanish on infinitely many distinct points.
Inserting an additional factor $u>0$ proves the second positive
definiteness assertion.  The logarithmic derivative follows from
the normally convergent product for $F$; for $x\geq0$ all its
denominators are positive.  On compact subintervals of $[0,\infty)$,
termwise differentiation is dominated by
$r!\|B\|^r\sum_jb_j$.  Thus
\[
 (-1)^rH^{(r)}(x)
 =r!\sum_j\frac{b_j^{r+1}}{(1+xb_j)^{r+1}}>0.
\]
This proves convergence, the identities, and strictness.
\end{proof}

\section{Weights, monodromy, and the additional geometric input}
\label{sec:root-weights}

Deligne (1980), \emph{La conjecture de Weil II}, Theorem 1.8.4,
states that for a punctually pure lisse sheaf of weight $w$ on the
punctured smooth curve over a finite field, the $j$th graded piece
of its local monodromy filtration is pure of weight $w+j$.
This is a theorem with a geometric-purity hypothesis, not an
implication from a list of zeros of $\zeta(s)$.
The mixed-characteristic question for $H^i$ of a smooth proper
variety over a local field is the corresponding assertion that
$\operatorname{gr}_j^M H^i$ has weight $i+j$.
An exact modern statement attributing it to Deligne is Scholze (2012),
\emph{Perfectoid Spaces}, Conjecture 9.3; Theorem 9.6 proves the
set-theoretic complete-intersection case in smooth projective toric
varieties.  No general mixed-characteristic proof is contained in
the supplied RH declarations.

\paragraph{Finding BR-WM-01 (a determinant and the Tate covariance do not
determine monodromy purity).}
For every prime power $q>1$ there are two Weil--Deligne linear
algebra data with the same semisimple Frobenius and the same
Frobenius determinant, both satisfying $FNF^{-1}=q^{-1}N$, for
which the monodromy-purity statement centered at weight $1$ has
different truth values.

\begin{proof}
Take a two-dimensional vector space with basis $e_0,e_2$ and
$Fe_0=e_0$, $Fe_2=qe_2$.  In the first datum put
$Ne_2=e_0$, $Ne_0=0$; in the second put $N=0$.
The covariance is checked on the basis and holds in both cases.
Their common Frobenius determinant is $(1-T)(1-qT)$.
For the first datum the monodromy filtration centered at zero is
$M_{-2}=0$, $M_{-1}=M_0=\langle e_0\rangle$, $M_1=V$.
Its graded pieces in degrees $-1$ and $1$ have eigenvalues $1$
and $q$, hence weights $0$ and $2$, which are $1-1$ and $1+1$.
For the second datum the monodromy filtration is concentrated in
degree zero, so its only graded piece is $V$, with weights $0,2$;
it is not pure of weight $1$.  This proves the claim.
\end{proof}

These data refute the proposed \emph{linear-algebra deduction} from
Frobenius roots and covariance.  They are not a counterexample to
Deligne's geometric conjecture: no realization as the cohomology
of a smooth proper variety of the required weight is asserted for
the second datum.  Also, the local Euler factor defined on
$\ker N$ differs between the two examples.  It would be incorrect
to call these examples indistinguishable by every ramified local
$L$-factor.  What fails is precisely the attempted passage from
a determinant on $V$ to its monodromy-graded purity.

\subsection{The tame triangle and the nodal universal curve}

The intake SA20, row C.008, attributes its triangle to cote 122;
the original scan for that row is not among the supplied files.
Its finite $\mathbb Z/\ell^a$ coefficients and twists type the
duality maps.  Numerical Frobenius weights require a rational
$\ell$-adic realization and a purity theorem.  After that passage,
$H^1(I_t,\mathbb Q_\ell)=\mathbb Q_\ell(-1)$ has weight $2$;
the monodromy logarithm lowers weights by $2$.  These facts explain
the appearance of adjacent twists, but do not prove the
mixed-characteristic monodromy-purity assertion.

There is a further distinction in the RESIDUE intake, H.3--H.6:
the ordinary cohomology of a singular nodal fibre is not the same
filtered object as the nearby smooth cohomology with its full
monodromy filtration.  The following elementary calculation makes
the distinction exact.

\paragraph{Finding BR-NODE-01 (the weight-zero part belongs to the dual graph).}
Let $C/\mathbb F_q$ be a geometrically connected proper nodal curve,
with $\ell\ne\operatorname{char}\mathbb F_q$, whose geometric
components, nodes, and two branches at every node are all defined over
$\mathbb F_q$.  Write $\widetilde C=\coprod_v C_v$ for the
normalization and $\Gamma$ for the dual graph.  There is a
Frobenius-equivariant exact sequence
\[
 0\longrightarrow H^1(\Gamma,\mathbb Q_\ell)
 \longrightarrow H^1(C_{\overline{\mathbb F}_q},\mathbb Q_\ell)
 \longrightarrow\bigoplus_vH^1(C_v,\mathbb Q_\ell)
 \longrightarrow0.
\]
The left term has weight zero and the right term weight one.
In particular a split rational curve with one ordinary node has
$H^1\simeq\mathbb Q_\ell$ of weight zero.  Its smooth genus-one
nearby fibre, if chosen in a semistable degeneration, also has the
weight-two monodromy quotient; that quotient is not an additional
piece of the singular fibre's ordinary $H^1$.

\begin{proof}
For the finite normalization map $\nu$, the exact sequence of
\'etale sheaves is
\[
 0\longrightarrow\mathbb Q_\ell
 \longrightarrow\nu_*\mathbb Q_\ell
 \longrightarrow\bigoplus_{e\in E(\Gamma)}\mathbb Q_{\ell,e}
 \longrightarrow0.
\]
Away from the nodes the first map is an isomorphism.  At a node
its stalk is the diagonal inclusion $\mathbb Q_\ell\to
\mathbb Q_\ell^2$, whose cokernel is one-dimensional; choosing
an orientation identifies the quotient with the difference of
branch values.  Finite pushforward and cohomology give the usual
normalization cohomology.  The first part of the long exact
sequence is consequently
\[
 0\to\mathbb Q_\ell\to\mathbb Q_\ell^{V(\Gamma)}
 \xrightarrow{\partial}\mathbb Q_\ell^{E(\Gamma)}
 \to H^1(C)\to\bigoplus_vH^1(C_v)\to0.
\]
The cokernel of the incidence map $\partial$ is graph cohomology.
All vertices and branches are rational, so Frobenius acts trivially
on this graph complex, giving weight zero.  Each $C_v$ is a smooth
proper curve; its $H^1$ has weight one by the finite-field Weil
theorem.  If $C$ is rational with one node, its graph has one
loop and its normalization has no $H^1$, proving the special case.
For a semistable genus-one smoothing, use the cited Tate-uniformization
calculation (Ochiai, 2010, \S4, p.~4, displayed Tate-module and
$H^1$ exact sequences; Berger, 2004, \S II.4.1--II.4.3 gives the
period calculation): the nonzero rank-one monodromy pairs the
weight-zero subspace with a weight-two quotient, exactly as in the
first datum of BR-WM-01.
The nearby fibre has dimension two in $H^1$, whereas the singular
fibre has dimension one.  Thus the extra quotient cannot be
ordinary $H^1$ of the singular curve.
\end{proof}

The already established decomposition theorem concerns the pure
\emph{complex} pushed forward by a proper morphism.  It does not
assert that every ordinary stalk cohomology group of that complex
is pure of its cohomological degree at a singular stratum.  This
distinction prevents a false purity claim for the boundary curves
of $\overline{\mathcal M}_{g,n}$.  The supplied entrance-category
claims H.3--H.6 remain background at the level at which their intake
checked them; their missing original proofs have not been silently
promoted to checks made in this run.

\subsection{The standard conjectures: which additional gate remains}

Grothendieck's \emph{Standard conjectures on algebraic cycles}
(1969), \S2, p.~195, asks for algebraicity of the K\"unneth
projectors (C); \S3, pp.~195--197, asks for the Lefschetz
standard conjectures (A and B); \S4, pp.~197--198, asks for
positivity of the primitive algebraic intersection form (Hdg).
The implication to the finite-field Weil conjectures is stated in
\S4, Remark (2), p.~197; the relation to numerical equivalence
is in Remarks (3)--(4), pp.~197--198.
There is an essential positive consequence of finite-field purity:
Katz--Messing (1974), Theorem 2(1), pp.~76--77, constructs algebraic
K\"unneth projectors over finite fields.  The distinct weights make the
Frobenius characteristic polynomials in different degrees coprime;
Chinese-remainder polynomials in the algebraic Frobenius correspondence
then act as the cohomological projectors.  Theorem 1 of that paper
transfers the required Frobenius characteristic polynomials to crystalline
cohomology under its stated axioms.  These are already established
consequences of the finite-field theorem, and a complete new proof of
that theorem would give another route to them.  They do not, by
themselves, construct the inverse Lefschetz correspondence or prove
positivity on algebraic primitive cycles in general.  The classical
number-field Dirichlet assertion does not add those correspondences.

For the abelian Artin--Tate subcategory the relevant structures are
already explicit: Artin projectors are obtained from finite
correspondences and character idempotents, and Tate objects have
their standard duality.  BR-AT-01 settles the global zero-location
assertion for their $L$-functions.  For arbitrary smooth projective
varieties the current gate is still an algebraic correspondence
realizing the inverse Lefschetz operator, together with the
required positive pairing on algebraic primitive classes.  This
work supplies neither that correspondence nor a counterexample to
the standard conjectures.  The unresolved conjectures are recorded
as such, rather than marked ``refuted'' from the absence of a proof.

\section{D\'erivateurs, localization, and information forgotten}
\label{sec:root-derivators}

The supplied DER\_LOCALIZERS intake identifies the linear
``Thm(?)'' at cote 134-8/048--054 and its remaining acyclicity
condition C.10.  Its original scans and the full proof of its
general counterexample C.26 are not supplied.  A specific motivic
instance is nevertheless established in the literature:
Voevodsky (2000), Theorem 3.1.12, gives strict homotopy invariance
for Nisnevich cohomology of homotopy-invariant presheaves with
transfers over a perfect field; Proposition 3.2.3 identifies the
Suslin-complex localization; Theorem 3.2.6 gives the embedding.
Those are geometric inputs to the localization theorem, independent
of RH.

\paragraph{Finding BR-DER-01 (constant rational coefficients detect exactly
rational homology).}
Let $\mathcal C$ be a nonzero presentable stable
$\mathbb Q$-linear category, with a specified nonzero object
$\mathbf1$ and the colimit-preserving exact action
$D(\mathbb Q)\to\mathcal C$, $V\mapsto V\otimes\mathbf1$.
For a functor $f:A\to B$ between small categories, the induced map
on homotopy colimits of the constant diagram $\mathbf1$ is an
equivalence if and only if $Nf$ is a rational homology equivalence.
Thus this particular constant-object localizer does not change
when one equips $\mathcal C$ with additional Tate twists,
Frobenius actions, or $L$-functions while retaining the specified
rational action and a nonzero unit.

\begin{proof}
The tensoring of a presentable stable category over spaces, followed
by the $\mathbb Q$-linear action, identifies the homotopy colimit
of a constant diagram on $A$ with
$C_*(NA;\mathbb Q)\otimes\mathbf1$.  Let $K$ be the cone in
$D(\mathbb Q)$ of the chain map induced by $Nf$.  Exactness shows
that the homotopy-colimit map is an equivalence exactly when
$K\otimes\mathbf1=0$.
Every complex of vector spaces splits as its homology together
with a contractible complex: choose complements to boundaries
inside cycles and to cycles inside each chain group.  The
differential identifies the latter complement with the
boundaries one degree lower, giving the contractible summands.
This works for unbounded complexes degree by degree.  Hence
$K\simeq\bigoplus_i H_i(K)[i]$ in $D(\mathbb Q)$.
If some $H_i(K)\ne0$, choose a nonzero vector and a linear
projection onto its span.  This exhibits $\mathbb Q[i]$ as a
retract of $K$.  Its image $\mathbf1[i]$ is then a nonzero
retract of $K\otimes\mathbf1$, a contradiction.  Conversely,
$K=0$ plainly implies $K\otimes\mathbf1=0$.
Thus the latter vanishing is equivalent to zero rational homology
of the cone, which is exactly the asserted condition on $Nf$.
\end{proof}

The quantifiers are deliberately explicit.  This proves the
constant-unit statement needed for rational motivic categories;
it does not identify all coefficient-dependent localizers, and
does not assert that every reversible localizer is of this form.
It also avoids the intake's unverified blanket full-faithfulness
assertion for arbitrary coefficient rings.

\paragraph{Finding BR-DER-02 (the reduced Tate object is essential).}
In geometric motives over a field, with the rational or usual
integral coefficients for which the cited motivic identities hold,
\[
 M(\mathbb G_m)\simeq\mathbf1\oplus\mathbf1(1)[1],\qquad
 \widetilde M(\mathbb G_m)[-1]\simeq\mathbf1(1).
\]
The formula in the DER\_LOCALIZERS intake omits the reduction if
$\mathbb Z_{\mathrm{tr}}(\mathbb G_m)$ means the full representable
presheaf with transfers.  Its full localized motive shifted by
$-1$ has an additional $\mathbf1[-1]$ summand.

\begin{proof}
Cover $\mathbb P^1$ by its two affine lines.  The motivic
Mayer--Vietoris triangle is
\[
 M(\mathbb G_m)\longrightarrow\mathbf1\oplus\mathbf1
 \longrightarrow M(\mathbb P^1)\longrightarrow M(\mathbb G_m)[1].
\]
The projective-line decomposition is
$M(\mathbb P^1)=\mathbf1\oplus\mathbf1(1)[2]$.
Each affine-line inclusion has degree-one component to
$\mathbf1$ and zero component to the Tate summand; the latter
lies in $\operatorname{Hom}(\mathbf1,\mathbf1(1)[2])
=\operatorname{Pic}(\operatorname{Spec}k)=0$.
After an invertible change of basis in $\mathbf1\oplus\mathbf1$,
the middle map is the identity on one unit summand and zero on
the other, with zero map to $\mathbf1(1)[2]$.  Taking its fibre
gives $M(\mathbb G_m)=\mathbf1\oplus\mathbf1(1)[1]$.
The rational point $1\in\mathbb G_m$ splits the augmentation;
its reduced fibre is exactly the second summand.  Shifting proves
the assertion.  The projective-line definition and reduction are
Voevodsky (2000), \S2.1 following Proposition 2.1.3; the Picard
calculation is Corollary 3.4.3.  This proof spells out the corrected
form rather than reproving the construction of motives.
\end{proof}

\section{A concrete boundary for the anabelian and Teichm\"uller bridges}
\label{sec:root-anabelian}

\paragraph{Finding BR-ANA-01 (identical abelian cohomological data do not
reconstruct a four-punctured sphere).}
The two hyperbolic curves over $\mathbb Q$
\[
 U_2=\mathbb P^1\setminus\{0,1,\infty,2\},\qquad
 U_3=\mathbb P^1\setminus\{0,1,\infty,3\}
\]
are not isomorphic even over $\overline{\mathbb Q}$, although
their geometric $\ell$-adic cohomology representations are
isomorphic in every degree.  For every prime $p\notin\{2,3\}$
and integer $r\geq1$ they have the same point count and local zeta
function:
\[
 \#U_i(\mathbb F_{p^r})=p^r-3,\qquad
 Z(U_i/\mathbb F_p,T)=\frac{(1-T)^3}{1-pT}.
\]
Thus the global good-prime cohomological zeta functions agree,
even after imposing RH on their Dirichlet/Tate factors.

\begin{proof}
An isomorphism of the affine smooth curves extends uniquely to
their smooth projective compactifications and permutes their
four missing points.  The six possible cross-ratios of an
unordered quadruple with parameter $t$ are
\[
 t,\quad 1-t,\quad t^{-1},\quad(1-t)^{-1},\quad
 t/(t-1),\quad(t-1)/t.
\]
For $t=2$ the resulting set is $\{2,-1,1/2\}$; for $t=3$ it
is $\{3,-2,1/3,-1/2,3/2,2/3\}$.  They are disjoint, so no
projective-linear transformation carries one unordered boundary
to the other.  This proves nonisomorphism.

For either boundary $D$ the localization sequence on
$\mathbb P^1_{\overline{\mathbb Q}}$ is Galois equivariant.
Since its four points are rational, the residue sequence gives
\[
 H^0(U_i,\mathbb Q_\ell)=\mathbb Q_\ell,\qquad
 H^1(U_i,\mathbb Q_\ell)
 =\ker\bigl(\mathbb Q_\ell(-1)^4
             \xrightarrow{\text{sum}}\mathbb Q_\ell(-1)\bigr)
 \simeq\mathbb Q_\ell(-1)^3,
\]
and higher ordinary cohomology vanishes.  This determines the
representations independently of $t$.  Alternatively compactly
supported cohomology has $H_c^1=\mathbb Q_\ell^3$ and
$H_c^2=\mathbb Q_\ell(-1)$.
At each stated prime the four rational points remain distinct,
so subtracting them from $\mathbb P^1(\mathbb F_{p^r})$ gives
$p^r+1-4$.  Exponentiating
$\sum_{r\geq1}(p^r-3)T^r/r$ yields the claimed zeta function.
\end{proof}

This is an explicit counterexample to reconstructing the
Teichm\"uller parameter from the indicated cohomological or
$L$-function data.  It does not contradict Grothendieck's
anabelian reconstruction proposals, whose input is a nonabelian
arithmetic fundamental group with its Galois action and peripheral
structure.  It shows why the abelian Artin--Tate result, even with
effective recovery of all its finite characters, is not itself
that nonabelian input.  A comparison with $\widehat{\mathrm{GT}}$
would have to specify and prove compatibility of a construction
with the noncommutative fundamental groups, their inertia, and
the tower maps.  None is constructed by the coincident words
``hexagon'' and ``tower'' in the supplied proof sketch.


% END root_core.tex

% BEGIN root_weight_criterion.tex
\section{A finite exterior-power gate for weight--monodromy}
\label{sec:finite-weight-gate}

Scholze (2012), arXiv:1111.4914, Lemma 9.5, printed p.~49,
records Deligne's tensor-power criterion: purity of the monodromy
filtration is equivalent to upper weight bounds on the inertia
invariants of every tensor power and every dual tensor power.
The following finite exterior-power formulation is proved here in
full. It reduces that infinite collection of representations; it
does not assert the needed bounds for an arbitrary geometric
representation. No historical novelty claim is made.

\paragraph{Finding BR-WM-02 (finite exterior-power criterion).}
Let $V$ be a complex vector space of dimension $d\geq1$, let
$q>1$, let $w\in\mathbb R$, and let $F$ be invertible and $N$
nilpotent with $FNF^{-1}=q^{-1}N$. Write $M$ for the monodromy
filtration centered at zero. On exterior powers, $N$ denotes the
induced derivation; on the dual it is $-{}^tN$. The following
three conditions are equivalent.
\begin{enumerate}
\item On $\operatorname{gr}_j^M V$, every eigenvalue of $F$
has absolute value $q^{(w+j)/2}$.
\item For each $1\leq m\leq d$, all eigenvalues on
$\ker(N:\bigwedge^m V\to\bigwedge^m V)$ have absolute value
at most $q^{mw/2}$, and all eigenvalues on
$\ker(N:\bigwedge^m V^\vee\to\bigwedge^m V^\vee)$ have absolute
value at most $q^{-mw/2}$.
\item For every $1\leq m\leq d$ there are constants
$C_m,C'_m>0$ and $A_m,A'_m\geq0$ such that, for every
positive integer $r$,
\[
 \left|\operatorname{Tr}\left(F^r\mid
             \ker N_{\wedge^m V}\right)\right|
       \leq C_m r^{A_m}q^{rmw/2},
 \quad
 \left|\operatorname{Tr}\left(F^r\mid
             \ker N_{\wedge^m V^\vee}\right)\right|
       \leq C'_m r^{A'_m}q^{-rmw/2}.
\]
\end{enumerate}
When these conditions hold, condition (3) holds with $A_m=A'_m=0$
and the respective dimensions in place of $C_m,C'_m$ (or any
positive constants when a space is zero).

\begin{proof}
We first justify replacing $F$ by its semisimple part $F_s$.
If $V_\lambda$ is the generalized $\lambda$-eigenspace of $F$,
the covariance identity implies
$N(V_\lambda)\subset V_{\lambda/q}$.
Since $F_s$ is scalar $\lambda$ on $V_\lambda$, it follows
that $F_sNF_s^{-1}=q^{-1}N$. The unipotent part
$F_u=F_s^{-1}F$ commutes with $N$ and $F_s$.
The same statements hold on exterior powers and duals.
Their kernels of $N$ and all monodromy-graded spaces are
preserved by both parts. On any such subquotient, multiplying
the semisimple operator by its commuting unipotent part
does not change its eigenvalues or the traces of its powers.
Thus all three conditions are unchanged by this replacement.

There is a decomposition into $F_s,N$-stable Jordan strings.
Here is a direct proof, including the invariant complement.
Let $l$ be the largest Jordan length of $N$ on a nonzero
remaining summand. Choose an $F_s$-eigenvector $v$ of
eigenvalue $\mu$ with $N^{l-1}v\ne0$, and put
$e_j=N^jv$ for $0\leq j<l$.
Then $F_se_j=\mu q^{-j}e_j$. Choose an eigenfunctional
$a$ of $F_s$ with
$a(F_sx)=\mu q^{-(l-1)}a(x)$ and $a(e_{l-1})=1$.
Such a functional exists by semisimplicity. It vanishes
on $e_j$ for $j<l-1$, because those eigenvalues are distinct.
The map
\[
 P(x)=\sum_{j=0}^{l-1}a(N^{l-1-j}x)e_j
\]
is the identity on the span $W$ of the $e_j$.
The equality $N^l=0$ shows $PN=NP$, including the endpoint
terms. Also $N^rF_s=q^rF_sN^r$, so the $j$th coefficient
of $P(F_sx)$ is $\mu q^{-j}$ times that of $P(x)$.
Therefore $PF_s=F_sP$. Its kernel is an invariant
complement. Induction on dimension gives the asserted
decomposition.

For a string of length $l$, assign
$h(e_j)=l-1-2j$. This is its monodromy grading:
$N$ lowers it by two, and $N^a$ is an isomorphism from
grade $a$ to grade $-a$ whenever those grades occur.
Thus condition (1) on this string says
\[
       |\mu|=q^{(w+l-1)/2}.                         \tag{WM.1}
\]
The monodromy filtration of the direct sum is the sum
of these filtrations, by its defining properties.

Assume condition (2). The determinant line
$\bigwedge^lW$ is contained in $\ker N$ in
$\bigwedge^lV$, and its eigenvalue is
$\delta=\mu^l q^{-l(l-1)/2}$.
The dual determinant line is an invariant line in
$\bigwedge^lV^\vee$, because the direct-sum decomposition
also splits the dual representation. The two bounds with
$m=l\leq d$ give
$|\delta|\leq q^{lw/2}$ and
$|\delta|^{-1}\leq q^{-lw/2}$, hence equality.
Solving it gives (WM.1) for every string, proving (1).

Assume condition (1). On each string define
\[
 H e_j=(l-1-2j)e_j,\qquad
 E e_j=j(l-j)e_{j-1},\quad E e_0=0.
\]
Direct multiplication gives $[E,N]=H$.
Make the strings orthogonal, and choose their basis norms
recursively by
\[
 \|e_{j+1}\|^2=(j+1)(l-j-1)\|e_j\|^2
       \quad(0\leq j<l-1).
\]
All these factors are positive, and these norms give
$N^*=E$. On $\bigwedge^mV$ with its induced positive
Hermitian form, the induced derivations again satisfy
$N^*=E$ and $[E,N]=H$.
If $x$ is in $\ker N$ and has $H$-weight $h$, then
\[
 h\|x\|^2=\langle [E,N]x,x\rangle=-\|Ex\|^2\leq0.
\]
Thus $h\leq0$ for every nonzero such vector.
The operators $F_s$ and $H$ commute and are semisimple.
By (WM.1), on $\bigwedge^mV$ each simultaneous
$F_s,H$-eigenvector of $H$-weight $h$ has Frobenius
absolute value $q^{(mw+h)/2}$: multiply the eigenvalue
absolute values of the chosen wedge basis.
The kernel of $N$ is stable under $F_s$ and $H$, so it
has a simultaneous eigenbasis. Its eigenvalues therefore
obey the first bound of (2). The dual strings have the
same lengths and central weight $-w$. Applying the
same argument to them gives the second bound.

Condition (2) implies (3), with no power of $r$,
by summing the eigenvalues of $F^r$ with their
algebraic multiplicities. For the converse we give
the required elementary power-sum argument. For a
finite-dimensional invertible operator with distinct
eigenvalues $\alpha$ and multiplicities $b_\alpha>0$,
\[
 \sum_{r\geq1}\operatorname{Tr}(F^r)z^r
       =\sum_\alpha\frac{b_\alpha\alpha z}{1-\alpha z}.
\]
A bound $|\operatorname{Tr}(F^r)|\leq Cr^AR^r$
makes the left side holomorphic for $|z|<R^{-1}$.
If $|\alpha|>R$, the right side has a pole at
$z=\alpha^{-1}$ inside that disc. The other terms
are holomorphic there, and $b_\alpha\ne0$, so the
pole cannot cancel. The identity initially holds
in a smaller disc and hence contradicts analytic
continuation. Therefore $|\alpha|\leq R$.
Apply this to the two invariant spaces in (3),
with $R=q^{mw/2}$ and $R=q^{-mw/2}$ respectively.
Zero-dimensional invariant spaces require no argument.
This proves (2) and completes all implications.
\end{proof}

\paragraph{Geometric use and the remaining gate.}
For an $\ell$-adic representation with unipotent inertia,
inertia invariants equal $\ker N$. Apply BR-WM-02
after every complex embedding of the algebraic Frobenius
eigenvalues. For $V=H^i(X_{\overline K},\mathbb Q_\ell)$,
take $w=i$ and $d=\dim V$. Weight--monodromy for this
specified $V$ is therefore equivalent to the two
families of trace bounds in (3), for only
$1\leq m\leq d$. A finite extension of $K$ first
makes the inertia unipotent; the weight assertion
is invariant under that extension.

This is a strictly finite list of representations,
although the trace bounds still quantify over all
positive Frobenius powers $r$. It uses actual
inertia invariants and both dual directions. Bounds
only for alternating total traces, only for
$\det(1-TF\mid V)$, or only in one dual direction
are not substituted for these hypotheses.
In a geometric realization, exterior projectors
place these representations inside tensor powers
of the cohomology, so a prospective positivity
argument has explicit targets. The adopted
Dirichlet zero-location theorem alone does not
establish these local trace bounds.

% END root_weight_criterion.tex

% BEGIN consequences_criteria_WM_formality.tex
\section{Weight--monodromy, Hyodo--Kato cohomology, and rigid analytic
formality}\label{weightmonodromy-hyodokato-cohomology-and-rigid-analytic-formality}

\textbf{Scope and provenance.} Literature cutoff: 20 November 2026.
Checked on 21 November 2026 UTC by consequences\_criteria, GPT-6 Astra
Pro; reasoning effort is not exposed. The adopted RH and Dirichlet-GRH
inputs are not needed for these local geometric results. The sources
below are additional geometric inputs. The 2026 preprints are used as
cited theorems, with their statements and indicated dependency arguments
inspected; their complete foundational proofs have not been
independently reconstructed.

\subsection{CC-WM01 --- Exact known classes of
weight--monodromy}\label{cc-wm01-exact-known-classes-of-weightmonodromy}

Fix the convention that geometric Frobenius \(F\) and monodromy satisfy
\(FNF^{-1}=q^{-1}N\), where \(q\) is the residue-field cardinality.
Monodromy purity of weight \(i\) means that eigenvalues on
\(\operatorname{gr}^M_j\) are algebraic and have absolute value
\(q^{(i+j)/2}\) under every complex embedding. The monodromy filtration
is centered at zero.

\begin{longtable}[]{@{}
  >{\raggedright\arraybackslash}p{(\columnwidth - 4\tabcolsep) * \real{0.3333}}
  >{\raggedright\arraybackslash}p{(\columnwidth - 4\tabcolsep) * \real{0.3333}}
  >{\raggedright\arraybackslash}p{(\columnwidth - 4\tabcolsep) * \real{0.3333}}@{}}
\toprule\noalign{}
\begin{minipage}[b]{\linewidth}\raggedright
Primary result and inspected locator
\end{minipage} & \begin{minipage}[b]{\linewidth}\raggedright
Exact relevant hypotheses
\end{minipage} & \begin{minipage}[b]{\linewidth}\raggedright
Conclusion and boundary
\end{minipage} \\
\midrule\noalign{}
\endhead
\bottomrule\noalign{}
\endlastfoot
Ito and Zou (2026), arXiv:2610.07890v1, \textbf{Theorem B}, p.~2;
\textbf{Theorem 4.4}, pp.~12--13 & \(K/\mathbf Q_p\) finite; \(X/K\)
hyper-Kähler: smooth projective, geometrically simply connected, with
\(H^0(X,\Omega^2_{X/K})=K\sigma\) for a nowhere-degenerate two-form;
definition, p.~1 & Weight--monodromy for every \(0\le i\le2\dim X\) and
every prime \(\ell\), including \(p\), via the associated Weil--Deligne
representation. \textbf{No \(b_2\ge4\) hypothesis occurs in this
theorem}; that restriction concerns separate compatibility results. \\
Scholze (2012), arXiv:1111.4914v1, \textbf{Theorem 9.6}, p.~49;
Conjecture 9.3, p.~48 & \(k\) a characteristic-zero local field; \(Y/k\)
geometrically connected, proper and smooth; \(Y\) a
\textbf{set-theoretic} complete intersection in a projective smooth
toric variety & Weight--monodromy in all degrees for \(\ell\ne p\). \\
Binda, Kato, and Vezzani (2025), arXiv:2207.00369v2, \textbf{Theorem
1.2}, p.~2; \textbf{Theorem 5.6}, p.~28 & \(K/\mathbf Q_p\) finite;
\(Y/K\) smooth and a \textbf{scheme-theoretic} complete intersection in
a projective smooth toric variety & The p-adic conjecture in every
degree. Theorem 5.6 more generally treats such a variety over
\(\mathbf C_p\), asserting quasi-purity of
\(H^i_{\mathrm{HK}}(Y^{\mathrm{an}})\). The
set-theoretic/scheme-theoretic distinction is explicit on p.~2. \\
Wear (2023), arXiv:2303.05610v1, \textbf{Theorem 1.1.2}, p.~3;
coefficient convention, pp.~1--2 & \(k\) a local field of residue
characteristic \(p\); \(A/k\) abelian; \(Y\) geometrically connected,
proper and smooth, and a \textbf{set-theoretic} complete intersection in
\(A\) & Weight--monodromy in all degrees for \(\ell\ne p\). This permits
equal or mixed characteristic and does not provide the p-adic case. \\
\end{longtable}

Scholze's \textbf{Theorem 9.4}, p.~48, states Deligne's
equal-characteristic result for \(k\) the local field at
\(x\in C(\mathbf F_q)\), where \(C/\mathbf F_q\) is a curve and the
variety comes from a proper smooth scheme over \(C\setminus\{x\}\). This
exactly scoped globally defined case suffices in his reduction.

The inspected Ito--Zou dependency chain is \textbf{Lemma 3.8}, p.~8
(purity in degrees at most two), \textbf{Lemma 4.2}, p.~12 (the
higher-degree tensor square is a summand of tensors of \(H^2\) and its
dual), and \textbf{Lemma 4.3}, p.~12 (tensor-square detection), followed
by Theorem 4.4. The complete last argument was read; its earlier
semisimplicity and Mumford--Tate inputs remain cited results.

In a semistable Hyodo--Kato model with \(q=p^a\), use the linear
operator \(\Phi^a\), rather than semilinear \(\Phi\), for traces. The
relation \(N\Phi=p\Phi N\) yields \(\Phi^aN(\Phi^a)^{-1}=q^{-1}N\).
Binda--Kato--Vezzani's p.~2 comparison discussion identifies the
pertinent potentially semistable p-adic realization; their Theorem 5.6's
final splitting argument was read, retaining the preceding motivic
comparison theorems as inputs.

\subsection{CC-WM02 --- The cochain-algebra
consequence}\label{cc-wm02-the-cochain-algebra-consequence}

Petrov and Zavyalov (2026), arXiv:2607.14517v2, \textbf{Theorem 1.2},
p.~2, is \textbf{Theorem 5.1 and Lemma 5.2}, p.~18. For \(X/K\) smooth
proper rigid analytic, \(K/\mathbf Q_p\) finite, weight--monodromy in
all degrees for a fixed prime \(\ell\), possibly \(p\), implies

\[
 R\Gamma_{\mathrm{et}}(X_{\mathbf C_p},\mathbf Q_\ell)
 \simeq \bigoplus_i H^i_{\mathrm{et}}(X_{\mathbf C_p},\mathbf Q_\ell)[-i]
\]

as underlying \(E_\infty\)-algebras. For \(\ell=p\), de Rham cochains
are formal too. Definitions 5.5--5.6, pp.~19--20, specify monodromy
purity; Examples 5.4, p.~19, use \(D_{\mathrm{pst}}\) at \(p\). No
canonical or Galois-equivariant splitting is asserted. The final proof
is on pp.~25--26.

Smooth proper \textbf{algebraic} varieties already have formal cochains
unconditionally, as the source explains on p.~2. Thus the substantive
additional application is to genuinely rigid analytic spaces. Its
\textbf{Construction 3.1 and Proposition 3.2}, p.~5, provide a smooth
proper rigid surface with nonformal cochains at every prime and in de
Rham cohomology. Smoothness and properness alone therefore cannot give
the needed hypothesis.

\subsection{CC-WM03 --- A finite family of exterior trace conditions
implies rigid
formality}\label{cc-wm03-a-finite-family-of-exterior-trace-conditions-implies-rigid-formality}

\textbf{Proposition.} Let \(K/\mathbf Q_p\) be finite and \(X/K\)
smooth, proper and rigid analytic. Fix a prime \(\ell\), including the
possibility \(\ell=p\). For each nonzero cohomology group let

\[
 V_i=D\bigl(H^i_{\mathrm{et}}(X_{\mathbf C_p},\mathbf Q_\ell)\bigr),
 \qquad d_i=\dim V_i,
\]

be the Weil--Deligne representation in Petrov--Zavyalov's Examples 5.4.
Choose a finite extension \(L/K\) killing the finite inertial actions in
these finitely many representations. Put \(q=\#k_L\), and write
\((F_i,N_i)\) for their Frobenius and monodromy over \(L\). Extend
scalars to an algebraic closure of the coefficient field and fix an
embedding of \(\overline{\mathbf Q}\) into that closure. Suppose that
all \(F_i\)-eigenvalues lie in its image, so they can be regarded as
algebraic numbers. For \(1\le m\le d_i\), define

\[
 T_{i,m}(r)=\operatorname{Tr}\left(F_i^r\mid
       \ker N_{\wedge^m V_i}\right),\qquad
 T^\vee_{i,m}(r)=\operatorname{Tr}\left((F_i^\vee)^r\mid
       \ker N_{\wedge^m V_i^\vee}\right).
\]

Here \(F_i^\vee\) denotes inverse-transpose action, and \(N\) acts as a
derivation on exterior powers, with \(N^\vee=-{}^tN\) on the dual. For
every \(i,m\) and embedding
\(\sigma:\overline{\mathbf Q}\hookrightarrow\mathbf C\), suppose that
positive constants \(C_{i,m,\sigma},C^\vee_{i,m,\sigma}\) and
nonnegative exponents \(A_{i,m,\sigma},A^\vee_{i,m,\sigma}\) exist such
that, for \textbf{every integer \(r\ge1\)},

\[
 \begin{split}
 |\sigma(T_{i,m}(r))|&\le C_{i,m,\sigma}r^{A_{i,m,\sigma}}q^{rmi/2},\\
 |\sigma(T^\vee_{i,m}(r))|&\le C^\vee_{i,m,\sigma}r^{A^\vee_{i,m,\sigma}}q^{-rmi/2}.
 \end{split} \tag{CC-WM.1}
\]

Then weight--monodromy holds in every degree over \(K\), and the
underlying \(E_\infty\)-algebra
\(R\Gamma_{\mathrm{et}}(X_{\mathbf C_p},\mathbf Q_\ell)\) is formal. For
\(\ell=p\), \(R\Gamma_{\mathrm{dR}}(X/K)\) is formal as well.
Conversely, weight--monodromy implies algebraicity and (CC-WM.1), with
exponents zero and constants equal to the respective kernel dimensions
when these dimensions are positive.

\textbf{Complete proof of the deduction.} There are finitely many
nonzero finite-dimensional cohomology groups. Intersect their open
inertia kernels and pass to a finite extension to obtain \(L\). Each
pair satisfies \(F_iN_iF_i^{-1}=q^{-1}N_i\).

Fix \(i\). Replacing \(F_i\) by its semisimple part leaves all relevant
traces and monodromy-graded eigenvalues unchanged, by the first part of
BR-WM-02. The same proof decomposes the pair into invariant Jordan
strings, with bases \(e_0,\ldots,e_{s-1}\) such that \(Ne_j=e_{j+1}\),
\(Ne_{s-1}=0\), and \(Fe_j=\mu q^{-j}e_j\). Since every \(\mu\) is
algebraic, these matrices define a pair over a number field. Thus every
\(\sigma\) can be applied to the pair, without requiring an algebraic
model for the full original Weil action. Exterior powers, duals and
kernels commute with this characteristic-zero field extension; their
power traces are precisely the algebraic numbers in (CC-WM.1).

Apply BR-WM-02 to this complex pair with dimension \(d_i\) and weight
\(w=i\). Its trace criterion gives

\[
 |\sigma(\alpha)|=q^{(i+j)/2}
 \quad\text{for every eigenvalue }\alpha
 \text{ on }\operatorname{gr}^M_jV_i.
\]

Doing this for every \(\sigma\) proves monodromy purity over \(L\),
including algebraicity. This is the hypothesis in Petrov--Zavyalov's
Definition 5.6.

For completeness, purity descends from \(L\) to \(K\). Write
\(q_0=\#k_K\). For any chosen geometric Frobenius \(\phi\) over \(K\),
some positive power \(\phi^f\) belongs to the open subgroup \(W_L\);
replacing \(L\) by a further finite extension if necessary does not
affect purity, since its Frobenius is a positive power after the finite
inertial action has been killed. The element \(\phi^f\) has degree \(f\)
over \(K\), so its eigenvalues have absolute value \(q_0^{f(i+j)/2}\) on
\(\operatorname{gr}^M_j\). Indeed, over \(L\) the finite inertial action
is trivial and its degree is \(f/[k_L:k_K]\); hence it acts by that
power of an \(L\)-Frobenius. If \(\alpha\) is a \(\phi\)-eigenvalue,
\(\alpha^f\) is such an eigenvalue. Algebraicity descends from
\(\alpha^f\) by the equation \(t^f-\alpha^f=0\), and every complex
absolute value of \(\alpha\) is \(q_0^{(i+j)/2}\). Scaling \(N\) under a
change of monodromy normalization does not change its filtration. This
proves descent; ascent follows by the same power calculation after
killing finite inertia.

Petrov--Zavyalov's Theorem 5.1 now gives étale formality. At \(p\),
Lemma 5.2 gives de Rham formality. The \(p\)-adic pair is obtained from
\(D_{\mathrm{pst}}\); when a semistable Hyodo--Kato description is
available, use its linear Frobenius and \(N\). No equality with ordinary
\(p\)-adic Galois inertia invariants is needed.

Conversely, monodromy purity gives algebraicity and BR-WM-02's first
condition under every complex embedding. That theorem bounds eigenvalues
on the two invariant exterior spaces by \(q^{mi/2}\) and \(q^{-mi/2}\).
A power trace is the sum of the eigenvalue powers, with algebraic
multiplicities, also for a nonsemisimple Frobenius. The triangle
inequality proves (CC-WM.1) with zero exponents and the stated
constants. A zero cohomology group requires no conditions; a zero kernel
admits any positive constant. Every assertion follows. \(\square\)

\textbf{Exact gain and limit.} There are at most \(2\sum_i d_i\) tested
representations; the quantifier over Frobenius powers remains infinite.
This is a criterion and its formality consequence, not a proof of the
bounds for arbitrary \(X\). For \(\ell\ne p\), after unipotent Galois
inertia is achieved, the kernels are actual inertia invariants of the
exterior representations. At \(p\), they belong to the Weil--Deligne or
Hyodo--Kato realization. Alternating traces on total cohomology do not
supply the separate inequalities. Both dual bounds and all complex
embeddings are required. The adopted Dirichlet GRH does not, by itself,
establish these geometric bounds.

\subsection{Exact primary bibliography and version
record}\label{exact-primary-bibliography-and-version-record}

\begin{itemize}
\tightlist
\item
  Binda, F., Kato, H., \& Vezzani, A. (2025). On the \(p\)-adic
  weight-monodromy conjecture for complete intersections in toric
  varieties. \emph{Inventiones Mathematicae, 241}, 559--603.
  \url{https://doi.org/10.1007/s00222-025-01344-x}. Inspected version:
  \url{https://arxiv.org/pdf/2207.00369v2,} dated 10 June 2025, 36 pages;
  locators above refer to this version. Published online 17 June 2025.
\item
  Ito, K., \& Zou, H. (2026). \emph{Semisimplicity, purity and
  Mumford-Tate conjecture for hyper-Kähler varieties}
  (arXiv:2610.07890v1) {[}Preprint{]}.
  \url{https://arxiv.org/abs/2610.07890v1}. Submitted 6 October 2026, 07:38:56
  UTC. Inspected PDF: \url{https://arxiv.org/pdf/2610.07890v1,} 23 pages.
\item
  Petrov, A., \& Zavyalov, B. (2026). \emph{Formality for rigid-analytic
  spaces satisfying the weight-monodromy conjecture}
  (arXiv:2607.14517v2) {[}Preprint{]}.
  \url{https://arxiv.org/abs/2607.14517v2}. First submitted 16 July 2026; v2
  revised 17 August 2026, 21:05:01 UTC. Inspected PDF:
  \url{https://arxiv.org/pdf/2607.14517v2,} 28 pages.
\item
  Scholze, P. (2012). Perfectoid spaces. \emph{Publications
  Mathématiques de l'IHÉS, 116}, 245--313.
  \url{https://doi.org/10.1007/s10240-012-0042-x}. Inspected author version:
  \url{https://arxiv.org/pdf/1111.4914v1,} submitted 21 November 2011,
  16:14:23 UTC, 51 pages; numbered locators above refer to that author
  version.
\item
  Wear, P. (2023). \emph{Perfectoid covers of abelian varieties and the
  weight-monodromy conjecture} (arXiv:2303.05610v1) {[}Preprint{]}.
  \url{https://arxiv.org/abs/2303.05610v1}. Submitted 9 March 2023, 22:41:32
  UTC. Inspected PDF: \url{https://arxiv.org/pdf/2303.05610v1,} 54 pages.
\end{itemize}

% END consequences_criteria_WM_formality.tex

% BEGIN consequences_criteria_stable_curve.tex
\section{Universal stable curves: exact stack purity and
decomposition}\label{universal-stable-curves-exact-stack-purity-and-decomposition}

Research module \nolinkurl{consequences_criteria}; GPT-6 Astra Pro; effort
setting not exposed; 20 November 2026. This module uses established
finite-field weight theory independently of the adopted Alpöge RH and
Dirichlet GRH. New deductions below are supplied with complete proofs.
No claim is made about independent reconstruction of the foundational
BBD or stack six-functor theorems.

\subsection{CC-ST01. Verified stack theorem and its
application}\label{cc-st01.-verified-stack-theorem-and-its-application}

\textbf{Primary theorem.} Sun (2012), arXiv:1009.4398v3, Theorem 1.2,
p.~3, treats proper morphisms of finite diagonal between finite-type
stacks over a finite field with affine stabilizers. For a pure input, it
gives a decomposition into shifted semisimple perverse cohomology
\textbf{after extension to the algebraic closure}. The convention
removing the subscript 0 is explicit in §1.3.3, p.~3. The weight
calculation uses Proposition 3.9(iii), p.~8, and Theorem 3.5, p.~6;
geometric semisimplicity and splitting are Theorems 3.11--3.12, p.~8,
with the application proved in §3.13, p.~9. These are arXiv printed-page
locators, not guessed journal-page translations.
\href{https://arxiv.org/pdf/1009.4398v3}{Primary PDF}.

\textbf{Geometric hypotheses.} Let \(g,n\geq0\) satisfy \(2g-2+n>0\),
let \(q\) be a power of \(p\), and let \(\ell\ne p\). Set \[
Y_0=\overline{\mathcal M}_{g,n,\mathbf F_q},\qquad
X_0=\overline{\mathcal M}_{g,n+1,\mathbf F_q},\qquad
\pi_0:X_0\longrightarrow Y_0,
\] where the last marking is forgotten and the curve stabilized. Put
\(d=3g-3+n\) and \(D=d+1\). Knudsen (1983), Corollary 2.6 and Theorem
2.7, p.~179, identify contraction and stabilization as inverse
constructions and give the smooth proper moduli stacks with
normal-crossing boundary. Knudsen (2012), introduction, pp.~1--2,
explicitly verifies that \(\pi_0\) is representable and is the universal
stable curve, hence is proper and flat with geometrically connected
nodal curve fibres. The local expression \(t_i\mapsto xy\) is also
printed in the proof of Knudsen's Theorem 2.7 on p.~179.
\href{https://journals.msp.org/mscand/article/download/1622/1621/1653}{Original
journal PDF}; \href{https://arxiv.org/pdf/1106.1588v2}{author's later
clarification}.

With \(\Lambda=\overline{\mathbf Q}_\ell\) and a fixed embedding
\(\iota:\Lambda\hookrightarrow\mathbf C\), define \[
K_0=\Lambda_{X_0}[D]=\mathrm{IC}_{X_0},\quad
L_0=R\pi_{0*}K_0,\quad
P_{i,0}={}^p\!H^i(L_0).
\] Then \(L_0\) is pure of weight \(D\), and \(P_{i,0}\) is pure of
weight \(D+i\). After base change to \(\overline{\mathbf F}_q\), there
is a noncanonical isomorphism \[
L\simeq\bigoplus_{i\in\mathbf Z}P_i[-i],
\] and every \(P_i\) is a semisimple perverse sheaf, hence a finite
direct sum of intersection complexes of irreducible lisse sheaves on
smooth locally closed substacks.

\textbf{Complete application proof.} Stability makes all geometric
stabilizers finite, so they are affine. The moduli stacks are of finite
type. Since \(\pi_0\) is representable and proper, its relative diagonal
is a closed immersion and therefore finite. Thus the hypotheses on the
morphism in Sun's theorem hold; arbitrary Artin stacks would not
suffice.

The dimension formula may be checked on the dense smooth-curve locus.
The tangent space of a smooth pointed curve is
\(H^1(C,T_C(-\sum p_i))\). Stability gives \(H^0(C,T_C(-\sum p_i))=0\),
and Riemann--Roch gives its dimension \(3g-3+n\). The normal-crossing
boundary has empty interior, and the stack is smooth, so this is its
dimension on every component. The total space has dimension \(D=d+1\).

The constant sheaf on the smooth stack \(X_0\) has weight zero. With
cohomological shifts, \(K_0\) is perverse and pure of weight \(D\);
equivalently, its unique ordinary cohomology sheaf sits in degree \(-D\)
and its Verdier dual is \(K_0(D)\). Properness gives
\(R\pi_{0!}=R\pi_{0*}\) in this finite-diagonal setting. The upper and
lower weight bounds of Sun's Proposition 3.9(iii) therefore give both
inequalities needed for purity of \(L_0\). His Theorem 3.5 gives the
weights \(D+i\). Finally, Theorems 3.11--3.12 give the asserted
geometric decomposition and semisimplicity. The description of simple
perverse objects is precisely the one used in Sun's Theorem 3.11, citing
Laszlo--Olsson, Theorem 8.2(ii). This proves the stated application.
\(\square\)

\textbf{Scope of the conclusion.} A pure complex need not have every
ordinary stalk cohomology group pure of its cohomological degree. In
particular, the statement about \(R\pi_{0*}\Lambda\) being pure of
weight zero does not assert that the \(H^1\) of every nodal fibre is
pure of weight one. The latter generally has both graph and
normalization contributions. Neither arithmetic semisimplicity nor a
Frobenius-equivariant choice of the displayed splitting follows merely
by retaining Sun's conclusion over \(\mathbf F_q\); the theorem
explicitly takes the geometric base change. Decomposition is not a
smooth-local property that can simply be descended from scheme charts:
Sun gives a counterexample in his introduction, pp.~1--2.

\subsection{CC-ST02. Three perverse degrees and their endpoint
summands}\label{cc-st02.-three-perverse-degrees-and-their-endpoint-summands}

In the preceding situation, the only possibly nonzero \(P_{i,0}\) have
\(i=-1,0,1\). Their weights are respectively \(d,d+1,d+2\).
Geometrically, the endpoints are \[
P_{-1}\simeq\Lambda_Y[d],\qquad
P_1\simeq\Lambda_Y[d](-1).
\] Consequently the constant-coefficient direct image admits a
noncanonical geometric decomposition \[
R\pi_*\Lambda_X\simeq
\Lambda_Y\;\oplus\;P_0[-d-1]\;\oplus\;\Lambda_Y(-1)[-2].
\tag{CC-ST.1}
\] If \(U=\mathcal M_{g,n,\overline{\mathbf F}_q}\subset Y\) is the
smooth-curve locus, then \[
P_0|_U\simeq (R^1\pi_*\Lambda)|_U[d].
\] The middle perverse sheaf can include boundary-supported summands;
equation (CC-ST.1) does not discard them.

\textbf{Complete proof.} Proper base change and the cohomological
dimension of a proper curve give \(R^j\pi_*\Lambda=0\) for
\(j\notin\{0,1,2\}\). Thus the ordinary cohomology of \(L\) is zero
outside degrees \(-d-1,-d,1-d\). For each of these degrees \(k\), the
support has dimension at most \(d\leq1-k\). The support criterion for
the middle perverse \(t\)-structure therefore gives
\(L\in{}^pD^{\leq1}\). The criterion is checked on an étale atlas;
relative dimension zero means no additional stack shift occurs.

Because \(X\) is smooth of dimension \(D\), Verdier duality gives
\(\mathbb D_XK=K(D)\). Proper duality then gives \(\mathbb D_YL=L(D)\).
Perverse duality interchanges \({}^pD^{\leq1}\) and \({}^pD^{\geq-1}\),
and Tate twisting preserves these categories, so \(L\in{}^pD^{\geq-1}\).
Hence \(P_i=0\) for \(|i|>1\). The same vanishing holds before geometric
base change, since that pullback is conservative and perverse exact. The
weight assertions now follow from CC-ST01.

To identify \(P_1\), use its geometric semisimplicity. A simple
constituent with support of dimension \(s\) restricts on a dense smooth
part of that support to a nonzero lisse sheaf shifted by \(s\). Its
contribution through the direct summand \(P_1[-1]\subset L\) therefore
has nonzero ordinary cohomology in degree \(1-s\). But \(L\) has no
ordinary cohomology above \(1-d\). Consequently \(1-s\leq1-d\), so
\(s=d\): every simple constituent of \(P_1\) has full support on a
component of \(Y\).

Over \(U\), smooth proper base change identifies \(P_1|_U\) with
\(R^2\pi_*\Lambda[d]|_U=\Lambda_U[d](-1)\), since the fibres are
geometrically connected smooth proper curves. A semisimple perverse
sheaf whose constituents all have full support is the intermediate
extension of its restriction to a suitable dense open. Since \(Y\) is
smooth and the local system just obtained is constant, this intermediate
extension is \(\Lambda_Y[d](-1)\), component by component. This gives
\(P_1\).

Perverse duality also gives \(\mathbb D_YP_1=P_{-1}(D)\). On the smooth
\(d\)-dimensional stack \(Y\), \[
\mathbb D_Y(\Lambda_Y[d](-1))=\Lambda_Y[d](d+1)=\Lambda_Y[d](D),
\] so \(P_{-1}=\Lambda_Y[d]\). Substitute these two identifications in
\(L\simeq P_{-1}[1]\oplus P_0\oplus P_1[-1]\) and shift by \(-D\) to
obtain (CC-ST.1). The same calculation of perverse degree on \(U\)
identifies the middle term with \(R^1\pi_*\Lambda[d]\). Every asserted
step has now been proved. \(\square\)

\subsection{CC-ST03. Why geometric semisimplicity does not imply
arithmetic
semisimplicity}\label{cc-st03.-why-geometric-semisimplicity-does-not-imply-arithmetic-semisimplicity}

Over \(\operatorname{Spec}\mathbf F_q\), let Frobenius act on
\(\Lambda^2\) by \[
F=\begin{pmatrix}1&1\\0&1\end{pmatrix}.
\] This is a continuous \(\ell\)-adic representation: the homomorphism
\(\widehat{\mathbf Z}\to\mathbf Z_\ell\) followed by
\(a\mapsto\left(\begin{smallmatrix}1&a\\0&1\end{smallmatrix}\right)\)
provides it, with the generator chosen as geometric Frobenius. All
eigenvalues are one, so the associated sheaf is punctually pure of
weight zero. It is not arithmetically semisimple, since its only
Frobenius-fixed line has no invariant complement. After pullback to
\(\operatorname{Spec}\overline{\mathbf F}_q\), the Galois action
disappears and the sheaf is a sum of two constant rank-one sheaves. This
fully explicit example shows why purity alone does not license the
stronger arithmetic conclusion. It does not assert that the particular
universal-curve decomposition has no arithmetic refinement.

\subsection{CC-ST04. Complete geometric support formula, including the
separating
boundary}\label{cc-st04.-complete-geometric-support-formula-including-the-separating-boundary}

Write \(j:U=\mathcal M_{g,n,\overline{\mathbf F}_q}\hookrightarrow Y\)
and \(\mathbb V=(R^1\pi_*\Lambda)|_U\). For each irreducible separating
boundary divisor \(D\subset Y\), let \(D^\circ\) be its open locus of
curves with exactly one node. Let \(\mathcal C_D\to D^\circ\) be the
finite étale two-element cover that chooses one of the two irreducible
components of the fibre. Define the rank-one local system \[
\epsilon_D=\Lambda[\mathcal C_D]/\Lambda\mathbf 1.
\] Here \(\Lambda\mathbf1\) is the diagonal constant sub-local system of
the component permutation local system. For a locally closed immersion
\(k_D:D^\circ\hookrightarrow Y\), the notation \((k_D)_{!*}\) means
intermediate extension to its closure followed by the closed pushforward
to \(Y\).

There is a geometric isomorphism of semisimple perverse sheaves \[
P_0\simeq j_{!*}(\mathbb V[d])\;\oplus\!
\bigoplus_{D\text{ separating}}(k_D)_{!*}\bigl(\epsilon_D[d-1]\bigr)(-1).
\tag{CC-ST.2}
\] Thus substituting (CC-ST.2) into (CC-ST.1) determines every perverse
support and its coefficient system. There are no simple constituents
whose support has codimension at least two, and no constituent supported
on a nonseparating boundary divisor. Intersection complexes in the
formula can nevertheless have nonzero stalks on deeper boundary strata.

The local system \(\epsilon_D\) is the sign of the permutation of the
two components. If those components are distinguished by their genera or
by the labelled markings, it is constant. In the standard labelled
moduli problem, the only boundary type allowing exchange is \(n=0\),
\(g\) even, with genera \(g/2\) and \(g/2\); there one must retain the
component-exchange local system on the stack. A twist of weight two is
essential: the coefficient system on \(D^\circ\) is \(\epsilon_D(-1)\),
so its intermediate extension has perverse weight \(2+(d-1)=d+1\).

\textbf{Complete proof.} A simple constituent of \(P_0\) supported in
dimension \(s\) contributes ordinary cohomology to \(L\) in degree
\(-s\) on a dense smooth part of its support. The upper ordinary degree
bound \(1-d\) from CC-ST02 forces \(-s\leq1-d\), hence \(s\geq d-1\).
Every proper support is therefore a divisor. Since
\(P_0|_U=\mathbb V[d]\) is lisse in one ordinary degree, no such proper
support meets \(U\). The full-support sum is \(j_{!*}(\mathbb V[d])\);
all other possible constituents lie on boundary divisors.

Fix such a divisor and remove its intersections with other boundary
branches and any smaller loci needed to make its restrictions lisse. On
the resulting open set, write \(i:D^\circ\hookrightarrow Y\) for the
closed smooth divisor and \(A=j_{!*}(\mathbb V[d])\). The
intermediate-extension property gives \({}^pH^0(i^*A)=0\). Also
\(i^*A\in{}^pD^{\leq0}\). On the smooth \((d-1)\)-dimensional divisor
the perverse structure on lisse complexes is the ordinary structure
shifted by \(d-1\), so \[
\mathcal H^{1-d}(i^*A)=0.
\] The required intermediate-extension vanishing is stated for stacks in
Laszlo--Olsson (2009), Lemma 6.1, author-preprint p.~10; only this
vanishing implication is used here.

Taking ordinary cohomology in degree \(1-d\) in the geometric
decomposition of \(L\) now gives, on \(D^\circ\), \[
(R^2\pi_*\Lambda)|_{D^\circ}\simeq
\Lambda_{D^\circ}(-1)\oplus E_D,
\tag{CC-ST.3}
\] where \(E_D[d-1]\) is the restriction of the sum of the
divisor-supported constituents of \(P_0\). The first term comes from
\(P_1[-1]\), and \(P_{-1}[1]\) contributes nothing in this degree. No
full-support middle term contributes by the vanishing just proved.

For a connected proper nodal curve, normalization identifies \(H^2\)
with the sum of \(\Lambda(-1)\) over its irreducible components: the
difference between the curve and its normalization is supported at
finitely many nodes and has no cohomology in degree two. This
identification is functorial, so over the one-node locus
\(R^2\pi_*\Lambda\) is exactly the component permutation local system
with a Tate twist. A nonseparating node leaves the curve irreducible,
giving a rank-one system; (CC-ST.3) then forces \(E_D=0\).

A separating node gives precisely two components. Their monodromy
permutation group is either trivial or the order-two group. In the
trivial case, every rank-one quotient by a constant line is constant. In
the order-two case, the constant subrepresentation is the unique
invariant line, and the quotient is the sign representation. Since the
coefficients have characteristic zero, these permutation representations
are semisimple. Equation (CC-ST.3) therefore gives
\(E_D\simeq\epsilon_D(-1)\) in both cases. Intermediate extension
uniquely determines each semisimple divisor-supported sum from this
dense-open restriction, which proves (CC-ST.2).

Finally, a component exchange preserves genus and each labelled marking.
It is impossible when the genera differ or when any labelled marking
lies on either component. If \(n=0\) and the two genera agree, the two
components are unordered and their exchange is exactly the finite
permutation character retained in \(\epsilon_D\). This proves the last
description and completes the proof. \(\square\)

\subsection{References}

Knudsen, F. F. (1983). The projectivity of the moduli space of stable
curves, II: The stacks \(M_{g,n}\). \emph{Mathematica Scandinavica, 52},
161--199. \url{https://doi.org/10.7146/math.scand.a-12001} . Original article,
Corollary 2.6 and Theorem 2.7, p.~179, visually inspected at
\url{https://journals.msp.org/mscand/article/download/1622/1621/1653} .

Knudsen, F. F. (2012). \emph{A closer look at the stacks of stable
pointed curves} (arXiv:1106.1588v2). \url{https://arxiv.org/abs/1106.1588v2} .
Introduction, pp.~1--2, directly inspected. The version inspected
supplies the author's clarification of the contraction/stabilization
construction.

Laszlo, Y., \& Olsson, M. (2009). Perverse \(t\)-structure on Artin
stacks. \emph{Mathematische Zeitschrift, 261}(4), 737--748.
\url{https://doi.org/10.1007/s00209-008-0348-z} . Author preprint, Theorem
8.2(ii), p.~12:
\url{https://www.cmls.polytechnique.fr/perso/laszlo/articleweb/faisceaux-pervers.pdf}
. Statement directly inspected; its whole foundational proof is not
reconstructed here.

Sun, S. (2012). Decomposition theorem for perverse sheaves on Artin
stacks over finite fields. \emph{Duke Mathematical Journal, 161}(12),
2297--2310. \url{https://doi.org/10.1215/00127094-1723657} . Inspected
version: \url{https://arxiv.org/pdf/1009.4398v3} . Exact theorem and page
locators are given above.

% END consequences_criteria_stable_curve.tex

% BEGIN consequences_criteria_RESIDUE.tex
\section{RESIDUE H.3, H.5, H.6: independently established categorical
content}\label{residue-h.3-h.5-h.6-independently-established-categorical-content}

Research module \nolinkurl{consequences_criteria}; GPT-6 Astra Pro; effort
setting not exposed; 20 November 2026. The available file
\texttt{sources/INTAKE\_RESIDUE\_20261120.md} reports H.3, H.5 and H.6
at statement level and explicitly says their proofs and cited sources
were not opened by that intake. The underlying \texttt{sections/H.tex}
is absent from the supplied source tree. The results below therefore
have their own complete proofs and precise models; they are not
retrospective certifications of unseen proofs.

\subsection{CC-RES01. Direct image for the universal curve is right Kan
extension}\label{cc-res01.-direct-image-for-the-universal-curve-is-right-kan-extension}

Let \(g,n\geq0\) with \(2g-2+n>0\), and work with the complex analytic
topological Deligne--Mumford stacks \[
X=\overline{\mathcal M}_{g,n+1}^{\mathrm{top}},\qquad
Y=\overline{\mathcal M}_{g,n}^{\mathrm{top}},\qquad
\pi:X\longrightarrow Y.
\] Stratify both by stable dual graph. Let \(\mathcal E\) be a compactly
generated presentable \(\infty\)-category, such as spaces, spectra, or
\(D(R)\) for an associative ring \(R\). Write \(\Pi_X\) and \(\Pi_Y\)
for the actual exit-path \(\infty\)-categories of these stratified
stacks. Then:

\begin{enumerate}
\def\labelenumi{\arabic{enumi}.}
\tightlist
\item
  The forgetful map induces an exit functor
  \(F=\Pi(\pi):\Pi_X\to\Pi_Y\).
\item
  The sheaf direct image \(\pi_*\) takes graph-constructible
  \(\mathcal E\)-valued sheaves to graph-constructible sheaves.
\item
  Under exodromy equivalences \(E_X,E_Y\), there is a natural
  equivalence \[
  E_Y\pi_*\simeq\operatorname{Ran}_{F}E_X.
  \tag{CC-RES.1}
  \] In particular, \[
  (E_Y\pi_*\mathcal F)(y)\simeq
  \lim_{(y\downarrow F)} E_X\mathcal F.
  \tag{CC-RES.2}
  \] The indexing category consists of pairs \((x,\alpha:y\to F(x))\).
  It is not replaced by the strict fibre of \(F\) without a further
  cofinality theorem. For \(\mathcal E=D(R)\), \(\pi_*\) denotes the
  derived sheaf direct image.
\end{enumerate}

\textbf{Inspected primary inputs.} Jansen (2023), Corollary 6.6 and
Theorem 6.7, p.~34, identify these exit categories with Charney--Lee
categories and give their exodromy equivalences; Theorem 6.8, p.~35,
gives the derived functor-category formulation. Jansen (2024), Theorem
2.17, pp.~10--11, gives proper base change for étale topological stacks;
Proposition 3.24, p.~21, gives compactly generated coefficients. Haine,
Porta, and Teyssier (2026 version), Theorem 3.2.3, pp.~34--35, states
that every morphism between exodromic stratified \(\infty\)-topoi
induces the requisite exit functor; Observation 2.3.6, p.~20, identifies
pullback with precomposition. Their Proposition 5.2.11, p.~53, supplies
the product formula used below. These are explicit locators in
arXiv:2308.09551v2, arXiv:2308.09550v3, and arXiv:2401.12825v2
respectively.

\textbf{Complete proof.} Forgetting the last marked point induces a map
of dual-graph strata: remove its leg and stabilize every now-unstable
rational vertex. Contraction of edges, corresponding to smoothing nodes,
is compatible with this operation because stable contraction of families
commutes with base change. Thus \(\pi\) is a stratified map. Jansen's
theorem says that both stratified stacks are exodromic. The
functoriality theorem of Haine--Porta--Teyssier gives \(F\), and its
pullback compatibility gives \[
E_X\pi^*\simeq F^*E_Y,
\] first for space-valued sheaves and then, by the compactly generated
coefficient construction, for \(\mathcal E\)-valued sheaves.

We prove the constructibility claim, which is necessary to identify the
actual sheaf direct image with the right adjoint inside constructible
sheaves. Fix a dual-graph stratum \(Y_G\). On a sufficiently small
contractible analytic chart \(B\to Y_G\), label the finitely many
components, nodes, branches and old marked points. The family over \(B\)
has constant nodal type. Its normalization is a disjoint union of proper
smooth families of compact surfaces, with disjoint labelled sections for
the old marked points and the two preimages of every node. One can
trivialize these families preserving all those sections: choose a
horizontal distribution tangent to the section images, using disjoint
neighbourhoods of the sections and a partition of unity; lift the radial
paths in a small coordinate ball \(B\). Properness makes the lifts exist
throughout these paths after shrinking the ball. They give smooth
trivializations depending continuously on the endpoint. Gluing each pair
of node sections then gives a fibre-preserving homeomorphism \[
\pi^{-1}(B)\simeq C\times B,
\tag{CC-RES.3}
\] where \(C\) is a fixed compact nodal curve, with its nodes and old
marked points distinguished.

This is a stratified trivialization for a refinement of the induced
graph stratification: an extra point in a component's smooth unmarked
part adds a leg to that vertex; at an old marked point it inserts a
rational tail; at a node it inserts a rational bridge. Each of these
conditions is preserved by the trivialization. Thus a
graph-constructible sheaf on \(X\) restricts to a constructible sheaf
for this finite point-and-component stratification on \(C\times B\).

Both \(C\) with this stratification and the ball \(B\) are conically
stratified. The exit product formula gives \[
\Pi(C\times B)\simeq\Pi(C)\times\Pi(B)\simeq\Pi(C),
\] since \(B\) is contractible. Exodromy therefore shows that every such
constructible sheaf on \(C\times B\) is pulled back from a constructible
sheaf \(\mathcal G\) on \(C\). These conical, finite-dimensional spaces
use the usual constructible sheaf categories; the equivalent hypersheaf
formulation of the product theorem yields the same objects here.

The map \(C\to *\) is proper. Proper base change for the product square
now identifies the direct image to \(B\) of the pullback of
\(\mathcal G\) with the constant sheaf having value
\(\Gamma(C,\mathcal G)\). Proper base change for \(\pi\) along
\(Y_G\hookrightarrow Y\), followed by the chart \(B\to Y_G\), identifies
this with \((\pi_*\mathcal F)|_B\). The base-change theorem holds also
for compactly generated coefficients: express \(\mathcal E\)-valued
sheaves through finite-limit-preserving functors from the opposite of the category of compact objects to
space-valued sheaves; evaluation detects the required equivalence,
reducing it to the space-valued statement. Hence \(\pi_*\mathcal F\) is
locally constant on each graph stratum. This proves (2).

It remains to identify the functor. For every
\(A\in\operatorname{Fun}(\Pi_Y,\mathcal E)\) and
\(\mathcal F\in\operatorname{Shv}_{\mathrm{cbl}}(X;\mathcal E)\),
exodromy and the sheaf adjunction give naturally \[
\begin{aligned}
\operatorname{Map}(A,E_Y\pi_*\mathcal F)
&\simeq\operatorname{Map}(E_Y^{-1}A,\pi_*\mathcal F)\\
&\simeq\operatorname{Map}(\pi^*E_Y^{-1}A,\mathcal F)\\
&\simeq\operatorname{Map}(F^*A,E_X\mathcal F).
\end{aligned}
\] Thus \(E_Y\pi_*E_X^{-1}\) is right adjoint to precomposition by
\(F\). By uniqueness of right adjoints, it is right Kan extension. Its
usual pointwise limit formula gives (CC-RES.2). This proves all
assertions. \(\square\)

\textbf{Relation to the intake.} This independently establishes a
precise geometric version of H.6 for the canonical exit categories.
Jansen identifies \(\Pi_{g,n}\) with
\(\mathrm{SC}_{g,n}^{\mathrm{op}}\); consequently the statement can be
expressed using presheaves on entrance categories, with right Kan
extension along the opposite \(u^{\mathrm{op}}\) of the entrance functor
\(u\). What remains uninspected is the equality of the unseen archive's
chosen, explicitly presented forgetful functor with this canonical
functor. The assertion is complex topological. No Frobenius action, Tate
twist, or finite-field weight filtration is created by exodromy; the
separate finite-field statement is CC-ST01--CC-ST04.

\subsection{CC-RES02. Complete lifting-category calculation for the
nodal local
model}\label{cc-res02.-complete-lifting-category-calculation-for-the-nodal-local-model}

Let \(\mathcal C\) be the ordinary category with objects \(0,1\), with
\[
\operatorname{End}(0)=\{1_0\},\quad
\operatorname{End}(1)=\mathbf Z,\quad
\operatorname{Hom}(0,1)=\{a\},\quad
\operatorname{Hom}(1,0)=\varnothing,
\] where composition in \(\mathbf Z\) is addition and \(m\circ a=a\).
Define \[
F:\mathcal C\times\mathcal C\longrightarrow\mathcal C
\] by \(F(1,1)=1\) and \(F(x)=0\) on the other three objects. On the
endomorphisms of \((1,1)\), set \(F(r,s)=r+s\); all arrows between
objects mapped to 0 go to \(1_0\), and all arrows from those objects to
\((1,1)\) go to \(a\).

For \(x\in\mathcal C^2\) and \(\alpha:F(x)\to y\), define the strict
lifting category \(\mathcal L(x,\alpha)\) as follows. Its objects are
arrows \(\beta:x\to z\) such that \(F(z)=y\) and \(F(\beta)=\alpha\). A
morphism from \(\beta\) to \(\beta'\) is an arrow \(\delta:z\to z'\)
with \(F(\delta)=1_y\) and \(\delta\beta=\beta'\). Then \[
|N\mathcal L(x,\alpha)|\simeq
\begin{cases}
S^1,&x=(0,0),\ \alpha=a:0\to1,\\
*,&\text{all other possible pairs }(x,\alpha),
\end{cases}
\tag{CC-RES.4}
\] In the exceptional case the lifting category itself is the one-object
groupoid \(B\mathbf Z\). In the remaining groupoid cases it is
equivalent to the terminal category; the cases with an initial object
need only have contractible nerve.

\textbf{Complete proof.} If \(F(x)=0\) and \(\alpha=1_0\), the lifting
category is the undercategory of \(x\) in the full subcategory
\(F^{-1}(0)\). The identity arrow of \(x\) is initial, so its nerve is
contractible.

Suppose \(F(x)=0\) and \(\alpha=a\). The target must be \(z=(1,1)\). Its
endomorphisms over \(1_1=0\in\mathbf Z\) form \[
K=\ker(\mathbf Z^2\xrightarrow{+}\mathbf Z)
 =\{(r,-r):r\in\mathbf Z\}\simeq\mathbf Z.
\] For \(x=(0,0)\) there is exactly one arrow \(\beta:x\to(1,1)\), and
every element of \(K\) fixes it. Hence the lifting category is exactly
the one-object groupoid \(BK\).

For \(x=(1,0)\) the possible arrows \(\beta\) are indexed by
\(\mathbf Z\), coming from the first coordinate. The element \((r,-r)\)
acts by translation by \(r\). This action is free and transitive, so its
action groupoid is equivalent to the terminal category. For \(x=(0,1)\)
the same argument uses translation by \(-r\) in the second coordinate.

Finally, if \(F(x)=1\), then \(x=(1,1)\), \(y=1\), and \(\alpha\) is an
arbitrary integer \(m\). Its lifts are the pairs \((r,s)\) with
\(r+s=m\). The group \(K\) acts freely and transitively on this set by
addition, so the lifting category again is equivalent to the terminal
category. No arrow \(1\to0\) exists, so these cases exhaust every
possible \((x,\alpha)\). This proves (CC-RES.4). \(\square\)

\textbf{Geometric interpretation with its precise limit.} The exit
category of a complex disk stratified by its origin and punctured disk
is \(\mathcal C\): the punctured stratum has fundamental group
\(\mathbf Z\) and the origin is the cone's initial object. The product
stratification of a bidisk has exit category \(\mathcal C^2\). The
multiplication map \((u,v)\mapsto uv\) induces the displayed functor
because the winding numbers add on \((\mathbf C^*)^2\). Knudsen's
universal-curve local model is exactly \(t=uv\) (1983, proof of Theorem
2.7, p.~179); the analytic nodal coordinates are also used explicitly in
Jansen (2023), Appendix B.2, pp.~43--44. At \((u,v)=(0,0)\) the extra
marked point lies at the node of the contracted curve; stabilization
inserts a rational bridge. Thus the single \(B\mathbf Z\) case is
precisely the local bridge-contraction followed by smoothing described
in the intake. This verifies that local mechanism without asserting a
global equivalence of all the archive's lifting categories with these
local ones.

\subsection{CC-RES03. Exact remaining scope of
H.5}\label{cc-res03.-exact-remaining-scope-of-h.5}

The independent results above resolve the actual sheaf/right-Kan
identification and compute the exceptional nodal lifting category in a
fully specified model. They do not by themselves prove that an arbitrary
categorical base change of the archive's entrance functor preserves
\(W_\infty\)-equivalences, nor its claimed global smoothness
classification. Properness of \(\pi\) supplies geometric proper base
change; that statement quantifies over geometric base changes and cannot
be silently replaced by the universal quantifier over functors in a
definition of \(W_\infty\)-properness. Likewise, a noncontractible local
lifting category does not exclude a possible contraction caused by
additional global morphisms unless the local-to-global comparison is
proved. The missing input is the archive's full functor, the definitions
it uses, and its comparison of global lifting categories with the local
models. H.5 remains an uninspected archive assertion; no contradiction
has been established.

\subsection{References}

Haine, P. J., Porta, M., \& Teyssier, J.-B. (2026). \emph{Exodromy
beyond conicality} (arXiv:2401.12825v2; revised 24 June 2026).
\url{https://arxiv.org/abs/2401.12825v2} . Inspected: Observation 2.3.1,
p.~19; Observation 2.3.6, p.~20; Theorem 3.2.3 and its proof,
pp.~34--35; Proposition 5.2.11, p.~53.
\url{https://arxiv.org/pdf/2401.12825v2} .

Jansen, M. Ø. (2023). \emph{Moduli stack of stable curves from a
stratified homotopy viewpoint} (arXiv:2308.09551v2).
\url{https://arxiv.org/abs/2308.09551v2} . Inspected: Corollary 6.6 and
Theorem 6.7, p.~34; Theorem 6.8, p.~35; Appendix B, local coordinates
and trivializations, pp.~40--45. \url{https://arxiv.org/pdf/2308.09551v2} .

Jansen, M. Ø. (2024). Stratified homotopy theory of topological
\(\infty\)-stacks: A toolbox. \emph{Journal of Pure and Applied Algebra,
228}(11), Article 107710. \url{https://doi.org/10.1016/j.jpaa.2024.107710} .
Inspected preprint arXiv:2308.09550v3: Theorem 2.17 and proof,
pp.~10--11; Definition 3.12 and following discussion, pp.~16--17; Remark
3.15, p.~17; Proposition 3.24, p.~21. \url{https://arxiv.org/pdf/2308.09550v3}
.

Knudsen, F. F. (1983). The projectivity of the moduli space of stable
curves, II: The stacks \(M_{g,n}\). \emph{Mathematica Scandinavica, 52},
161--199. \url{https://doi.org/10.7146/math.scand.a-12001} . Theorem 2.7 and
its displayed formal local map, p.~179, visually inspected:
\url{https://journals.msp.org/mscand/article/download/1622/1621/1653} .

% END consequences_criteria_RESIDUE.tex

% BEGIN root_linear_archive.tex
\section{Cote 134-8: the linear lemmas, normalization, and a proved range of the proposed theorem}

\subsection{The recovered question and the scope of its predicates}
The primary source is Grothendieck's \emph{Pursuing stacks}, cote
134-8, manuscript footers 47--54, in the official Montpellier PDF
\url{https://grothendieck.umontpellier.fr/134-8.pdf}.  This PDF has a
cover: its physical page number is the manuscript footer plus one.
The displayed \emph{Thm(?)} is on footer 48, physical page 49.
The original uses a connective chain-complex target, denoted with a
subscript on \(D\); we do not replace it by the unbounded derived
category.  The statements proved below use the unambiguous modern
notation \(D_{\geq0}(k)\), with homological differentials of degree
\(-1\).

On footer 47 the scalar functor \(\kappa(U)=\mathbf1\otimes_k U\)
is required to commute with inverse limits and is given a left adjoint
\(\varphi\).  The separately mentioned right adjoint of \(\kappa\)
is \(\operatorname{Hom}_{\mathcal M}(\mathbf1,-)\).  The setup treats
\(\kappa\) as exact.  There are enough projectives, the canonical
map \(L\varphi(\mathbf1)\to k\) is an isomorphism, and a monomorphism
\(\alpha:\mathbf1\to J\) is sought with
\[
 X\longrightarrow J\otimes X\text{ monic},\qquad
 L\varphi(J\otimes X)=0\quad(X\in\mathcal M).
 \tag{134.1}
\]
The vanishing includes degree zero.  The proposed conclusion is that
localizing \(\mathcal M\) at the maps inverted by \(L\varphi\)
gives the connective derived category of scalar modules.  Footer 48
then adds projective generators of finite presentation and sketches a
chain-complex functor and its right adjoint.  Its right margin also
records a projectivity condition on the scalar modules of morphisms.
We state this extra condition explicitly in BR-13405.

The supplied DER\_LOCALIZERS intake reports results C.1--C.29,
including a counterexample outside the finitely presented-generator
case.  Their original proofs are absent.  The present proofs therefore
establish the claims they state directly; they do not certify the
unseen counterexample or silently identify its hypotheses with these
handwritten predicates.  No RH or GRH input is needed for this section.

\subsection{BR-13401. The weakened form of Lemme 1 is valid}
Let \(\mathcal M,\mathcal A\) be abelian categories, with enough
projectives in \(\mathcal M\), and let
\(\Lambda:\mathcal M\to\operatorname{Ch}_{\geq0}(\mathcal A)\)
be an exact additive functor.  Put \(\varphi=H_0\Lambda\).
Assume that, for every projective \(P\in\mathcal M\),
\[
 H_j(\Lambda P)=0\quad(j>0).
 \tag{134.2}
\]
Then \(\varphi\) is right exact and there is a natural isomorphism
\[
 L\varphi(X)\simeq\Lambda X
 \quad\text{in }D_{\geq0}(\mathcal A),\qquad X\in\mathcal M.
 \tag{134.3}
\]
It is unnecessary to require \(\Lambda P\) to be a projective object
of the chain-complex category.  This is the weakened condition written
beside Lemme 1 on manuscript footer 53, physical page 54.

\paragraph{Complete proof.}
Exactness of \(\Lambda\) and right exactness of \(H_0\) imply right
exactness of \(\varphi\).  Choose a projective resolution
\(P_\bullet\to X\) in nonnegative homological degrees.  Form the
first-quadrant double complex
\(B_{a,b}=(\Lambda P_a)_b\), with the usual sign in one
differential so that the differentials anticommute.  Its total complex
uses only the finite direct sum \(\bigoplus_{a+b=n}B_{a,b}\) in each
degree \(n\).

The augmentations \(P_\bullet\to X\) give a map
\(\operatorname{Tot}B\to\Lambda X\).  In every fixed internal degree
\(b\), its augmented horizontal complex is exact, since the degree
\(b\) functor of \(\Lambda\) is exact.  The first-quadrant filtration
therefore proves this map is a quasi-isomorphism.  In the other
direction, the natural augmentations
\(\Lambda P_a\to H_0(\Lambda P_a)[0]\) give a map
\(\operatorname{Tot}B\to\varphi(P_\bullet)\).  Its vertical homology
is an isomorphism in degree zero and zero in positive degree by
(134.2).  The other first-quadrant filtration proves that it too is a
quasi-isomorphism.  These assertions can equivalently be checked by
successive filtration quotients in each total degree; no convergence
of an infinite diagonal is being assumed.  The complex
\(\varphi(P_\bullet)\) computes \(L\varphi(X)\), giving (134.3).
Comparison maps between projective resolutions are unique up to chain
homotopy; applying the additive functor and totalization preserves
these homotopies.  The resulting isomorphism is consequently natural
and independent of the resolution. \(\square\)

\subsection{BR-13402. Projective objects of the connective chain category}
Suppose \(\mathcal A\) is abelian with enough projectives and
coproducts, and coproducts of projectives are projective.  For
\(n\geq1\), let \(D^n(Q)\) be the complex with \(Q\) in degrees
\(n,n-1\), identity differential, and zero elsewhere.  A complex
\(P\) is projective as an object of
\(\operatorname{Ch}_{\geq0}(\mathcal A)\) if and only if it is
isomorphic to
\[
 U[0]\oplus\bigoplus_{n\geq1}D^n(Q_{n-1}),
 \tag{134.4}
\]
where \(U,Q_j\) are projective in \(\mathcal A\).  The sum here is
locally finite in the chain degree.  In particular, a projective
complex is not the same as an arbitrary complex with projective terms.

\paragraph{Complete proof.}
There are natural identifications
\[
 \operatorname{Hom}_{\mathrm{Ch}}(Q[0],C)
   =\operatorname{Hom}_{\mathcal A}(Q,C_0),\qquad
 \operatorname{Hom}_{\mathrm{Ch}}(D^nQ,C)
   =\operatorname{Hom}_{\mathcal A}(Q,C_n).
\]
Consequently the indicated building blocks are projective.  Their
coproducts are projective, proving one implication.

For the converse, choose projective epimorphisms \(Q_n\to P_n\)
for all \(n\geq0\).  The preceding identifications give a degreewise
epimorphism
\[
 Q_0[0]\oplus\bigoplus_{n\geq1}D^nQ_n\longrightarrow P.
\]
It splits because \(P\) is projective.  Thus the terms of \(P\)
are projective, \(H_jP=0\) for \(j>0\), and \(H_0P\) is
projective, all by taking the corresponding direct summands of the
source.  Write \(B_j=\operatorname{im}(d:P_{j+1}\to P_j)\).
The sequence \(0\to B_0\to P_0\to H_0P\to0\) splits, so
\(B_0\) is projective.  Inductively,
\(0\to B_j\to P_j\to B_{j-1}\to0\) splits for every \(j\geq1\),
because \(H_jP=0\) and \(B_{j-1}\) is projective.  These splittings
identify \(P_0\) with \(H_0P\oplus B_0\) and \(P_j\) with
\(B_j\oplus B_{j-1}\), with differential the identity on the second
factor and zero on the first.  This is (134.4), with \(U=H_0P\).
\(\square\)

\subsection{BR-13403. Unit normalization and the exact remaining counit test}
Let \(k\) be a commutative ring.  Assume the following precise data:
\(\mathcal M\) is a cocomplete \(k\)-linear abelian category with
enough projectives; \(\kappa(U)=\mathbf1\otimes_k U\) is an exact
scalar functor with left adjoint \(\varphi\); and there is an exact
\(k\)-linear adjunction
\[
 \Lambda:\mathcal M\rightleftarrows
 \operatorname{Ch}_{\geq0}(k):G
 \tag{134.5}
\]
with both functors exact, \(H_0\Lambda=\varphi\), and (134.2).
Assume also that the canonical counit
\(\Lambda G(k[0])\to k[0]\) is a quasi-isomorphism.  Equivalently,
under the identification below, this is the canonical normalization
\(L\varphi(\mathbf1)\to k\), including degree zero.

Set \(W=\Lambda^{-1}(\mathrm{qis})\).  Then (134.5) induces
mutually inverse equivalences on the localizations, with the induced
unit and counit, if and only if
\[
 L\varphi(GP)=0
 \quad\text{for every acyclic projective object }
 P\in\operatorname{Ch}_{\geq0}(k).
 \tag{134.6}
\]
Thus (134.6) is an exact residual test for this specified adjunction,
not merely a sufficient list of examples.  It is the question left at
the bottom of manuscript footer 54, physical page 55.

\paragraph{Complete proof.}
For every scalar module \(U\) and every \(X\in\mathcal M\),
the adjunctions give natural isomorphisms
\[
 \operatorname{Hom}_{\mathcal M}(X,G(U[0]))
 =\operatorname{Hom}_{\mathrm{Ch}}(\Lambda X,U[0])
 =\operatorname{Hom}_k(\varphi X,U)
 =\operatorname{Hom}_{\mathcal M}(X,\kappa U).
\]
Yoneda identifies \(G(U[0])\) with \(\kappa U\), and identifies
the two counits.  If \(U\) is projective, it is a summand of a free
module.  Both \(\kappa\) and \(\Lambda\) preserve coproducts:
the first by its scalar-tensor construction, the second because it is
a left adjoint.  Therefore the normalization for \(k\) implies that
\(\Lambda G(U[0])\to U[0]\) is a quasi-isomorphism for every
projective \(U\), by coproducts and retracts.  The coproduct of these
quasi-isomorphisms is a quasi-isomorphism in modules, since direct
sums are exact.

Every projective chain object splits by BR-13402 as \(U[0]\oplus P\)
with \(U\) projective and \(P\) acyclic projective.  Additivity of
both adjoints means that its counit is a quasi-isomorphism exactly when
the counit on \(P\) is one.  Since \(P\) is acyclic, the latter
condition is precisely \(\Lambda GP\simeq0\), equivalently (134.6)
by BR-13401.  The exact-adjunction criterion F-SM17 proves that counits
on all projective chain objects are quasi-isomorphisms if and only if
all counits are, and that this is equivalent to the stated localized
adjunction equivalence.  This proves both implications. \(\square\)

\subsection{BR-13404. The proposed equivalence holds over fields and over \(\mathbf Z\)}
Retain all the data of BR-13403.  Assume \(\mathcal M\) has enough
injectives and the tensor interval (134.1) exists.  Suppose there is
an integer \(r\geq0\) such that every projective \(k\)-module has
injective dimension at most \(r\).  Then (134.6) holds and the exact
adjunction induces an equivalence
\[
 W^{-1}\mathcal M\simeq D_{\geq0}(k).
 \tag{134.7}
\]
In particular this conclusion holds for a field \(k\) and for
\(k=\mathbf Z\), once the specified exact adjunction has been
constructed.  These are the rational and integral scalar rings used
in the programme.  The result does not assert an unbounded derived
equivalence or eliminate the construction hypotheses of BR-13403.

\paragraph{Complete proof.}
If \(I\) is injective in \(\mathcal M\), its monomorphism
\(I\to J\otimes I\) splits.  Applying \(L\varphi\) and (134.1)
shows \(L\varphi(I)=0\), since it is a retract of zero.  Any object
\(X\) having a finite injective resolution consequently satisfies
\(L\varphi(X)=0\).  Indeed, in a short exact sequence
\(0\to X\to I\to X'\to0\) with \(I\) injective, the triangle for
the left derived functor gives
\(L\varphi(X')\simeq L\varphi(X)[1]\).  Starting at the final
injective term of a finite resolution and working backwards proves
the assertion, including degree-zero vanishing.

Let \(P\) be an acyclic projective chain object.  BR-13402 gives
\(P=\bigoplus_{n\geq1}D^n(Q_{n-1})\) for projective scalar
modules \(Q_j\).  Choose length-\(r\) injective resolutions of the
\(Q_j\).  Apply each disk functor and take the locally finite sum
over \(j\).  In each chain degree only two of these resolutions
occur, so the resulting sequence is exact; it is a length-\(r\)
resolution of \(P\).

Each of its terms is injective as a chain object.  To see this without
an unsupported assertion about infinite sums of injectives, note that
\(D^{j+1}(I)\) is right adjoint, in the appropriate variable, to
the exact evaluation \(C\mapsto C_j\):
\[
 \operatorname{Hom}_{\mathrm{Ch}}(C,D^{j+1}I)
    =\operatorname{Hom}_k(C_j,I).
\]
It is therefore injective if \(I\) is injective.  The locally finite
sum of these disks equals their product degree by degree.  Products
of injective objects are injective, since maps into the product are
products of maps and products of surjections of abelian groups are
surjective.  This proves the asserted injectivity.

Because \(\Lambda\) is exact, its right adjoint \(G\) sends
injectives to injectives.  Because \(G\) is exact, applying it to
the length-\(r\) resolution gives a finite injective resolution of
\(GP\) in \(\mathcal M\).  The first paragraph now gives
\(L\varphi(GP)=0\).  BR-13403 proves (134.7).

Over a field, every module is injective, so \(r=0\) works.  Over
\(\mathbf Z\), every free module \(F\) has
\(0\to F\to F\otimes\mathbf Q\to(F\otimes\mathbf Q)/F\to0\),
whose two latter groups are divisible and hence injective.  Every
projective abelian group is a summand of a free one; taking summands
shows injective dimension at most one.  For completeness, the
divisibility assertion uses the elementary injectivity criterion:
a homomorphism \(n\mathbf Z\to D\) extends to \(\mathbf Z\to D\)
exactly when the prescribed image of \(n\) is divisible by \(n\).
The usual ideal criterion for injectivity then applies.  Thus \(r=1\)
works over \(\mathbf Z\). \(\square\)

\subsection{BR-13405. An explicit construction of the exact adjunction in the corollary's locally projective case}
Let \(R\) be a small \(k\)-linear category such that every
\(R(x,y)\) is a projective \(k\)-module, and let
\(\mathcal M=\operatorname{Fun}_k(R^{\mathrm{op}},\operatorname{Mod}k)\).
Assume \(\kappa(U)=T\otimes_k U\), for a fixed right \(R\)-module
\(T\), is exact and has a left adjoint \(\varphi\).  Then there
is an exact \(k\)-linear functor \(\Lambda\) representing
\(L\varphi\) on objects and an exact right adjoint \(G\), as in
(134.5).  This supplies a fully specified construction when the
finite-presentation projective generators and the additional scalar
projectivity condition on footer 48 are interpreted by such a ringoid.

\paragraph{Complete proof.}
Write \(P_x=R(-,x)\) for the right representable and set
\(D(x)=\varphi(P_x)\).  The dependence on \(x\) makes \(D\) a
left \(R\)-module.  Since \(\kappa\) is exact, its left adjoint
\(\varphi\) sends projectives to projectives, so each \(D(x)\)
is a projective scalar module.  In fact it is finitely generated:
\[
 \operatorname{Hom}_k(D(x),U)
   =\operatorname{Hom}_{\mathcal M}(P_x,\kappa U)
   =T(x)\otimes_k U
\]
commutes with arbitrary coproducts.  A projective module with this
property is finitely generated.  To verify this last assertion, write
it as a summand of \(k^{(I)}\).  Its inclusion in that direct sum
has, by the coproduct property, only finitely many nonzero coordinate
maps, so it is a summand of a finite free module.

There is a natural isomorphism \(\varphi(X)=D\otimes_R X\).
Both sides preserve coproducts and cokernels, and they agree naturally
on each representable.  Every right module has a presentation by
coproducts of representables; applying the two functors to such a
presentation proves the isomorphism and its naturality.

Take a projective resolution \(B_\bullet\to D\) by left
\(R\)-modules, choosing each \(B_a\) to be a coproduct of left
representables \(R(x,-)\).  Put
\[
 \Lambda(X)=B_\bullet\otimes_R X.
 \tag{134.8}
\]
Every degree of this functor is a sum of evaluation functors, and is
therefore exact; hence \(\Lambda\) is exact.  Its zeroth homology
is \(D\otimes_R X=\varphi X\).  For \(X=P_x\), the complex is
\(B_\bullet(x)\), whose positive homology vanishes and whose zeroth
homology is \(D(x)\).  Coproducts and retracts give the same
positive-homology vanishing for every projective \(X\).
BR-13401 proves that (134.8) represents \(L\varphi\).

Define a right \(R\)-module by
\[
 (GC)(x)=\operatorname{Hom}_{\operatorname{Ch}_{\geq0}(k)}
                    (B_\bullet(x),C).
 \tag{134.9}
\]
The ordinary tensor--Hom adjunction, degree by degree with the chain
map compatibility, gives \(\Lambda\dashv G\).  Every term
\(B_a(x)\) is a coproduct of projective scalar modules \(R(y,x)\),
and is projective.  The complex \(B_\bullet(x)\) resolves the
projective module \(D(x)\).  Its successive boundary sequences
therefore split, exactly as in BR-13402; it is a projective object of
the chain category.  It follows that (134.9) is an exact functor of
\(C\), for every \(x\).  Exactness in a functor category is
objectwise, so \(G\) is exact.  This completes the construction.
\(\square\)

\subsection{BR-13406. Positive-degree interval acyclicity is insufficient}
The degree-zero clause in (134.1) cannot be removed.  More precisely,
replace its vanishing by
\(L_i\varphi(J\otimes X)=0\) only for \(i>0\), keeping the scalar
unit normalization and all the exact-adjunction data.  The proposed
canonical equivalence need no longer hold.

\paragraph{Complete counterexample.}
Take a nonzero field \(k\), \(\mathcal M=\operatorname{Mod}k\)
with its ordinary tensor product, \(\kappa=\varphi=\mathrm{id}\),
and \(J=\mathbf1=k\), \(\alpha=\mathrm{id}\).  The interval map
\(X\to J\otimes X\) is the identity, all positive derived functors
of \(\varphi\) vanish, and \(L\varphi(k)=k\).  Let
\(\Lambda(X)=X[0]\) and \(G(C)=C_0\).  These are exact adjoint
functors, since a chain map from \(X[0]\) into a nonnegative complex
is exactly a module map \(X\to C_0\).  They satisfy (134.2).
The class \(W\) consists precisely of isomorphisms of modules, so
the canonical localized functor still has only complexes concentrated
in degree zero in its essential image.  The object \(k[1]\) of
\(D_{\geq0}(k)\) has nonzero first homology and cannot lie in that
image.  Thus the asserted canonical equivalence fails.  Exactly the
missing degree-zero condition fails, since
\(\varphi(J\otimes X)=X\) need not be zero. \(\square\)

\paragraph{What is settled and what remains.}
The original Lemme 1 is proved under its weaker, explicitly annotated
hypothesis.  F-SM17 proves the exact-adjunction/counit lemma without an
infinite-coproduct assumption on its right adjoint.  BR-13403 identifies
the residual acyclic-projective test, and BR-13404 discharges it over
fields and \(\mathbf Z\) under the specified exact-adjunction data.
BR-13405 constructs those data in the locally projective ringoid case.
For a general scalar ring without the finite-injective-dimension
condition, (134.6) remains the precise test for the raw adjunction.
An equivalence constructed by a different derived or model-category
adjunction would require its own comparison with (134.5).  The much
stronger assertion that every ringoid presentation automatically
satisfies the scalar projectivity hypothesis has not been made.


% END root_linear_archive.tex

% BEGIN source_cote1348_counit.tex
\section{Cote 134-8, Lemma 2: a complete counit criterion}
\label{sec:SM-cote1348-counit}

The manuscript at handwritten pages 52--53 considers an adjunction from
an abelian category to a category of complexes, the inverse image of
quasi-isomorphisms, and tests on the adjunction counit.  Its conditions
labelled $(ii')$ and $(ii'')$ have revisions that should not be
silently identified with a modern assertion.  The following gives a
precise version of the clear counit/projective-object criterion.  All
complexes here are homological and concentrated in degrees at least
zero.  The exactness hypothesis on the right adjoint is explicit.

\paragraph{Finding F-SM-17 (localization equivalence and projective counit tests).}
Let $\mathcal A$ and $\mathcal M$ be abelian categories, with
$\mathcal A$ having enough projectives.  Put
$\mathcal B=\operatorname{Ch}_{\geq0}(\mathcal A)$ and let
\[
 \Lambda:\mathcal M\rightleftarrows\mathcal B:G
 \qquad(\Lambda\dashv G)
\]
be an additive adjunction such that both functors are exact.  Write
$\eta$ and $\varepsilon$ for its unit and counit, let
$\mathcal Q$ be the quasi-isomorphisms in $\mathcal B$, and put
$W=\Lambda^{-1}(\mathcal Q)$.  Fix universes in which the following
localizations exist.  Then these conditions are equivalent:
\begin{enumerate}
\item For every $C\in\mathcal B$, the counit
      $\varepsilon_C:\Lambda G(C)\to C$ is a quasi-isomorphism.
\item The same assertion holds for every projective object $P$ of
      the abelian category $\mathcal B$.
\item $G$ sends $\mathcal Q$ into $W$, and the induced functor
      \[
       \overline\Lambda:\mathcal M[W^{-1}]
          \longrightarrow \mathcal B[\mathcal Q^{-1}]
          =D_{\geq0}(\mathcal A)
      \]
      is an equivalence.
\end{enumerate}
Whenever these hold, $G^{-1}(W)=\mathcal Q$, the induced functors
$\overline\Lambda\dashv\overline G$ are mutually inverse equivalences,
and every $\eta_X$ belongs to $W$.

The equivalence of the first and third conditions does not require
abelian categories or exactness; it uses the displayed adjunction and
the saturation of quasi-isomorphisms.  For the implication from the
second condition to the first, exactness of the composite
$H=\Lambda G$ suffices in place of separate exactness of both functors.
No preservation of infinite coproducts by $G$ is assumed.

\begin{proof}
\emph{Enough projectives in the category of complexes.}
For a projective $P\in\mathcal A$, let $S^0(P)$ be $P$ in degree zero.
For $n\geq1$, let $D^n(P)$ have $P$ in degrees $n$ and $n-1$,
with the identity as its only nonzero differential.  There are natural
identifications
\[
 \operatorname{Hom}_{\mathcal B}(S^0(P),C)
       =\operatorname{Hom}_{\mathcal A}(P,C_0),
 \qquad
 \operatorname{Hom}_{\mathcal B}(D^n(P),C)
       =\operatorname{Hom}_{\mathcal A}(P,C_n).
\]
An epimorphism of complexes is an epimorphism in every degree.
Consequently these stalk and disk complexes are projective objects of
$\mathcal B$.

For any $C$, choose epimorphisms $P_n\to C_n$ with each $P_n$
projective.  The associated chain maps give an epimorphism
\[
 Q=S^0(P_0)\oplus\bigoplus_{n\geq1}D^n(P_n)\longrightarrow C.
\]
This coproduct exists without any assumption about arbitrary
coproducts in $\mathcal A$: in degree $j$ its value is the finite
biproduct $P_j\oplus P_{j+1}$.  A map from $Q$ lifts across any
epimorphism of complexes by lifting separately its maps from the
$P_n$ into degree $n$.  Hence $Q$ is projective.  Iterating on kernels
gives a projective resolution of every object of $\mathcal B$.

\emph{The finite-diagonal total-complex argument.}
We record the elementary convergence fact used below.  For a first
quadrant bicomplex $X_{r,s}$ with commuting differentials $d_h,d_v$,
give
\[
 \operatorname{Tot}(X)_N=\bigoplus_{r+s=N}X_{r,s},
 \qquad d=d_h+(-1)^r d_v.
\]
If every column is acyclic, the total complex is acyclic.  Indeed, for
a total cycle choose its component with largest first index $r$.
The cycle equation makes that component a vertical cycle, since there
is no component with first index $r+1$.  A vertical primitive, with
the appropriate sign, removes it by subtracting a total boundary;
the only newly introduced component has first index $r-1$.  Repeating
terminates after finitely many steps.  The same proof, starting with
the largest second index and using horizontal primitives, works if
every row is acyclic.  This argument uses only the finitely many
positions on a diagonal.  In a general $\mathcal A$ one may perform
this cycle calculation after applying
$\operatorname{Hom}_{\mathcal A}(P,-)$ for each projective $P$.
That functor is exact and commutes with the finite sums in question;
enough projectives makes these functors detect zero homology.

Two consequences will be used.  A morphism of first quadrant
bicomplexes that is a quasi-isomorphism in each column induces a
quasi-isomorphism of total complexes: take its columnwise mapping
cone, apply the column argument, and identify its total complex with
the mapping cone of the total map, using the usual signs.  Also, if
$X_{\bullet,s}\to C_s$ is an exact augmented horizontal complex for
each $s$, the augmentation
$\operatorname{Tot}(X)\to C$ is a quasi-isomorphism.  For this, insert
$C_s$ as the zeroth column and shift the old columns right by one;
the resulting first quadrant bicomplex has exact rows, and its total
complex is the cone of the augmentation, up to the corresponding
sign convention.  Thus its cone is acyclic.  Neither consequence
uses a totalization with infinite sums in a fixed degree.

\emph{Projective tests imply all tests.}
Assume the second condition, and write $H=\Lambda G$.  Choose a
projective resolution in the abelian category $\mathcal B$,
\[
 \cdots\longrightarrow P_2\longrightarrow P_1
       \longrightarrow P_0\longrightarrow C\longrightarrow0 .
\]
Regard $(P_r)_s$ as a first quadrant bicomplex.  Since $H$ is exact,
the augmented horizontal complex
\[
 \cdots\longrightarrow H(P_1)\longrightarrow H(P_0)
       \longrightarrow H(C)\longrightarrow0
\]
is also exact in $\mathcal B$, hence in each internal degree.
The preceding augmentation argument gives quasi-isomorphisms
$\operatorname{Tot}(P_\bullet)\to C$ and
$\operatorname{Tot}(H(P_\bullet))\to H(C)$.

Naturality of the counit gives a morphism between these bicomplexes.
On its $r$th column it is $\varepsilon_{P_r}$, a quasi-isomorphism by
assumption. The column argument therefore makes the top arrow in
the commutative square
\[
\begin{array}{ccc}
 \operatorname{Tot}(H(P_\bullet))&
       \xrightarrow{\operatorname{Tot}(\varepsilon)}&
                         \operatorname{Tot}(P_\bullet)\\
 \downarrow&&\downarrow\\
 H(C)&\xrightarrow{\varepsilon_C}&C
\end{array}
\]
a quasi-isomorphism.  Both vertical arrows are quasi-isomorphisms.
On homology the square then shows that $\varepsilon_C$ is an
isomorphism in every degree.  This proves the first condition.  The
converse implication to the second is immediate.

Notice that $H$ was applied separately to the resolution objects
$P_r$.  We neither equated $H(\operatorname{Tot}P_\bullet)$ with
$\operatorname{Tot}H(P_\bullet)$ nor applied $G$ to any infinite
coproduct.

\emph{All counits imply the localization equivalence.}
Suppose the first condition holds.  For a morphism $v:C\to C'$,
naturality gives
\[
 \varepsilon_{C'}\circ\Lambda G(v)=v\circ\varepsilon_C.
\]
Since both counits are quasi-isomorphisms, the two-out-of-three
property gives
\[
 \Lambda G(v)\in\mathcal Q
       \quad\Longleftrightarrow\quad v\in\mathcal Q.
\]
This is $G^{-1}(W)=\mathcal Q$ and, in particular, allows $G$ to
induce $\overline G$ on the localizations.  The triangular identity
\[
 \varepsilon_{\Lambda X}\circ\Lambda(\eta_X)
       =1_{\Lambda X}
\]
gives $\Lambda(\eta_X)\in\mathcal Q$, hence $\eta_X\in W$.
The unit and counit therefore descend to natural isomorphisms on the
localized categories.  Their triangular identities remain valid,
so $\overline\Lambda$ and $\overline G$ are inverse equivalences.

\emph{The converse with the localized adjunction made explicit.}
Assume the third condition.  Because both functors preserve the
respective classes being inverted, their original unit and counit
descend to natural transformations on the localizations: naturality
on original arrows implies naturality on their formal inverses.
Their triangular identities give an adjunction
$\overline\Lambda\dashv\overline G$.

The unit of an adjunction whose left adjoint is fully faithful is an
isomorphism: for each $Z$, its map on $\operatorname{Hom}(Z,-)$ is,
under adjunction, the fully faithful map on morphisms induced by the
left adjoint, so the conclusion follows from Yoneda.  Apply this to
$\overline\Lambda$.  The triangular identity makes the counit an
isomorphism on objects in its essential image, and essential
surjectivity and naturality make it an isomorphism on every object.
Thus the image of each $\varepsilon_C$ in
$\mathcal B[\mathcal Q^{-1}]$ is invertible.

Finally, a chain map whose image is invertible in this localization
is a quasi-isomorphism.  Each homology functor factors through the
localization, so it sends that map to an isomorphism.  This proves
the first condition and completes all implications.
\end{proof}

\paragraph{Scope of the projective test and of the source reading.}
The phrase \emph{projective object of
$\operatorname{Ch}_{\geq0}(\mathcal A)$} means projectivity in the
abelian category.  It is stronger than being projective in every
degree: a nonzero projective object placed only in degree one is
termwise projective but is not a projective complex.  For example the
evident epimorphism $D^1(P)\to S^1(P)$ admits no chain section when
$P\ne0$.  Testing the counit on \emph{all} termwise projective
complexes would be a stronger assumption than condition two and
therefore also sufficient; the two classes must nevertheless not be
identified.  Testing only projectives of $\mathcal A$ placed in
degree zero is not what has been proved here.

The source's arrow from $\mathcal M$ to the category of complexes is
the left adjoint; its reverse arrow is the right adjoint.
Handwritten page~53 explicitly discusses the right adjoint preserving
epimorphisms, which, together with its automatic left exactness,
gives the exactness used here.  Page~52 displays the all-counits
quasi-isomorphism condition and an equivalence after localization.
F-SM-17 establishes this precise modern formulation completely.
It does not assert that every overwritten formulation of
$(ii')$ or $(ii'')$ has been deciphered, nor does it infer a localization
equivalence from the left functor alone without the stated compatibility
of its right adjoint.

The primary locator is Grothendieck, A. (n.d.),
\emph{Pursuing stacks: The Modelizing Story, Chapter VII,
Linearization of homotopy types}, cote~134-8, handwritten
pp.~52--53, physical PDF pp.~53--54:
\url{https://grothendieck.umontpellier.fr/134-8.pdf}.
The argument above is a new complete presentation of the stated
criterion, with no claim that its categorical mechanism is historically
new or that the overwritten variants have been recovered verbatim.

% END source_cote1348_counit.tex

% BEGIN source_cote1348_purity.tex
\section{Cote 134-8: two extensions using purity}
\label{sec:SM-cote1348-purity}

Retain the precise connective setup of BR-13403 and the ringoid
construction of BR-13405. Thus $k$ is commutative,
$R$ is a small $k$-linear category with projective scalar modules
$R(x,y)$, $\mathcal M=\operatorname{Mod}\!-\!R$, and
\[
 \varphi(X)=D\otimes_R X,\qquad
 \Lambda(X)=B_\bullet\otimes_R X,\qquad
 (GC)(x)=\operatorname{Hom}_{\mathrm{Ch}_{\geq0}(k)}(B_\bullet(x),C).
 \tag{SM24.1}
\]
Here $B_\bullet\to D$ is the chosen resolution by coproducts of
left representables. The canonical unit normalization and the full
interval condition, including degree zero, remain in force.
For precision we also write the scalar compatibility used below:
\[
 J\otimes_{\mathcal M}(X\otimes_k U)
   \simeq (J\otimes_{\mathcal M}X)\otimes_k U,\qquad
 \alpha_{X\otimes_k U}=\alpha_X\otimes_k U.
 \tag{SM24.2}
\]
It holds for a $k$-linear tensor product preserving colimits in each
variable: present $U$ by free modules and use coproducts and cokernels.
We make it explicit rather than infer it from an unspecified tensor
operation. It is the scalar compatibility of the tensor setup being
used.

\paragraph{Finding F-SM-24 (the ordinary injective-dimension hypothesis can be replaced).}
Under these hypotheses, either one of the following two conditions
suffices for the specified adjunction to induce
\[
 \mathcal M[W^{-1}]\simeq D_{\geq0}(k),
 \qquad W=\Lambda^{-1}(\mathrm{qis}).
 \tag{SM24.3}
\]
\begin{enumerate}
\item There is a finite $r$ such that every projective $k$-module
      has a pure-injective resolution of length at most $r$.
\item Every $R(x,y)$ is finitely generated projective over $k$.
\end{enumerate}
The second condition imposes no injective or pure-injective dimension
condition on $k$. The first, by the known cardinality theorem stated
below, applies to every countable scalar ring and more generally to
$|k|=\aleph_t$ for a finite integer $t\geq0$.
In the first condition, $R(x,y)$ may be infinitely generated.

\begin{proof}
\emph{Purity conventions and elementary facts.}
A monomorphism of right modules is pure if tensoring with every
left module keeps it injective. Equivalently, the associated short
exact sequence remains exact after applying $\operatorname{Hom}(F,-)$
for every finitely presented right module $F$.
This equivalence also holds for ringoid modules. Indeed, a presentation
of $F$ by finite sums of representables makes a lift of a map
$F\to C$ precisely a solution of a finite system of module equations
whose constants are in the submodule. Tensoring with the cokernel
of the transposed finite matrix tests solvability of that same system.
Arbitrary tensor tests reduce to finitely presented tests because
every module is a filtered colimit of finitely presented modules.
In particular, products of pure short exact sequences are pure:
maps from a finitely presented module into a product are tuples of
maps, and products of surjections of abelian groups are surjective.

A module $E$ is pure injective if maps into $E$ extend across pure
monomorphisms. Products of pure-injective modules are pure injective,
by extending the coordinate maps. Every scalar module $Q$ has a
pure monomorphism into a pure-injective module. Here is the needed
construction, avoiding any envelope assumption. Write
$Q^+=\operatorname{Hom}_{\mathbb Z}(Q,\mathbb Q/\mathbb Z)$.
For every module $V$, the character module $V^+$ is pure injective:
use
\[
 \operatorname{Hom}_k(M,V^+)
   =\operatorname{Hom}_{\mathbb Z}(V\otimes_k M,\mathbb Q/\mathbb Z)
\]
and injectivity of $\mathbb Q/\mathbb Z$.
The evaluation map $e_Q:Q\to Q^{++}$ is pure.
Indeed $e_Q^+:Q^{+++}\to Q^+$ has the section $e_{Q^+}$.
For every $V$, the character dual of $V\otimes e_Q$ is therefore
surjective. Since $\mathbb Q/\mathbb Z$ is an injective cogenerator
of abelian groups, $V\otimes e_Q$ is injective. Thus $Q^{++}$ is
the required pure-injective target.

We will also use the fact that all functors
$\operatorname{Tor}_n^R(D,-)$ preserve pure short exact sequences.
For ringoid modules, the character dual is taken at every object
and reverses the variance; the Ext groups below are in left
$R$-modules.
For completeness, suppose $0\to X\to Y\to Z\to0$ is pure.
The evaluation pairing
$X^+\otimes_R X\to\mathbb Q/\mathbb Z$ extends across
$X^+\otimes_R X\to X^+\otimes_R Y$, which is injective by purity.
The resulting map $X^+\to Y^+$ splits restriction
$Y^+\to X^+$. Thus the character-dual sequence of left
$R$-modules splits.
The natural identity
\[
 \operatorname{Tor}_n^R(D,X)^+
   \simeq \operatorname{Ext}_R^n(D,X^+)
 \tag{SM24.4}
\]
follows by using the projective resolution $B_\bullet$ and the
exactness of character duality.
Applying $\operatorname{Ext}_R^n(D,-)$ to the split dual sequence
and dualizing back proves the assertion. Exactness is reflected
by the character dual because it is an injective cogenerator.
In particular, an $L\varphi$-acyclic module has only
$L\varphi$-acyclic pure submodules.

\emph{Relative pure injectives are killed by the actual interval.}
For $x\in R$ and a scalar module $E$, put
\[
 C_x(E)(y)=\operatorname{Hom}_k(R(x,y),E).
 \tag{SM24.5}
\]
This is the right adjoint to evaluation at $x$:
\[
 \operatorname{Hom}_{\mathcal M}(X,C_x(E))
     =\operatorname{Hom}_k(X(x),E).
 \tag{SM24.6}
\]
Call a monomorphism objectwise pure if its evaluation at every
$x$ is pure over $k$. When $E$ is pure injective, (SM24.6)
shows that $C_x(E)$ is injective with respect to objectwise pure
monomorphisms. So is every product of such modules.

By (SM24.2), the interval monomorphism
$\alpha_X:X\to J\otimes_{\mathcal M}X$ is objectwise pure:
after tensoring its evaluation with any $U$, it is the evaluation
of the assumed monomorphism $\alpha_{X\otimes_k U}$.
Consequently, if
\[
 I=\prod_{a\in A} C_{x_a}(E_a),\qquad E_a\text{ pure injective},
\]
the identity on $I$ extends across $\alpha_I$.
Thus $I$ is a retract of $J\otimes_{\mathcal M}I$.
The full interval acyclicity gives
\[
 L\varphi(I)=0.
 \tag{SM24.7}
\]
No assertion that $I$ is an ordinary injective object is needed.

\emph{The counit test is a product of coinduced scalar coefficients.}
Let $P$ be an acyclic projective object of
$\mathrm{Ch}_{\geq0}(k)$. By BR-13402,
\[
 P=\bigoplus_{n\geq0}D^{n+1}(Q_n)
   =\prod_{n\geq0}D^{n+1}(Q_n),
 \qquad Q_n\text{ projective over }k.
\]
The equality with the product holds degreewise, since only two
disks contribute in any degree. Since $G$ is a right adjoint,
\[
 GP=\prod_{n\geq0}\operatorname{Hom}_k(B_n(-),Q_n).
 \tag{SM24.8}
\]
Writing $B_n=\bigoplus_{a\in A_n}R(x_a,-)$ gives
\[
 GP=\prod_{n\geq0}\prod_{a\in A_n} C_{x_a}(Q_n).
 \tag{SM24.9}
\]
The chain-map formula behind (SM24.8) is
$\operatorname{Hom}_{\mathrm{Ch}}(B,D^{n+1}Q)
=\operatorname{Hom}_k(B_n,Q)$.
In particular, no preservation of coproducts by $G$ was assumed.

\emph{First sufficient condition.}
Choose pure-injective resolutions of length $r$ for all $Q_n$.
The functor $C_x$ is exact, since $R(x,y)$ is projective over $k$
for each $y$. Apply these functors and take the products in
(SM24.9). Products are exact in this module category, so the
result is a finite exact sequence
\[
 0\longrightarrow GP\longrightarrow I^0\longrightarrow
       I^1\longrightarrow\cdots\longrightarrow I^r
       \longrightarrow0
\]
whose terms have the form appearing in (SM24.7).
Every $L\varphi(I^j)$ vanishes.
The long exact sequences of the derived functors, or successive
distinguished triangles, now give $L\varphi(GP)=0$.
This induction is finite; no infinite totalization is involved.

\emph{Second sufficient condition.}
Embed each $Q_n$ purely into $E_n=Q_n^{++}$.
For any finitely presented right $R$-module $F$, the scalar
module $F(x)$ is finitely presented over $k$: evaluate a finite
presentation of $F$ by representables and use the present
finitely generated projectivity of all $R(x,y)$.
If $0\to Q_n\to E_n\to V_n\to0$ is the resulting pure sequence,
(SM24.6) shows that
\[
 0\to C_x(Q_n)\to C_x(E_n)\to C_x(V_n)\to0
\]
is pure as a sequence of right $R$-modules. Exactness uses
projectivity of $R(x,y)$; the purity test uses finite presentation
of $F(x)$ and scalar purity.
Taking products gives a pure monomorphism
\[
 GP\longrightarrow I=
       \prod_{n\geq0}\prod_{a\in A_n}C_{x_a}(E_n).
\]
By (SM24.7), $L\varphi(I)=0$.
By (SM24.4), every $\operatorname{Tor}_n^R(D,GP)$ injects
into the corresponding zero group for $I$.
Thus $L\varphi(GP)=0$ in this case also.

In either case the exact residual condition BR-13403 is satisfied.
Together with the unchanged canonical unit normalization,
F-SM-17 gives (SM24.3) for the specified connective adjunction.
\end{proof}

\paragraph{The countable-ring consequence uses a known theorem.}
Chach\'olski, Neeman, Pitsch, and Scherer (2018),
Theorem~9.2, p.~926, explicitly state the
Kie{\l}pi\'nski--Simson/Gruson--Jensen theorem:
a ring of cardinality $\aleph_t$, $t\in\mathbb N$, has pure
global dimension at most $t+1$. Their preceding paragraph
states this in terms of finite pure-injective dimensions.
Thus F-SM-24(1) applies with $r=t+1$, and in particular with
$r=1$ for every countably infinite scalar ring.
Finite rings are covered as well: apply that theorem to
$k\times\mathbb Z$ and take the first idempotent component of
a pure-injective resolution. A pure sequence and a pure-injective
module over a product ring are exactly pairs of such sequences
and modules over the factors, as follows from the componentwise
tensor and Hom formulas. Hence the resulting resolution over
$k$ has the same length. This proves the stated countable-ring
consequence without claiming a new proof of the cardinality theorem.

\paragraph{Finding F-SM-25 (why the second proof does not automatically extend to arbitrary projective morphism modules).}
Coinduction need not carry a scalar pure short exact sequence to a
pure sequence of ringoid modules when its morphism modules are
infinitely generated projective. The following is a complete
counterexample to that particular inference. It is not a
counterexample to the interval theorem.

\begin{proof}
Take $k=\mathbb Z$ and the one-object ringoid $R=\mathbb Z[t]$,
which is free, but infinitely generated, over $k$.
Let
$E=\widehat{\mathbb Z}=\prod_p\mathbb Z_p$ and $V=E/\mathbb Z$,
where $\mathbb Z$ embeds diagonally.
The exact sequence
\[
 0\longrightarrow\mathbb Z\longrightarrow E
       \longrightarrow V\longrightarrow0
 \tag{SM25.1}
\]
is pure over $\mathbb Z$. For every $n>0$, the map
$\mathbb Z/n\to E/nE$ is the Chinese-remainder isomorphism;
tensoring with $\mathbb Z$ also preserves the injection.
The structure theorem for finitely presented abelian groups,
followed by filtered colimits, gives all tensor tests for purity.

Choose $a=\sum_{j\geq2}2^{j!}\in\mathbb Z_2$.
It is not an integer: its binary expansion has infinitely many
zeros and ones, whereas an ordinary integer has an eventually
zero or eventually one $2$-adic expansion.
Let $e\in E$ have $2$-adic component $a$ and all other components
zero. For each $n\geq0$, choose
$0\leq r_n<2^n$ with $r_n\equiv a\pmod{2^n}$, taking $r_0=0$.
Then
\[
 q_n=(e-r_n)/2^n\in E,\qquad
 2q_{n+1}-q_n=(r_n-r_{n+1})/2^n\in\mathbb Z.
\]
Division is legitimate in the $2$-adic component by the chosen
congruence and in every other component because $2$ is a unit.
Thus $2^{-n}\mapsto q_n+\mathbb Z$ defines a homomorphism
\[
 h:\mathbb Z[1/2]\longrightarrow V.
\]
It has no lift to $E$. The image of $1$ under any such lift would
belong to $\bigcap_n2^n E$, so its $2$-adic component would be
zero. But being a lift would make its difference from $e$ a
diagonal integer, forcing $a$ to be an integer, a contradiction.

Now $F=R/(2t-1)$ is a finitely presented $R$-module whose
underlying abelian group is $\mathbb Z[1/2]$.
For $C(U)=\operatorname{Hom}_{\mathbb Z}(R,U)$, the coinduction
adjunction identifies
\[
 \operatorname{Hom}_R(F,C(U))
       =\operatorname{Hom}_{\mathbb Z}(\mathbb Z[1/2],U).
\]
The coinduced version of (SM25.1) is short exact because
$R$ is free over $\mathbb Z$. It is not pure over $R$:
the displayed Hom map from $C(E)$ to $C(V)$ is not surjective,
by the constructed $h$. This proves the claim.
\end{proof}

\paragraph{The residual generality is stated exactly.}
F-SM-24 removes the ordinary scalar injective-dimension condition
for all countable scalar rings, for the other rings of finite
pure global dimension just described, and for every scalar ring
when the ringoid is locally finitely generated projective.
The remaining case has neither available finite pure-injective
dimension nor the finite-presentation property used in the second
proof. For that case the exact unresolved assertion is still
\[
 \operatorname{Tor}_j^R
 \left(D,\prod_{n\geq0}\prod_{a\in A_n}C_{x_a}(Q_n)\right)=0
 \quad(j\geq0)
 \tag{SM25.2}
\]
for every sequence of projective scalar modules $Q_n$.
F-SM-25 blocks one proposed automatic purity argument for
(SM25.2); it does not provide data satisfying the original
interval and normalization assumptions and violating (SM25.2).
The ordinary finite-injective-dimension proof, the relative-purity
proofs here, and a hypothetical infinite resolution must therefore
not be conflated.

\subsection*{Source for the cardinality input}
Chach\'olski, W., Neeman, A., Pitsch, W., \& Scherer, J. (2018).
Relative homological algebra via truncations.
\emph{Documenta Mathematica, 23}, 895--937.
\url{https://ems.press/content/serial-article-files/26419?nt=1}.
Exact input: \S9.2 and Theorem~9.2, p.~926.
The source attributes this result to
Kie{\l}pi\'nski--Simson (1975) and Gruson--Jensen (1973);
their original proofs are not claimed to have been reconstructed here.


% END source_cote1348_purity.tex

% BEGIN source_cote1348_fp.tex
\section{Further progress on the arbitrary-ring case of cote 134-8}

\subsection{F-SM-26: the full interval forces infinite relative Tor amplitude}
Assume the data of BR-13403 and the full interval condition (134.1).
Write \(T=\mathbf1\), and suppose \(k\ne0\).  For every object \(X\), put
\[
 S(X)=\operatorname{coker}\bigl(X\xrightarrow{\alpha_X}J\otimes X\bigr).
\]
Then
\[
 L\varphi(S^nT)\simeq k[n]\qquad(n\ge0).
 \tag{SM26.1}
\]
In the ringoid situation of BR-13405, where
\(\varphi(X)=D\otimes_R X\), this implies
\[
 \operatorname{Tor}^{R}_n(D,S^nT)\cong k\ne0
 \qquad(n\ge0).
 \tag{SM26.2}
\]
Consequently \(D\) has infinite flat dimension, and hence infinite
projective dimension, as a left \(R\)-module.  A finite flat-dimension
hypothesis on \(D\) cannot close the remaining counit test while retaining
the actual nonzero normalized interval hypotheses.

\paragraph{Complete proof.}
The defining sequence
\(0\to X\to J\otimes X\to S(X)\to0\) is exact.
The exact complex-valued representative \(\Lambda\) sends it to a short
exact sequence of complexes.  Its associated triangle, together with
\(\Lambda(J\otimes X)\simeq0\), gives
\[
 L\varphi(SX)\simeq L\varphi(X)[1].
\]
Here homological shift \([1]\) raises the homology degree by one.
Iteration and the canonical normalization \(L\varphi(T)\simeq k[0]\)
give (SM26.1).  Taking degree \(n\) homology gives (SM26.2).
Finite flat dimension \(r\) would make every such Tor group zero for
\(n>r\), a contradiction.  Finite projective dimension implies finite
flat dimension, completing the proof. \(\square\)

\subsection{F-SM-27: an arbitrary-scalar result when \(D\) is of type \(\mathrm{FP}_\infty\)}
Retain the ringoid, scalar, exact-adjunction and full interval hypotheses
of BR-13403--05.  Thus \(R\) is a small \(k\)-linear category, each
\(R(x,y)\) is a projective scalar module,
\(\mathcal M=\operatorname{Fun}_k(R^{\mathrm{op}},\operatorname{Mod}k)\),
\(\varphi=D\otimes_R-\), and the specified adjunction is formed from
a resolution \(B_\bullet\to D\) by coproducts of left representables.
Retain the canonical unit normalization, including degree zero.

Assume the scalar compatibility used in F-SM-24:
\[
 J\otimes(X\otimes_kU)\cong(J\otimes X)\otimes_kU,\qquad
 \alpha_{X\otimes_kU}=\alpha_X\otimes_kU.
 \tag{SM27.1}
\]
This holds for a \(k\)-linear tensor product preserving colimits in each
variable.  Finally suppose that \(D\), as a left \(R\)-module, has a
resolution
\[
 \cdots\longrightarrow A_2\longrightarrow A_1
 \longrightarrow A_0\longrightarrow D\longrightarrow0
 \tag{SM27.2}
\]
in which every \(A_n\) is a finitely generated projective left
\(R\)-module: equivalently, a retract of a finite coproduct of left
representables.  This is the additional condition
\(D\in\mathrm{FP}_\infty(R)\); no finite length for (SM27.2) is required.

Then the exact adjunction specified in BR-13405 induces
\[
 W^{-1}\mathcal M\simeq D_{\ge0}(k).
 \tag{SM27.3}
\]
There is no cardinality, ordinary injective-dimension or pure
injective-dimension restriction on \(k\), and the scalar modules
\(R(x,y)\) need not be finitely generated.  More strongly, with
\[
 C_x(U)(y)=\operatorname{Hom}_k(R(x,y),U),
 \tag{SM27.4}
\]
\(L\varphi\) vanishes on every \(C_x(U)\) and on every product of such
objects, for arbitrary scalar modules \(U\).

\paragraph{Exact known input.}
We use Benson and Goodearl's (2000) periodic-flat-module theorem:
a flat module \(V\) in an exact sequence
\(0\to V\to P\to V\to0\), with \(P\) projective, is projective.
The precise primary locator is Theorem 2.5, printed p.~50
(also stated as Theorem 1.1 on p.~46), in
\emph{Periodic flat modules, and flat modules for finite groups}.
Its proof occupies pp.~50--51, after Lemmas 2.1--2.4.
The complete deduction from that known theorem to the contractibility
needed below is given here.  In particular, ordinary acyclicity of a
complex of projectives, unbounded to the right cohomologically, is
not substituted for pure acyclicity.

\paragraph{Complete proof.}
We use the convention that a pure monomorphism of scalar modules
remains injective after tensoring with every scalar module.
A scalar module \(E\) is pure-injective if maps to \(E\) extend across
all such monomorphisms.  For every scalar module \(V\), its character
module
\[
 V^+=\operatorname{Hom}_{\mathbf Z}(V,\mathbf Q/\mathbf Z)
\]
is pure-injective: tensor--Hom adjunction identifies extension across
a pure monomorphism with extension to the injective abelian group
\(\mathbf Q/\mathbf Z\).  We will use this elementary fact both to
produce tests and to detect purity.

\emph{1. Coinduction with pure-injective coefficients is annihilated.}
There is a natural adjunction
\[
 \operatorname{Hom}_{\mathcal M}(X,C_x(E))
   =\operatorname{Hom}_k(X(x),E).
 \tag{SM27.5}
\]
Consequently, when \(E\) is pure-injective, \(C_x(E)\) is injective
with respect to all monomorphisms that are scalar-pure at every
object of \(R\).  Condition (SM27.1) and the interval's monomorphism
condition show that \(\alpha_X:X\to J\otimes X\) is objectwise
scalar-pure: tensoring it with any \(U\) gives the monomorphism
\(\alpha_{X\otimes_k U}\).  Thus \(\alpha_{C_x(E)}\) splits in
\(\mathcal M\), by extending the identity using (SM27.5).
The interval condition gives
\[
 L\varphi(C_x(E))=0\qquad(E\text{ scalar pure-injective}).
 \tag{SM27.6}
\]
This argument does not declare \(C_x(E)\) ordinarily injective.

\emph{2. Dualizing the finite-projective resolution.}
Fix \(x\), and form the cochain complex
\[
 K_x^n=\operatorname{Hom}_{R\text{-}\mathrm{Mod}}
                (A_n,R(x,-)),\qquad
 \delta^n(f)=f\circ d_{n+1}.
 \tag{SM27.7}
\]
Every \(K_x^n\) is a projective scalar module: for
\(A_n=R(y,-)\) it is \(R(x,y)\), and the general case follows by
finite coproducts and retracts.  For every scalar module \(U\),
evaluation and finite-projective duality give natural isomorphisms
\[
 A_n\otimes_R C_x(U)
   \cong\operatorname{Hom}_k(K_x^n,U).
 \tag{SM27.8}
\]
For a representable this is the defining formula (SM27.4).
Finite coproducts convert to finite products, which coincide, and
retracts preserve the isomorphism.  This also proves naturality in
\(A_n\); hence the isomorphisms commute with the differentials.
The homological complex \(A_\bullet\otimes_R C_x(U)\), computing
\(L\varphi(C_x(U))\), is therefore the transpose complex
\(\operatorname{Hom}_k(K_x^\bullet,U)\).

By (SM27.6), this transpose complex is acyclic for every pure-injective
\(U\).  In particular take the injective cogenerator
\(E_0=\operatorname{Hom}_{\mathbf Z}(k,\mathbf Q/\mathbf Z)\).
Then
\(\operatorname{Hom}_k(K_x^\bullet,E_0)
 =\operatorname{Hom}_{\mathbf Z}(K_x^\bullet,\mathbf Q/\mathbf Z)\);
exactness and the cogenerator property imply that \(K_x^\bullet\)
is acyclic.

Let \(Z^n=\ker(\delta^n)\).  As \(K_x^\bullet\) starts in degree zero,
acyclicity implies \(Z^0=0\) and gives exact sequences
\[
 0\longrightarrow Z^n\longrightarrow K_x^n
   \xrightarrow{\delta^n}Z^{n+1}\longrightarrow0.
 \tag{SM27.9}
\]
These sequences are pure.  To see this directly, fix a pure-injective
\(E\) and a map \(f:Z^n\to E\), with \(n\ge1\).
The map \(f\delta^{n-1}:K_x^{n-1}\to E\) is a cycle in the
transpose complex.  Its acyclicity gives \(g:K_x^n\to E\) with
\(g\delta^{n-1}=f\delta^{n-1}\).  Surjectivity onto \(Z^n\)
implies \(g|_{Z^n}=f\).  The assertion for \(n=0\) is immediate.
In particular this extension property holds for every \(E=V^+\).
Tensor--Hom adjunction and the fact that
\(\mathbf Q/\mathbf Z\) detects nonzero elements imply that
\(V\otimes_k Z^n\to V\otimes_k K_x^n\) is injective for every
\(V\).  This proves purity of (SM27.9).

\emph{3. Contractibility, using the genuine periodicity input.}
A pure submodule of a flat module is flat.  In the present instance
one can check this without another theorem: for an injection
\(V'\to V\), put \(Z^n\to K_x^n\) in the tensor square.
Both horizontal maps are injective by purity, and the right vertical
map is injective by flatness of \(K_x^n\).  Therefore the left
vertical map is injective.  This is the flatness criterion for
\(Z^n\).

Take the direct sum of (SM27.9) for \(n\ge0\).
Since \(Z^0=0\), the two end sums are canonically isomorphic after
shifting their indices.  We obtain
\[
 0\longrightarrow Z\longrightarrow P\longrightarrow Z
 \longrightarrow0,\qquad
 Z=\bigoplus_{n\ge0}Z^n,\quad P=\bigoplus_{n\ge0}K_x^n.
 \tag{SM27.10}
\]
Here \(Z\) is flat and \(P\) is projective.
The cited Benson--Goodearl Theorem 2.5 makes \(Z\) projective.
Each \(Z^n\), as a direct summand of \(Z\), is consequently
projective.  Thus every sequence (SM27.9) splits.
Choosing splittings identifies
\(K_x^n=Z^n\oplus Z^{n+1}\) and
\(\delta^n(z,w)=(w,0)\); the maps
\(h^n(z,w)=(0,z)\) give a contracting homotopy
(with \(h^0=0\)).  This proves contractibility of \(K_x^\bullet\),
including its unbounded right-hand end.

For every scalar module \(U\), transpose this homotopy.
Then \(\operatorname{Hom}_k(K_x^\bullet,U)\) is contractible.
By (SM27.8),
\[
 L\varphi(C_x(U))=0\qquad\text{for every }x,U.
 \tag{SM27.11}
\]

\emph{4. Products and the specified adjunction.}
Each functor \(A_n\otimes_R-\) commutes with arbitrary products,
because \(A_n\) is a retract of a finite sum of representables,
whose tensor functors are evaluation functors.
Therefore for any family \(X_i\) there is a natural isomorphism of
complexes
\[
 A_\bullet\otimes_R\prod_iX_i
   \cong\prod_i(A_\bullet\otimes_R X_i).
 \tag{SM27.12}
\]
Products of scalar modules are exact, so a product of acyclic
complexes is acyclic.  Formula (SM27.11) now proves the stronger
vanishing assertion for every product of coinduced objects.
Notice that this product argument computes \(L\varphi\) with \(A\);
it does not assume that the originally chosen \(B_n\) is finite.

Let \(P\) be an acyclic projective object of
\(\operatorname{Ch}_{\ge0}(k)\).  BR-13402 writes
\[
 P=\bigoplus_{n\ge0}D^{n+1}(Q_n)
   =\prod_{n\ge0}D^{n+1}(Q_n),
 \tag{SM27.13}
\]
where each \(Q_n\) is projective; the equality of sum and product is
degreewise, since only two summands occur in any chain degree.
The right adjoint \(G\) preserves products.  For every \(y\),
\[
 G(D^{n+1}Q_n)(y)
   =\operatorname{Hom}_{\mathrm{Ch}}(B_\bullet(y),D^{n+1}Q_n)
   =\operatorname{Hom}_k(B_n(y),Q_n).
\]
The latter formula is the disk adjunction, obtained by sending the
degree-\(n\) map freely to \(Q_n\); its degree-\(n+1\) companion is
determined by the differential.  Write the chosen resolution term as
\(B_n=\bigoplus_{a\in I_n}R(x_{n,a},-)\).
It follows naturally as right \(R\)-modules that
\[
 GP\cong\prod_{n\ge0}\prod_{a\in I_n}C_{x_{n,a}}(Q_n).
 \tag{SM27.14}
\]
The preceding product vanishing gives \(L\varphi(GP)=0\).
BR-13403 and F-SM-17 conclude (SM27.3).  This last invocation uses
the canonical unit normalization to handle projective complexes'
degree-zero summands; the coinduction vanishing itself did not need
that normalization. \(\square\)

\paragraph{Exact remaining scope.}
Condition \(\mathrm{FP}_\infty\) is an added hypothesis on \(D\),
not a consequence asserted here of the existence of finitely
presented projective generators of \(\mathcal M\).
It supplies finite-projective duality in (SM27.8) and product
commutation in (SM27.12).  Neither identity has been asserted for an
arbitrary infinite coproduct of representables.
Thus F-SM-27 adds a third proved range to F-SM-24:
arbitrary scalar rings and arbitrary projective scalar morphism
modules, provided \(D\) has this finite-generation property in every
resolution degree.  The case without any of these extra hypotheses
remains the exact test (134.6).

A countable-subring argument has not removed that last condition.
Although a tensor element uses finitely many tensors, its coefficients
in \(GP\) are global linear maps on the modules \(B_n(y)\).
A boundary found after restriction to a smaller ringoid can require
extending such maps back to all of \(B_n(y)\), compatibly with the
ringoid action.  Finite support of the outer tensor expression does
not provide that extension.  No automatic descent or compatible
extension assertion is used in the proved result.

\paragraph{Primary reference.}
Benson, D. J., \& Goodearl, K. R. (2000).
Periodic flat modules, and flat modules for finite groups.
\emph{Pacific Journal of Mathematics, 196}(1), 45--67.
\url{https://doi.org/10.2140/pjm.2000.196.45}.
Theorem 2.5, p.~50; proof pp.~50--51.
Primary publisher PDF:
\url{https://msp.org/pjm/2000/196-1/pjm-v196-n1-p03-s.pdf}.


% END source_cote1348_fp.tex

% BEGIN source_cote1348_compact_unit.tex
\subsection{F-SM-28: a compact unit and compact interval do not force \(D\in\mathrm{FP}_\infty\)}
For every nonzero commutative ring \(k\), there is an example satisfying
the full unit normalization and interval hypotheses of BR-13403--05,
with the following additional properties:
\begin{itemize}
\item the monoidal unit and the interval are finitely presented
projective objects;
\item the category has a countable set of finitely presented projective
generators, and all scalar modules of morphisms between these generators
are either \(0\) or \(k\);
\item the left ringoid module \(D\) representing \(\varphi\) is not
finitely generated, so in particular is not of type
\(\mathrm{FP}_\infty\).
\end{itemize}
The desired connective localization equivalence nevertheless holds
in this example.  The example therefore disproves the proposed
automatic-\(\mathrm{FP}_\infty\) inference, not the original
equivalence.

\paragraph{Complete construction and proof.}
Let \(\mathbb N\) denote the category with one morphism \(i\to j\)
when \(i\le j\).  Put
\[
 \mathcal M=\operatorname{Fun}
     \bigl(\mathbb N,\operatorname{Ch}_{\ge0}(k)\bigr),
 \qquad
 T=\operatorname{const}_{\mathbb N}(k[0]).
 \tag{SM28.1}
\]
This is a cocomplete \(k\)-linear abelian category.  Give it the
pointwise tensor product of chain complexes, with the usual signed
differential; its monoidal unit is \(T\).  Define
\[
 \kappa(U)=\operatorname{const}_{\mathbb N}(U[0]),\qquad
 \Lambda(X)=\operatorname*{colim}_{i\in\mathbb N}X_i,\qquad
 \varphi(X)=H_0\Lambda(X).
 \tag{SM28.2}
\]
The constant functor \(G(C)=\operatorname{const}_{\mathbb N}(C)\)
is right adjoint to \(\Lambda\), by the universal property of colimits.
Both are exact: exactness of \(G\) is pointwise, and filtered colimits
of modules are exact, so \(\Lambda\) is exact in every chain degree.
The latter familiar exactness can also be read elementwise: a kernel
or relation in a filtered colimit is witnessed at one sufficiently
large index; a finite collection of such witnesses has a common
upper index.  Thus kernels and cokernels of short exact sequences
commute with this colimit.

For a scalar module \(U\), a chain map \(X_i\to U[0]\) is precisely
a map \(H_0X_i\to U\).  Taking compatible families gives
\[
 \operatorname{Hom}_{\mathcal M}(X,\kappa U)
   \cong\operatorname{Hom}_k
       \left(\operatorname*{colim}_{i}H_0X_i,U\right)
   =\operatorname{Hom}_k(\varphi X,U).
 \tag{SM28.3}
\]
Hence \(\varphi\dashv\kappa\).
The scalar functor \(\kappa\) is exact, preserves all inverse limits,
and is compatible with the tensor products.  It also has the required
right adjoint: since \(0\) is initial in \(\mathbb N\), a map
\(T\to X\) is determined by an arbitrary element of \((X_0)_0\);
thus \(\operatorname{Hom}_{\mathcal M}(T,X)=(X_0)_0\), and
scalar tensor--Hom gives \(\kappa\dashv\operatorname{Hom}_{\mathcal M}(T,-)\).

For \(a\ge1\), let \(C_a=D^a(k)\) be the disk with identity
differential from degree \(a\) to degree \(a-1\), and set
\(C_0=k[0]\).  For each \(i,a\ge0\), define
\[
 (P_{i,a})_j=
 \begin{cases}0,&j<i,\\ C_a,&j\ge i,\end{cases}
 \tag{SM28.4}
\]
with identity transition maps after \(i\).  These free extensions
satisfy
\[
 \operatorname{Hom}_{\mathcal M}(P_{i,a},X)
   =\operatorname{Hom}_{\mathrm{Ch}}(C_a,X_i)
   =(X_i)_a.
 \tag{SM28.5}
\]
The formula for \(a=0\) is the sphere adjunction and the formula for
\(a\ge1\) is the left disk adjunction.  Evaluation is exact and
commutes with filtered colimits.  Therefore every \(P_{i,a}\) is
projective and finitely presented.  They generate: for every element
in every \((X_i)_a\), the corresponding map \(P_{i,a}\to X\) contributes
that element, and the sum of these maps is a degreewise epimorphism.
In particular
\[
 T=P_{0,0}
\]
is finitely presented and projective.

We have
\[
 \Lambda(P_{i,0})=k[0],\qquad
 \Lambda(P_{i,a})=D^a(k)\quad(a\ge1).
\]
Thus \(\Lambda\) has zero positive homology on projectives, first on
the generators and then on their coproducts and retracts.
The finite-diagonal proof BR-13401 gives
\[
 L\varphi(X)\simeq\Lambda(X).
 \tag{SM28.6}
\]
In particular, \(L\varphi(T)\to k[0]\) is the canonical identity
normalization.

Take
\[
 J=P_{0,1}=\operatorname{const}_{\mathbb N}(D^1(k)),
 \qquad
 \alpha:T\longrightarrow J
\]
to be the identity inclusion in degree zero.  The interval \(J\)
is also finitely presented and projective.
For any \(X\), at every diagram index the tensor complex has
\[
 ((J\otimes X)_i)_n=(X_i)_n\oplus(X_i)_{n-1},\qquad
 d(x,y)=(d_Xx+y,-d_Xy),
 \tag{SM28.7}
\]
where the second summand is zero for \(n=0\).
The map \(\alpha_X\) is \(x\mapsto(x,0)\), hence monic.
The maps \(h_n(x,y)=(0,x)\) satisfy \(dh+hd=1\).
They commute with the diagram structure maps, so give a natural
contraction in \(\mathcal M\).  Applying the exact colimit functor
preserves this contraction.  Equation (SM28.6) therefore gives
\(L\varphi(J\otimes X)=0\) in every degree.
The scalar compatibility (SM27.1) also holds directly.
This verifies the full interval, including its degree-zero clause.

Now take \(R\) to be the full small \(k\)-linear category of the
objects \(P_{i,a}\), with
\[
 R((i,a),(j,b))=
 \operatorname{Hom}_{\mathcal M}(P_{i,a},P_{j,b}).
\]
Restricted Yoneda identifies \(\mathcal M\) with right \(R\)-modules.
For completeness, any object has a presentation by sums of the
\(P_{i,a}\), as just proved; projectivity and finite presentation make
restricted Yoneda exact and make it preserve those sums.  Such a
presentation proves full faithfulness, and a presentation of any
right \(R\)-module by representables gives essential surjectivity
by taking the corresponding cokernel in \(\mathcal M\).

By (SM28.5), a scalar morphism module is zero when \(i<j\), and
when \(i\ge j\) it is
\(\operatorname{Hom}_{\mathrm{Ch}}(C_a,C_b)=(C_b)_a\).
For \(b=0\) this is \(k\) only when \(a=0\).
For \(b\ge1\) it is \(k\) precisely when \(a=b\) or \(a=b-1\).
Thus each scalar morphism module is \(0\) or \(k\).
All hypotheses required for the exact-adjunction construction
BR-13405 are satisfied.

The left \(R\)-module \(D\) representing \(\varphi\) is
\[
 D(i,a)=\varphi(P_{i,a})=
 \begin{cases}k,&a=0,\\0,&a\ge1.\end{cases}
 \tag{SM28.8}
\]
Furthermore
\[
 R((i,0),(j,0))=
 \begin{cases}k,&j\le i,\\0,&j>i,\end{cases}
 \tag{SM28.9}
\]
and the nonzero arrows act on (SM28.8) by their scalar coefficients.
If \(D\) were finitely generated, it would be a quotient of a finite
sum of left representables.  A map from a representable based at
\((i,a)\) to \(D\) is specified by an element of \(D(i,a)\).
For \(a\ge1\) it is zero.  For the remaining finitely many indices
\((i,0)\), choose an upper bound \(N\).  At every \((j,0)\) with
\(j>N\), all those representables evaluate to zero by (SM28.9).
Their image cannot equal \(D(j,0)=k\ne0\).
This contradiction proves that \(D\) is not finitely generated.

Finally, this construction still models the desired connective
derived category.  The counit
\(\Lambda G(C)\to C\) is an identity for every complex \(C\).
The unit \(X\to G\Lambda X\) becomes an identity after \(\Lambda\),
hence belongs to \(W=\Lambda^{-1}(\mathrm{qis})\).
Also \(G\) sends quasi-isomorphisms to \(W\), since \(\Lambda G\)
is the identity.  The ordinary adjunction consequently descends
with invertible unit and counit, and gives
\[
 W^{-1}\mathcal M\simeq D_{\ge0}(k).
\]
For an adjunction formed from any representable resolution as in
BR-13405, the same conclusion, including the raw counit criterion,
also follows from F-SM-24's finitely generated scalar-morphism case.
Thus the failure of finite generation for \(D\) creates no
counterexample to the original claim in this example. \(\square\)

\paragraph{Source scope and consequence for the search.}
A fresh visual reading of manuscript footers 47 and 48 found the
conditions (a)--(d), followed by the corollary's finitely presented
projective generators and its scalar-projectivity marginal condition.
Those lines do not state finite presentation of the monoidal unit.
The example above supplies that stronger property anyway, together
with finite presentation and projectivity of the interval itself.
Therefore neither adding a compact unit nor adding this compact
interval makes F-SM-27 an automatic solution of the full general-ring
question.  This is an explicit construction, not an inference from
an unread part of the manuscript.


% END source_cote1348_compact_unit.tex

% BEGIN root_methods.tex
\section{The geometric method: a usable arithmetic determinant reduction}
\label{sec:root-method}

There are two distinct multiplicities in the prompt.  An automorphic
restriction multiplicity is the dimension of a space of equivariant
functionals.  A Stepanov multiplicity is an order of vanishing or an
integral divisibility order.  A theorem relating them must include an
integral lattice, a valuation, and compatible maps; a complex
representation-theoretic dimension alone does not supply those data.
The following finite-dimensional form isolates the actual numerical
inequality needed by the arithmetic determinant method.

\paragraph{Finding BR-DET-01 (divisibility versus archimedean size).}
Let $K$ be a number field with ring of integers $\mathcal O_K$.
Let $A$ be an $r\times r$ matrix over $\mathcal O_K$.  Suppose its
$j$th column has every entry in a nonzero integral ideal $\mathfrak a_j$,
and for each embedding $\sigma:K\hookrightarrow\mathbb C$ the
Euclidean norm of that column is at most $C_{\sigma,j}>0$.
Count all embeddings, including both members of a complex pair.
If
\[
 \sum_{j=1}^r\log N\mathfrak a_j
   >\sum_{\sigma:K\hookrightarrow\mathbb C}
                     \sum_{j=1}^r\log C_{\sigma,j},             \tag{D.1}
\]
then $\det A=0$.  If $\det A\ne0$, the reverse weak inequality
is necessary.

\begin{proof}
Expansion by permutations shows
$\det A\in\mathfrak a_1\cdots\mathfrak a_r$.
For $a=\det A\ne0$, inclusion of integral ideals
$(a)\subseteq\mathfrak a_1\cdots\mathfrak a_r$ gives
\[
 \prod_jN\mathfrak a_j
 \leq N(a)=|N_{K/\mathbb Q}(a)|
 =\prod_\sigma|\sigma(a)|.
\]
Hadamard's determinant inequality, applied separately at every
embedding, gives
$|\sigma(a)|\leq\prod_j C_{\sigma,j}$.
Multiplying and taking logarithms proves the reverse weak
inequality.  Its contradiction with (D.1) proves vanishing.
\end{proof}

\paragraph{Finding BR-DET-02 (the exact accumulation law for iterated lifts).}
Let $L_j$ be free finite-rank $\mathcal O_K$-lattices with Hermitian
norms at the embeddings, and let
$\theta_j:L_{j-1}\to L_j$ be $\mathcal O_K$-linear maps.
Suppose
\[
 \theta_j(L_{j-1})\subseteq\mathfrak b_jL_j,
 \qquad \|\theta_{j,\sigma}\|\leq c_{j,\sigma}.
\]
Then
\[
 \theta_m\cdots\theta_1(L_0)
   \subseteq(\mathfrak b_1\cdots\mathfrak b_m)L_m,
 \qquad
 \|\theta_{m,\sigma}\cdots\theta_{1,\sigma}\|
   \leq\prod_{j=1}^m c_{j,\sigma}.                            \tag{D.2}
\]
After an integral evaluation map $e:L_m\to\mathcal O_K^r$, columns
obtained from chosen vectors $v_k\in L_0$ obey BR-DET-01 with
the bound
\[
 C_{\sigma,k}=\|e_\sigma\|
       \left(\prod_{j=1}^m c_{j,\sigma}\right)\|v_{k,\sigma}\|,
\]
or any strictly positive upper bound if a column is zero.
Thus a valid application of ``one order
per lift'' requires the ideal-inclusion assertion in (D.2), and
a valid use of arbitrarily many lifts requires their accumulated
logarithmic gain to exceed their accumulated norm cost.

\begin{proof}
For any ideal $\mathfrak b$ and any $\mathcal O_K$-linear map
$\theta$, one has $\theta(\mathfrak bL)=\mathfrak b\theta(L)$:
both sides consist of finite sums of coefficients from
$\mathfrak b$ times images of elements of $L$.
Induction proves the lattice inclusion.  Submultiplicativity of
operator norms proves the other inequality.  An integral
evaluation map preserves the ideal inclusions; applying the
composite operator to $v_k$ bounds its norm by the product of
the operator norms and $\|v_{k,\sigma}\|$. The preceding
finding applies to the resulting matrix.
\end{proof}

These findings are elementary reductions, not a replacement for
Alp\"oge's arithmetic estimates.  They explain exactly which
content the four-page method would have to export in order to
transfer to another motivic family.  CM-period divisibility and
the integral theta structures supply the ideals.  A Rallis
formula and period estimates may control the archimedean norms.
The Iwasawa averaging and relative trace formula must preserve
the necessary independence and integrality while improving the
cost.  The order-$m$ cancellation in the sketch matters because
it changes the sign of the strict inequality (D.1).  Neither
that cancellation nor a uniform tower estimate is a formal
consequence of the five exported RH declarations.

The exact external comparison points are Gan--Gross--Prasad
(2009), Conjectures 17.1 and 17.3 (local distinguished members),
Theorem 14.1 and the reductions in Theorems 15.1 and 16.1
(multiplicity-one statements in their stated cases), and
Conjecture 24.1 (global periods and central values).  These are
the statuses and numbering of the supplied work; we do not call
all of its conjectures theorems from its presence in the ZIP.
Gan--Qiu--Takeda (2014), Theorem 11.4, is a precise Rallis
inner-product formula, with range, cuspidality, and local
conditions.  Applying it repeatedly requires checking those
conditions at every step.  Its special values generally belong
to higher-rank automorphic $L$-functions, outside the stated
Dirichlet scope.  GRH also allows zeros at the centre; it does
not assert that a central period is nonzero.

\subsection{What this does for crystals and for towers}

For a rank-one abelian Artin--Tate object unramified at a
realization prime $p$, its crystalline realization has Frobenius
eigenvalue $p^{-n}$ times the unramified finite character, in the
convention of BR-AT-01.  Its Newton slope is $-n$ and its
Hodge--Tate weight is $-n$.  The period $t$ with $\varphi(t)=pt$ gives this normalization
(Berger, 2004, \S II.3.1 and II.4.3); an unramified finite twist
changes the eigenvalue only by a root of unity.  This local rank-one
compatibility is already known; GRH is a global assertion about the continued
Euler product.  For a ramified finite character one must first
state a potentially crystalline or other appropriate local
realization, rather than silently use the unramified formula.
For the CM curve $y^2=x^3+1$, the separate argument below does
give a genuine new use of Dirichlet GRH: it controls the supply
of inert, trace-zero primes used by the divisibility mechanism.
Its two-dimensional $H^1$ remains outside the diagonal
Artin--Tate Hodge types.

In a cyclotomic tower of conductor $p^a$, the finite Frobenius
recognition theorem below applies at every fixed level, with
$Q=\operatorname{lcm}(p^a,4)$ and bound
$(\varphi(Q)\log Q)^2$.  This is an explicit effective
consequence for the tower's finite abelian layers.  An inverse
limit of integral Selmer groups, a statement about its
characteristic ideal, or a Vandiver-type vanishing assertion
requires additional integral information.  BR-MIX-01 supplies
an explicit example of information lost on passing to the
weight-graded $L$-function; it prevents substituting a
semisimple spectral calculation for that missing integral
calculation.

\subsection{Arithmetic duality is a precise analogy, not a missing construction}

Artin--Verdier duality is a pairing for constructible torsion
\'etale coefficients on number rings, with compact support and
the real-place conventions.  A standard exact source is Milne
(2006), \emph{Arithmetic Duality Theorems}, Chapter II,
Theorem 3.1(b), p.~177.  Its degree-three duality explains the arithmetic
three-dimensional analogy.  It does not equip
$\operatorname{Spec}\mathbb Z$ with a positive Hilbert-space
cohomology whose determinant is $\xi$.  In fact the infinitely
many zeros of $\xi$ already exclude an ordinary characteristic
polynomial of a single finite-dimensional complex operator.
If such a finite-dimensional operator $A$ existed, its
determinant $\det(sI-A)$ would be a polynomial with at most
$\dim A$ zeros, contradicting the known zero count.
An infinite-dimensional space, a trace class transformation
of the variable as in BR-TR-01, or an appropriate regularized
determinant is additional structure.  This is the precise
analytic point at which the arithmetic-duality analogy and a
Hilbert--P\'olya or Deninger construction must be connected by
a new theorem.


% END root_methods.tex

% BEGIN consequences_criteria.tex
\section{Consequences and criteria under the adopted RH and Dirichlet
GRH}\label{consequences-and-criteria-under-the-adopted-rh-and-dirichlet-grh}

Research module \nolinkurl{consequences_criteria}. Author/checker: GPT-6
Astra Pro; effort setting not exposed. Date: 20 November 2026 UTC. The
mathematical input is the RH and Dirichlet-GRH return adopted by the
task. This module does not claim an independent proof or a kernel check
of that return. Known conditional results are used with their actual
hypotheses; an implication asserted here is relative to the adopted
input, and hence has no remaining RH condition within the task's adopted
theory.

The propositions labelled CC-F are deductions and counterexamples
established in this run; no claim of historical priority or exhaustive
novelty is made.

The main distinction is between \textbf{theorems whose hypotheses are
now discharged}, \textbf{equivalent criteria whose asserted positivity
or approximation now follows}, and \textbf{additional geometric or
statistical structures that an equivalence does not supply}. An
input-scope obstruction does not establish logical independence: it
identifies the extra hypothesis in the available theorem and prevents an
unjustified application of that theorem.

\subsection{1. Highest-priority
outcomes}\label{highest-priority-outcomes}

\begin{longtable}[]{@{}
  >{\raggedright\arraybackslash}p{(\columnwidth - 6\tabcolsep) * \real{0.2500}}
  >{\raggedright\arraybackslash}p{(\columnwidth - 6\tabcolsep) * \real{0.2500}}
  >{\raggedright\arraybackslash}p{(\columnwidth - 6\tabcolsep) * \real{0.2500}}
  >{\raggedright\arraybackslash}p{(\columnwidth - 6\tabcolsep) * \real{0.2500}}@{}}
\toprule\noalign{}
\begin{minipage}[b]{\linewidth}\raggedright
ID
\end{minipage} & \begin{minipage}[b]{\linewidth}\raggedright
Outcome
\end{minipage} & \begin{minipage}[b]{\linewidth}\raggedright
Exact operative source
\end{minipage} & \begin{minipage}[b]{\linewidth}\raggedright
What is discharged or remains
\end{minipage} \\
\midrule\noalign{}
\endhead
\bottomrule\noalign{}
\endlastfoot
CC-F01 & Every smooth proper Borger \(\Lambda\)-scheme over
\(\mathbb Z[1/M]\) has cohomological Dirichlet-completed factors
satisfying RH on their weight-shifted lines & Borger (2009), Theorem
0.1, p.~3; Theorem 5.3, p.~17; Corollary 6.2 and (6.2.1), p.~20 &
Dirichlet GRH suffices. This is a statement about realizations and their
\(L\)-data; no Chow-motive decomposition or cycle-surjectivity assertion
is inferred \\
CC-F01b & The rational compactly supported cohomology of every separated
finite-type \(\Lambda\)-scheme over \(\mathbb Z\) actually splits into
finite abelian characters and Tate twists & Borger (2009), §6 setup and
Corollary 6.2, equation (6.2.1), followed by the weight-space descent
proved here & No torsion-freeness hypothesis or semisimplification is
inserted. Dirichlet GRH supplies the resulting primitive factors;
bad-fiber model zeta functions require separate treatment \\
CC-C01 & Effective Chebotarev with the Lagarias--Odlyzko error term for
\(L/K\) whenever the top field \(L/\mathbb Q\) is abelian &
Lagarias--Odlyzko (1977), Theorem 1.1, equation (1.3) & Their exact
analytic hypothesis is RH for \(\zeta_L\), discharged by Dirichlet
factorization. Merely assuming \(L/K\) abelian does not do this \\
CC-C02 & A prime in every reduced class \(a\pmod q\), \(q>3\), below
\((\varphi(q)\log q)^2\) & Lamzouri--Li--Soundararajan (LLS), corrected
arXiv v3 (2016), Corollary 1.2, p.~3 & Dirichlet GRH suffices; a
representative for every class supports finite character tests,
including higher-rank semisimple representations \\
CC-F02 & Siegel's lower bound for primitive real Dirichlet \(L(1,\chi)\)
has a computable constant for every positive rational exponent & LLS,
Theorem 1.5, pp.~3--4, plus the full effective deduction below &
Dirichlet GRH suffices. This is the analytic Siegel theorem; it is not
an asserted effective integral-point theorem \\
CC-F03 & The Hilbert-class-field hypothesis is a real scope boundary
even for an imaginary quadratic base & Explicit splitting field of
\(x^3-x-1\), proved below & The required ideal-class-character
\(L\)-function can be an irreducible degree-two Artin \(L\)-function
over \(\mathbb Q\), outside Dirichlet factorization \\
CC-P01--CC-P04 & Weil positivity, Li positivity, the strong
Nyman--Beurling approximation, and the bare spectral realization all
hold in their precise senses & Exact locators in §4 & Stronger
constructions, support ranges, approximation asymptotics, and geometric
functoriality are separate questions \\
CC-P07 & The concrete de Branges sufficient positivity conditions
analyzed by Conrey--Li are false and remain false & Conrey--Li (2000),
§§3.1--3.2 and §4 Remark & A sufficient condition is not automatically
an equivalent criterion. RH does not imply these conditions \\
\end{longtable}

\subsection{2. The Borger bridge, with a full
proof}\label{the-borger-bridge-with-a-full-proof}

\subsubsection{\texorpdfstring{CC-F01. Smooth proper
\(\Lambda\)-cohomology satisfies its shifted analytic
RH}{CC-F01. Smooth proper \textbackslash Lambda-cohomology satisfies its shifted analytic RH}}\label{cc-f01.-smooth-proper-lambda-cohomology-satisfies-its-shifted-analytic-rh}

\textbf{Statement.} Let \(X\) be smooth and proper over
\(\mathbb Z[1/M]\), of finite type, and equipped with a
\(\Lambda\)-structure in Borger's sense. Fix a prime \(\ell\), an
algebraic closure \(\overline{\mathbb Q}_\ell\), and an identification
of the algebraic roots of unity in it with complex roots of unity. Put
\[
 V^j=H^j_{\text{ét}}(X_{\overline{\mathbb Q}},\mathbb Q_\ell).
\] Then \(V^{2m+1}=0\), and there is a finite family of primitive
Dirichlet characters, counted with nonnegative integer multiplicities
\(r_{m,\chi}\), such that \[
 V^{2m}\otimes\overline{\mathbb Q}_\ell
 \simeq\bigoplus_\chi(\chi_\ell\otimes\overline{\mathbb Q}_\ell(-m))^{\oplus r_{m,\chi}}.
 \tag{CC-F01.1}
\] The convention for \(\chi_\ell\) is that its value on geometric
Frobenius at an unramified prime \(p\) is \(\chi(p)\). Define the
primitive Dirichlet completion of the cohomological factor by \[
 L^{\mathrm{Dir}}_{2m}(X,s)=\prod_\chi L(s-m,\chi)^{r_{m,\chi}}.
 \tag{CC-F01.2}
\] Every nontrivial zero of this factor lies on \(\Re s=m+\tfrac12\).
The product has a pole at \(s=m+1\) of order precisely \(r_{m,1}\) and
has no other poles. The corresponding completed factors have the usual
gamma and conductor factors of these primitive Dirichlet characters. At
all sufficiently good primes, (CC-F01.2) has exactly the cohomological
Euler factor of \(X\).

If \(S\) contains \(\ell\), the bad primes of \(X\), and every character
conductor prime, then the incomplete factor satisfies the exact identity
\[
 L^S_{2m}(X,s)
 =L^{\mathrm{Dir}}_{2m}(X,s)
   \prod_{p\in S}\prod_\chi
       (1-\chi(p)p^{m-s})^{r_{m,\chi}},
 \tag{CC-F01.3}
\] where a primitive character is assigned value zero at a prime
dividing its conductor. The explicitly deleted Euler factors can
introduce zeros on \(\Re s=m\). They are excluded from the
nontrivial-zero assertion. It would be incorrect to assert that every
zero of an incomplete Euler product lies on the shifted critical line.

\textbf{Proof.} Borger's Theorem 0.1(c) says that the Hodge
decomposition of \(H^j(X(\mathbb C),\mathbb C)\) has no components of
type \((a,b)\) with \(a\ne b\). Smoothness and properness give purity of
this Hodge structure of weight \(j\); its nonzero components therefore
satisfy \(a+b=j\). Consequently odd cohomology vanishes. Comparison with
étale cohomology proves \(V^{2m+1}=0\).

Borger's Corollary 6.2, equation (6.2.1), supplies a cyclotomic field
\(F=\mathbb Q(\zeta_N)\) over which the representation is a sum of Tate
twists, with multiplicities the diagonal Hodge numbers. In degree \(2m\)
the only possible Hodge type is \((m,m)\). Therefore \[
 V^{2m}|_{G_F}\simeq\mathbb Q_\ell(-m)^{b_{2m}}.
\] The twist \(W_m=V^{2m}(m)\) is thus trivial on \(G_F\) and factors
through the finite abelian group \(\mathrm{Gal}(F/\mathbb Q)\).
Characteristic-zero representations of a finite group are semisimple, by
averaging a projection over the group. Every irreducible representation
of a finite abelian group over an algebraically closed field is
one-dimensional: a common eigenspace of the commuting operators is
invariant, and irreducibility forces it to be the whole space. It
follows that \(W_m\) splits into finite-order characters. Characters of
\(\mathrm{Gal}(\mathbb Q(\zeta_N)/\mathbb Q)\) are Dirichlet characters;
replacing an imprimitive character by its primitive character changes
the modulus description, not the Galois character. This proves
(CC-F01.1).

On \(\mathbb Q_\ell(-m)\), geometric Frobenius has eigenvalue \(p^m\).
Hence the good-prime determinant is \[
 \det(1-\mathrm{Frob}_p p^{-s}\mid V^{2m})
   =\prod_\chi(1-\chi(p)p^{m-s})^{r_{m,\chi}}.
\] Its inverse is the Euler factor of (CC-F01.2). Euler products agree
in a right half-plane of absolute convergence, and their analytic
continuations consequently agree wherever defined. Removing the finitely
many factors in \(S\) gives (CC-F01.3).

By the adopted Dirichlet GRH, every nontrivial zero of \(L(s-m,\chi)\)
has real part \(m+1/2\). A finite product with nonnegative
multiplicities has no other nontrivial zeros. Only the primitive trivial
character contributes a pole, and it contributes \(\zeta(s-m)\). At
\(s=m+1\), all nonprincipal Dirichlet factors are nonzero by the
classical nonvanishing of \(L(1,\chi)\); thus the pole order is exactly
\(r_{m,1}\). Finally, if \(\chi(p)\ne0\), it is a root of unity, so
\(1-\chi(p)p^{m-s}=0\) implies \(p^{m-\Re s}=1\), or \(\Re s=m\). If
\(\chi(p)=0\), the factor is identically one. This proves every
assertion. \(\square\)

\textbf{Scope and significance.} This supplies analytic RH to an entire
geometric class identified inside the supplied \(\mathbf F_1\)
literature. The argument proves a decomposition of realizations and
their Dirichlet-completed \(L\)-functions. It does not identify a Chow
motive with a direct sum of Artin--Tate Chow motives, prove a Hodge or
Tate cycle map surjective, or construct the infinite-dimensional
cohomology envisaged in Borger's introduction, p.~3. Borger's Theorem
6.1 applies much more generally to separated finite-type
\(\Lambda\)-schemes; Corollary 6.2 gives their rational compactly
supported cohomology after cyclotomic restriction. For arbitrary mixed
objects one must separately control the actual extensions and bad-prime
inertia. CC-F01b proves that Borger's particular rational cohomology
class actually splits, so no such extension is being silently discarded
there; its asserted arithmetic-model comparison is restricted to
good-prime data.

\subsubsection{\texorpdfstring{CC-F01b. Actual rational splitting for
every separated finite-type
\(\Lambda\)-scheme}{CC-F01b. Actual rational splitting for every separated finite-type \textbackslash Lambda-scheme}}\label{cc-f01b.-actual-rational-splitting-for-every-separated-finite-type-lambda-scheme}

\textbf{Statement.} Let \(X\) be a separated \(\Lambda\)-scheme of
finite type over \(\mathbb Z\), let \(n\ge0\), and let \(\ell\) be
prime. For \[
 V=H^n_{\text{ét},c}(X_{\overline{\mathbb Q}},\mathbb Q_\ell),
\] there is an actual isomorphism, after extension of coefficient field,
\[
 V\otimes\overline{\mathbb Q}_\ell
 \simeq\bigoplus_{m,\chi}
 (\chi_\ell\otimes\overline{\mathbb Q}_\ell(-m))^{\oplus r_{m,n,\chi}}.
 \tag{CC-F01b.1}
\] The characters have finite abelian image, the multiplicities are
nonnegative integers, and only finitely many occur. In particular this
rational representation is semisimple. This assertion is about
cohomology with \(\mathbb Q_\ell\) coefficients: it does not assume that
integral or finite-coefficient cohomology is torsion-free, and asserts
neither integral splitting nor absence of torsion. The primitive
completion of its good-prime \(L\)-data is \[
 L_n^{\mathrm{Dir}}(X,s)=\prod_{m,\chi}L(s-m,\chi)^{r_{m,n,\chi}},
 \tag{CC-F01b.2}
\] whose nontrivial zeros lie on the lines \(\Re s=m+1/2\) for the
weights present. Deleted Euler factors are explicit finite products and
can add zeros on \(\Re s=m\). For a nonproper or singular model, no
equality with the zeta function of its bad fibers is asserted.

\textbf{Proof.} In the setup of Borger §6, \(X\) is exactly a separated
finite-type \(\Lambda\)-scheme over \(\mathbb Z\). Corollary 6.2, p.~20,
equation (6.2.1), gives an isomorphism of the actual representation \[
 V|_H\simeq\bigoplus_m\mathbb Q_\ell(-m)^{r_{m,n}},
 \qquad H=G_{\mathbb Q(\zeta_N)}.
\] The source does not insert a semisimplification superscript here. Let
\(\epsilon_\ell:G_\mathbb Q\to\mathbb Z_\ell^\times\) be the arithmetic
cyclotomic character, so that \(\mathbb Q_\ell(-m)\) carries
\(\epsilon_\ell^{-m}\). Its restriction to \(H\) has infinite image:
\(H\) has finite index in \(G_\mathbb Q\), while the image of the full
cyclotomic character is infinite, and the image of a finite-index
subgroup has finite index in that image. The powers
\(\epsilon_\ell^{-m}|_H\) are distinct for distinct integers \(m\),
since the torsion subgroup of \(\mathbb Z_\ell^\times\) is finite.

Define the corresponding isotypic subspace \[
 V_m=\{v\in V:\rho(h)v=\epsilon_\ell(h)^{-m}v\text{ for every }h\in H\}.
\] The given direct-sum decomposition gives \(V=\bigoplus_m V_m\). Since
\(H\) is normal and \(\epsilon_\ell\) is a character of the whole group,
for \(g\in G_\mathbb Q\), \(h\in H\), and \(v\in V_m\), \[
 \rho(h)\rho(g)v
 =\rho(g)\rho(g^{-1}hg)v
 =\epsilon_\ell(g^{-1}hg)^{-m}\rho(g)v
 =\epsilon_\ell(h)^{-m}\rho(g)v.
\] Thus every \(V_m\) is \(G_\mathbb Q\)-stable. The twist \(V_m(m)\) is
trivial on \(H\), hence factors through the finite abelian quotient
\(G_\mathbb Q/H\). By characteristic-zero semisimplicity and the
character decomposition of a finite abelian group, it splits over
\(\overline{\mathbb Q}_\ell\) into one-dimensional finite-order
characters. Undoing the twists proves (CC-F01b.1). The determinant
computation and primitive-character identification in CC-F01 apply to
every summand and give (CC-F01b.2) at good primes. The adopted Dirichlet
GRH and the same elementary computation of zeros of deleted Euler
factors prove the analytic assertions. No passage through a nonsplit
extension or an exactness claim for inertia invariants occurs.
\(\square\)

The nontrivial zero contributions in (CC-F01b.2) cannot be canceled by
poles of other factors: their real parts are half-integers, whereas the
only possible poles occur at the integers \(m+1\). Trivial zeros and
possible cancellation among trivial zeros and poles are separate from
this assertion. The importance of this refinement is that Borger's
actual rational splitting eliminates an extension obstruction inside
this specific geometric class. General mixed Artin--Tate objects outside
this class can still have nontrivial extensions. Borger's Remark 6.3
explicitly cautions against an integral-coefficient analogue.

\subsubsection{CC-F01a. The hexagonal CM curve does not itself furnish
Borger's finite-type
descent}\label{cc-f01a.-the-hexagonal-cm-curve-does-not-itself-furnish-borgers-finite-type-descent}

Borger (2009), §2.9, pp.~10--11, proves that a connected smooth proper
curve of genus at least one over \(\mathbb Q\), with a smooth proper
model over some \(\mathbb Z[1/M]\), has no \(\Lambda\)-structure. In
particular the genus-one CM curve with hexagonal period lattice in the
adopted note cannot be identified with a smooth proper
\(\Lambda\)-scheme over \(\mathbb Z[1/M]\). CM endomorphisms at split
primes do not supply the commuting Frobenius lifts at all good primes
required in Borger's definition. This refutes that particular
identification; it does not refute the note's use of a CM curve or its
separate reference to \(\mathbf F_1\).

\subsection{3. Classical consequences and exact hypothesis
map}\label{classical-consequences-and-exact-hypothesis-map}

\subsubsection{CC-C01. Lagarias--Odlyzko for abelian top
fields}\label{cc-c01.-lagariasodlyzko-for-abelian-top-fields}

For a finite Galois extension \(L/K\), let \(G=\mathrm{Gal}(L/K)\), let
\(C\) be a conjugacy class, let \(\pi_C(x,L/K)\) count unramified prime
ideals of \(K\) of norm at most \(x\) with Frobenius in \(C\), and let
\(n_L=[L:\mathbb Q]\), \(\Delta_L=|\mathrm{disc}(L)|\). Lagarias and
Odlyzko (1977), Theorem 1.1, equation (1.3), require \textbf{RH for
\(\zeta_L\)} and conclude, for \(x>2\), \[
 \left|\pi_C(x,L/K)-\frac{|C|}{|G|}\operatorname{Li}(x)\right|
 \le c_1\left\{\frac{|C|}{|G|}\sqrt{x}\log(\Delta_Lx^{n_L})+\log\Delta_L\right\},
\] with an absolute effectively computable \(c_1\). The inspected
witness is the full retypeset primary text, Theorem 1.1 on its p.~4; the
original article occupies pp.~409--464. Their Corollary 1.2 gives an
effective least-prime norm bound
\(c_2(\log\Delta_L)^2(\log\log\Delta_L)^4\); later sharpenings should
not be silently substituted for this displayed corollary.

If \(L/\mathbb Q\) is abelian, its Dedekind zeta function is a product
of primitive Dirichlet \(L\)-functions, including \(\zeta\); the adopted
statements therefore discharge the analytic hypothesis. Any intermediate
\(K\) is allowed. If only \(L/K\) is abelian, this factorization into
Dirichlet \(L\)-functions over \(\mathbb Q\) is unavailable in general.
For unrestricted \(L\), the missing input is RH for that top-field
Dedekind zeta function, not Artin holomorphy as an additional hypothesis
of Theorem 1.1.

\subsubsection{CC-C02. Least primes and finite character
tests}\label{cc-c02.-least-primes-and-finite-character-tests}

LLS, corrected arXiv v3 (2016), Corollary 1.2, p.~3, gives \[
 P(a,q)\le (\varphi(q)\log q)^2\qquad(q>3,\ (a,q)=1)
\] under Dirichlet GRH. The small moduli are finite explicit cases.
Theorem 1.1 gives stronger subgroup-escape bounds for \(q\ge3000\);
Corollary 1.1 gives least quadratic nonresidue \(<(\log q)^2\) for prime
\(q\ge5\). These need only Dirichlet GRH. The corrected v3 constants
must be used: Theorem 1.2 has \(\alpha(2)=0.8\), \(\alpha(3)=0.7\), and
\(\alpha(h)=0.66\) for \(h>3\). Earlier-version constants are not
interchangeable.

For a rank-one character, its values on group generators determine it.
For an arbitrary representation, the trace at the chosen generators need
not determine the character. The safe finite certificate tests every
residue class (or every class of the actual quotient), as in the
separate archive module.

The earlier explicit bound is also now checked against its original
source: Bach and Sorenson (1996), Theorem 5.3 (ERH), printed p.~1731,
gives \(P(a,q)<2(q\log q)^2\) for positive moduli \(q\ge2\) and
\((a,q)=1\). Their §2, p.~1720, specifies RH for \(\zeta_E\) and
\(\zeta_{\mathbb Q}\); in this application \(E=\mathbb Q(\zeta_q)\), so
the adopted input suffices. Their Corollary 3.5, pp.~1722--1723,
concerns a prime in every ideal class of a quadratic field and
explicitly sets \(E\) equal to its Hilbert class field. That application
has the wider top-field hypothesis discussed in CC-F03, even though its
base is quadratic. The stronger LLS bound is the one used for the finite
certificate.

\subsubsection{CC-C03. Miller's deterministic
algorithm}\label{cc-c03.-millers-deterministic-algorithm}

Miller (1976), Theorem 2, p.~301, proves under his ERH a deterministic
primality algorithm taking \(O(b^4\log\log b)\) Turing-machine steps for
an integer of binary length \(b\). His Appendix, pp.~316--317,
formulates the relevant ERH for Dirichlet \(L\)-functions. The adopted
Dirichlet GRH therefore discharges it. This validates that algorithm and
its complexity analysis; deterministic primality testing was already
unconditional by other algorithms, so one must not describe the
consequence as the first proof that primality belongs to P. The exact
bound on p.~301 was inspected visually as well as in extracted text.

\subsubsection{CC-C04. Bach's explicit witnesses and the class-group
boundary}\label{cc-c04.-bachs-explicit-witnesses-and-the-class-group-boundary}

Bach (1990), \emph{Explicit bounds for primality testing and related
problems}, has distinct applications. Its abstract gives a Miller or
Solovay--Strassen compositeness witness at most \(2(\log n)^2\), under
the relevant ERH. LLS §1.1, p.~1, attributes to Bach a prime outside a
proper subgroup below \(2(\log q)^2\), but that prime is allowed to
divide \(q\); imposing \(p\nmid q\) changes that cited bound to
\(3(\log q)^2\). This distinction must be preserved when applying a
generator theorem.

Belabas, Diaz y Diaz, and Friedman (2008), \emph{Small generators of the
ideal class group}, introduction pp.~2--3 of the inspected author
preprint, explicitly distinguish the two zero sums used by Bach: one for
a nontrivial class character and one for the Dedekind zeta function.
Their own Theorem 1.1, preprint p.~3, assumes RH for every nontrivial
ideal-class-character \(L\)-function. This supplies a directly inspected
primary locator for the hypothesis boundary.

The commonly cited class-group corollary is generation by prime ideals
of norm at most \(12(\log|\Delta_K|)^2\), under GRH for the relevant
ideal-class-character Hecke \(L\)-functions. This is identified as a
consequence of Bach's Theorem 4, p.~376, by subsequent primary
algorithmic literature. \textbf{Access/check status:} the original AMS
article was not successfully retrieved in this run (HTTP 403); its
theorem number is not represented here as directly checked against the
original. Its exact original wording is a bibliographical gate, while
the stronger operative LLS bounds and the mathematical hypothesis
boundary are directly checked. The following explicit example explains
why ``\(K/\mathbb Q\) abelian'' does not by itself discharge the
class-character hypothesis.

\subsubsection{CC-F03. An abelian base with a non-Dirichlet
class-character
factor}\label{cc-f03.-an-abelian-base-with-a-non-dirichlet-class-character-factor}

\textbf{Statement.} Put \(K=\mathbb Q(\sqrt{-23})\), and let \(H\) be
the splitting field of \(f(x)=x^3-x-1\). Then \(H/K\) is a nontrivial
unramified cyclic cubic extension. A nontrivial character of
\(\mathrm{Gal}(H/K)\) is an unramified class character of \(K\), but its
induction to \(G_\mathbb Q\) is an irreducible two-dimensional
representation. It is not a direct sum, or a virtual integer
combination, of abelian one-dimensional representations. Thus its Hecke
\(L\)-function is not supplied by an abelian-Artin factorization into
Dirichlet \(L\)-functions.

\textbf{Proof.} A reducible monic cubic in \(\mathbb Z[x]\) has an
integer root dividing its constant term. Neither \(1\) nor \(-1\) is a
root of \(f\), so \(f\) is irreducible. Its discriminant is
\(-4(-1)^3-27(-1)^2=-23\), a nonsquare. A transitive subgroup of \(S_3\)
is \(A_3\) or \(S_3\); the discriminant being a nonsquare forces
\(\mathrm{Gal}(H/\mathbb Q)=S_3\). The fixed field of \(A_3\) is \(K\),
since the product of the root differences squares to \(-23\).
Consequently \(\mathrm{Gal}(H/K)=A_3\cong C_3\).

For \(p\ne23\), the polynomial has separable reduction. Its roots lie in
an unramified extension of \(\mathbb Q_p\): choose a finite residue
extension splitting the reduction, lift it to the corresponding
unramified local extension, and lift each simple root by Hensel's lemma.
Hence the splitting field is unramified at all such primes.

At \(p=23\), direct multiplication gives \[
 x^3-x-1\equiv(x-10)^2(x-3)\pmod{23}.
\] The simple root \(3\) lifts to a root \(a\in\mathbb Z_{23}\), and
polynomial division writes \(f=(x-a)g\) with \(g\) monic quadratic over
\(\mathbb Z_{23}\). The resultant of \(x-a\) and \(g\) is a unit,
because \(3\) differs from \(10\) modulo 23. The discriminant product
formula \(\mathrm{disc}(f)=\mathrm{disc}(g)\mathrm{Res}(x-a,g)^2\)
therefore gives valuation one for \(\mathrm{disc}(g)\). The splitting
field of \(g\) is a ramified quadratic extension of \(\mathbb Q_{23}\);
the linear factor adds nothing. Its inertia group acts as a
transposition on the three roots. The relative inertia group for \(H/K\)
is its intersection with \(A_3\), which is trivial. Thus \(H/K\) is
unramified also above 23. There are no real places of \(K\), so there is
no archimedean ramification. Global class field theory identifies this
cyclic unramified extension with a quotient of the ideal class group.

Let \(\eta\) be a nontrivial complex character of \(A_3\). Its conjugate
by a transposition is \(\eta^{-1}\ne\eta\). The induced representation
has character values \(2,0,-1\) on the identity, a transposition, and a
three-cycle, respectively. Its character norm is
\((4+3\cdot0+2\cdot1)/6=1\), so it is irreducible. Every one-dimensional
character of \(S_3\) kills \(A_3=[S_3,S_3]\), and every virtual
combination of them has equal character values on the identity and on a
three-cycle. The induced character does not. More generally, every
abelian character of \(G_\mathbb Q\) is trivial on its commutator
subgroup; the image of this commutator subgroup in the quotient \(S_3\)
is \(A_3\). A virtual combination of abelian \(G_\mathbb Q\)-characters
likewise cannot equal this induced character. Artin induction of
\(L\)-functions identifies \(L_K(s,\eta)\) with the degree-two induced
Artin \(L\)-function, so the claimed abelian factorization is
unavailable. \(\square\)

This is a counterexample to the proposed \emph{factorization route}, not
a proof that Dirichlet GRH could never imply this Hecke GRH by a
different theorem. In this example the missing hypothesis is the
zero-line theorem for that class-character factor, or sufficiently
strong Dedekind GRH for the corresponding Hilbert class extension. One
need not identify \(H\) with the entire Hilbert class field for this
obstruction.

\subsubsection{CC-F03b. A sufficient condition that does discharge the
class-character
hypothesis}\label{cc-f03b.-a-sufficient-condition-that-does-discharge-the-class-character-hypothesis}

\textbf{Statement.} Let \(K/\mathbb Q\) be a number field whose Hilbert
class field \(H_K\) is abelian over \(\mathbb Q\). Then every
ideal-class-character Hecke \(L\)-function of \(K\) is a product of
primitive Dirichlet \(L\)-functions, and its GRH follows from the
adopted input. In particular this applies when \(K\) is imaginary
quadratic and its ideal class group has exponent dividing two.

\textbf{Proof.} Write \(G=\mathrm{Gal}(H_K/\mathbb Q)\),
\(U=\mathrm{Gal}(H_K/K)\), and let \(\eta\) be a character of \(U\). The
Hilbert class field correspondence identifies \(\eta\) with the chosen
ideal-class character. Because \(G\) is finite abelian,
\(\mathrm{Ind}_U^G\eta\) is a direct sum of one-dimensional characters
of \(G\). Induction compatibility of Artin \(L\)-functions therefore
writes \(L_K(s,\eta)\) as the product of the corresponding finite
abelian \(L\)-functions over \(\mathbb Q\), which are primitive
Dirichlet factors. This proves the first claim.

For imaginary quadratic \(K\), its Hilbert class field is stable under
complex conjugation, and conjugation acts on the class group by
inversion: an ideal \(\mathfrak a\) satisfies
\([\overline{\mathfrak a}]=[\mathfrak a]^{-1}\), since
\(\mathfrak a\overline{\mathfrak a}\) is generated by the ideal norm.
The exact sequence for \(\mathrm{Gal}(H_K/\mathbb Q)\) splits by complex
conjugation, which has order two and restricts nontrivially to \(K\).
Thus this group is the semidirect product of the class group by the
inversion action. If the class group has exponent dividing two,
inversion is trivial and the semidirect product is abelian. The first
part applies. \(\square\)

\subsubsection{CC-C05. Littlewood's
consequences}\label{cc-c05.-littlewoods-consequences}

Titchmarsh, revised by Heath-Brown (1986), Theorem 14.2 and equations
(14.2.5)--(14.2.6), pp.~336--337, gives the Littlewood bounds leading
from RH to the Lindelöf hypothesis. Theorem 14.25(A)--(C), pp.~369--370,
gives convergence of \(\sum\mu(n)n^{-s}=1/\zeta(s)\) for \(\Re s>1/2\),
and the equivalence \[
 \mathrm{RH}\quad\Longleftrightarrow\quad
 M(x)=O_\varepsilon(x^{1/2+\varepsilon})
 \quad\text{for every }\varepsilon>0.
\] These are now available with the RH hypothesis discharged. This does
not yield the much stronger inequality \(|M(x)|\le\sqrt{x}\).

For the edge of the Dirichlet critical strip, the exact effective form
used here is LLS Theorem 1.5. Write \[
 B_1(q)=\log\log q-\log2+\tfrac12+\frac1{\log\log q}
       +\frac{14\log\log q}{\log q}.
\] For a primitive character of conductor \(q\ge10^{10}\), \[
 |L(1,\chi)|\ge\frac{\pi^2}{12e^\gamma B_1(q)}.
\] LLS Corollary 1.3 then gives, for negative fundamental discriminant
\(-q\), \[
 h(-q)\ge\frac{\pi}{12e^\gamma}\frac{\sqrt q}{B_1(q)}.
\] These require only Dirichlet GRH. Bounds for Hecke or higher-rank
automorphic \(L\)-functions require their corresponding analytic
hypotheses.

\subsubsection{CC-F02. A fully effective analytic Siegel
bound}\label{cc-f02.-a-fully-effective-analytic-siegel-bound}

\textbf{Statement.} For every positive rational \(\varepsilon\), there
is an algorithm producing a positive rational \(c_\varepsilon\) such
that every primitive real nonprincipal Dirichlet character \(\chi\) of
conductor \(q\) satisfies \[
 L(1,\chi)\ge c_\varepsilon q^{-\varepsilon}.
\]

\textbf{Proof.} For \(q\ge10^{10}\), put \(t=\log q\) and \(u=\log t\).
Then \(t>23\) and \(u>3\). Since \(-\log2+1/2<0\), \(1/u<1/3\), and
\(14/t<1\), \[
 B_1(q)\le2u+\tfrac13<3u.
\] The displayed LLS bound yields \[
 L(1,\chi)\ge \frac{\pi^2}{36e^\gamma\log\log q}.
\] The value is positive: the real Euler product is positive on \(s>1\),
continuity and nonvanishing at \(s=1\) give positivity at 1. Choose a
computable integer \(Q\ge10^{10}\) with
\(\log Q\ge\max(2,2/\varepsilon^2)\). For \(q\ge Q\), the elementary
estimate \(e^v\ge v^2/2\) gives \[
 q^\varepsilon=e^{\varepsilon t}
 \ge\frac{\varepsilon^2t^2}{2}\ge t\ge\log t=\log\log q.
\] Consequently \(q^\varepsilon L(1,\chi)\ge\pi^2/(36e^\gamma)\)
throughout that range.

Only finitely many conductors remain. Primitive real characters of
conductor below \(Q\) can be enumerated by finite arithmetic. For each
such nonprincipal \(\chi\), periodicity and zero sum on a period imply
\(A(x)=\sum_{n\le x}\chi(n)\) has \(|A(x)|\le q\). Abel summation gives
the certified tail estimate \[
 \left|L(1,\chi)-\sum_{n\le N}\frac{\chi(n)}n\right|
 \le \frac{2q}{N}.
\] Increase \(N\) until the rational number
\(b_\chi=\sum_{n\le N}\chi(n)/n-2q/N\) is positive. This terminates
because the partial sums converge to the positive number \(L(1,\chi)\).
Since \(q^\varepsilon\ge1\), one may choose a positive rational
\(c_\varepsilon\) below both \(\pi^2/(36e^\gamma)\) and every resulting
\(b_\chi\). The finite minimum and the computable constant permit such a
choice. For \(q<Q\),
\(L(1,\chi)\ge b_\chi\ge c_\varepsilon\ge c_\varepsilon q^{-\varepsilon}\);
for \(q\ge Q\), the preceding tail-range estimate applies. This
algorithm establishes effectivity, without claiming practical
efficiency. \(\square\)

\subsubsection{CC-C06. Cramér's RH theorem, as distinct from his
stronger
conjecture}\label{cc-c06.-cramuxe9rs-rh-theorem-as-distinct-from-his-stronger-conjecture}

Cramér's classical RH consequence is
\(p_{n+1}-p_n=O(\sqrt{p_n}\log p_n)\). An inspected primary modern
source is Dudek (2014), \emph{On the Riemann hypothesis and the
difference between primes}, arXiv:1402.6417v2: Theorem 1, p.~2, states
Cramér's result; Theorem 2, p.~2, proves the explicit version \[
 (x-\tfrac4\pi\sqrt x\log x,x]\quad\text{contains a prime for every }x\ge2;
\] Theorem 3 replaces \(4/\pi\) by \(1+\varepsilon\) for sufficiently
large \(x\). RH is the sole unproved analytic hypothesis. Cramér's
conjectural \(O((\log p_n)^2)\) gap bound is not the theorem being
discharged. The original 1920 theorem number was not directly inspected;
the exact operative locators here are Dudek's numbered statements, with
the original work retained in the bibliography.

\subsubsection{CC-C07. Selberg's mean square for short
intervals}\label{cc-c07.-selbergs-mean-square-for-short-intervals}

The usable RH consequence is \[
 \int_1^X|\psi(t+h)-\psi(t)-h|^2\,dt
 \ll Xh(\log X)^2\qquad(1\le h\le X).
\] An inspected modern primary statement is Bhowmik--Schlage-Puchta,
\emph{Mean representation number of integers as the sum of primes},
Lemma 1, which explicitly attributes the input to Selberg (1943),
equation (13). The original Selberg article is \emph{On the normal
density of primes in small intervals, and the difference between
consecutive primes}, pp.~87--105; it is paper 12 in \emph{Collected
papers}, vol.~I. The original equation was not directly inspected in
this run, so that locator is marked as transmitted by the later primary
paper.

The elementary consequence needed here requires no further zero
hypothesis. Apply the mean square with upper endpoint \(2X\). For fixed
\(\delta>0\), the measure of \(t\in[X,2X]\) for which
\(|\psi(t+h)-\psi(t)-h|>\delta h\) is at most
\(O(X(\log X)^2/(\delta^2h))\), by integrating over that set. Thus if
\(h/(\log X)^2\to\infty\), the relative exceptional measure tends to
zero. This is an almost-all-interval statement, not a bound in every
interval and not the Selberg eigenvalue conjecture.

\subsubsection{CC-C08. Montgomery's restricted pair
correlation}\label{cc-c08.-montgomerys-restricted-pair-correlation}

Montgomery (1973), §1, unnumbered Theorem, p.~181, equations (1)--(2),
defines \[
 F(\alpha,T)=\left(\frac{T}{2\pi}\log T\right)^{-1}
 \sum_{0<\gamma,\gamma'\le T}
 T^{i\alpha(\gamma-\gamma')}
 \frac4{4+(\gamma-\gamma')^2}.
\] Under RH, \[
 F(\alpha,T)=(1+o(1))T^{-2\alpha}\log T+\alpha+o(1)
\] uniformly on \(0\le\alpha\le1-\varepsilon\), for each fixed
\(\varepsilon>0\). RH suffices for this theorem. The unrestricted
pair-correlation prediction, discussed separately in the same paper at
equations (11)--(12), is not obtained by dropping a hypothesis from the
displayed theorem: its Fourier-support obstruction remains. RH
determines the horizontal locations, not all vertical correlations or
simplicity.

\subsection{4. Criteria and programmes}\label{criteria-and-programmes}

\begin{longtable}[]{@{}
  >{\raggedright\arraybackslash}p{(\columnwidth - 6\tabcolsep) * \real{0.2500}}
  >{\raggedright\arraybackslash}p{(\columnwidth - 6\tabcolsep) * \real{0.2500}}
  >{\raggedright\arraybackslash}p{(\columnwidth - 6\tabcolsep) * \real{0.2500}}
  >{\raggedright\arraybackslash}p{(\columnwidth - 6\tabcolsep) * \real{0.2500}}@{}}
\toprule\noalign{}
\begin{minipage}[b]{\linewidth}\raggedright
ID
\end{minipage} & \begin{minipage}[b]{\linewidth}\raggedright
Precise conclusion now available
\end{minipage} & \begin{minipage}[b]{\linewidth}\raggedright
Exact locator
\end{minipage} & \begin{minipage}[b]{\linewidth}\raggedright
Remaining mathematical task
\end{minipage} \\
\midrule\noalign{}
\endhead
\bottomrule\noalign{}
\endlastfoot
CC-P01 & The global Weil quadratic form is nonnegative on its specified
test-function space & Connes--Consani, supplied source, Appendix C,
Proposition C.1, p.~51; explicit formula Appendix B & The local
archimedean or semilocal comparison with a particular cutoff/Sonin
projection is an additional identity or inequality, not an automatic
termwise consequence of the global sum \\
CC-P02 & Every classical Li coefficient is nonnegative; indeed CC-F04
proves strict positivity for the nonzero Dirichlet systems & Li (1997),
Theorem 1, p.~325; generalized criterion Lagarias (2007), Theorem 2.4,
p.~1702 (preprint Theorem 2.2, p.~11) & Higher-rank automorphic
positivity requires RH for the corresponding \(L\)-function \\
CC-P03 & \(\mathbf1_{(0,1]}\) belongs to the closed span of
\(x\mapsto\{1/(nx)\}\), \(n\ge1\), in \(L^2(0,\infty)\) & Báez-Duarte
(2002), Theorem 1.1, p.~1 & Convergence of a particular proposed
coefficient sequence or optimal distance asymptotics needs its own
theorem; the ordinary Möbius truncations diverge, as discussed at (1.1),
p.~2 \\
CC-P04 & There is an abstract self-adjoint operator whose spectral
ordinates are the zeros, counting multiplicity & Elementary construction
CC-F05 below & Arithmetic origin, a prescribed trace formula, canonical
domain, and geometric functoriality do not follow from the diagonal
construction \\
CC-P05 & Connes's RH-equivalent trace assertion for the \(\mathbb Q\)
version has its zero hypothesis discharged & Connes (1999), §VIII,
Theorem 5 in positive characteristic, and the explicit \(\mathbb Q\)
analogue on preprint pp.~44--47 & The literally numbered theorem
concerns positive characteristic. General number-field Hecke
\(L\)-functions exceed the adopted Dirichlet statement \\
CC-P06 & Deninger's proposed eigenvalue location is correct for the
adopted abelian Artin--Tate factors & Deninger, \emph{The Hilbert--Pólya
strategy and height pairings}, §2, Conjectures 1--2, pp.~2--5 & The
functorial cohomology, determinant formula, Hodge star, and positive
height pairing remain construction/positivity questions; CC-F06 refutes
a general converse from eigenvalue location to the adjoint identity \\
CC-P07 & The specific sufficient de Branges conditions tested by
Conrey--Li remain refuted & Conrey--Li (2000), Theorems 1--2;
§§3.1--3.2, equations (3.1)--(3.4); §4 Remark & Any successful revised
Hilbert-space construction must avoid those disproved positivity
assumptions \\
CC-P08 & The \(\Lambda\)-geometric class in CC-F01 has its analytic RH &
Borger Theorem 0.1 and Corollary 6.2 & Infinite-dimensional
\(\mathbf F_1\) cohomology, an archimedean theory, and a geometric proof
of RH are separate parts of the programme \\
CC-P09 & The already proved restricted Katz--Sarnak density results can
now be interpreted with real ordinates in the Dirichlet families &
Rubinstein (2001), Theorem 3.1, pp.~150--151 & Full support and
higher-rank families do not follow. Theorem 3.1 itself was already
unconditional, allowing complex ordinates \\
\end{longtable}

\subsubsection{CC-P01. Exact global positivity and the archimedean
boundary}\label{cc-p01.-exact-global-positivity-and-the-archimedean-boundary}

For \(g\in C_c^\infty(\mathbb R_{>0})\), use
\(\widetilde g(s)=\int_0^\infty g(x)x^{s-1}dx\) and impose
\(\widetilde g(0)=\widetilde g(1)=0\). On RH the zero-side form is \[
 Q(g)=\sum_\rho\widetilde g(\rho)
       \overline{\widetilde g(1-\overline\rho)}
     =\sum_\rho|\widetilde g(\rho)|^2\ge0.
\] The sums converge because the Mellin transform decays rapidly on a
vertical strip and the zero count is \(O(T\log T)\). Proposition C.1 of
the supplied Connes--Consani paper expresses the equivalent
prime-and-archimedean inequality with their sign convention as
\(\sum_v W_v(g*\overline g^{\sharp})\le0\); it allows finitely many
additional Mellin vanishing conditions as well. The global inequality is
now a theorem in the adopted theory. Their Theorem 1 and Theorem 6.11
concern a more specific comparison of the archimedean term with a
Sonin-space trace, with support and vanishing restrictions. Those
results remain their existing proved theorems; RH alone does not justify
enlarging their support class or declaring each separate local term
positive.

\subsubsection{CC-F04. Strict real-part positivity of Dirichlet Li
coefficients}\label{cc-f04.-strict-real-part-positivity-of-dirichlet-li-coefficients}

\textbf{Statement.} For the zeros of \(\zeta\), or of any primitive
Dirichlet \(L\)-function after the usual completion and removal of the
trivial-character poles, define the Li coefficients by the standard
symmetrically ordered zero sum \[
 \lambda_n=\sum_\rho\left[1-\left(1-\frac1\rho\right)^n\right],\qquad n\ge1.
\] Under the adopted statements, \(\Re\lambda_n>0\) for every \(n\ge1\).
For a real character, or for \(\zeta\), \(\lambda_n\) itself is real and
strictly positive. The criterion is the known theorem of Li and its
generalized form; the short proof here records the strictness and the
treatment of complex characters.

\textbf{Proof.} Every zero has the form \(\rho=1/2+i\gamma\), and \[
 z_\rho=1-\frac1\rho=\frac{-1/2+i\gamma}{1/2+i\gamma}
\] has modulus one. Write \(z_\rho=e^{i\theta_\rho}\). Each real-part
summand is \(1-\cos(n\theta_\rho)\ge0\). For \(|\gamma|\to\infty\),
\(\theta_\rho=O(1/|\gamma|)\); hence these real parts are
\(O_n(1/\gamma^2)\). The usual zero count \(O(T\log(q(T+2)))\), with
\(q\) fixed, makes the real-part series absolutely convergent, so taking
real parts of the prescribed symmetric limit gives the sum of these
nonnegative terms.

There are infinitely many zeros with arbitrarily large \(|\gamma|\). For
finite \(\gamma\), \(z_\rho\ne1\), while \(z_\rho\to1\) as
\(|\gamma|\to\infty\). For a fixed \(n\), choose a sufficiently large
zero so that the nonzero principal argument satisfies
\(0<|\theta_\rho|<2\pi/n\). Then \(1-\cos(n\theta_\rho)>0\), so the
entire sum is strictly positive. In the real-character and zeta cases,
conjugate zeros occur with equal multiplicities; their contributions are
conjugates, making the standard sum real. \(\square\)

No assertion of reality is made for the coefficient attached to a
complex Dirichlet character. The generalized Li criterion uses its real
part (Lagarias, published Theorem 2.4; preprint Theorem 2.2).

\subsubsection{CC-F05. The bare Hilbert--Pólya existence
statement}\label{cc-f05.-the-bare-hilbertpuxf3lya-existence-statement}

\textbf{Statement and proof.} List the nontrivial zeros of \(\zeta\),
with multiplicity, as \(1/2+i\gamma_j\), and choose an enumeration with
\(|\gamma_j|\to\infty\). On \(\mathcal H=\ell^2(\mathbb N)\), put \[
 D(a_j)=(\gamma_j a_j),\qquad
 \operatorname{Dom}(D)=\{a\in\ell^2:\sum_j\gamma_j^2|a_j|^2<\infty\}.
\] Finite sequences give a dense domain. Since the \(\gamma_j\) are
real, \(D\) is symmetric. If \(b\in\operatorname{Dom}(D^*)\), testing
the adjoint relation against the coordinate vectors forces
\((D^*b)_j=\gamma_jb_j\); membership of \(D^*b\) in \(\ell^2\) then
forces \(b\in\operatorname{Dom}(D)\). Conversely every such \(b\)
satisfies the adjoint relation by Cauchy--Schwarz. Thus \(D^*=D\). The
resolvent \((D-i)^{-1}\) is diagonal with entries
\((\gamma_j-i)^{-1}\to0\), and its truncations are finite-rank operators
converging in operator norm. Therefore it is compact. Each coordinate
vector is an eigenvector with the prescribed ordinate, with
multiplicities retained. If a complex number \(z\) is not an ordinate,
discreteness and escape to infinity give \(\inf_j|\gamma_j-z|>0\), so
the diagonal operator \(R_z\) with entries \((\gamma_j-z)^{-1}\) is
bounded. The coordinate identity \(DR_z=I+zR_z\) shows that its range is
in \(\operatorname{Dom}(D)\), bounded in the graph norm; also
\((D-z)R_z=I\) and \(R_z(D-z)=I\) on their respective domains. Hence
there is no additional spectrum. This proves the abstract existence
statement. \(\square\)

This construction uses the zeros themselves as input. It supplies none
of the independent arithmetic explanation sought in a stronger
Hilbert--Pólya programme. The same construction works for each primitive
Dirichlet function. No zero-simplicity assumption is required.

\subsubsection{CC-P05. Connes: exact theorem
scope}\label{cc-p05.-connes-exact-theorem-scope}

Connes's \emph{Trace formula in noncommutative geometry and the zeros of
the Riemann zeta function} has several assertions that should not be
merged. Its §III Theorem 1, preprint p.~13, gives a spectral realization
of critical zeros (with the multiplicity cutoff stated there), with
noncritical zeros entering a resonance description. The adopted RH
removes the noncritical alternatives for \(\zeta\) and the relevant
Dirichlet twists. In §VIII, Theorem 5 is explicitly formulated for
global fields of positive characteristic. The text then constructs and
discusses the corresponding \(\mathbb Q\) cutoff using the real place
and prolate spaces on preprint pp.~44--47, ending with the stated
analogues of Lemma 1, Theorem 5 and Lemma 3. For \(\mathbb Q\), unitary
Hecke characters are Dirichlet characters times imaginary norm powers;
these shift ordinates without changing the real zero line. Thus the
analytic hypothesis in that \(\mathbb Q\) criterion is covered. The
statement for general number-field Hecke characters requires additional
GRH outside the adopted family. A proof of an equivalent trace assertion
does not, by itself, provide a new geometric derivation of its required
cutoff limit.

\subsubsection{CC-F06. Eigenvalue purity does not imply Deninger's
positive adjoint
structure}\label{cc-f06.-eigenvalue-purity-does-not-imply-deningers-positive-adjoint-structure}

\textbf{Statement.} For every real \(w\), there is a complex linear
endomorphism with spectrum on \(\Re s=w/2\) that admits no
positive-definite Hermitian form for which \(A+A^*=wI\).

\textbf{Proof.} On \(\mathbb C^2\), put \(A=(w/2)I+N\), where
\(Ne_2=e_1\) and \(Ne_1=0\). Its only eigenvalue is \(w/2\), and
\(N^2=0\ne N\). If a positive Hermitian form made \(A+A^*=wI\), then
\(N^*=-N\). Taking \(x=e_2\) gives \[
 \|Nx\|^2=\langle x,N^*Nx\rangle
          =-\langle x,N^2x\rangle=0,
\] contradicting positive definiteness and \(Nx=e_1\ne0\). \(\square\)

Deninger's proposed identity in §2, pp.~4--5, is
\(\nu(f,f')=(\theta f,f')+(f,\theta f')\), obtained from a trace, cup
product and Hodge star. It implies the desired line. The example proves
that a general converse from line location to that positive adjoint
identity is false. It does not refute the existence of the particular
geometric objects Deninger conjectures. Semisimplicity, a pairing and
the functorial geometry are additional content.

\subsubsection{CC-P07. What the de Branges counterexamples
settle}\label{cc-p07.-what-the-de-branges-counterexamples-settle}

Conrey and Li's Theorems 1--2 formulate sufficient positivity conditions
in de Branges spaces. In §3.1 they show that the proposed conditions
fail for \(E(z)=\xi(1-iz)\), and in §3.2 they give analogous failures
for the character modulo 4. Their §4 Remark attributes to Sarnak a
nonnumerical proof of failure of the kernel-shift condition (3.3), via
the value distribution of \(\log\zeta\), with a Dirichlet analogue for
(3.8). The derivative-test counterexamples to (3.1) in §3 are cited as
the published numerical results, not as computations independently
certified in this run. The nonnumerical kernel-shift refutation already
suffices to disprove the blanket assertion that every proposed
sufficient positivity condition follows from RH. The accurate outcome is
a refutation of those particular sufficient conditions. It would be a
logical error to deduce them from RH.

\subsubsection{CC-P09. Katz--Sarnak and the restriction that
survives}\label{cc-p09.-katzsarnak-and-the-restriction-that-survives}

Rubinstein's Theorem 3.1, pp.~150--151, treats primitive quadratic
characters \(\chi_d\), averaged over fundamental discriminants with
\(X/2\le|d|<X\). For product Schwartz tests
\(f(x_1,\ldots,x_n)=\prod_i f_i(x_i)\), each \(f_i\) even, with Fourier
support \(\sum_i|u_i|<1\), the normalized \(n\)-level density is the
symplectic determinant with kernel \[
 K_{-1}(x,y)=\frac{\sin\pi(x-y)}{\pi(x-y)}
            -\frac{\sin\pi(x+y)}{\pi(x+y)}.
\] The sentence immediately preceding the theorem explicitly allows
complex ordinates and says RH is not assumed. Thus the theorem itself
was already unconditional. Dirichlet GRH now supplies real ordinates but
does not remove its Fourier-support restriction. Full number-field
family statistics, and analogues involving higher-rank \(L\)-functions,
are not consequences of just the adopted zero-line statements. The
finite-field monodromy and equidistribution part of Katz--Sarnak
likewise was not waiting for number-field RH.

\subsection{5. Explicit boundary
catalogue}\label{explicit-boundary-catalogue}

\begin{longtable}[]{@{}
  >{\raggedright\arraybackslash}p{(\columnwidth - 2\tabcolsep) * \real{0.5000}}
  >{\raggedright\arraybackslash}p{(\columnwidth - 2\tabcolsep) * \real{0.5000}}@{}}
\toprule\noalign{}
\begin{minipage}[b]{\linewidth}\raggedright
Proposed extension
\end{minipage} & \begin{minipage}[b]{\linewidth}\raggedright
Status and actual additional gate
\end{minipage} \\
\midrule\noalign{}
\endhead
\bottomrule\noalign{}
\endlastfoot
General effective Chebotarev from Dirichlet GRH & Theorem 1.1 of
Lagarias--Odlyzko needs RH for \(\zeta_L\). It is discharged when
\(L/\mathbb Q\) is abelian; arbitrary top fields need more \\
Bach class-group generators for every abelian \(K/\mathbb Q\) & The
class-character Hecke \(L\)-functions need their own zero theorem, or a
suitable Hilbert-class-extension Dedekind GRH; CC-F03 shows why abelian
base is insufficient for the offered factorization argument \\
Full generalized Li positivity & Lagarias Theorem 2.4 (preprint 2.2) is
family by family. General unitary cuspidal \(GL_N\) functions require
their own RH; Dirichlet GRH covers the finite-order \(GL_1/\mathbb Q\)
family and its norm shifts \\
Full GUE / Katz--Sarnak statistics / simplicity of all zeros & The cited
restricted statistics keep their test-function restrictions. No theorem
here upgrades zero location to arbitrary correlations or simplicity \\
Cramér's logarithmic-square gap conjecture & Distinct from Cramér's
RH-conditional square-root-times-log theorem \\
Universal local positivity in an explicit formula & Positivity of a sum
of local terms does not imply positivity of each summand.
Connes--Consani's extra local comparison has its own hypotheses \\
Deninger's positive cohomology or a canonical Hilbert--Pólya operator &
The analytic target is supplied for the adopted factors; functorial
spaces, adjoints, pairings and trace compatibility remain. CC-F06
disproves an unrestricted linear-algebra converse \\
A \(\Lambda\)-structure on the hexagonal elliptic curve & Refuted by
Borger §2.9; split-prime CM Frobenius lifts are weaker data \\
A Chow-motive Artin--Tate decomposition from the cohomological
\(L\)-factorization & Not proved by the factorization; a realization is
not automatically fully faithful or conservative in the required motivic
sense \\
Effective integral points merely from the phrase ``effective Siegel'' &
No such implication is asserted. CC-F02 proves only the analytic lower
bound for real Dirichlet \(L(1,\chi)\) and its class-number
consequences \\
\end{longtable}

\subsection{6. Source register and exact-reading
status}\label{source-register-and-exact-reading-status}

The following are primary works or mirrors of primary texts. Web
references in the companion TSV are provenance aids; the paper should
cite authors, dates and the locators above. No original statement number
is guessed when an original was not inspected.

\begin{enumerate}
\def\labelenumi{\arabic{enumi}.}
\item
  \textbf{Borger, J. (2009).} \emph{\(\Lambda\)-rings and the field with
  one element} (arXiv:0906.3146v1). Theorems 0.1, 5.3, 6.1; Corollary
  6.2, equation (6.2.1); Corollaries 6.5--6.6; §2.9.
  \url{https://arxiv.org/pdf/0906.3146} . Supplied TeX and external PDF both
  inspected. The PDF's regenerated date is not used as the publication
  date.
\item
  \textbf{Lamzouri, Y., Li, X., \& Soundararajan, K. (2015; corrected
  preprint 2016).} Conditional bounds for the least quadratic
  non-residue and related problems. \emph{Mathematics of Computation,
  84}(295), 2391--2412. \url{https://doi.org/10.1090/S0025-5718-2015-02925-1}
  . Operative corrected text: \url{https://arxiv.org/pdf/1309.3595} , v3, 8
  December 2016, internal date 12 December 2016. Theorems 1.1, 1.4, 1.5;
  Corollaries 1.1--1.3. PDF and extracted text inspected.
\item
  \textbf{Lagarias, J. C., \& Odlyzko, A. M. (1977).} Effective versions
  of the Chebotarev density theorem. In A. Fröhlich (Ed.),
  \emph{Algebraic number fields: L-functions and Galois properties}
  (pp.~409--464). Academic Press. Theorem 1.1, equation (1.3); Corollary
  1.2. Inspected retypeset primary text:
  \url{https://aareyanmanzoor.github.io/assets/articles/lagarias-odlyzko.pdf}
  .
\item
  \textbf{Miller, G. L. (1976).} Riemann's hypothesis and tests for
  primality. \emph{Journal of Computer and System Sciences, 13}(3),
  300--317. Theorem 2, p.~301; Appendix, pp.~316--317.
  \url{https://www.cs.cmu.edu/\textasciitilde} glmiller/Publications/Papers/Mi76.pdf
  . Original PDF inspected, including rendered p.~301.
\item
  \textbf{Bach, E. (1990).} Explicit bounds for primality testing and
  related problems. \emph{Mathematics of Computation, 55}(191),
  355--380. \url{https://doi.org/10.1090/S0025-5718-1990-1023756-8} . Original
  PDF access failed; author abstract and LLS attribution inspected.
  Theorem 4 / p.~376 class-group locator remains to be directly checked.
\item
  \textbf{Titchmarsh, E. C. (1986).} \emph{The theory of the Riemann
  zeta-function} (2nd ed., D. R. Heath-Brown, rev.). Oxford University
  Press. Theorems 14.2 and 14.25(A)--(C). Original book scan:
  \url{https://sites.math.rutgers.edu/\textasciitilde} zeilberg/EM18/TitchmarshZeta.pdf
  . Inspected extracted text; equations crosschecked against theorem
  context.
\item
  \textbf{Cramér, H. (1920).} Some theorems concerning prime numbers.
  \emph{Arkiv för Matematik, Astronomi och Fysik, 15}(5), 1--33.
  Historical attribution; original statement number not inspected.
  Operative source follows.
\item
  \textbf{Dudek, A. W. (2014/2015).} On the Riemann hypothesis and the
  difference between primes. Preprint arXiv:1402.6417v2; published in
  \emph{International Journal of Number Theory, 11}(3), 771--778.
  Theorems 1--3 in the inspected preprint (published numbering
  1.1--1.3). \url{https://arxiv.org/pdf/1402.6417} .
\item
  \textbf{Selberg, A. (1943).} On the normal density of primes in small
  intervals, and the difference between consecutive primes. \emph{Archiv
  for Mathematik og Naturvidenskab, 47}(6), 87--105. Equation (13), as
  cited by the next primary work; original not directly inspected.
\item
  \textbf{Bhowmik, G., \& Schlage-Puchta, J.-C. (2010).} Mean
  representation number of integers as the sum of primes. \emph{Nagoya
  Mathematical Journal, 200}, 27--33.
  \url{https://doi.org/10.1215/00277630-2010-010} . Lemma 1, p.~2 in arXiv v3,
  with explicit attribution to Selberg equation (13). Primary PDF
  inspected: \url{https://arxiv.org/pdf/0806.3295} ; initial mirror:
  \url{https://paperzz.com/doc/9064999/mean-representation-number-of-integers-as-the-sum}
  .
\item
  \textbf{Montgomery, H. L. (1973).} The pair correlation of zeros of
  the zeta function. In \emph{Analytic number theory} (Proceedings of
  Symposia in Pure Mathematics, Vol. 24, pp.~181--193). American
  Mathematical Society. \url{https://doi.org/10.1090/pspum/024/9944} . §1
  unnumbered Theorem, p.~181, equations (1)--(2); conjectural equations
  (11)--(12). Primary-text mirror inspected:
  \url{https://studylib.net/doc/7952799/the-pair-correlation-of-zeros-of-the-zeta-function}
  .
\item
  \textbf{Li, X.-J. (1997).} The positivity of a sequence of numbers and
  the Riemann hypothesis. \emph{Journal of Number Theory, 65}(2),
  325--333. \url{https://doi.org/10.1006/jnth.1997.2137} . Theorem 1, p.~325;
  equation (1.4). Primary-text mirror:
  \url{https://paperzz.com/doc/7086946/--science-direct} .
\item
  \textbf{Lagarias, J. C. (2007).} Li coefficients for automorphic
  \(L\)-functions. \emph{Annales de l'Institut Fourier, 57}(5),
  1689--1740. \url{https://doi.org/10.5802/aif.2311} . Published Theorem 2.4,
  p.~1702:
  \url{https://www.numdam.org/article/AIF\_2007\_\_57\_5\_1689\_0.pdf} . In
  arXiv v4 (2005), \url{https://arxiv.org/pdf/math/0404394} , the
  corresponding statement is Theorem 2.2, p.~11. Both primary statements
  inspected; their different numbering is retained.
\item
  \textbf{Báez-Duarte, L. (2002/2003).} A strengthening of the
  Nyman--Beurling criterion for the Riemann hypothesis. \emph{Rendiconti
  Lincei. Matematica e Applicazioni, 14}, 5--11. Theorem 1.1 in
  arXiv:math/0202141v2: \url{https://arxiv.org/pdf/math/0202141} . PDF
  inspected. The quantitative continuation is arXiv:math/0205003; not
  used here without checking its exact numbered statement.
\item
  \textbf{Connes, A. (1999).} Trace formula in noncommutative geometry
  and the zeros of the Riemann zeta function. \emph{Selecta Mathematica,
  5}, 29--106. \url{https://arxiv.org/pdf/math/9811068} . §VIII, Theorem 5,
  and preprint pp.~44--47. PDF inspected; characteristic scope checked
  against the theorem's opening words.
\item
  \textbf{Connes, A., \& Consani, C. (2020).} Weil positivity and trace
  formula, the archimedean place. arXiv:2006.13771.
  \url{https://arxiv.org/pdf/2006.13771} . Supplied TeX
  \nolinkurl{weil-compo.tex}; Appendix C Proposition C.1; Theorem 1;
  Theorem 6.11; Appendix B. Supplied text and PDF locators inspected.
\item
  \textbf{Deninger, C. (2010).} The Hilbert--Pólya strategy and height
  pairings. In \emph{Casimir force, Casimir operators and the Riemann
  hypothesis} (pp.~275--283). Walter de Gruyter. §2, Conjectures 1--2
  and the adjoint calculation, preprint pp.~2--5. Author/institute PDF:
  \url{https://www.uni-muenster.de/SFB878/publications/files/php7aMxKR3029.pdf}
  ; arXiv: \url{https://arxiv.org/pdf/1001.1621} . PDF inspected; publication
  metadata checked on the author's publications page.
\item
  \textbf{Conrey, J. B., \& Li, X.-J. (2000).} A note on some positivity
  conditions related to zeta- and \(L\)-functions. \emph{International
  Mathematics Research Notices, 2000}(18), 929--940.
  \url{https://arxiv.org/pdf/math/9812166} . Theorems 1--2; §§3.1--3.2; §4
  Remark. Primary PDF inspected.
\item
  \textbf{Rubinstein, M. (2001).} Low-lying zeros of \(L\)-functions and
  random matrix theory. \emph{Duke Mathematical Journal, 109}(1),
  147--181. Theorem 3.1, pp.~150--151.
  \url{https://web.williams.edu/Mathematics/sjmiller/public\_html/ntprob17/handouts/Lfns/Rubinstein\_ThesisPaper\_Duke.pdf}
  . Original PDF inspected. The theorem's explicit absence of an RH
  assumption was checked.
\item
  \textbf{Belabas, K., Diaz y Diaz, F., \& Friedman, E. (2008).} Small
  generators of the ideal class group. \emph{Mathematics of Computation,
  77}(262), 1185--1197. \url{https://doi.org/10.1090/S0025-5718-07-02003-0} .
  Inspected author preprint dated 5 December 2005:
  \url{https://www.math.u-bordeaux.fr/\textasciitilde} kbelabas/research/OnBach.pdf
  . Introduction pp.~2--3 and Theorem 1.1, p.~3.
\item
  \textbf{Bach, E., \& Sorenson, J. (1996).} Explicit bounds for primes
  in residue classes. \emph{Mathematics of Computation, 65}(216),
  1717--1735. \url{https://doi.org/10.1090/S0025-5718-96-00763-6} . Original
  author-hosted PDF inspected:
  \url{https://www.researchgate.net/profile/Jonathan-Sorenson-2/publication/220576951\_Explicit\_bounds\_for\_primes\_in\_residue\_classes/links/55100ea70cf2ac2905afbceb/Explicit-bounds-for-primes-in-residue-classes.pdf}
  . Theorem 5.3, p.~1731; §2, p.~1720; Corollaries 3.4--3.5,
  pp.~1722--1723. The numerical computations reported by the authors
  were not rerun here.
\end{enumerate}

\subsection{7. Check record and unresolved bibliographical
gates}\label{check-record-and-unresolved-bibliographical-gates}

Every new proof in this module must be read as a whole after
composition. The companion
\nolinkurl{consequences_criteria_CHECK_EVENTS.tsv} records the actual
completion of those reads and any correction; it does not call a
same-author proof read an independent re-proof. The adopted Alpöge
statements receive an event for their use, not for a claimed
verification of the unseen full archive.

The unresolved bibliographical items are narrowly delimited: direct
inspection of Bach's original Theorem 4 and the original Selberg
equation (13). These are not presented as completed direct-source
checks. The operative least-prime, cohomological, effective-Siegel and
criterion results above have explicit inspected primary locators. No
mathematical claim of global logical independence is made from a failure
of the current factorization route.

% END consequences_criteria.tex

% BEGIN consequences_criteria_GT_addendum.tex
\section{Current Grothendieck--Teichmüller scope: 20 November
2026}\label{current-grothendieckteichmuxfcller-scope-20-november-2026}

Research addendum \nolinkurl{consequences_criteria}; GPT-6 Astra Pro;
effort setting not exposed. This updates the literature status without
enlarging the adopted Alpöge input or claiming a kernel check of another
proof archive.

\subsection{CC-GT01. The 2026 rational claim has substantive primary
support}\label{cc-gt01.-the-2026-rational-claim-has-substantive-primary-support}

It would be outdated to label rational Deligne--Drinfeld freeness simply
open. OpenAI's 23 September 2026 preprint, Theorem 1.1, asserts freeness
of the rational graded solution Lie algebra on one generator in each odd
weight at least three, together with its weight-completed version. Its
mirrored introduction, the statement of its all-weight upper bound, and
parts of the argument were inspected; the entire 41-page original proof
and its Lean archive were not independently checked here. Plectis Atlas
itself is an index of other parties' reports, not our verification
method. The Valency page could not be retrieved directly.

Two subsequent primary mathematical texts supply stronger evidence than
an index entry:

\begin{longtable}[]{@{}
  >{\raggedright\arraybackslash}p{(\columnwidth - 4\tabcolsep) * \real{0.3333}}
  >{\raggedright\arraybackslash}p{(\columnwidth - 4\tabcolsep) * \real{0.3333}}
  >{\raggedright\arraybackslash}p{(\columnwidth - 4\tabcolsep) * \real{0.3333}}@{}}
\toprule\noalign{}
\begin{minipage}[b]{\linewidth}\raggedright
Source
\end{minipage} & \begin{minipage}[b]{\linewidth}\raggedright
Exact locator inspected
\end{minipage} & \begin{minipage}[b]{\linewidth}\raggedright
Mathematical contribution
\end{minipage} \\
\midrule\noalign{}
\endhead
\bottomrule\noalign{}
\endlastfoot
Willwacher (2026b), arXiv:2610.15592v4, 28 October & Theorem 1.1, p.~2;
Theorem 2.2, p.~3, proved in §§3.1--3.3, pp.~4--6 & The required upper
bound \(\dim_{\mathbb Q}\mathfrak{grt}_1(\mathbb Q)_n\le d_n\), where
\(d_n\) is the dimension of the free Lie algebra on odd generators \\
Hirose (2026), arXiv:2610.23719v1, 16 October & Theorem 4 and §1.2, p.~2
& An integral Hoffman spanning theorem modulo confluence relations,
combined there with Brown and Furusho to identify the motivic group with
\(\mathrm{GRT}\) \\
\end{longtable}

For the first paper, the full chain proving the upper bound was read:
integral rank comparison, the characteristic-two projection, its
explicit reconstruction, the image restriction, and the free-Lie
dimension calculation. The key derivation intertwining and square-zero
identities were checked algebraically. This is a targeted reading of an
existing proof, not an independent reproof or a formal verification.
Theorem 1.3 and §4 give further characteristic-two structure, but they
are unnecessary when Brown's injection is retained. Hirose's theorem and
stated motivic consequence were inspected; his full technical proof was
not checked in this run.

Neither proof uses RH or Dirichlet GRH. They are external 2026
mathematical results, not achievements attributable to the adopted
analytic input.

\subsection{CC-GT02. What Drinfeld and Brown had already
supplied}\label{cc-gt02.-what-drinfeld-and-brown-had-already-supplied}

Drinfeld (1991), Proposition 6.3, p.~859, gives a homogeneous element of
\(\mathfrak{grt}_1\) in every odd degree \(n\ge3\) with nonzero
depth-one term. The unnumbered ``Questions'' immediately following ask
about precisely one generator in each odd degree and freeness. The
primary proposition, proof and questions were visually inspected on the
rendered original scan; its existence assertion should not be conflated
with freeness or exhaustion.

Brown (2012), Theorem 1.1, p.~950, gives the motivic Hoffman basis; the
paragraph immediately following explains why it proves Conjecture 1 on
p.~949, namely faithfulness of the motivic action. The structure of the
motivic prounipotent Lie algebra is recalled on p.~950. An exact
numbered modern formulation of its free-subalgebra consequence is
Willwacher (2026a), Theorem 1.1, p.~1. His Theorem 1.2, also p.~1,
extends the assertion to any homogeneous odd-degree family with nonzero
depth-one coefficients. These are injectivity statements. The upper
bound in CC-GT01 supplies the distinct exhaustion step.

\subsection{CC-GT03. Combining injection and upper bound, including
completion}\label{cc-gt03.-combining-injection-and-upper-bound-including-completion}

\textbf{Statement.} Let \[
 \mathfrak f=\operatorname{Lie}_{\mathbb Q}\langle e_3,e_5,e_7,\ldots\rangle,
 \qquad \deg e_j=j.
\] Write
\(\mathfrak g^{\mathrm{gr}}=\bigoplus_n\mathfrak{grt}_{1,n}(\mathbb Q)\)
for the graded algebra of polynomial solutions, and
\(\widehat{\mathfrak g}=\prod_n\mathfrak{grt}_{1,n}(\mathbb Q)\) for its
weight completion; some sources use \(\mathfrak{grt}_1\) for the latter.
Take the weight-preserving injection
\(\iota:\mathfrak f\hookrightarrow\mathfrak g^{\mathrm{gr}}\) supplied
by Brown's result and the upper bound of Willwacher (2026b), Theorem
1.1. Then \(\iota\) is an isomorphism, and its extension
\(\prod_n\mathfrak f_n\to\widehat{\mathfrak g}\) is an isomorphism of
complete Lie algebras.

\textbf{Proof.} In every weight \(n\), the space \(\mathfrak f_n\) is
finite-dimensional: only finitely many generators can occur, and every
generator has degree at least three. Injectivity gives \[
 \dim\mathfrak f_n\le\dim\mathfrak{grt}_{1,n}.
\] The cited upper bound reverses this inequality. Therefore the
injective linear map \(\iota_n\) is bijective in every weight. Their
direct sum is a bijective graded Lie homomorphism, hence a Lie algebra
isomorphism. Their product is a continuous bijection for the product,
equivalently weight-filtration, topologies; its inverse is the product
of the degreewise inverses and is continuous. A bracket component of
fixed weight is a finite sum of brackets of lower-weight components, so
the product map and its inverse preserve the completed bracket. This
proves the completed assertion. \(\square\)

This deduction identifies exactly which result completed the rational
theorem. It supplies no theorem about the full profinite absolute Galois
group.

\subsection{CC-GT04. The full profinite question is a different
assertion}\label{cc-gt04.-the-full-profinite-question-is-a-different-assertion}

Fix Drinfeld's original profinite group \(\widehat{\mathrm{GT}}\), with
its commutator, invertibility, symmetry, hexagon and pentagon
conditions. Bortnovskyi, Dolgushev, Holikov, and Pashkovskyi (2026
revision), §1.3, equations (1.5)--(1.8), pp.~4--5, gives the fixed Ihara
embedding and cyclotomic compatibility. In that paper's coordinates, \[
 \operatorname{Ih}(g)=\bigl((\chi_{\mathrm{cyc}}(g)-1)/2,f_g\bigr),
 \qquad
 x\mapsto x^{\chi_{\mathrm{cyc}}(g)},\quad
 y\mapsto f_g^{-1}y^{\chi_{\mathrm{cyc}}(g)}f_g.
\] Its cyclotomic projection is \((m,f)\mapsto2m+1\), so the composite
equals \(\chi_{\mathrm{cyc}}\). The paper's introduction distinguishes
surjectivity of this whole embedding from its proved surjectivity onto
particular quotients; its Theorems 5.3 and 5.7 concern such quotient
results. Rational Lie-algebra freeness does not establish full profinite
surjectivity.

The safe status statement for this report is precise: \textbf{the full
profinite identification is not settled by the adopted RH/Dirichlet GRH
or by the rational Deligne--Drinfeld theorem; the inspected 2026 primary
profinite work treats full surjectivity as unresolved.} No exhaustive
claim about all unpublished assertions is needed. The gentle,
pro-\(\ell\), prounipotent, and higher-genus variants must not be
substituted for the fixed original profinite embedding.

\subsection{CC-GT05. The exact positive role of the cyclotomic
quotient}\label{cc-gt05.-the-exact-positive-role-of-the-cyclotomic-quotient}

The cyclotomic character and Kronecker--Weber identify finite abelian
characters of \(G_{\mathbb Q}\) with Dirichlet characters. Consequently
the adopted GRH controls the analytic factors of finite abelian Artin
representations and their Tate twists, including the Borger realizations
of CC-F01 and CC-F01b. Those results also apply to arithmetic
representations known to factor through the cyclotomic projection in the
preceding diagram.

This analytic statement concerns those representations and their
\(L\)-functions. It does not determine the kernel of the cyclotomic
projection, the associator component \(f_g\), or surjectivity of
\(\operatorname{Ih}\). Abstractly, surjectivity after projection cannot
establish surjectivity before projection: for any profinite abelian
group \(A\), the map \(A\to A\times C_2\), \(a\mapsto(a,1)\), has
surjective composite with the projection to \(A\) and is a proper
embedding. This elementary observation is a counterexample to that
general inference, not a counterexample to any conjecture about the
actual group \(\widehat{\mathrm{GT}}\).

\subsection{Exact bibliography and access
record}\label{exact-bibliography-and-access-record}

\begin{enumerate}
\def\labelenumi{\arabic{enumi}.}
\tightlist
\item
  \textbf{Brown, F. (2012).} Mixed Tate motives over \(\mathbb Z\).
  \emph{Annals of Mathematics, 175}(2), 949--976.
  \url{https://doi.org/10.4007/annals.2012.175.2.10} . Primary PDF:
  \url{https://annals.math.princeton.edu/wp-content/uploads/annals-v175-n2-p10-p.pdf}
  . Theorem 1.1 and its stated consequence directly inspected.
\item
  \textbf{Drinfeld, V. G. (1991).} On quasitriangular quasi-Hopf
  algebras and on a group that is closely connected with
  \(\operatorname{Gal}(\overline{\mathbb Q}/\mathbb Q)\).
  \emph{Leningrad Mathematical Journal, 2}(4), 829--860. Original
  Russian publication: \emph{Algebra i Analiz, 2}(4), 149--181 (1990).
  Scan: \url{https://www.lamfa.u-picardie.fr/marin/docs/drinfeld.pdf} .
  Proposition 6.3 and the following unnumbered Questions, p.~859,
  visually inspected on the rendered original page.
\item
  \textbf{Willwacher, T. (2026a).} \emph{On free generators in the
  Grothendieck--Teichmüller Lie algebra} (arXiv:2607.18793v2, 6 August
  2026). \url{https://arxiv.org/pdf/2607.18793v2} . Theorems 1.1--1.2, p.~1,
  directly inspected.
\item
  \textbf{Willwacher, T. (2026b).} \emph{The Deligne--Drinfeld
  conjecture after OpenAI} (arXiv:2610.15592v4, 28 October 2026).
  \url{https://arxiv.org/pdf/2610.15592v4} . Theorem 1.1 and its entire proof
  chain in §§1--3 read; key algebraic identities checked. Theorem 1.3
  and §4 are additional, not required inputs to CC-GT03.
\item
  \textbf{Hirose, M. (2026).} \emph{Confluence relations, Hoffman basis
  over \(\mathbb Z_{(2)}\), and the Deligne--Drinfeld conjecture}
  (arXiv:2610.23719v1, 16 October 2026).
  \url{https://arxiv.org/pdf/2610.23719v1} . Theorem 4 and §1.2, p.~2,
  directly inspected; full technical proof not independently checked
  here.
\item
  \textbf{Bortnovskyi, I., Dolgushev, V. A., Holikov, B., \&
  Pashkovskyi, V. (2024/2026).} \emph{Accessing non-abelian quotients of
  the Grothendieck--Teichmueller group via elementary tools}
  (arXiv:2405.11725v3, 6 August 2026).
  \url{https://arxiv.org/pdf/2405.11725v3} . Introduction, §1.3, equations
  (1.5)--(1.8), pp.~4--5, and the described quotient scope directly
  inspected.
\item
  \textbf{OpenAI. (2026, September 23).} \emph{The Deligne--Drinfeld
  conjecture}. OpenAI Math Release preprint.
  \url{https://github.com/openai/math/blob/fd4aeeb2ee4fc729c18d98444fed42fd0529eeeb/preprints/The-Deligne-Drinfeld-conjecture-September-23-2026/paper.pdf}
  . Inspected mirror:
  \url{https://wcook04.github.io/plectis-atlas/papers/The-Deligne-Drinfeld-conjecture-September-23-2026.html}
  . Theorem 1.1 and dimension-bound route inspected; the complete
  original proof and Lean archive were not checked in this run.
\end{enumerate}

% END consequences_criteria_GT_addendum.tex

% BEGIN source_and_masullo.tex
% Candidate section. This file is intended for inclusion in the root paper.
% IDs F-SM-* identify results proved in this run, not claims of historical novelty.
\section{The two supplied routes: exact scope and mathematical bridges}
\label{sec:source-masullo}

\subsection{What is adopted, what has been checked, and what is not exported}

We use Alp\"oge's five declarations in \nolinkurl{ALPOGE_Main.lean} as the research
inputs specified in the task.  The check performed here is an inspection of the
declarations and of \nolinkurl{ALPOGE_Statements.lean}, not a fresh Lean kernel
verification: the import \nolinkurl{Supplements.Kernels.MainTheoremsAll} is absent
from the supplied ZIP, as is the remainder of the formal development.  Neither
\nolinkurl{lean} nor \nolinkurl{lake} is installed in the working environment.  The
supplied scan is a report about another source tree; no claim is made that the
175,187 files were rescanned here.  This distinction does not change the adopted
mathematical premises, but it determines the correct check-event label.

The declaration names and their exact targets are as follows.
\begin{enumerate}
\item \texttt{riemannHypothesis : RiemannHypothesis := rhOfAll}; the target is
Mathlib's named proposition.
\item \texttt{riemannZeta\_eq\_zero\_re : RiemannZetaZerosOnLine := rhOfAll};
the supplied definition excludes the negative even integers and $s=1$.
\item \texttt{LFunction\_ne\_zero\_of\_half\_lt\_re :
DirichletLNonvanishingRightOfHalf := allChi}; this quantifies over every positive
modulus, every complex Dirichlet character, and $\Re s>1/2$, with the exception
$(\chi,s)=(1,1)$.
\item \texttt{LFunction\_eq\_zero\_re : DirichletLZerosOnLine := claimAll.2.1};
this concerns zeros in the open strip $0<\Re s<1$.
\item \texttt{LFunction\_eq\_zero\_re\_of\_isPrimitive :
PrimitiveDirichletLZerosOnLine := critAll}; this is the primitive-character
specialization, again in the open strip.
\end{enumerate}
The pinned Mathlib revision is
\nolinkurl{0df444a360eaa60ab8c11dca51a86af692955474}, in
\nolinkurl{ALPOGE_lakefile.toml}.  None of these five declarations states a
weight-monodromy theorem, a theorem on all Hecke characters of all number
fields, a nonabelian Artin holomorphy theorem, or a theorem specifying an
integral algebraic $K$-group of $\mathbb Z$.

Two corrections to the prompt's reading of the four-page note are material.
The theta-lift chain has length of order $\sqrt X$, both in the paragraph on
page~1 and in step~8 on page~4.  The radical is visible in the PDF.  The claim
about $K$-groups in step~17, footnote~6, is explicitly tentative and gives no
degrees or quantifiers.  It cannot be turned into a specified $K$-group
vanishing theorem merely by citing the note.  The relevant export, if one
exists in the full development, must first be identified.

\subsection{Masullo's different target and its exact remaining gate}

Masullo, \emph{The Music of the Primes}, page~6, specifies all number fields
and all unitary Hecke characters.  This is a larger target than the five
Alp\"oge declarations above.  The analytic route can be located precisely:
\begin{enumerate}
\item Section~1.2.1, pages~13--15: a divisor exact sequence of nuclear
Fr\'echet spaces, then Hardy--Sobolev Hilbert completions whose spectral
quotients retain zeros in a prescribed vertical strip.
\item Sections~1.2.2--1.2.3, pages~15--16, and Chapter~11: orbital bases
and a decomposition of a lifted geometric kernel into local terms.
\item Theorem~10.4.2.8, pages~463--471: the spectral trace after the
relevant truncation/projection operations.
\item Theorem~11.2.2.1, pages~490--491: orbital nuclearity on the
zero line for Sobolev degree $s\geq 2[F:\mathbb Q]$.
\item Theorem~11.3.2.1, pages~495--496: orbital trace calculations;
Theorem~12.4.1.1, pages~514--519: the $S$-local trace calculation.
\item Theorem~12.5.1.1, pages~519--520: the proposed global trace
identity with only the zeros inside the chosen strip on its spectral side.
\item Sections~1.4.2--1.4.5, pages~28--30: the proposed deduction of
zero location by holomorphic families of test functions.
\end{enumerate}
The end of Chapter~13 is not a usable theorem statement: on page~522,
Theorem~13.3.1.2 consists of a heading with instructions to insert the
statement and prove the main results.  Theorem~13.3.1.1 on the same page is
marked for correction.  More substantially, Section~12.3 on page~514, devoted
to the global-to-$S$-local deconstruction, consists only of its heading and an
instruction to include the section.  The proof of Theorem~12.5.1.1 on
page~520 uses the resulting decomposition and exchange of diagonal integration
with the local sum.  These are concrete missing links in the supplied written
argument.  The following proposition supplies the zero-location deduction
once the appropriate trace identity has been established; it does not supply
the missing kernel identity.

\paragraph{Finding F-SM-01 (a shrinking-strip trace criterion).}
Let $\mathcal Z$ be a locally finite multiset in
$\{\lambda\in\mathbb C:|\Re\lambda|\leq B\}$, with positive integral
multiplicities $m(\lambda)$, and suppose
\[
 \sum_{\lambda\in\mathcal Z}
       \frac{m(\lambda)}{1+|\Im\lambda|^2}<\infty.
\]
For $f\in C_c^\infty(0,\infty)$ put
$\widehat f(\lambda)=\int_0^\infty f(t)e^{\lambda t}\,dt$.
Assume that for some sequence $a_j\downarrow0$,
\begin{equation}\label{eq:sm-strip}
 \sum_{\lambda\in\mathcal Z}m(\lambda)\widehat f(\lambda)
  =\sum_{\substack{\lambda\in\mathcal Z\\|\Re\lambda|<a_j}}
           m(\lambda)\widehat f(\lambda)
 \qquad(f\in C_c^\infty(0,\infty)).
\end{equation}
Then every $\lambda\in\mathcal Z$ is imaginary.  Conversely, if every
$\lambda$ is imaginary, \eqref{eq:sm-strip} holds for every $a_j>0$.

\begin{proof}
Fix $a=a_j$ and let $\mathcal E$ be the submultiset with
$|\Re\lambda|\geq a$.  The hypothesis gives
$\sum_{\mathcal E}m(\lambda)\widehat f(\lambda)=0$ for every such test
function.  Fix a real number $w>B$ and a complex number $z$ with
$\Re z>B$.  On $t\geq0$ define
\[
 h_{z,w}(t)=e^{-zt}-e^{-wt}.
\]
Choose smooth functions $\eta,\zeta$ on $[0,\infty)$ with $0\leq\eta,\zeta
\leq1$, $\eta=0$ on $[0,1]$, $\eta=1$ on $[2,\infty)$, $\zeta=1$ on
$[0,1]$, and $\zeta=0$ on $[2,\infty)$.  Set
$f_{\epsilon,R}(t)=\eta(t/\epsilon)\zeta(t/R)h_{z,w}(t)$, where
$0<\epsilon\leq1$ and $R\geq1$.  These are admissible compactly supported
smooth test functions on $(0,\infty)$.

For $\beta\in[-B,B]$, the $L^1$ norms of
$g_{\epsilon,R,\beta}=e^{\beta t}f_{\epsilon,R}(t)$ and its second
derivative are bounded independently of $\epsilon,R,\beta$.  Here is the
estimate at the only potentially singular endpoint.  Since $h_{z,w}(0)=0$,
$h_{z,w}(t)=O(t)$, $h'_{z,w}(t)=O(1)$ and $h''_{z,w}(t)=O(1)$ for
$0\leq t\leq2$.  On $[\epsilon,2\epsilon]$ the terms containing
$\eta''(t/\epsilon)/\epsilon^2$ therefore have $L^1$ norm
$O(\epsilon^{-2}\epsilon\epsilon)=O(1)$, and those containing
$\eta'(t/\epsilon)/\epsilon$ also have $L^1$ norm $O(1)$.
The remaining terms near zero are bounded.  At infinity all terms are
bounded by a polynomial in $t$ times $e^{-ct}$, with
$c=\min(\Re z,w)-B>0$; the derivatives of $\zeta(t/R)$ have size
$O(R^{-1})$ and $O(R^{-2})$.  This proves the uniform bounds.
Twice integrating by parts in the Fourier variable, together with the
$L^1$ bound when $|\gamma|\leq1$, now gives
\[
 |\widehat f_{\epsilon,R}(\beta+i\gamma)|
 \leq\frac{C}{1+\gamma^2}
 \qquad(\beta\in[-B,B]).
\]
Dominated convergence in $t$, followed by dominated convergence in the
absolutely summable zero series, gives
\[
 0=\sum_{\lambda\in\mathcal E}m(\lambda)
       \left(\frac1{z-\lambda}-\frac1{w-\lambda}\right)
       \qquad(\Re z>B).
\]
The series on the right is normally convergent on every compact subset of
$\mathbb C\setminus\mathcal E$: outside a finite set each summand is bounded
by $C_Km(\lambda)/(1+|\Im\lambda|^2)$.  It therefore defines a
meromorphic function on $\mathbb C$.  At any $\lambda_0\in\mathcal E$
its residue is the strictly positive integer $m(\lambda_0)$; local finiteness
ensures that the other terms are holomorphic near $\lambda_0$.
The identity theorem, applied to the connected complement of the discrete
set $\mathcal E$, makes the meromorphic function identically zero.  A zero
meromorphic function has no pole, so $\mathcal E$ is empty.  This holds for
every $a_j$, proving $\Re\lambda=0$ for every $\lambda$.  The converse
is immediate because no terms are omitted.
\end{proof}

\paragraph{Finding F-SM-02 (the precise theorem-level comparison).}
For one fixed unitary Hecke character, suppose the established full explicit formula
and a proposed shrinking-strip trace formula have exactly the same
geometric distribution on test functions supported at idelic norm $>1$.
After taking its character isotype and setting $t=\log|g|$, their equality
is \eqref{eq:sm-strip} for $\lambda=\rho-1/2$.  The usual zero-counting
bound $N(T)=O(T\log(T+2))$ implies the summability required in
F-SM-01, so the trace formula implies GRH for this character.  Conversely,
GRH implies this numerical trace identity by the full explicit formula.
For the Dirichlet characters covered by Alp\"oge's adopted theorem, the
numerical identity is therefore available without assuming Masullo's
operator-theoretic construction.

\begin{proof}
The trace distributions subtract to exactly \eqref{eq:sm-strip}.  By a
dyadic decomposition,
\[
 \sum_{2^k<|\Im\lambda|\leq2^{k+1}}
 \frac{m(\lambda)}{1+|\Im\lambda|^2}
 \ll \frac{2^k(k+1)}{2^{2k}},
\]
whose sum is finite; local finiteness handles bounded ordinates.
Apply F-SM-01.  Conversely, all zeros lie in every open strip of positive
width centered on the critical line, so restricting the spectral sum removes
no terms.  This proves the equality of distributions.  No assertion about
the existence of a particular kernel, basis, or geometric realization is
used or obtained in this converse.
\end{proof}

The full idelic explicit formula is recorded in Masullo's Theorem~B.4.4.7,
page~623, whereas the restricted one is Theorem~12.5.1.1,
pages~519--520.  To apply F-SM-02 to the manuscript one must first align the
Haar-measure and multiplicity conventions in those statements, and prove
the missing deconstruction.  F-SM-01 is independent of those normalization
choices.  It also avoids passing separately to divergent archimedean
integrals: the difference of exponentials vanishes at the identity.

\subsection{A larger Artin bridge, conditional on Masullo's larger target}

\paragraph{Finding F-SM-03 (off-line Artin holomorphy from Hecke GRH).}
Assume GRH for every finite-order Hecke character of every number field.
Let $K$ be a number field and $V$ a finite-dimensional complex
representation of $\operatorname{Gal}(\overline K/K)$ with finite image.
Then its Artin $L$-function is holomorphic and nonzero on
\[
 U=\{s:0<\Re s<1,\quad\Re s\ne1/2\}.
\]
The same assertion for $V(n)$ holds on $U-n$ under the convention
$L(V(n),s)=L(V,s+n)$.

\begin{proof}
Choose a finite Galois extension $E/K$ through which $V$ factors and put
$G=\operatorname{Gal}(E/K)$.  Brauer induction gives
\[
 [V]=\sum_{j=1}^r n_j\operatorname{Ind}_{H_j}^{G}\psi_j,
 \qquad n_j\in\mathbb Z,\qquad\dim\psi_j=1.
\]
The $\psi_j$ have finite order.  With $K_j=E^{H_j}$, Artin reciprocity
and inductivity of Artin $L$-functions give the meromorphic identity
\[
 L_K(V,s)=\prod_{j=1}^r L_{K_j}(\psi_j,s)^{n_j}.
\]
Each Hecke factor is holomorphic in $0<\Re s<1$, including the trivial
character's Dedekind zeta function, whose possible pole is at $1$.
The assumed GRH says that no such factor vanishes on $U$.  An integral
power, positive, zero, or negative, of a nowhere-zero holomorphic function
is again nowhere-zero and holomorphic.  Their finite product has the same
property.  Substituting $s+n$ proves the twisted assertion.
\end{proof}

The exact representation-theoretic input is Serre (1977),
\emph{Linear Representations of Finite Groups}, \S10.5, Theorem~20,
page~78.  The $L$-function factorization is equation~(2) in Cogdell,
\emph{On Artin $L$-functions}, \S2.1, page~10; the compatibility of the
ramified local factors is explained in \S2.2, page~12.  This is a conditional
bridge to all finite-image Artin motives, substantially beyond the
abelian-over-$\mathbb Q$ scope of the supplied Alp\"oge theorem.
It does not prove holomorphy on the critical line: factors in the denominator
can vanish there.  Nor does the proof imply general automorphy, the Artin
conjecture, or the existence of a motivic realization functor.

\subsection{Two explicit corrections to the geometric analogies}

\paragraph{Finding F-SM-04 (the norm does not recover the unit strata).}
The assertion on Masullo page~493 that a saturated all-local-units stratum
in $\mathbb A_F^\times$ is the inverse image of the corresponding unit
stratum of $\mathbb A_{\mathbb Q}^\times$ under the norm map is false.
Here ``saturated'' has its explicit meaning from page~492: a set of places
contains all places over a rational prime whenever it contains one.
The displayed stratum on page~493 is a subset of the ideles themselves,
defined by $g_v\in\mathcal O_{F_v}^{\times}$ at each listed place;
it is not a subset enlarged by multiplication by $F^{\times}$.
The correct condition at a rational prime $p$ for the norm to be a unit is
\[
 \sum_{v\mid p} f(v/p)v_v(g_v)=0,
\]
where $v_v$ assigns value $1$ to a uniformizer and $f(v/p)$ is the residue
degree.  The stronger all-local-units condition is $v_v(g_v)=0$ for
every $v\mid p$.

\begin{proof}
For a finite extension $F_v/\mathbb Q_p$, a uniformizer has norm of
$p$-adic valuation $f(v/p)$, and the norm of a unit is a unit; by factoring
an arbitrary element into a uniformizer power times a unit,
$v_p(N_{F_v/\mathbb Q_p}(g_v))=f(v/p)v_v(g_v)$.  Multiplication of the
local norm factors proves the displayed formula.
Take $F=\mathbb Q(i)$ and $p=5$, which splits into two degree-one places
$v_1,v_2$.  Thus $F_{v_j}\simeq\mathbb Q_5$.  Let the idele $g$ have
$g_{v_1}=5$, $g_{v_2}=5^{-1}$ and every other component $1$.
Its norm is the identity idele because the two $5$-adic components multiply
to $1$.  Neither of these components is a local unit, so $g$ does not lie
in the asserted unit stratum.  This is the required counterexample.
The all-local-units stratum remains well-defined directly; the error in
identifying it with a norm inverse image does not by itself disprove
the later trace identity.
\end{proof}

\paragraph{Finding F-SM-05 (a maximal curve need not have gonality $g+1$).}
The equality for maximal curves asserted on Masullo page~535 is false.
Over $\mathbb F_9$, the smooth plane curve
\[
 H:\quad Y^3Z+YZ^3=X^4
\]
is maximal of genus $3$ and has gonality at most $3$, hence strictly
less than $g+1=4$.

\begin{proof}
In characteristic $3$, the three partial derivatives of
$Y^3Z+YZ^3-X^4$ are $-X^3,Z^3,Y^3$, so they have no common projective
zero.  The curve is smooth.  Over an algebraic closure, any two positive-degree
plane-curve components intersect by B\'ezout's theorem, creating a singular
point; a repeated component would also contradict smoothness.  Thus this
smooth plane quartic is geometrically integral.  Its genus is
$(4-1)(4-2)/2=3$ (the plane-curve genus formula, Stacks Project,
Lemma~53.9.3).  For each $x\in\mathbb F_9$, $x^4$
is the norm of $x$ and lies in $\mathbb F_3$; the trace
$y\mapsto y^3+y$ from $\mathbb F_9$ onto $\mathbb F_3$ has three
elements in each fibre.  There are therefore $9\cdot3=27$ affine
points.  At $Z=0$ the equation forces $X=0$, leaving the unique point
$[0:1:0]$.  Thus $\#H(\mathbb F_9)=28=9+1+2\cdot3\cdot3$,
which is the maximal value.  The nonconstant rational function $x=X/Z$
has function-field degree at most $3$, because $y=Y/Z$ satisfies
$y^3+y=x^4$ over $\mathbb F_9(x)$.  It extends to a morphism to
$\mathbb P^1$ on the smooth projective curve; consequently the gonality
is at most $3$.
\end{proof}

The correction strengthens rather than defeats the point being made in
Masullo \S14.2, pages~534--536: the proposed number-field multiplicity
mechanism cannot simply be transported to function fields.  For this
curve the Frobenius eigenvalue multiplicity in $H^1$ is $2g=6$, whereas
the gonality is at most $3$: the finite-field Weil theorem gives modulus
$3$ for each of its six eigenvalues, and their sum is
$9+1-28=-18$, so equality in the triangle inequality makes all six
equal to $-3$.

\subsection{The categorical and positivity information still required}

The stack diagrams are especially relevant to Grothendieck's programme,
but their stated status matters.  Masullo's equations~(8.4.0.6) and
(8.4.0.7), page~390, and (8.4.0.8)--(8.4.0.9), page~391, relate
Fr\'echet-space divisor sequences to periodic-orbit quotient stacks and
boundary divisors.  Remark~8.4.0.1 on page~391 explicitly says that the
vertical arrows are not morphisms.  They express an interpretation, not
a constructed sheaf, realization functor, or equivalence of categories.
Consequently these diagrams do not supply an arithmetic cohomology theory
of $\operatorname{Spec}\mathbb Z$, a six-functor formalism, or a comparison
with the \emph{D\'erivateurs}.  A literal categorical bridge would require
those arrows to be constructed and their compatibility with the trace,
tensor products, Tate twists, and local terms proved.

\paragraph{Finding F-SM-06 (spectral location alone does not supply
semisimplicity or a positive unitary realization).}
There exists a complex finite-dimensional operator with purely imaginary
spectrum that is not skew-adjoint for any positive definite Hermitian
inner product.

\begin{proof}
On $\mathbb C^2$ take $A=i\gamma I+N$, where $\gamma\in\mathbb R$,
$Ne_2=e_1$ and $Ne_1=0$.  Its spectrum is the singleton $\{i\gamma\}$.
If $A$ were skew-adjoint for a positive definite Hermitian form, then
$e^{tA}$ would preserve its norm for every real $t$: differentiating
$\langle e^{tA}v,e^{tA}v\rangle$ gives zero.  But
$e^{tA}e_2=e^{i\gamma t}(e_2+te_1)$, whose norm is at least
$|t|\|e_1\|-\|e_2\|$ by the triangle inequality and therefore tends
to infinity.  This contradicts norm preservation.
\end{proof}

This refutes a general inference, not RH and not the possibility of a
particular Hilbert--P\'olya realization.  Masullo's spectral discussion
expressly permits Jordan blocks (\S1.2.1, page~14), so zero location
must be separated from a theorem killing nilpotent parts or proving
zero simplicity.

\paragraph{Finding F-SM-07 (a bounded biorthogonal lift is not automatically
a Schauder basis).}
The proof step in Masullo's Lemma~10.4.2.5, page~462, that obtains
the Schauder basis property from the displayed boundedness and
biorthogonality alone is invalid.

\begin{proof}
Let $E=\ell^2(\{0,1,2,\ldots\})$, $E_0=\mathbb Ce_0$, $F=E$,
and let $x_n$ be the image of $e_n$ in $E/E_0$ for $n\geq1$.
These $x_n$ form the usual orthonormal basis.  Set $u_n=e_0+e_n$
and let $\phi_n$ be the $n$th coordinate functional.  Then
$\phi_n(u_m)=\delta_{nm}$ and $\|u_n\|=\sqrt2$.
The collar conditions hold with $r=1$, $r'=11/10$ and $\epsilon=1/2$.
Yet
\[
 \frac1N\sum_{n=1}^{N}u_n
   =e_0+\frac1N\sum_{n=1}^{N}e_n\longrightarrow e_0
\]
in $E$.  Thus $e_0$ belongs to the closed span of the $u_n$, while
every $\phi_n(e_0)=0$.  If the $u_n$ were a Schauder basis of that
closed span, applying the continuous biorthogonal coordinate functionals
to the expansion of $e_0$ would give every coefficient zero, forcing
$e_0=0$, a contradiction.
\end{proof}

This counterexample concerns the inference used in the written proof.
It does not refute the existence of suitable lifts in the Hilbert-space
case.  The following construction supplies that repair in the form relevant
to the Hardy--Sobolev Hilbert spaces.

\paragraph{Finding F-SM-08 (Hilbert-space basis lifting with a controlled
graph).}
Let $E$ be a Hilbert space, $E_0$ a closed subspace, $q:E\to E/E_0$ the
quotient, and $F\subset E$ a linear subspace such that $F\cap E_0$ is
dense in $E_0$.  Let $(x_n)$ be a Schauder basis of $E/E_0$ with
coefficient functionals $(\phi_n)$, or a Markushevich basis with
its total biorthogonal functionals $(\phi_n)$, and assume each
$x_n\in q(F)$.
For any positive numbers $\delta_n$ satisfying
$\sum_n\|\phi_n\|\delta_n<\infty$, there are $u_n\in F$ with
$q(u_n)=x_n$, $\|u_n-Sx_n\|<\delta_n$, where
$S:E/E_0\to E_0^\perp$ is the isometric orthogonal section, such that
$(u_n)$ is respectively a Schauder or Markushevich basis of its closed span.
That span maps
isomorphically to $E/E_0$ under $q$.

\begin{proof}
Choose $v_n\in F$ with $q(v_n)=x_n$.  Then $Sx_n-v_n\in E_0$.
Density of $F\cap E_0$ provides $w_n\in F\cap E_0$ with
$\|v_n+w_n-Sx_n\|<\delta_n$; put $u_n=v_n+w_n$ and
$d_n=u_n-Sx_n\in E_0$.  The series
\[
 Kx=\sum_n\phi_n(x)d_n
\]
converges absolutely for every $x\in E/E_0$ and defines a bounded
operator to $E_0$, with norm at most
$\sum_n\|\phi_n\|\delta_n$.  Set $T=S+K$.
Then $qT=\operatorname{id}$, so $\|Tx\|\geq\|x\|$ and $T$
is a bounded linear isomorphism onto a closed subspace of $E$.
The inverse is the restriction of $q$.  Since $Tx_n=u_n$, in the
Schauder case the image under $T$ of the basis expansion of $x$ is a
convergent basis expansion in $(u_n)$.  In the Markushevich case,
density of the span and totality of the biorthogonal functionals
$\phi_n\circ q$ on the image are transported by the isomorphism $T$.
Their closed span is $T(E/E_0)$ in both cases, proving every assertion.
\end{proof}

Thus one lift-related gap has an explicit repair.  This does not repair
all of Theorem~10.4.2.8 or the missing Section~12.3: the estimates for
the resulting integral kernels and the interchange of diagonal
integration with deconstruction are separate requirements.

\subsection{Comparison of methods and what each can actually transfer}

\paragraph{Finding F-SM-09 (the exact power-sum bridge to purity).}
Let $\alpha_1,\ldots,\alpha_r$ be distinct nonzero complex numbers,
let $c_1,\ldots,c_r$ be nonzero complex coefficients, and let $R>0$.
Suppose there are constants $C>0$ and $A\geq0$ such that
\[
 \left|\sum_{j=1}^r c_j\alpha_j^n\right|\leq C n^A R^n
 \qquad(n\geq1).
\]
Then $|\alpha_j|\leq R$ for every $j$.  If additionally the set of
$\alpha_j$ is invariant under $\alpha\mapsto R^2/\alpha$, then
$|\alpha_j|=R$ for every $j$.

\begin{proof}
The power-series
\[
 F(z)=\sum_{n\geq1}\left(\sum_{j=1}^r c_j\alpha_j^n\right)z^n
\]
converges absolutely on $|z|<R^{-1}$ by the assumed bound.  Near $0$ it
is the rational function
\[
 \sum_{j=1}^r\frac{c_j\alpha_jz}{1-\alpha_jz}.
\]
For any $j$, the latter has a genuine pole at $z=1/\alpha_j$:
the $j$th summand has residue $-c_j/\alpha_j\ne0$ and the other
summands are holomorphic there because the $\alpha_j$ are distinct.
If $|\alpha_j|>R$, this pole would lie in the disk on which the power
series is holomorphic, contradicting equality by analytic continuation.
Thus $|\alpha_j|\leq R$.  Under the stated reciprocal symmetry,
$R^2/\alpha_j$ is another listed number, so
$R^2/|\alpha_j|\leq R$, which gives the reverse inequality.
\end{proof}

Taking $R=q^{w/2}$ identifies the precise missing inputs for a purity
argument of the form suggested by Alp\"oge's footnote~3: bounds over
\emph{all} finite extensions, a nonzero net coefficient at each eigenvalue,
and the appropriate duality.  A bound for a single finite field is
insufficient for this deduction.  For the trace of an honest finite-dimensional
space the coefficients are the positive algebraic multiplicities, so
they cannot cancel.  For alternating cohomological traces this must
be proved separately.

\paragraph{Finding F-SM-10 (why alternating traces alone do not give purity).}
For any $R>0$ there is a graded vector space with Frobenius whose
alternating traces satisfy every bound in F-SM-09, while a Frobenius
eigenvalue has modulus greater than $R$.

\begin{proof}
Take $V^0=V^1=\mathbb C$, all other graded pieces zero, and let
Frobenius act by the same scalar $a>R$ on both nonzero pieces.
Then $\sum_i(-1)^i\operatorname{tr}(F^n\mid V^i)=a^n-a^n=0$
for every $n\geq1$, while both pieces have eigenvalue $a$.
\end{proof}

This is the exact point of a no-cancellation or positivity input;
it supplies spectral information that a virtual determinant can hide.
Even F-SM-09 determines eigenvalue moduli, not the nilpotent part of
Frobenius, as F-SM-06 illustrates.

Alp\"oge's step~5 (\emph{Remark}, page~3 and footnote~3) seeks a
noncancelling sign in cohomological power sums; his steps~7--10 and~12--16
(page~4) obtain arithmetic multiplicity and norm estimates from unitary
theta lifts, periods, and integral structures.  Masullo's primary gain is
different: a coarser completion retains only part of the zero divisor,
while the geometric trace is asserted to remain unchanged.  F-SM-01
shows precisely why preservation of that trace forces the omitted part
to be empty.  The first route needs a positivity/divisibility estimate;
the second needs trace compatibility through completion, lifting,
localization, and a limiting process.

Neither gain automatically gives algebraic cycles, crystalline
Frobenius, a realization functor, a weight filtration on an arbitrary
motive, or reconstruction of nonabelian fundamental groups.  In the
supplied text the Iwasawa tower is an arithmetic amplification device
(Alp\"oge steps~11 and~19), and the hexagonal CM curve supplies a
specific theta-geometric setting (step~3).  Neither is identified by a
theorem with the Teichm\"uller tower or the profinite GT group.
The testable bridges are the explicit $L$-function identities and the
trace/norm compatibilities just described.  The correspondence between
the words ``tower'', ``hexagon'', or ``multiplicity'' by itself does not
constitute any such compatibility.

The Masullo route also does not eliminate the logical possibility of a
geometric or foliated realization of a different construction.  Its
Sections~1.2.7 and~1.5 compare its own analytic objects with Deninger's
expectations; the absence of a three-manifold in this construction is
not a nonexistence theorem for all other constructions.

% END source_and_masullo.tex

% BEGIN source_masullo_numeric.tex
% Supplement to source_and_masullo.tex. Original proofs of the stated deductions;
% no claim of historical novelty. Background inputs are explicitly identified.
\section{The actual Hilbert-space sheets and a numerical completion over \texorpdfstring{$\mathbb Q$}{Q}}
\label{sec:masullo-numeric}

\subsection{What cote 1 actually says}

The supplied transcription of cote~1, scan~079, gives Corollaries~4 and~5:
coefficientwise domination of the Fredholm determinant of $T$ by that of
$|T|$, and hence $|\det(1+zT)|\leq\det(1+|z||T|)$.  The following
eigenvalue/product theorem and its trace corollary are crossed out, as is the
draft proof continuing on scan~080.  Scans~111--112 state and attempt to prove
the genus-zero lemma used by that draft.  These supplied sheets contain no
Riemann-$\xi$ function, arithmetic spectral realization, or unfinished question
asserting that such a realization should exist.  Thus BR-TR-01 and BR-TR-02
are arithmetic applications of this analytic framework.  The next two
findings also give complete proofs of the actual analytic assertions on the
supplied sheets, including the crossed-out product assertion.  ``Crossed
out'' is a description of this source, not a claim that the classical theorem
is false or remains open.

\paragraph{Finding F-SM-11 (the proposition on scans 111--112, corrected).}
Let $f(z)=\sum_{n\geq0}a_nz^n\not\equiv0$ and
$h(z)=\sum_{n\geq0}b_nz^n\not\equiv0$ be entire, with
$|a_n|\leq|b_n|$ for all $n$.  Suppose that $h$ has genus zero, and that the
nonzero zeros $z_j$ of $f$, repeated with multiplicity, satisfy
$\sum_j|z_j|^{-1}<\infty$.  Then $f$ has genus zero.  Here genus zero means
a constant times a power of $z$ times a normally convergent product of
factors $1-z/z_j$, with no nonconstant exponential factor.

\begin{proof}
Write $h(z)=c z^r\prod_j(1-z/w_j)$, where
$\sum_j|w_j|^{-1}<\infty$.  The function
$H(z)=|c|z^r\prod_j(1+z/|w_j|)$ has nonnegative coefficients that majorize
the absolute values of the coefficients of $h$.  To see this, first use the
finite products, then pass to their locally uniform limits; Cauchy's
coefficient formula preserves each coefficient inequality.  Therefore
$|f(z)|\leq H(|z|)$.

Every genus-zero product $P$ satisfies, for each $\varepsilon>0$,
$|P(z)|\leq C_\varepsilon e^{\varepsilon|z|}$.  Indeed, split its factors
after an index $N$ for which the sum of the reciprocal moduli in the tail is
less than $\varepsilon/2$.  The tail is bounded by
$\exp((\varepsilon/2)|z|)$, using $\log(1+x)\leq x$; the finite product
and the power of $z$ are bounded by a constant times
$\exp((\varepsilon/2)|z|)$.  Applying this to $H$ proves the same bound
for $f$.  Hadamard factorization for an entire function of order at most one
now gives
\[
 f(z)=e^{A+bz}P(z),\qquad P(z)=z^m\prod_j(1-z/z_j).
\]
More explicitly, the genus-one factors in the usual Hadamard formula may
be replaced by genus-zero factors, because $\sum_j1/z_j$ converges
absolutely; the resulting linear exponential is absorbed into $b$.

Suppose $b\ne0$.  Replace $z$ by $\overline b z/|b|$ to make $b>0$;
the preceding radial bounds are unchanged.  Fix
$\pi/3<\theta<\pi/2$, and put $\delta=b\cos\theta/2>0$.  On the sector
$|\arg z|\leq\theta$, the bound for $f$ with $\varepsilon=\delta$ gives
\[
 |P(z)|\leq C e^{-\delta|z|}.
\]
Globally the genus-zero bound for $P$ gives
$|P(z)|\leq C'e^{\delta|z|/3}$.  Set $\omega=e^{2\pi i/3}$.
For every nonzero $z$, at least one of $z,\omega z,\omega^2z$ belongs
to this sector, because its half-width exceeds $\pi/3$.  Hence the entire
function $G(z)=P(z)P(\omega z)P(\omega^2z)$ satisfies
\[
 |G(z)|\leq C(C')^2e^{-\delta|z|/3}.
\]
Liouville's theorem makes $G$ constant, and the displayed bound makes that
constant zero.  A finite product of nonzero entire functions cannot vanish
identically.  This contradicts $P\not\equiv0$, so $b=0$.
\end{proof}

Three transcription-level points matter.  Scan~112 literally specifies
$\pi/3<\theta<\pi$ while subsequently requiring
$\Re z\geq c|z|$ throughout the sector, for $c>0$.
The latter assertion requires $\theta<\pi/2$; the proof above makes the
compatible choice $\pi/3<\theta<\pi/2$.  Scan~111 writes a rotation by
$b/|b|$; with the substitution $f(\Theta z)$ the correcting rotation is
$\Theta=\overline b/|b|$.  Finally, coefficient domination gives
$|f(z)|\leq H(|z|)$, not generally $|f(z)|\leq|h(z)|$ at the same complex
argument, as written on scan~112.  For example, $f=1$ and $h=1+z$
satisfy the coefficient inequality but violate the latter inequality at
$z=-1$.  These corrections repair the proof of the stated proposition.
The uncertain sum/product sign in the transcription of scan~111 is not
used as a premise: the canonical-product hypothesis is stated explicitly.

\paragraph{Finding F-SM-12 (the Fredholm assertions on scans 079--080).}
Let $T$ be a trace-class operator on a complex Hilbert space.  Put
$D_T(z)=\sum_{n\geq0}\alpha_n(T)z^n$, where
$\alpha_n(T)=\operatorname{Tr}(\bigwedge^nT)$ and $\alpha_0(T)=1$.
If $s_j$ are the singular values and $\lambda_j$ the nonzero eigenvalues,
each repeated with its multiplicity (algebraic multiplicity for
$\lambda_j$), then
\[
 |\alpha_n(T)|\leq\alpha_n(|T|),\quad
 |D_T(z)|\leq D_{|T|}(|z|),\quad
 \sum_j|\lambda_j|\leq\|T\|_1,
\]
\[
 D_T(z)=\prod_j(1+z\lambda_j),\qquad
 \operatorname{Tr}T=\sum_j\lambda_j.
\]
In particular $D_T$ has genus zero.  The elementary trace-class ideal,
exterior-power, and compact-operator spectral theorems are the analytic
background here; no arithmetic hypothesis is used.

\begin{proof}
The singular values of $\bigwedge^nT$ are the products
$s_{j_1}\cdots s_{j_n}$ with $j_1<\cdots<j_n$.  Thus
\[
 |\operatorname{Tr}(\bigwedge^nT)|
 \leq\|\bigwedge^nT\|_1
 =\sum_{j_1<\cdots<j_n}s_{j_1}\cdots s_{j_n}
 =\alpha_n(|T|).
\]
The last quantity is at most $\|T\|_1^n/n!$, proving entire convergence
and the two domination inequalities.  Diagonalizing the positive compact
operator $|T|$ gives directly
$D_{|T|}(z)=\prod_j(1+zs_j)$, a genus-zero product.

For any finite collection of nonzero eigenvalues, include their full
generalized eigenspaces and let $E$ be their finite-dimensional invariant
sum.  A unitary triangularization of $T|_E$ gives an orthonormal basis
$e_1,\ldots,e_d$ with diagonal entries the selected eigenvalues.  There is
a diagonal unitary on $E$, extended by zero on $E^\perp$, whose product
with $T$ has trace $\sum_{j=1}^d|\lambda_j|$.  Its operator norm is at
most one.  The trace ideal inequality therefore bounds this sum by
$\|T\|_1$.  Increasing the finite collection proves the summability
assertion.

For completeness, the determinant zero information needed next does not
assume the desired product formula.  The exterior-power definition is
continuous under trace-norm approximation, uniformly for $z$ in compact
sets.  For fixed $n$, this follows by telescoping the tensor product
$A^{\otimes n}-B^{\otimes n}$ and restricting by the bounded
antisymmetrizing projection.  The bounds
$|\alpha_n(A)|\leq\|A\|_1^n/n!$ make the tails uniformly small on
bounded trace-norm sets, proving the asserted locally uniform continuity.
Finite-rank determinant identities consequently imply similarity
invariance and
$\det((1+A)(1+B))=\det(1+A)\det(1+B)$ for trace-class $A,B$.
In particular, if $1+A$ is invertible, its inverse has the form $1+B$
with $B$ trace class, so $\det(1+A)\ne0$.
For a nonzero eigenvalue $\lambda$, the compact spectral theorem gives
the decomposition by its finite-dimensional Riesz subspace and its
invariant closed complement.  Similarity invariance and the direct-sum
identity factor $D_T$ into the ordinary determinant on the Riesz subspace
and the determinant on the complement.  At $z=-1/\lambda$, the latter
factor is nonzero, and the former has zero of order exactly the algebraic
multiplicity of $\lambda$.  At every other point where $1+zT$ is
invertible the determinant is nonzero.  These are all its zeros by the
compact spectral theorem.

Consequently the nonzero zeros of $D_T$ have summable reciprocal moduli.
Its coefficients are dominated by those of the genus-zero function
$D_{|T|}$.  Finding F-SM-11 gives genus zero for $D_T$.  Its zero divisor
and $D_T(0)=1$ identify it as $\prod_j(1+z\lambda_j)$.  Normal convergence
and $\sum_j|\lambda_j|<\infty$ allow differentiation at zero, giving
$\operatorname{Tr}T=\sum_j\lambda_j$.
\end{proof}

This proves the classical Hilbert-space product and trace assertions
written in the deleted passage without using them to prove the
genus-zero lemma.  It should not be described as an RH-dependent solution
of an arithmetic question on those sheets.

\subsection{A fully normalized numerical form of Masullo's main identity}

We now separate a numerical consequence that can be completed from the
adopted Dirichlet theorem from the geometric deconstruction not supplied
in Masullo's Section~12.3.  No kernel-decomposition claim is needed for
the following result.

\paragraph{Finding F-SM-13 (Masullo 12.5.1.1 over $\mathbb Q$, as a numerical distribution).}
Adopt Alp\"oge's Dirichlet zero-location theorem.  Let $\chi$ be a
primitive Dirichlet character of conductor $q$, including the trivial
character of conductor one; write $\chi(-1)=(-1)^a$ with $a\in\{0,1\}$.
Let $h\in C_c^\infty(\mathbb R_{>0})$ satisfy $h(1)=0$, and put
\[
 \widehat h(z)=\int_0^\infty h(t)t^z\,\frac{dt}{t},
 \qquad\delta_\chi=\begin{cases}1,&q=1,\\0,&q>1.\end{cases}
\]
Define the following concrete linear functionals:
\[
 A_a(h)=\frac12\int_0^\infty h(t)t^{1/2}
 \left(\frac1{|1-t|}+\frac{(-1)^a}{1+t}\right)\frac{dt}{t}
 =\int_0^\infty
 \frac{e^{-(a+1/2)x}}{1-e^{-2x}}
 \bigl(h(e^x)+h(e^{-x})\bigr)\,dx,
\]
\[
 B_\chi(h)=\sum_{n\geq1}\frac{\Lambda(n)}{\sqrt n}
 \bigl(\chi(n)h(n)+\overline\chi(n)h(n^{-1})\bigr).
\]
The integrals are ordinary absolutely convergent integrals, and the prime
power sum is finite.  For every $\alpha>0$ one has
\begin{equation}\label{eq:SM-numeric-strip}
 A_a(h)+B_\chi(h)
 =\delta_\chi\bigl(\widehat h(1/2)+\widehat h(-1/2)\bigr)
 -\!\sum_{\substack{\rho:\,L(\rho,\chi)=0,\ 0<\Re\rho<1\\
                         |\Re\rho-1/2|<\alpha}}
       m_\chi(\rho)\widehat h(\rho-1/2).
\end{equation}
The zero sum is absolutely convergent.  Thus this is an equality of
continuous linear functionals on the closed test-function subspace
$\{h\in C_c^\infty(\mathbb R_{>0}):h(1)=0\}$, with its usual inductive
limit topology.

Under the canonical identification
$C_{\mathbb Q}\simeq\widehat{\mathbb Z}^{\times}\times\mathbb R_{>0}$,
take Haar measure $du\,dt/t$, where $du$ has total mass one, and set
\[
 \phi_{\chi,h}(u,t)=\widehat\chi(u)h(t),\qquad
 c_{\chi,s}(u,t)=\widehat\chi(u)^{-1}t^s.
\]
Here $\widehat\chi$ denotes the finite character on
$\widehat{\mathbb Z}^{\times}$ induced by $\chi$.
Here and above $L(s,\chi)$ denotes the finite Dirichlet series.  Masullo's
notation $L(c)$ includes the archimedean factor (Definition~3.2.1.2,
page~123).  With these conventions his finite factor satisfies
$L_f(c_{\chi,s})=L(s,\chi)$, and his completed function is
\[
 L_{\mathrm M}(c_{\chi,s})=
 \pi^{-(s+a)/2}\Gamma((s+a)/2)L(s,\chi).
\]
It has exactly the nontrivial zero divisor appearing above.  Moreover
\[
 M(\phi_{\chi,h})(c_{\chi,s}|\cdot|^{-1/2})
       =\widehat h(s-1/2).
\]
Equation~\eqref{eq:SM-numeric-strip} is exactly the $\chi$-component of
Masullo's Theorem~12.5.1.1, pages~519--520, for $F=\mathbb Q$ and
$I=[-\alpha,\alpha]$.  Finite Fourier decomposition consequently proves
that numerical theorem for every $\phi\in\mathcal D(C_{\mathbb Q})$
with $\phi(C_{\mathbb Q}^1)=0$, in Masullo's own Haar normalization.

\begin{proof}
We give the contour calculation as well as the local normalization checks.
The classical analytic inputs are the continuation and functional equation
of primitive Dirichlet $L$-functions, the order-one completed function,
and its zero count $N_\chi(T)=O_\chi(T\log(2+T))$.  These are background
theorems; this proof does not claim to re-prove the analytic continuation
or the functional equation.  They are the standard inputs behind the
explicit formula (Iwaniec and Kowalski, 2004, Sections~4.6, 5.3, and~5.5).

Let
\[
 \mathcal L_\chi(s)=(q/\pi)^{(s+a)/2}
            \Gamma((s+a)/2)L(s,\chi),\qquad H(s)=\widehat h(s-1/2).
\]
The completed function has precisely the nontrivial zeros, and has simple
poles at $0,1$ exactly when $\delta_\chi=1$.  Its functional equation is
$\mathcal L_\chi(s)=\varepsilon_\chi\mathcal L_{\bar\chi}(1-s)$,
where the nonzero constant $\varepsilon_\chi$ disappears after logarithmic
differentiation.  Integration by parts after $t=e^x$ shows that $H(s)$
decays faster than any inverse power of $|\Im s|$, uniformly on every
fixed vertical strip.  This and the zero count prove absolute convergence
of $\sum_\rho m_\chi(\rho)H(\rho)$.

Integrate $(\mathcal L'_\chi/\mathcal L_\chi)(s)H(s)$ around rectangles
with vertical sides $c=3/2$ and $1-c=-1/2$.  Horizontal sides can be chosen
to contribute zero in the limit.  One explicit justification uses the
Hadamard logarithmic derivative of
$[s(s-1)]^{\delta_\chi}\mathcal L_\chi(s)$: in each interval $[j,j+1]$
choose a height at distance at least $c_0/(j\log(2+j))$ from the ordinates
of all zeros of modulus at most $2j+4$; the zero count leaves such a height
for a sufficiently small fixed $c_0>0$.  Do the same below the real axis.
The nearby terms of the logarithmic derivative are then polynomially
bounded in $j$, and its tail is bounded using
$\sum_{|\rho|>R}|\rho|^{-2}=O_\chi(\log(2+R)/R)$.
The rapid decay of $H$ kills both horizontal integrals.

The residue theorem and the functional equation therefore give
\begin{align*}
 \sum_\rho m_\chi(\rho)H(\rho)-\delta_\chi(H(0)+H(1))
 =\frac1{2\pi i}\int_{(c)}\biggl[
 \frac{\mathcal L'_\chi}{\mathcal L_\chi}(s)H(s)
 +\frac{\mathcal L'_{\bar\chi}}{\mathcal L_{\bar\chi}}(s)H(1-s)
 \biggr]\,ds.
\end{align*}
The constant $\tfrac12\log(q/\pi)$ in each logarithmic derivative
contributes $\log(q/\pi)h(1)=0$, by Mellin inversion.  On this line the
Dirichlet series for $L'/L$ converges absolutely.  Mellin inversion gives
\[
 \frac1{2\pi i}\int_{(c)}H(s)n^{-s}\,ds=n^{-1/2}h(n),\qquad
 \frac1{2\pi i}\int_{(c)}H(1-s)n^{-s}\,ds=n^{-1/2}h(n^{-1}).
\]
Thus the two $L'/L$ terms contribute $-B_\chi(h)$.

For the gamma terms use, for $\Re z>0$,
\[
 \psi(z)=-\gamma+\int_0^\infty
                 \frac{e^{-v}-e^{-zv}}{1-e^{-v}}\,dv.
\]
First truncate the $v$-integral to $[\epsilon,R]$.  All constant terms
pair with $h(1)=0$.  Fubini and Mellin inversion give for the remaining
term
\[
 -\frac12\int_\epsilon^R
 \frac{e^{-(a/2+1/4)v}}{1-e^{-v}}
       \bigl(h(e^{v/2})+h(e^{-v/2})\bigr)\,dv.
\]
The truncated digamma integrals on the line $\Re s=c$ are bounded by
$C(1+|\Im s|)$ uniformly in $\epsilon,R$: near zero use the mean-value
bound on $e^{-v}-e^{-(s+a)v/2}$, and away from zero use exponential
decay.  This justifies passage to the limit against rapidly decreasing
$H(s)+H(1-s)$.  The resulting real integral is absolutely convergent,
because $h(1)=0$ and $h$ is smooth, and equals $-A_a(h)$ after $v=2x$.
The expression for $A_a$ in the statement follows by splitting its
$t$-integral at one and substituting $t=e^x,e^{-x}$; the elementary
identity
\[
 \frac12\left(\frac1{1-r}+\frac{(-1)^a}{1+r}\right)
       =\frac{r^a}{1-r^2},\qquad 0<r<1,
\]
gives exactly its displayed coefficient.  We have proved the full
explicit formula, with the sum over all nontrivial zeros.
The adopted Dirichlet theorem puts every such zero on $\Re\rho=1/2$.
For every $\alpha>0$ its restriction to the indicated open strip therefore
omits no zero, proving~\eqref{eq:SM-numeric-strip}.

We now identify the idelic terms.  For an idele $b=(b_\infty,(b_p)_p)$,
put $r=\operatorname{sgn}(b_\infty)\prod_p p^{v_p(b_p)}\in\mathbb Q^\times$.
Its class has the unique representative
$(t,(u_p)_p)=(b_\infty/r,(b_p/r)_p)$ with $t>0$ and all $u_p$ units;
moreover $t=|b|$.  This proves the stated group identification.
Masullo's Definitions~A.3.2.1--A.3.2.2, page~560, give
$\mu_\infty=dx/|x|$, $\mu_p(\mathbb Z_p^\times)=1$, and
\[
 \mu_1(C_{\mathbb Q}^1)=\frac{2^{r_1}(2\pi)^{r_2}h_{\mathbb Q}R_{\mathbb Q}}
 {\sqrt{|d_{\mathbb Q}|}w_{\mathbb Q}}=\frac{2}{2}=1.
\]
Hence the quotient measure is precisely $du\,dt/t$, without an omitted
volume factor.

For $p\nmid q$, the local idele $b_p=p^kz$ with
$z\in\mathbb Z_p^\times$ has
$\phi_{\chi,h}(b_p)=\chi(p)^{-k}h(p^{-k})$.  The $k=0$ term is zero;
on $k>0$ one has $|1-b_p|_p=1$, and on $k=-m<0$ one has
$|1-b_p|_p=p^m$.  Each shell has multiplicative measure one.  Its local
integral, multiplied by $\log p$, is therefore
\[
 (\log p)\sum_{m\geq1}p^{-m/2}
 \bigl(\chi(p)^m h(p^m)+\overline\chi(p)^m h(p^{-m})\bigr).
\]
For $p\mid q$, primitivity makes the local unit character nontrivial;
unit averaging annihilates each shell with $k\ne0$, while $k=0$ is again
identically zero since $h(1)=0$.  This gives exactly $B_\chi(h)$.
At the real place,
$\phi_{\chi,h}(x)=\operatorname{sgn}(x)^a h(|x|)$; splitting
$\mathbb R^\times$ into its two components and retaining Masullo's
factor $1/2$ gives exactly $A_a(h)$.  The potential singularity at $x=1$
is integrable because $h(1)=0$, so Weil's principal value agrees with the
ordinary integral here.

Orthogonality on the mass-one compact group gives the asserted Mellin
transform, gives zero for all other finite-character components, and gives
$M(\phi_{\chi,h})(|\cdot|^{\pm1/2})=
\delta_\chi\widehat h(\pm1/2)$.  Every Bruhat test function on
$C_{\mathbb Q}$ factors through a finite quotient of its compact factor
and is smooth with compact support on its positive-real factor.  Its
finite Fourier expansion is therefore
$\phi(u,t)=\sum_\chi\widehat\chi(u)h_\chi(t)$, with each character
indexed once by its primitive conductor.  The condition
$\phi(u,1)=0$ for all $u$ is equivalent to $h_\chi(1)=0$ for every
component.  Adding the finitely many component identities proves the
stated assertion for all $\phi$.  Continuity follows on each fixed
compact-support/finite-quotient test space from the local integral bounds,
the finite prime-power sum, and the rapid-decay estimate for the absolutely
convergent zero sum; these estimates are compatible with the inductive
limit topology.
\end{proof}

\paragraph{What is completed and what remains geometric.}
Finding F-SM-13 identifies the numerical distribution in Masullo's
restricted-strip formula for every allowed test function over
$\mathbb Q$, and gives all conductor, parity, Haar, and Mellin
normalizations.  The conductor term vanishes here for the precise reason
$h(1)=0$; it has not been silently discarded in a formula for unrestricted
test functions.  The statement indexes each idelic character by its
primitive conductor; treating an imprimitive Dirichlet series as though it
were that same completed idelic $L$-function would introduce extra deleted
Euler factors.

This numerical identity does not prove Masullo's geometric assertion
\[
 \int_{C_F}k(x,x)\,d\mu(x)
 =\int_{C_F}k_S(x,x)\,d\mu(x)
   +\sum_{p\nmid S}\int_{C_F}k_p(x,x)\,d\mu(x),
\]
nor identify the kernels in the missing Section~12.3, justify an exchange
of the local sum with diagonal restriction/integration, or establish its
Sobolev, nuclearity, and quotient compatibility hypotheses.  Equality of
two already identified numerical distributions does not establish that
a proposed kernel construction realizes them.  It does, however, complete
the exact numerical statement of Theorem~12.5.1.1 for $\mathbb Q$ from
the adopted Alp\"oge theorem and classical analytic inputs.  No extension
to arbitrary number fields or their Hecke characters is made by this
argument.

% END source_masullo_numeric.tex

% BEGIN source_masullo_kernel.tex
\section{The actual Masullo kernel: a diagonal bound, an indentation defect, and a concrete repair}
\label{sec:SM-actual-kernel}

We use the operators actually named in the manuscript:
\[
 H=W_I^s(C_F)=L^2(C_F,w_{I,s}\,d\mu),\quad
 S_\phi=A(\phi)P_a^{\mathrm{spec}},\quad
 T_{\rho,\phi}=M_\rho A(\phi)P_a^{\mathrm{geom}}.
\]
The projections are Definition~10.4.1.1, p.~456; nuclearity for
$S_\phi$ is asserted in Theorem~10.4.2.8, pp.~463--464.
Definition~10.4.3.7, pp.~473--474, gives
$k^{\mathrm{geom}}_\phi(x,y)=\phi(xy^{-1})-k^{\mathrm{spec}}_\phi(x,y)$.
Equation~(10.4.3.48), p.~473, assigns to $T_{\rho,\phi}$ the
weighted-measure kernel
$w_{I,s}(y)^{-1}\rho(x)k^{\mathrm{geom}}_\phi(x,y)$.
The first result is conditional on the indicated actual $S_\phi$
being trace class and on this kernel identification.  It does not
replace that operator by a separately constructed zero spectrum or
recheck the entire earlier nuclearity argument.

\subsection{An explicit diagonal convergence estimate}

\paragraph{Finding F-SM-18 (trace-density domination and valid-indentation shell convergence).}
Let $\nu=w\mu$, where $\mu$ is a positive Radon measure and $w>0$.
Every trace-class operator $S$ on $L^2(X,\nu)$ has a unique
$d_S^\nu\in L^1(X,\nu)$ with
\[
 \operatorname{Tr}(M_fS)=\int_X f\,d_S^\nu\,d\nu
 \quad(f\in L^\infty),\qquad \|d_S^\nu\|_1\leq\|S\|_1.
 \tag{SM.18.1}
\]
Its density relative to $\mu$ is $d_S^\mu=w\,d_S^\nu$, with
the same $L^1$ norm.  For the actual $S_\phi$, if $\phi(1)=0$,
the trace-density geometric diagonal is $d=-d_{S_\phi}^\mu$, so
\[
 \|d\|_{L^1(\mu)}\leq\|S_\phi\|_1.
 \tag{SM.18.2}
\]
When the manuscript's averaging procedure gives the trace
diagonal of its spectral kernel, this agrees with its smoothened
geometric diagonal, since the convolution kernel has continuous
diagonal $\phi(1)$.  Its cutoff diagonal is $\rho d$;
for a sequence $|\rho_j|\leq1$ with $\rho_j\to1$ almost
everywhere, $\rho_jd\to d$ in $L^1$.

Suppose now $X=G/H$, with $G$ second countable, locally compact
abelian and $H$ closed and discrete with counting measure.
Take $\mu$ to be the quotient Haar measure and $D$ a measurable
fundamental domain carrying its corresponding measure $\mu_D$.
For a nonnegative measurable normalized indentation $b$, put
\[
 \sum_{r\in H}b(r^{-1}x)=1,\qquad
 \omega_b(x,h)=\sum_{r\in H}b(r^{-1}x)b(r^{-1}h^{-1}x).
 \tag{SM.18.3}
\]
Let measurable subsets $G_N$ increase to $G$, set
$\alpha_N=1_{G_N}$, and define on $D$
\[
 \beta_N(x)=\sum_{h\in H}\omega_b(x,h)
                       (1-\alpha_N(h^{-1}x)).
 \tag{SM.18.4}
\]
Then $0\leq\beta_N\leq1$, $\beta_N\to0$, and the omitted
absolute diagonal mass is bounded, uniformly for $|\rho|\leq1$, by
\[
 \int_D\sum_h|\rho(\pi x)d(\pi x)\omega_b(x,h)|
                 (1-\alpha_N(h^{-1}x))\,d\mu_D(x)
 \leq \int_D|d(\pi x)|\beta_N(x)\,d\mu_D(x)\longrightarrow0.
 \tag{SM.18.5}
\]
The initial-stratum and successive-shell diagonal densities
are absolutely summable in $L^1$, with sum of their norms equal
to $\|d\|_1$.  Thus diagonal integration and bounded cutoff
limits commute with this shell sum.  This conclusion does not
identify the pieces with independently specified local operators.

\begin{proof}
Write a singular-value expansion
$S=\sum_j\sigma_j u_j\otimes v_j^*$ with unit vectors $u_j,v_j$.
The series $\sum_j\sigma_j u_j\overline{v_j}$ converges in $L^1$,
since $\|u_j\overline{v_j}\|_1\leq1$ and
$\sum_j\sigma_j=\|S\|_1$.  Its sum is $d_S^\nu$.
The rank-one trace identity, followed by trace-norm and $L^1$
limits, proves (SM.18.1).  Testing against all bounded functions
gives uniqueness; the same estimate proves contractivity of
$S\mapsto d_S^\nu$.  Change of measure gives $d_S^\mu$.
The defining difference for the geometric kernel, with $\phi(1)=0$,
gives (SM.18.2), and dominated convergence gives the cutoff limit.

The discrete subgroup $H$ is countable.  All terms in $\omega_b$
are nonnegative, and Tonelli gives
\[
 \sum_h\omega_b(x,h)
 =\sum_r b(r^{-1}x)\sum_h b(r^{-1}h^{-1}x)=1.
 \tag{SM.18.6}
\]
Consequently $\beta_N$ lies in $[0,1]$.  For fixed $x,h$,
$\alpha_N(h^{-1}x)\to1$; dominated convergence with the
probability weights $\omega_b(x,h)$ gives $\beta_N(x)\to0$.
Another dominated convergence argument against $|d|$ proves
(SM.18.5), including uniformity in $\rho$.

Explicitly, let $c_0=\sum_h\omega_b(x,h)\alpha_0(h^{-1}x)$
and $c_j=\sum_h\omega_b(x,h)(\alpha_j-\alpha_{j-1})(h^{-1}x)$
for $j\geq1$.  Then $c_j\geq0$ and $\sum_jc_j=1$.
The shell densities are $d(\pi x)c_j(x)$, and Tonelli yields
$\sum_j\|dc_j\|_1=\|d\|_1$.  This proves every stated interchange.
\end{proof}

For a common-Hilbert-space realization of a proposed trace-class
operator remainder $E_N$, an equivalent trace-norm target is
\[
 \|d_{E_N}\|_1\leq\|E_N\|_1
   =\inf_{E_N=U_NV_N}\|U_N\|_2\|V_N\|_2\longrightarrow0.
 \tag{SM.18.7}
\]
The infimum runs over Hilbert--Schmidt factorizations.
The inequality $\|UV\|_1\leq\|U\|_2\|V\|_2$ is
Cauchy--Schwarz with trace duality.  Conversely, for
$E=W|E|$, take $U=W|E|^{1/2}$ and $V=|E|^{1/2}$;
both squared Hilbert--Schmidt norms equal $\operatorname{Tr}|E|$.
This proves the infimum formula.  Its vanishing is equivalent
to trace-norm convergence, and is sufficient for diagonal
convergence; it is not necessary for a scalar trace equality alone.
Hilbert--Schmidt convergence alone is insufficient: $N^{-1}$
times a rank-$N$ orthogonal projection has Hilbert--Schmidt norm
$N^{-1/2}$, trace norm one, and trace one.

Theorem~11.1.2.1, pp.~480--483, states a cofinal intersection
of closed carving subspaces and an identification by special
theta functions.  It does not state a common-space realization
and remainder estimate (SM.18.7).  On the other hand,
(SM.18.5) solves diagonal summation once the actual spectral
trace-class input and a normalized indentation are available.
Identifying those diagonal pieces with the local kernels remains
an additional assertion in the unwritten Section~12.3.

\subsection{The displayed orbital indentation vanishes}

\paragraph{Finding F-SM-19 (literal formula defect and nonzero test).}
Interpret equation~(11.3.2.2), p.~495, as the scalar-valued
function it is called, with the ordinary Haar measure
$d\mu_1(g)$ of its adjacent equations:
\[
 U(x,y)=\sum_{\xi\in F^\times}\int_{E_F^1}
   1_{\tau(\mathbb R_{>0})}(\xi gx)
   1_{\tau(\mathbb R_{>0})}(\xi gy)\,d\mu_1(g).
 \tag{SM.19.1}
\]
For $F=\mathbb Q$ this is identically zero, is not a normalized
indentation, and cannot yield the stated orbital trace formula
(11.3.2.1) under this literal interpretation.

\begin{proof}
Definition~A.2.3.6, p.~549, equation~(A.2.3.14), defines
$E_F^1$ as an actual subset of $\mathbb A_F^{\times,1}$.
For $\mathbb Q$ it is
$E_{\mathbb Q}^1=\{1\}_\infty\times\widehat{\mathbb Z}^{\,\times}$.
Definitions~A.3.2.1--2, pp.~559--560, use ordinary Haar measure
here, of volume one, with counting measure assigned only to
$\mathbb Q^\times$.  This infinite compact group has atomless
Haar measure: a positive singleton mass $c$ would give mass
$Nc$ to $N$ distinct points for every $N$, a contradiction.

Fix $x$ and $\xi$.  If $\xi gx=\tau(t)$ for $g\in E_{\mathbb Q}^1$,
the product formula and $|\tau(t)|=t$ force $t=|x|$.
Thus $g=\xi^{-1}\tau(|x|)x^{-1}$ is the only possible value.
The first indicator in (SM.19.1) is supported on at most one
point of $E_{\mathbb Q}^1$.  Each integral is zero and the
countable sum is zero, for all $x,y$.

For a nonzero right side of the orbital formula, choose
$F=\mathbb Q$, $\chi=1$, $S=\{\infty\}$, $a=1/2$, $s=4$,
and any $\alpha>0$.  Then $\operatorname{Pic}(\mathbb Z)=1$
and all stated ramification assumptions hold.
Choose $h\in C_c^\infty((3/2,5/2))$ with $h(2)=1$ and set
$\phi(u,t)=h(t)$ on
$C_{\mathbb Q}\simeq\widehat{\mathbb Z}^{\,\times}\times\mathbb R_{>0}$.
This is a Bruhat test function, finite-unit invariant and
supported over norms greater than one.  For the finite uniformizer
idele $\pi_2$, $|\pi_2|=1/2$, whence
\[
 \sum_{m\in\mathbb Z}\phi(\pi_2^m)
       =\sum_mh(2^{-m})=1.
\]
For the shell containing the place $2$, the stated right side is
$\log2/\mu_2(\mathbb Z_2^\times)>0$; the stated normalization
makes this $\log2$.  The literal indented kernel on the left
is zero, including after cutoffs, smoothening, and any integration
in the remaining norm variable.  The claimed equality therefore
fails with formula (SM.19.1) interpreted as printed.
\end{proof}

The introductory discussion, pp.~21--22, explicitly speaks of
scalar-valued measurable functions and allows indicators of
fundamental domains.  It introduces no Dirac interpretation
for the indicators in (SM.19.1).  A distributional replacement
or independently normalized disintegration would be a new
formula requiring its own operator and diagonal estimates.
This finding does not rule out such a corrected realization.

\subsection{A valid rational indentation and compact averaging}

\paragraph{Finding F-SM-20 (two explicit repairs and the remaining identification).}
For $\mathbb Q$, put
$Z=\mathbb R_{>0}\times\widehat{\mathbb Z}^{\,\times}
\subset\mathbb A_{\mathbb Q}^\times$ and $b=1_Z$.
Then $b$ is a continuous normalized indentation and
$b^\Delta(x,y)=\sum_\xi b(\xi^{-1}x)b(\xi^{-1}y)$ is valid.
For $x\in Z$ and $q\in\mathbb Q^\times$,
\[
 b^\Delta(x,qx)=1_{\{q=1\}}.
 \tag{SM.20.1}
\]
Hence this lift preserves the global diagonal but does not itself
recover the nontrivial prime-orbital formulas.

There is also a legitimate compact averaging for the actual operators.
Write $K=\widehat{\mathbb Z}^{\,\times}$ with probability Haar
measure, and assume $\phi(u,t)=h(t)$, $\rho(u,t)=\rho(t)$.
For $H_0=L^2(\mathbb R_{>0},w(t)\,dt/t)$, put
\[
 (Jf)(u,t)=f(t),\quad (J^*F)(t)=\int_KF(u,t)\,du,\quad E=JJ^*.
\]
Then $E$ is the orthogonal projection onto $K$-invariants, and
\[
 T_{\rho,\phi}=ET_{\rho,\phi}E,\qquad S_\phi=ES_\phi E.
 \tag{SM.20.2}
\]
For either operator $T$ in (SM.20.2), if $T$ is trace class,
\[
 \operatorname{Tr}_H T=\operatorname{Tr}_{H_0}(J^*TJ),
 \qquad \|T\|_1=\|J^*TJ\|_1.
 \tag{SM.20.3}
\]
If $K_T$ is its weighted Hilbert--Schmidt kernel, the compressed
kernel is the actual compact average
\[
 K_0(t,s)=\int_K\int_KK_T((u,t),(v,s))\,du\,dv,\qquad
 \|K_0\|_2\leq\|K_T\|_2.
 \tag{SM.20.4}
\]

\begin{proof}
For an idele $x$, let $r=\prod_p p^{v_p(x_p)}>0$, a finite
product, and take $q=\operatorname{sgn}(x_\infty)r$.
Then $q^{-1}x\in Z$.  Two such rational multipliers differ
by a rational number with all finite valuations zero and
positive sign, hence by one.  Thus the translates of $Z$
are disjoint and cover the ideles.  This subset is open
and closed and maps homeomorphically onto $C_{\mathbb Q}$.
Its intersection with the inverse image of a compact subset
of the quotient is compact.  Therefore $b$ has the stated
indentation properties and $\sum_qb(q^{-1}x)=1$.
For $x\in Z$ the first factor in $b^\Delta(x,qx)$ forces
$\xi=1$, and the second then forces $q=1$, proving (SM.20.1).

In particular, for the initial subgroup
$G_0=\mathbb R^\times\times\widehat{\mathbb Z}^{\,\times}$,
all outside-shell diagonal pieces from this lift are zero:
the only surviving $q$ is one and $x\in Z\subset G_0$.
The part labelled $S$ still contains the original global
spectral diagonal.  Its identification with a separately
defined $S$-local Paley--Wiener operator has not followed.

For averaging, Fubini and the volume-one normalization give
the displayed adjoint and $J^*J=1$, so $E$ is the claimed
orthogonal projection.  The actual weight $w_{I,s}$ depends
only on the norm.  Thus $K$ acts unitarily.
The closed geometric theta-function subspace of
Definition~10.4.1.1 is $C_F$-invariant, as explicitly used
in Lemma~10.4.1.2, p.~456.  That subspace and its orthogonal
complement are therefore $K$-invariant; their projections
commute with every $K$-translation and hence with $E$.
Convolution by $\phi$ satisfies
$EA(\phi)=A(\phi)=A(\phi)E$, by integrating the compact
variable, and $M_\rho$ commutes with $E$.  This proves
(SM.20.2).

Relative to $H=JH_0\oplus(JH_0)^\perp$, such a $T$ is
$(J^*TJ)\oplus0$.  The traces and singular values of this
block operator prove (SM.20.3).  Fubini first for simple
tensors and then $L^2$ density proves (SM.20.4); Jensen's
inequality for probability Haar measure on $K\times K$
gives its norm bound.
\end{proof}

These repairs permit a genuine trace calculation on the norm
coordinate for the indicated invariant part.  They do not
identify the compressed operator with the $S$-local operator
in Theorem~12.4.1.1, pp.~514--519.  That needs a specified
intertwiner carrying the actual kernels, projections and cutoff
limits to those local operators.  The zero indentation in
(SM.19.1) supplies none, and the closed-subspace carving
statement supplies no trace-norm estimate for that identification.

\subsection{Two precise obstructions to naive replacements}

\paragraph{Finding F-SM-21 (the unchanged archimedean local kernel misses a nonzero prime term).}
There exists a real $h\in C_c^\infty((3/2,5/2))$ such that
$h(2)=1$ and
\[
 P(h):=\widehat h(1/2)+\widehat h(-1/2)=0,\qquad A_0(h)=0,
 \tag{SM.21.1}
\]
where $A_0$ is the real, even archimedean distribution in F-SM-13.
Under the adopted Dirichlet theorem, its full zero distribution is
\[
 -\sum_\rho m(\rho)\widehat h(\rho-1/2)
       =\frac{\log2}{\sqrt2}>0.
 \tag{SM.21.2}
\]
Consequently the following two manuscript identifications cannot
be joined by simply replacing the zero indentation with $1_Z$
and leaving the $S=\{\infty\}$ local operator unchanged:
the global spectral trace identification of
Theorem~10.4.2.8 and equation~(12.5.1.4), and the local
Paley--Wiener trace limit of Theorem~12.4.1.1.
Under these explicitly stated operator inputs, the former
geometric cutoff trace tends to $\log2/\sqrt2$, whereas the
latter trace tends to zero, including the combined removed
moment correction $P(h)$.

More precisely, if trace-preserving transports put those
trace-class cutoffs on a common Hilbert space, their difference
$E_\lambda$ must satisfy
\[
 \liminf_{\lambda\to\infty}\|E_\lambda\|_1
       \geq \frac{\log2}{\sqrt2}.
 \tag{SM.21.3}
\]
This rules out that unchanged-kernel identification; it does not
rule out an assembly with the required nonzero prime contribution.

\begin{proof}
On the indicated interval, write
\[
 P(h)=\int h(t)p(t)\,dt,\quad
 A_0(h)=\int h(t)a(t)\,dt,\quad
 p(t)=\frac{t+1}{t^{3/2}},\quad
 a(t)=\frac{\sqrt t}{t^2-1}.
\]
The ratio
$R(t)=a(t)/p(t)=t^2/((t-1)(t+1)^2)$ is strictly decreasing for
$t>1$: its logarithmic derivative is
$(-t^2+t-2)/(t(t^2-1))<0$.
Choose nonzero nonnegative bumps $v_1,v_2$ supported,
respectively, in $(1.55,1.7)$ and $(2.3,2.45)$.
The quotients $A_0(v_i)/P(v_i)$ are weighted averages of
$R$ over these disjoint ordered intervals, so the first
is strictly larger than the second.  The matrix with
columns $(P(v_i),A_0(v_i))$ is therefore invertible.
Choose a real bump $u$ supported in $(1.9,2.1)$ with
$u(2)=1$, solve that two-by-two system for real $c_1,c_2$,
and set $h=u-c_1v_1-c_2v_2$.  This proves (SM.21.1) and
$h(2)=1$.

The only nonzero term in the finite distribution of F-SM-13 is
$\Lambda(2)h(2)/\sqrt2=\log2/\sqrt2$; the reciprocal terms
are outside the support.  Its explicit formula reads
$A_0(h)+B_1(h)=P(h)-\sum_\rho m(\rho)\widehat h(\rho-1/2)$,
so (SM.21.2) follows.  All zero sums converge absolutely.
The adopted zero-location input makes the same statement
hold in every positive-width strip used by the manuscript.

The source's global operator identification is exactly the
left side of (SM.21.2).  For $S=\{\infty\}$, its stated
local trace limit is $A_0(h)$, since $h(1)=0$ removes the
cutoff logarithm.  The combined moment correction also
vanishes by (SM.21.1).  Thus their trace difference tends
to $\log2/\sqrt2$.  The elementary inequality
$|\operatorname{Tr}E_\lambda|\leq\|E_\lambda\|_1$
proves (SM.21.3).
\end{proof}

\paragraph{Finding F-SM-22 (multiplying by a compact diagonal Dirac kernel destroys compactness).}
For the actual compact-invariant operators of F-SM-20, write
$T_{\rho,\phi}=E_K\otimes B_\rho$ on
$L^2(K)\otimes H_0$, where $E_K$ projects onto constants
and $B_\rho=J^*T_{\rho,\phi}J$.
If its radial kernel is $K_0(t,s)$, the proposed
multiplicative replacement
\[
 K_0(t,s)\longmapsto \delta_K(uv^{-1})K_0(t,s)
 \tag{SM.22.1}
\]
defines $1_{L^2(K)}\otimes B_\rho$ on tensor test functions.
If $B_\rho\ne0$, this operator is not compact and hence
is not trace class.  Under the operator inputs of F-SM-21,
$B_\rho\ne0$ for all sufficiently advanced cutoffs in that
example.  Thus (SM.22.1) cannot serve as a trace-class repair
of the ordinary-function formula.

\begin{proof}
The normalized compact diagonal distribution acts by
$\int_K\delta_K(uv^{-1})f(v)\,dv=f(u)$, so integrating
the product kernel gives $1\otimes B_\rho$ on simple
tensors and then on the Hilbert completion.
The group $K=\widehat{\mathbb Z}^{\,\times}$ has infinitely
many distinct continuous characters, forming an infinite
orthonormal set $(\chi_j)$ in $L^2(K)$.
Choose a unit vector $v\in H_0$ with $B_\rho v\ne0$.
The unit vectors $\chi_j\otimes v$ have images
$\chi_j\otimes B_\rho v$, which are mutually orthogonal
and have the same positive norm.  Their images have no
convergent subsequence.  The operator is not compact.
For the example in F-SM-21, its asserted global trace
limit is positive, so the trace-class $B_\rho$ has
nonzero trace and is nonzero for all sufficiently
advanced cutoffs.  This proves the final assertion.
\end{proof}

This is an obstruction to the particular \emph{multiplicative}
Dirac interpretation.  A subsequent compact-isotypic
projection changes the operator again: compression by $E_K$
restores $E_K\otimes B_\rho$.  Such a repair is possible
as an operation, but it returns to the global operator
and still does not supply the missing local identification.

\paragraph{Primary source and verification scope.}
Masullo, A. (2026). \emph{The music of the primes}
[supplied manuscript].
\url{https://math216.wordpress.com/wp-content/uploads/2026/11/the-music-of-the-primes.pdf}.
Additional exact locators: pp.~21--22; Lemma~12.2.1.1,
pp.~504--505; Proposition~12.2.2.2, pp.~508--510;
the missing Section~12.3, p.~514;
Theorem~11.3.2.1 and equation~(11.3.2.2), p.~495;
Definition~A.2.3.6, p.~549; Definitions~A.3.2.1--2,
pp.~559--560.  The p.~495 formula was inspected visually.
F-SM-18--22 prove the statements written here, without
certifying the earlier full trace-class construction or the
absent local-operator identification.

% END source_masullo_kernel.tex

% BEGIN source_K_bridge.tex
\section{Integral $K$-groups: concrete known consequences and the missing export}
\label{sec:source-K-bridge}

The final footnote in Alp\"oge's four-page remark suggests consequences for
$K$-groups of $\mathbb Z$, but the supplied declarations do not identify a
degree, a torsion prime, an eigenspace, or a comparison map.  There are
nevertheless precise existing integral results with which such an export
can be compared.  We keep ordinary integral $K$-groups, their reductions
modulo $p$, and their $p$-primary subgroups distinct.

\subsection{The exact cyclotomic comparison}

\paragraph{Finding F-SM-14 (the mod-$p$ comparison and its vanishing consequence).}
Let $p$ be an odd prime, let $n\geq1$, put
$F_p=\mathbb Q(\zeta_p)$, and let $C_p$ be the $p$-primary subgroup of
$\operatorname{Cl}(F_p)$.  Let $\omega$ be the Teichm\"uller character of
$\Delta_p=\operatorname{Gal}(F_p/\mathbb Q)$.  There is an isomorphism of
$\mathbb F_p$-vector spaces
\begin{equation}\label{eq:SM-Kclass}
 K_{2n}(\mathbb Z)/p
 \simeq H^2_{\text{ét}}
       (\mathbb Z[1/p],\mu_p^{\otimes(n+1)})
 \simeq (C_p/pC_p)(\omega^{-n}).
\end{equation}
Here the last expression denotes a character eigenspace.  In particular,
\[
 K_{2n}(\mathbb Z)\{p\}=0
 \quad\Longleftrightarrow\quad C_p(\omega^{-n})=0.
\]
The comparison with \emph{mod-$p$} \'etale cohomology in the first
isomorphism is a known theorem, not a consequence of RH: see Weibel
(2013), Chapter~VI, Corollary~8.3, chapter page~48, following
Theorem~8.2, chapter page~47.  The following proof spells out its
class-group specialization and the vanishing equivalence.

\begin{proof}
Let $R=\mathbb Z[\zeta_p,1/p]$.  The map
$\operatorname{Spec}R\to\operatorname{Spec}\mathbb Z[1/p]$ is finite
\'etale Galois with group $\Delta_p$, of order $p-1$.  Taking invariants
on $\mathbb F_p[\Delta_p]$-modules is exact, by averaging; the
Hochschild--Serre spectral sequence therefore identifies cohomology over
$\mathbb Z[1/p]$ with the invariants of cohomology over $R$.

The Kummer sequence gives
\[
 0\longrightarrow\operatorname{Pic}(R)/p
 \longrightarrow H^2_{\text{ét}}(R,\mu_p)
 \longrightarrow\operatorname{Br}(R)[p]\longrightarrow0.
\]
There is exactly one prime of $F_p$ above $p$, generated by
$1-\zeta_p$.  Thus localizing at $p$ leaves the ideal class group
unchanged.  The field $F_p$ is totally imaginary, and the local-invariant
sequence for the Brauer group of its ring of $S$-integers has only this
one finite place; the sum-of-invariants map is then injective.  Hence
$\operatorname{Br}(R)[p]=0$.  The Kummer identification is
$\Delta_p$-equivariant and reads
\[
 H^2_{\text{ét}}(R,\mu_p)
       \simeq C_p/pC_p.
\]
Because $R$ contains $\mu_p$, the higher coefficient twist is constant
over $R$.  With its descent action retained, it gives
\[
 H^2_{\text{ét}}(R,\mu_p^{\otimes(n+1)})
 \simeq (C_p/pC_p)\otimes_{\mathbb F_p}\mu_p^{\otimes n}.
\]
The last tensor factor has character $\omega^n$ modulo $p$.
Its invariant subspace is exactly the $\omega^{-n}$ eigenspace of
$C_p/pC_p$, proving~\eqref{eq:SM-Kclass} from the cited comparison.

The finite $p$-group $C_p$ splits into its Teichm\"uller eigenspaces,
since $p-1$ is invertible in $\mathbb Z_p$; reduction modulo $p$ commutes
with the corresponding idempotents.  A finite $p$-group $A$ has
$A/pA=0$ if and only if $A=0$: if $A=pA$, iteration gives
$A=p^rA=0$ for $p^r$ annihilating $A$.  Apply this first to each class
component, and then to the $p$-primary component of the finite group
$K_{2n}(\mathbb Z)$.  The latter finiteness is the usual integral
$K$-theory finiteness theorem, also used in the cited comparison.
\end{proof}

The assertion is a comparison of reductions modulo $p$ and hence a
vanishing criterion.  It does not identify the entire finite
$p$-primary groups with one another.  In positive-rank odd degrees the
correct comparison uses $p$-adic coefficients/completion; an
uncompleted finitely generated $K$-group cannot surject onto a copy of
$\mathbb Z_p$.  This distinction is material when reading informal
Chern-map summaries.

\subsection{A complete known vanishing and its class-group consequence}

Dutour Sikiri\'c, Elbaz-Vincent, Kupers, and Martinet (2026),
Theorem~6.4, prove
\[
 K_8(\mathbb Z)=0.
\]
The theorem is on page~19 of the author-posted updated manuscript;
Corollary~7.1 on page~20 concludes
$C_p(\omega^{p-5})=0$ for every prime $p>3$.  The journal article is in
\emph{Journal of the Institute of Mathematics of Jussieu},
\emph{25}(1), 565--590, DOI~10.1017/S1474748025101394;
its online publication date is 24 November 2025.  This is a known result
independent of the adopted Alp\"oge input.

The final deduction has a short, exact audit trail.  Their
Proposition~6.1 puts $H_9(BQ;\mathbb Z)$ in the class of finite groups
whose prime divisors are at most seven, and Proposition~6.3 puts the
kernel of the relevant Hurewicz map in the analogous class for five.
Since $K_8(\mathbb Z)=\pi_9(BQ)$, its possible prime divisors are at
most seven.  The odd primes $3,5,7$ are regular (their relevant Bernoulli
numerators are those of $B_2=1/6$ and $B_4=-1/30$), so Weibel VI,
Proposition~10.5, chapter page~63, removes their torsion.  The
$2$-primary part is zero by Weibel VI, Corollary~9.8, chapter page~55.
This last input makes explicit the prime-two step that is not written in
the updated manuscript's very short proof of Theorem~6.4.  The Voronoi
cell enumerations and the preceding homology computations are carried
as published inputs; they have not been rerun here.

Applying F-SM-14 with $n=4$ gives
$(C_p/pC_p)(\omega^{-4})=0$, hence
$C_p(\omega^{-4})=0$.  As $-4\equiv p-5\pmod{p-1}$, this is the stated
cyclotomic class-group consequence.  The argument works at the level of
finite quotients without an incorrect uncompleted Chern map.  It also
does not need the assertion, as printed in the manuscript's Section~7,
that $H^1_{\text{ét}}(\mathbb Z[1/p],\mathbb Z_p(j))$
always vanishes when $p>j+1$.  For the positive twists relevant here,
that assertion requires even $j$: for odd $j>1$ this group has a
rank-one summand, as
Weibel VI, Exercise~8.3, chapter page~50, specifies.

\subsection{A density-one family with all degrees specified}

Chen et al. (2026), arXiv:2602.05090v1, Theorem~1.1 on page~2, prove
the following quantitative partial regularity statement.  For each
$\alpha>1/2$, let
\[
 M_\alpha(p)=\left\lfloor\frac{\sqrt p}{(\log p)^\alpha}\right\rfloor.
\]
The number of primes $p\leq X$ for which some even index
$2\leq2k\leq\min(M_\alpha(p),p-3)$ has
$p\mid\operatorname{num}(B_{2k})$ is
$O_\alpha(X/(\log X)^{2\alpha})$.  Their Corollary~1.4, page~3,
gives the corresponding vanishing of the even class-group components;
Corollary~1.7(4), page~5, gives the $K_{4k-2}$ consequence.
The formal verification reported by that paper concerns its main
elementary estimate; no Lean build of that archive has been performed
in this run, and the deep comparison inputs remain cited mathematics.

\paragraph{Finding F-SM-15 (both families of integral torsion consequences).}
For a fixed $\alpha>1/2$, outside
$O_\alpha(X/(\log X)^{2\alpha})$ odd primes $p\leq X$, the following
hold simultaneously for every integer $k$ with
$2\leq2k\leq\min(M_\alpha(p),p-3)$ and every integer $b\geq0$:
\begin{align*}
 K_{4k-2+2b(p-1)}(\mathbb Z)\{p\}&=0,\\
 K_{2p-2-4k+2b(p-1)}(\mathbb Z)\{p\}&=0.
\end{align*}
In particular, each fixed $\alpha$ gives a density-one set of primes
with these simultaneous vanishings.  The first line for $b=0$ is the
published consequence just cited; the second line makes explicit the
$4r$-degree consequence of its partial Vandiver components.

\begin{proof}
For a nonexceptional prime the partial regularity theorem gives
$p\nmid\operatorname{num}(B_{2k})$ throughout the range.  The
Herbrand--Ribet theorem identifies this with
$C_p(\omega^{1-2k})=0$; its exponent is also written $p-2k$ modulo
$p-1$.  Leopoldt reflection implies
$C_p(\omega^{2k})=0$.  These are exactly the two cited cyclotomic
inputs in Chen et al., equation~(4) and the deduction of
Corollary~1.4, page~3.

For the first line of the conclusion put
$n=2k-1+b(p-1)\geq1$ in F-SM-14.  Then
$-n\equiv1-2k\pmod{p-1}$.  For the second put
$n=p-1-2k+b(p-1)$.  The bound $2k\leq p-3$ ensures $n\geq2$, and
$-n\equiv2k\pmod{p-1}$.  Thus F-SM-14 proves both assertions, for
every $b$ at once.  They introduce no additional exceptional primes.
Finally, the prime number theorem gives
\[
 \frac{O_\alpha(X/(\log X)^{2\alpha})}{\pi(X)}
       =O_\alpha((\log X)^{1-2\alpha})\longrightarrow0.
\]
This proves the density assertion.
\end{proof}

The degrees in the second line are divisible by four because $p$ is
odd.  The degrees in the first line are congruent to two modulo four.
Both are precise $p$-primary assertions; neither line is an integral
vanishing assertion for an entire $K$-group at every prime.  Also, these
are pre-existing partial-regularity and comparison inputs, independent
of complex RH; their proof does not manufacture a missing exported
theorem from Alp\"oge's final footnote.

\subsection{The degree-one boundary and what an export must state}

\paragraph{Finding F-SM-16 (the elementary degree-one calculation).}
$K_1(\mathbb Z)\simeq\{\pm1\}\simeq\mathbb Z/2$.

\begin{proof}
Use the classical definition
$K_1(\mathbb Z)=\operatorname{GL}_\infty(\mathbb Z)^{\mathrm{ab}}$.
The determinant is a surjection onto $\{\pm1\}$.  Its kernel is
$\operatorname{SL}_\infty(\mathbb Z)$, which is generated by elementary
matrices: the Euclidean algorithm reduces the first column of a
unimodular matrix to a unit vector using integer row additions and signed
row interchanges; column additions clear the rest of its first row, and
induction applies to the remaining unimodular block.  Sign corrections
are allowed because
$\left(\begin{smallmatrix}0&1\\-1&0\end{smallmatrix}\right)
=e_{12}(1)e_{21}(-1)e_{12}(1)$ and its square is
$\operatorname{diag}(-1,-1)$.  In determinant one the final block also
has determinant one.  Stabilization supplies an extra index when needed.

For distinct indices $i,j,r$,
$e_{ij}(a)=[e_{ir}(a),e_{rj}(1)]$, so every elementary matrix is a
commutator in the stable group.  Conversely every commutator has
determinant one.  Thus the determinant kernel is precisely the
commutator subgroup, proving the asserted abelianization.
\end{proof}
The exact standard locator is Weibel III, Example~1.3.5, chapter
page~4, following Whitehead's Lemma~1.3.3.  This elementary nonzero
group already prevents a reading of the tentative footnote as blanket
vanishing in all positive degrees.

To make a further integral consequence of the supplied analytic method
reviewable, an export must specify its degree set, its coefficient or
torsion convention, its prime range, and the exact arithmetic map it
controls.  For example, an exported vanishing of the component
$C_p(\omega^{-n})$ would give the corresponding $p$-primary vanishing
in $K_{2n}(\mathbb Z)$ by F-SM-14.  To conclude an entire finite
$K_{2n}(\mathbb Z)$ is zero, one must cover every prime, including the
separate prime-two calculation.  Zero location for complex Dirichlet
$L$-functions alone supplies no such integral eigenspace or
divisibility statement in the supplied five declarations.

\subsection{Primary references and checked locators}

Chen, E., Cummins, C., Eltschig, B., Grubisic, D., Haller, L.,
Hong, L., Kurghinyan, A., Lau, K., Leather, H., Lee, S.,
Markosyan, A., Ono, K., Patel, M., Pendharkar, G., Rathi, V.,
Schneidman, A., Seeker, V., Sengupta, S., Sinha, I., \ldots\
Zhang, J. (2026).
\emph{Almost all primes are partially regular}
(arXiv:2602.05090v1).
\url{https://arxiv.org/abs/2602.05090}.
The used locators are Theorem~1.1, p.~2; equation~(4),
the reflection implication and Corollary~1.4, p.~3; and
Corollary~1.7(4), p.~5.  The natural-language proof of the main
quantitative theorem in Sections~2--3, pp.~6--10, was read throughout;
the formalization scope is stated in Section~4, pp.~10--11.

Dutour Sikiri\'c, M., Elbaz-Vincent, P., Kupers, A., \& Martinet, J.
(2026). Voronoi complexes in higher dimensions, cohomology of
$\operatorname{GL}_N(\mathbb Z)$ for $N\geq8$ and the triviality of
$K_8(\mathbb Z)$.
\emph{Journal of the Institute of Mathematics of Jussieu, 25}(1),
565--590.
\url{https://doi.org/10.1017/S1474748025101394}.
Updated author manuscript:
\url{https://jamartin.perso.math.cnrs.fr/Publications/DEKM.pdf}.
The used locators in this manuscript are Proposition~6.1, p.~18;
Proposition~6.3 and Theorem~6.4, p.~19; and Corollary~7.1, p.~20.
The complete final deductions were read; upstream Voronoi
computations were not reproduced.

Weibel, C. A. (2013).
\emph{The K-book: An introduction to algebraic K-theory}
(Graduate Studies in Mathematics, Vol.~145).
American Mathematical Society.
\url{https://doi.org/10.1090/gsm/145}.
Author's chapter files:
\url{https://sites.math.rutgers.edu/~weibel/Kbook/Kbook.III.pdf}
and
\url{https://sites.math.rutgers.edu/~weibel/Kbook/Kbook.VI.pdf}.
The exact locators are III, Example~1.3.5, chapter p.~4;
VI, equation~(8.1.1) and Theorem~8.2, chapter p.~47;
VI, Corollary~8.3, chapter p.~48; VI, Exercise~8.3, chapter p.~50;
VI, Corollary~9.8, chapter p.~55; VI, Proposition~10.5,
chapter p.~63; and VI, Theorem~10.9, chapter p.~65.
Chapter pagination differs from the printed volume; theorem
numbering is unchanged.  These primary statements and their
adjacent arguments were read directly.  The underlying motivic
comparison theorems are carried as established inputs.

% END source_K_bridge.tex
\clearpage
\part{Unresolved questions and unchecked source assertions}
\section{Current gates and approaches still in progress}
No mathematical assertion in this part is promoted to a theorem by
its inclusion. The executable question register is \nolinkurl{QUESTIONS.tsv};
the dated approach record is \nolinkurl{LEDGER.md}. A missing historical
proof package and an open mathematical conjecture are different statuses.

\begin{description}[style=nextline,leftmargin=0pt,labelwidth=0pt,labelsep=0pt]
\item[Cote 14, scans 25--38: original quotient identification.]
The generated quotient has full proofs in AQ-14-FORWARD and
AQ-14-CONVERSE. Original Definition1.1, Proposition1.2 and
Corollary2.3 now identify that quotient under the explicitly bracketed
projective-parameter convention. AQ-14-QUASIPROJECTIVE now removes the projective qualification
in characteristic zero for smooth quasiprojective parameters, by an
integral equality of the generated relations. AQ-14-NONPROPER now
proves the cohomological implication in every characteristic.
AQ-14-ALTERATION proves that the relation quotient in characteristic
$p>0$ is $p$-primary torsion, hence proves equality after inverting $p$
and on the prime-to-$p$ torsion in the resulting cycle quotients.
The possible residual $p$-primary relation difference is the remaining
integral question. Supplied scan38 is physical PDFpage38,
handwritten page37; both locators are retained.
\item[Cote 15, scans 100--106: full simple-motive classification.]
Weil-number occurrence is proved. The cited 2026 CM Hodge theorem,
with the inspected Milne implications, discharges the rational Tate and
standard-conjecture inputs for abelian varieties and their products.
AQ-ABELIAN-RECOGNITION proves full faithfulness after scalar extension
and Frobenius-polynomial recognition in the abelian--Artin numerical
category. CC-AG01--03 further classify its idempotent completion over
$\overline{\mathbf Q}_\ell$ by finite Weil-number gradings and
compute its diagonalizable group, including its Artin torsion quotient. CC-XC01--04 now reduce membership of a specified smooth projective
variety to a determinant condition on one projected finite family
of cross-cycle maps. This condition forces semisimplicity and yields
fixed-prime Tate and numerical--homological comparison on the stated
products. Extending the classification to arbitrary varieties still
requires producing those actual cross-correspondences. The
quaternion example AQ-RATIONAL-FIBRE also shows why a rational neutral
fibre functor and arbitrary rational multiplicities cannot be silently
inserted. AQ-WEIL-NOTATION recovers the historical totally real and CM
closures on the earlier original sheets and proves the norm-condition
identification; that historical notation is no longer a gate.
The hypotheses of Milne (1994), \S2, have therefore been
discharged in a specified sector, not deleted globally.
\item[Cote 18, scans 11--13 and 59: period density.]
The overly strong formulation is refuted, and generation for one
elliptic curve has a precise known answer. Non-CM transcendence degree
four and the non-CM algebraicity question for the Hodge-splitting
scalar remain the specified transcendence gates. Group generation
does not imply density of a particular period locus; AQ-COCYCLE
gives the exact distinction.
\item[Cote 134-8, scans 47--54: the linear theorem.]
The original predicates and adjunction directions have been read.
BR-13401--06 and F-SM-17 prove the precise modern comparison and
counit criteria, the equivalence over fields and \(\mathbf Z\) under
the stated hypotheses, and a locally projective ringoid construction.
F-SM-24 extends this to all countable rings and all scalar rings of
cardinality $\aleph_t$ for finite $t$, under the explicit scalar
compatibility, and to arbitrary scalar rings when the ringoid morphism
modules are finitely generated projective. F-SM-27 gives a third range:
arbitrary scalar rings and projective morphism modules if the left
module $D$ has a resolution by finitely generated projectives.
Outside these ranges the remaining test is
\(L\varphi(GP)=0\) on acyclic projective chain objects, equivalently
the explicit Tor condition SM25.2. F-SM-26 shows why finite flat
dimension of $D$ cannot furnish another case: it contradicts the
nonzero full-interval normalization. The intake's
three-sort counterexample remains an uninspected earlier construction;
no claim about its exact compatibility with the original hypotheses
has been promoted to a new proof here.
\item[SA08, cote 136/117--121: strict PD sources and localization.]
AQ-SA08-LIFT refutes a specific degreewise-free lifting input.
The original admissibility conditions and common-source and strict
equalization conjectures have now been retrieved on physical pages
111, 113, 115, 117, 119 and 121, footers one smaller. Their exact
transcription and a proof or counterexample under those conditions
are in progress. The stronger intake all-arrow terminality and SA09
comparison still require their distinct category definitions and proof
context. The lifting counterexample is not a refutation of those
different existence questions.
\item[Cote 123: apposition outside type A.]
The quotient-dimension, rational-singularity and Lie-rank questions
are answered. CC-RK01--05 supplies the remaining classical and
exceptional ranks, subregular uniqueness and the analytic transfer
to unipotent group elements. BR-12305--08 treats all central forms
of the group quotient. CC-AP01--07 completely classifies type A
apposition quadruples, their generators, opposition and outer action.
Explicit classifications in the other simple types remain in progress;
type A already refutes a universal distinguished outer-invariant class.
\item[Weight--monodromy and universal-curve categories.]
The local monodromy operator and its invariant exterior powers must
be retained. BR-WM-02 supplies a finite gate and has passed a
recorded independent whole-proof reading. CC-ST01--04 give the geometric
stack decomposition with its explicit separating-divisor summands.
CC-RES01--02 prove the canonical exit/right-Kan statement and the nodal
lifting-category calculation. Identifying an unseen chosen archive
functor with that canonical functor, or deducing all categorical base
changes from geometric proper base change, requires another argument.
Purity of a proper pushforward
complex is not purity of every ordinary singular-fibre cohomology group.
\item[The analytic method as a transferable theorem.]
The lattice ideals, repeated theta-lift estimates, period nonvanishing,
and tower bounds must be exported with exact hypotheses from the
full argument. BR-DET-01/02 identify the numerical determinant gate.
The footnote on $K$-groups contains no degree-indexed theorem to apply.
F-SM-14--16 instead give exact comparisons and explicit families of
torsion consequences of separately cited integral arithmetic theorems.
\item[Masullo's operator realization.]
The fully normalized numerical shrinking-strip identity is proved in
the Dirichlet case in F-SM-13, including the finite Fourier expansion
of the tests and the archimedean parity term. F-SM-18 proves diagonal summation with a valid normalized indentation,
and F-SM-20 gives the actual compact-invariant compression. F-SM-19
refutes the printed indentation: it vanishes for $\mathbf Q$.
F-SM-21--22 rule out two specific naive replacements, by a nonzero
prime-trace discrepancy and by noncompactness. The remaining
global-to-$S$-local kernel assembly in \S12.3, p.~514, must supply
a valid intertwiner and the required prime corrections; the
diagonal domination step alone no longer describes the gate. All-number-field Hecke scope and
zero-multiplicity assertions require additional proofs.
\item[Profinite GT, anabelian reconstruction and geometric spectra.]
The full fixed Ihara embedding's surjectivity, arithmetic associator
data, and a spectrum constructed geometrically from primes are distinct
targets. The current rational GT theorem and the adopted Dirichlet
theorem do not identify those objects. Explicit counterexamples here
refute only the stated routes through abelian cohomology or bare spectra.
\end{description}

\section{Checks and current status}
Every completed new proof is associated with its actual reading events
in \nolinkurl{CHECK_EVENTS.tsv}. No unproved candidate in this part is
used as an input to a completed proof. The timed checkpoints and the
three-hour, then hourly movement reviews are recorded in
\nolinkurl{LEDGER.md}; the two-no-movement stopping rule remains in force.
\clearpage
\section{References for the root geometric and analytic sections}
References in the consequences and GT sections additionally give their
own exact source records. All archive locators refer to the numbered
scan, excluding the unnumbered cover, at the Montpellier catalogue:
\url{https://grothendieck.umontpellier.fr/archives-grothendieck/}.
\begin{enumerate}[leftmargin=0pt,labelsep=0pt,label={}]
\item Beilinson, A. A., Bernstein, J., \& Deligne, P. (1982).
\emph{Faisceaux pervers}. Ast\'erisque, 100. Operative locators:
Stabilit\'es 5.1.14, pp.~128--129; Corollary 5.4.4,
Theorem 5.4.5 and Corollary 5.4.6, p.~142.
\url{https://publications.ias.edu/sites/default/files/Faisceaux%20pervers.pdf}
\item Berger, L. (2004). \emph{An introduction to the theory of
$p$-adic representations} [Author preprint, arXiv:math/0210184].
Operative locators: \S II.3.1 and II.4.1--II.4.3.
\url{https://arxiv.org/abs/math/0210184}
\item Deligne, P. (1980). La conjecture de Weil II.
\emph{Publications Math\'ematiques de l'IH\'ES, 52}, 137--252.
Operative locators: Theorem 1.8.4; Problem 1.8.15,
pp.~178--179; 1.9.5, p.~181.
\url{https://publications.ias.edu/sites/default/files/Number40.pdf}
\item Deligne, P., \& Goncharov, A. B. (2005). Groupes fondamentaux
motiviques de Tate mixte. \emph{Annales Scientifiques de l'\'Ecole
Normale Sup\'erieure, 38}(1), 1--56.
Operative locators in supplied source: \S1.1, equations
(1.1.2), (1.1.4), Proposition 1.7 and its proof.
\url{https://arxiv.org/abs/math/0302267}
\item Gan, W. T., Gross, B. H., \& Prasad, D. (2009).
\emph{Symplectic local root numbers, central critical $L$-values,
and restriction problems} [Preprint, arXiv:0909.2999].
Operative locators: Theorems 14.1, 15.1, 16.1;
Conjectures 17.1, 17.3, 24.1. These labels refer to the supplied version.
\url{https://arxiv.org/abs/0909.2999}
\item Gan, W. T., Qiu, Y., \& Takeda, S. (2014).
\emph{The regularized Siegel--Weil formula (the second term identity)
and the Rallis inner product formula} [Preprint, arXiv:1207.4709].
Operative locator: Theorem 11.4.
\url{https://arxiv.org/abs/1207.4709}
\item Grothendieck, A. (1969). Standard conjectures on algebraic
cycles. In \emph{Algebraic geometry, Bombay 1968} (pp.~193--199).
Oxford University Press. Operative locators: \S2, p.~195;
\S3, pp.~195--197; \S4 and Remarks 2--4, pp.~197--198.
\url{https://download.uni-mainz.de/mathematik/Algebraische%20Geometrie/Lehre/WS23.Padische.1968.Grothendieck.StandardConjs.pdf}
\item Kashiwara, M. (1986). A study of variation of mixed Hodge
structure. \emph{Publications of the Research Institute for Mathematical
Sciences, 22}(5), 991--1024.
\url{https://doi.org/10.2977/PRIMS/1195177264}.
Operative locators: \S1.8--1.9; Theorems 4.4.1 and 4.5.2;
Propositions 4.5.1 and 5.2.1.
\item Katz, N. M., \& Messing, W. (1974). Some consequences of the
Riemann hypothesis for varieties over finite fields.
\emph{Inventiones Mathematicae, 23}, 73--77.
Operative locators: Theorem 1; Theorem 2(1), pp.~76--77.
\url{https://web.math.princeton.edu/~nmk/old/katzmessing.pdf}
\item Milne, J. S. (2006). \emph{Arithmetic duality theorems} (2nd ed.).
BookSurge. Operative locator: Chapter II, Theorem 3.1(b), p.~177;
Corollaries 3.2--3.3 specify the ordinary and compact-support variants.
\url{https://jmilne.org/math/Books/ADTnot.pdf}
\item Pearlstein, G. (2000). \emph{Degenerations of mixed Hodge structure}
(arXiv:math/0002030v3, 6 October). Operative locator: Theorem 5.1(i)--(iii),
p.~30 of this preprint. \url{https://arxiv.org/abs/math/0002030v3}
\item Scholze, P. (2012). \emph{Perfectoid spaces}
[Preprint, arXiv:1111.4914]. Operative locators in the inspected
preprint: Conjecture 9.3 and Theorem 9.4, p.~48;
Lemma 9.5 and Theorem 9.6, p.~49.
\url{https://arxiv.org/abs/1111.4914}
\item Titchmarsh, E. C. (1986). \emph{The theory of the Riemann
zeta-function} (2nd ed., revised by D. R. Heath-Brown).
Oxford University Press. Operative locators: \S2.12,
pp.~29--30; Theorems 9.3--9.4 and equation (9.4.3), p.~214.
\item Voevodsky, V. (2000). Triangulated categories of motives
over a field. In V. Voevodsky, A. Suslin, \& E. M. Friedlander,
\emph{Cycles, transfers, and motivic homology theories}
(Annals of Mathematics Studies, Vol.~143, pp.~188--238).
Princeton University Press. Operative locators in the author PDF:
\S2.1 after Proposition 2.1.3; Theorem 3.1.12; Corollary 3.4.3.
\url{https://www.math.ias.edu/vladimir/sites/math.ias.edu.vladimir/files/s5_0.pdf}
\item Deligne, P., \& Katz, N. (Eds.). (1973).
\emph{Groupes de monodromie en g\'eom\'etrie alg\'ebrique,
SGA 7 II} (Lecture Notes in Mathematics, Vol.~340). Springer.
Operative locator: Expos\'e XV, Th\'eor\`eme 3.4(ii)--(iii),
pp.~195--196 (specialization, vanishing cycles and Picard--Lefschetz).
\end{enumerate}

\section{Primary bibliography and source audit for the AQ
module}\label{primary-bibliography-and-source-audit-for-the-aq-module}

All mathematical attributions below distinguish the cited theorem from
the deductions proved in \nolinkurl{archive_questions.tex}. Page numbers
for Milne, Achter--Casalaina-Martin--Vial, KSV and Caradot--Lin refer to
the specified author/arXiv version, not a different published
pagination. The original archive quotations come from the supplied
transcriptions or the explicitly identified official scans recovered in
this run; the intake reports are secondary records of earlier work and
do not substitute for their absent proof files.

\subsection{References in APA style}\label{references-in-apa-style}

Achter, J. D., Casalaina-Martin, S., \& Vial, C. (2021). \emph{On the
image of the second ℓ-adic Bloch map} (arXiv:2007.07180v2).
\url{https://arxiv.org/pdf/2007.07180v2}

\begin{itemize}
\tightlist
\item
  Relevant primary statements: Appendix A, Propositions A.23--A.26
  (functoriality and specialization), A.28--A.29 (codimension one and
  zero-cycles/Roitman). The appendix explicitly treats arbitrary
  codimension where used. The entire paper is not claimed independently
  reproved.
\end{itemize}

Bloch, S. (1979). Torsion algebraic cycles and a theorem of Roitman.
\emph{Compositio Mathematica, 39}(1), 107--127.
\url{https://www.numdam.org/article/CM_1979__39_1_107_0.pdf}

\begin{itemize}
\tightlist
\item
  Used: Proposition 3.5, p. 115, correspondence compatibility;
  Proposition 3.8, p. 117, specialization; Theorem 4.2, p. 119,
  prime-to-characteristic Roitman. The modern Appendix A above supplies
  a convenient explicit formulation of the same operations.
\end{itemize}

Boavida de Brito, P., Horel, G., \& Robertson, M. (2019). Operads of
genus zero curves and the Grothendieck--Teichmüller group.
\emph{Geometry \& Topology, 23}, 299--346.
\url{https://doi.org/10.2140/gt.2019.23.299}

\begin{itemize}
\tightlist
\item
  Primary PDF: \url{https://msp.org/gt/2019/23-1/gt-v23-n1-p07-s.pdf}
\item
  Used: Definition 7.2(II), p. 324, the profinite
  Grothendieck--Teichmüller hexagon. This does not define the real
  hexagonal elliptic lattice of the archive.
\end{itemize}

Broe, M. (2026). \emph{The Tate conjecture for abelian fourfolds over
finite fields} (arXiv:2608.28651v1, August 19).
\url{https://arxiv.org/pdf/2608.28651v1}

\begin{itemize}
\tightlist
\item
  Theorem 1.1, p. 1, gives the codimension-two cycle-map surjection for
  every abelian fourfold over a finite field. Corollary 1.2, p. 2, and
  Corollary 8.2, p. 20, give standard conjecture D in dimension four
  over arbitrary fields. Corollary 8.1(i), p. 20, gives rational equals
  numerical equivalence with rational coefficients over a finite field.
  Exact statements and their scope were inspected; the complete preprint
  proof was not independently checked here. The later all-abelian result
  below is stronger for the present realization application.
\end{itemize}

Caradot, A., \& Lin, Z. (2026). \emph{Differential graded algebras with
divided powers and homotopy Lie algebras} (arXiv:2511.22614v4, June 20).
\url{https://arxiv.org/pdf/2511.22614v4}

\begin{itemize}
\tightlist
\item
  Checked directly: §2.1, pp. 5--6, the termwise-free convention;
  Proposition 3.10 and its proof, pp. 19--20; Theorem 3.15, pp. 23--24.
  AQ-SA08-LIFT refutes the degreewise-surjective lifting statement in
  this exact version. Its map is not a quasi-isomorphism, so it does not
  refute Theorem 3.15. No assertion is made about changes in a later
  journal version.
\end{itemize}

Chudnovsky, G. V. (1980). Algebraic independence of values of
exponential and elliptic functions. In \emph{Proceedings of the
International Congress of Mathematicians (Helsinki, 1978)} (Vol. 1, pp.
339--350). Academia Scientiarum Fennica.

\begin{itemize}
\tightlist
\item
  Official proceedings:
  \url{https://www.mathunion.org/fileadmin/ICM/Proceedings/ICM1978.1/ICM1978.1.ocr.pdf}
\item
  Exact primary locator: §2A, Theorem 1 and Corollary 2, printed p. 343.
  The algebraic-invariants setting is fixed on p. 339. Theorem 1 gives a
  transcendence-degree lower bound for periods and quasi-periods;
  Corollary 2 gives algebraic independence of a nonzero CM elliptic
  period and π. The meeting year is 1978; publication is 1980. The
  relevant twelve-page extract was read locally.
\end{itemize}

Duren, P. L. (1983). \emph{Univalent functions} (Grundlehren der
mathematischen Wissenschaften, Vol. 259). Springer.
\url{https://link.springer.com/book/9780387907956}

\begin{itemize}
\tightlist
\item
  Standard locator: Theorem 2.3, p. 31, Koebe quarter theorem. The book
  locator was corroborated indirectly; the AQ module supplies the entire
  needed area-theorem/odd-function/Möbius proof, so the bound does not
  depend on an unexamined book proof.
\end{itemize}

Grothendieck, A. (n.d.). \emph{Théorie des cycles algébriques:
Conjectures de Hodge, Tate, Lefschetz, Weil} {[}Unpublished manuscripts,
cote 14{]}. Archives Grothendieck, Université de Montpellier.
\url{https://grothendieck.umontpellier.fr/14.pdf}

\begin{itemize}
\tightlist
\item
  Supplied source: scan 38, manuscript p. 37. It states the rational
  ℓ-adic conclusion, labels Théorème 3.1 ``dém. pas faite,'' and asks
  for a converse. The original quotient is now identified:
  Definition1.1/Proposition1.2(i), physical PDF page24/footer23;
  Proposition1.2(v), physical25/footer24; Theorem2.1,
  physical31/footer30; Corollary2.3, physical34/footer33. The word
  projectif is bracketed in Definition1.1; the proofs cover its
  projective-parameter reading. Physical page38 corresponds to
  footer/manuscript37, matching the supplied scan038 convention.
\item
  Definition 1.1 says only a field k. No explicit algebraically closed
  convention was found in the preceding relevant handwritten pages
  inspected. Our theorem retains algebraic closure as an explicit
  hypothesis; the finite-field rational-point counterexample shows why
  it matters.
\end{itemize}

Grothendieck, A. (n.d.). \emph{Théorie arithmétique. Théorie de Galois
des motifs} {[}Unpublished manuscripts, cote 15{]}. Archives
Grothendieck, Université de Montpellier.
\url{https://grothendieck.umontpellier.fr/15.pdf}

\begin{itemize}
\tightlist
\item
  Supplied sources: scans 100, 102, 104 and 106, manuscript pp. 99, 101,
  103 and 105. Official physical pages 90--106 were additionally
  retrieved where relevant, and the defining sheets visually inspected:
  physical94/footer93 identifies \texttt{Q′} as the totally real
  algebraic closure of Q; physical98/footer97 writes the norm and
  norm-one normalization; physical100/footer99 gives the boxed
  all-conjugate purity correction. Physical102/footer101,104/footer103
  and106/footer105 were also visually compared with the supplied
  transcriptions. The exhaustion question on physical104 remains in a
  deleted passage. AQ-WEIL-NOTATION proves the exact field/group
  identification with those norm/purity conditions retained. The
  counterexample to a valuation-only paraphrase does not attribute that
  reading to the manuscript term \emph{unité absolue}. Uncertain
  superscripts and crossed revisions remain source qualifications.
\end{itemize}

Grothendieck, A. (n.d.). \emph{Motivic, period and realization notes}
{[}Unpublished manuscripts, cote 18{]}. Archives Grothendieck,
Université de Montpellier.
\url{https://grothendieck.umontpellier.fr/18.pdf}

\begin{itemize}
\tightlist
\item
  Supplied sources: scans 9, 11, 13, 38 and 59; manuscript pp. 8, 10,
  12, 37 and 58. These contain respectively the Tate-filtration
  conventions, strong period triangle, three elliptic questions,
  real-involution claim and cocycle-image genericity reformulation. The
  title here is descriptive, not asserted to be a manuscript title.
\end{itemize}

Grothendieck, A. (1983). \emph{Courbes elliptiques réelles; capping and
nodal-propagation notes} {[}Unpublished manuscripts, cote 67{]}.
Archives Grothendieck, Université de Montpellier.
\url{https://grothendieck.umontpellier.fr/67.pdf}

\begin{itemize}
\tightlist
\item
  Supplied sources: scans 74--82, archive foliation 76--84, and scan 91,
  archive foliation 93. Scan 91 is explicitly dated ``Déc 83.'' This
  date is not assigned to every sheet of the cote. Scan 75 specifies a
  nonzero tangent tensor; scan 78 uses its inverse dual pairing and
  mutually polar tangent disks; scans 76--77 impose no norm restriction
  on that tensor. These are the precise source conventions used in
  AQ-67-RANGE.
\end{itemize}

Grothendieck, A. (n.d.). \emph{Pseudoline-relative-position surface
notes} {[}Unpublished manuscripts, cote 154{]}. Archives Grothendieck,
Université de Montpellier.
\url{https://grothendieck.umontpellier.fr/154.pdf}

\begin{itemize}
\tightlist
\item
  Supplied source: scan 120. Its Euler characteristic is
  cellular/topological. No finite-field scheme or Frobenius structure on
  that surface is specified. The title is descriptive.
\end{itemize}

Habegger, P. (2014). \emph{Singular moduli that are algebraic units}
(arXiv:1402.1632v1). \url{https://arxiv.org/pdf/1402.1632v1}

\begin{itemize}
\tightlist
\item
  Used: §1, p. 1, the classical integrality of CM j-invariants, to
  certify that j = −6912/23 is not CM. The AQ proof does not use the
  paper's finiteness theorem for singular units.
\end{itemize}

Honda, T. (1968). Isogeny classes of abelian varieties over finite
fields. \emph{Journal of the Mathematical Society of Japan, 20}(1--2),
83--95. \url{https://doi.org/10.2969/jmsj/02010083}

\begin{itemize}
\tightlist
\item
  Main Theorem, p. 83: existence for weight-one Weil integers. Primary
  scan:
  \url{https://people.math.rochester.edu/faculty/doug/otherpapers/Honda-Ab.pdf}
\item
  AQ-WEIL-EXIST derives the specified weight-i cohomological occurrence
  using an i-th root and Künneth. This is an unconditional known
  consequence, not a new RH application.
\end{itemize}

Kreutz, T., Shen, M., \& Vial, C. (2026). \emph{De Rham--Betti classes
with coefficients} (arXiv:2206.08618v3).
\url{https://arxiv.org/pdf/2206.08618v3}

\begin{itemize}
\tightlist
\item
  Also: \url{https://arxiv.org/html/2206.08618v3}
\item
  Read: Proposition 3.1, Theorem 5.13(ii), Proposition 5.14. The
  coefficient/CM-field restriction in Theorem 5.13(ii) is satisfied for
  one elliptic curve. The result identifies the generated period torsor,
  answering 18/013(b); it does not settle maximal transcendence degree
  for a non-CM period point.
\end{itemize}

Lagarias, J. C., \& Odlyzko, A. M. (1977). Effective versions of the
Chebotarev density theorem. In A. Fröhlich (Ed.), \emph{Algebraic number
fields: L-functions and Galois properties} (pp. 409--464). Academic
Press.

\begin{itemize}
\tightlist
\item
  Full retype of the original article:
  \url{https://aareyanmanzoor.github.io/assets/articles/lagarias-odlyzko.pdf}
\item
  Used: Theorem 1.1, equation (1.3), retype p. 4. For the quadratic
  extension Q(√−3)/Q, degree 2 and discriminant 3, its hypothesis
  follows from ζ\_K(s)=ζ(s)L(s,χ\_−3) and the adopted Dirichlet GRH. Its
  effective absolute constant is left as c1, not invented numerically.
\end{itemize}

Lamzouri, Y., Li, X., \& Soundararajan, K. (2015). Conditional bounds
for the least quadratic non-residue and related problems.
\emph{Mathematics of Computation, 84}(295), 2391--2412.
\url{https://arxiv.org/pdf/1309.3595v3}

\begin{itemize}
\tightlist
\item
  Operative source: corrected arXiv v3, dated December 8, 2016;
  Corollary 1.2, p. 3. Under Dirichlet GRH, q\textgreater3 and (a,q)=1,
  the least prime in the class a mod q is at most (φ(q) log q)².
  AQ-FINITE uses Q=lcm(N,4), so the q\textgreater3 condition is
  automatic.
\end{itemize}

Lamzouri, Y., Li, X., \& Soundararajan, K. (2017). Corrigendum to
``Conditional bounds for the least quadratic non-residue and related
problems.'' \emph{Mathematics of Computation, 86}(307), 2551--2554.
\url{https://doi.org/10.1090/mcom/3261}

\begin{itemize}
\tightlist
\item
  Bibliographic/version control: the operative theorem was checked in
  corrected v3. The standalone corrigendum's complete text was not
  independently opened in this AQ run; no unverified early-version
  constants are used.
\end{itemize}

Milne, J. S. (1994). Motives over finite fields. In U. Jannsen, S. L.
Kleiman, \& J.-P. Serre (Eds.), \emph{Motives} (Proceedings of Symposia
in Pure Mathematics, Vol. 55, Part 1, pp. 401--459). American
Mathematical Society. \url{https://jmilne.org/math/articles/1994aP.pdf}

\begin{itemize}
\tightlist
\item
  Exact author-PDF locators: Conjecture 1.14 and Proposition 1.15, p.
  11; Corollary 1.16, pp. 11--12; start of §2, p. 12; Proposition 2.6
  and Remark 2.7, p. 15.
\item
  The injectivity/classification and abelian/Artin generation claims use
  Tate. The eigenvalue-existence argument is unconditional via Honda.
  These two uses are separated explicitly in AQ-WEIL-EXIST.
\item
  Author's corrections page:
  \url{https://jmilne.org/math/xnotes/Seattle1.html}
\end{itemize}

Milne, J. S. (1999). Lefschetz motives and the Tate conjecture.
\emph{Compositio Mathematica, 117}(1), 45--76.
\url{https://doi.org/10.1023/A:1000776613765}

\begin{itemize}
\tightlist
\item
  Author manuscript:
  \url{https://www.jmilne.org/math/articles/1999bP.pdf}
\item
  Theorem 7.1, journal p. 72, proves CM Hodge implies the Tate statement
  (0.1) for all abelian varieties over an algebraic closure of a finite
  field. Introduction (0.1), journal p. 46, includes numerical equals
  homological and algebraicity of all Tate classes; the paragraph
  immediately following explains finite-model descent and the full
  pole-order statement. Remark 7.6(b), journal p. 76, extends the
  consequence to homological motives of abelian type. The CM-Hodge
  premise is supplied by the 2026 result below.
\end{itemize}

Milne, J. S. (2002). Polarizations and Grothendieck's standard
conjectures. \emph{Annals of Mathematics, 155}(2), 599--610.
\url{https://annals.math.princeton.edu/2002/155-2/p07}

\begin{itemize}
\tightlist
\item
  Operative final author version, September 19, 2001:
  \url{https://www.jmilne.org/math/articles/2002aP.pdf}
\item
  Theorem 3.3(a),(b), author pp. 8--9, gives numerical versus ℓ-adic
  homological equivalence and Hodge standard for all abelian varieties
  over algebraically closed fields, conditional on CM Hodge. Remark 3.5,
  author p. 9, treats good Weil cohomology theories. Earlier
  long-version Theorem 4.6 is not the numbering of this final author
  PDF. The primitive Hodge pairing has the sign (−1)\^{}r; the
  abbreviated 2026 report\textquotesingle s displayed pairing omits that
  sign.
\end{itemize}

Milne, J. S. (2026). \emph{On the Hodge, Tate, and standard conjectures
for abelian varieties} (Hexagon:2611.00244v1, November 9).
\url{https://www.jmilne.org/math/xnotes/HTS.pdf}

\begin{itemize}
\tightlist
\item
  Theorem 2.1 states CM Hodge; Theorem 3.1, p. 6, states the full Tate
  conjecture for all abelian varieties over finite fields; Theorem 3.5,
  p. 8, states Hodge standard for all abelian varieties. Section 3.7
  explains reduction to finite fields. This author report confirms the
  current status and cites the underlying proofs; it is not substituted
  for a new full proof check of those papers.
\end{itemize}

Mohajer, M. (2026). \emph{Rational elliptic period relations across
realizations: Complex lattices, local p-adic periods, and Drinfeld
modules} (arXiv:2611.25399v1, November 19).
\url{https://arxiv.org/pdf/2611.25399v1}

\begin{itemize}
\tightlist
\item
  Section 1.1, equation (1.3), pp. 2--3, retains the classical full
  complex-period conjecture; Proposition 2.1 concerns holomorphic E2.
  Theorems 4.2 and 4.4 concern local B\_dR and normalized C\_p periods.
  These exact scopes were inspected; the complete 39-page preprint proof
  is not certified here. Its local statements do not settle the complex
  E2* splitting-scalar problem.
\end{itemize}

Sebbar, A., \& Besrour, K. (2026). \emph{Zeros of Katz's p-adic
Eisenstein series E2: Hypergeometric congruences and isogenies}
(arXiv:2611.14502v1, November 10).
\url{https://arxiv.org/pdf/2611.14502v1}

\begin{itemize}
\tightlist
\item
  Section 8.4, Remark 8.11, Proposition 8.12, and Question 8.13, pp.
  32--33, leave algebraic non-CM zeros of the p-adic Katz form as a
  question and obtain nonvanishing under the crystalline period
  conjecture. Algebraicity over Q\_p is explicitly distinguished from
  algebraicity over Q. The abstract and these exact statements were
  read; the whole preprint is not claimed checked.
\end{itemize}

Tate, J. (1966). Endomorphisms of abelian varieties over finite fields.
\emph{Inventiones Mathematicae, 2}, 134--144.
\url{https://doi.org/10.1007/BF01404549}

\begin{itemize}
\tightlist
\item
  Primary scan:
  \url{https://pazuki.perso.math.cnrs.fr/index_fichiers/Tate66.pdf}
\item
  Main Theorem, p. 134, is the Tate-module endomorphism theorem. Theorem
  2(a), p. 140, identifies Q{[}π{]} with the centre of the semisimple
  endomorphism algebra, giving the semisimple Frobenius input. Theorem
  2(d), p. 140, identifies the quaternion algebra in the
  scalar-Frobenius supersingular case. AQ-RATIONAL-FIBRE verifies the
  explicit model over F\_9 before invoking this theorem.
\end{itemize}

van Hoften, P., \& Sempliner, J. (2026). \emph{The Hodge conjecture for
CM abelian varieties} (Hexagon:2610.00259v2).
\url{https://hexagonmath.org/2610.00259v2}

\begin{itemize}
\tightlist
\item
  Author-hosted PDF:
  \url{https://polvanhoften2.github.io/Pol\%20van\%20Hoften\%20and\%20Jack\%20Sempliner\%20-\%20The\%20Hodge\%20conjecture\%20for\%20CM\%20abelian\%20varieties.pdf}
\item
  Theorem I, p. 1, proves rational Hodge for every complex CM abelian
  variety. Sections 1.1.2--1.1.3, pp. 1--2, state the consequences for
  standard conjectures and finite-field Tate using Milne. The exact
  theorem and downstream hypotheses were inspected. The complete proof
  is a cited research input, not independently certified in the AQ
  module. Version 2 is retained explicitly to avoid relying on an
  unspecified first posting.
\end{itemize}

Waldschmidt, M. (2004). Open Diophantine problems. \emph{Moscow
Mathematical Journal, 4}(1), 245--305.
\url{https://webusers.imj-prg.fr/~michel.waldschmidt/articles/pdf/ODP-MoscowMathJ.pdf}

\begin{itemize}
\tightlist
\item
  Exact locator: p. 277, question immediately after Conjecture 3.22
  (Katz), about the nonholomorphic weight-two Eisenstein invariant for a
  non-CM lattice with algebraic invariants. The AQ module writes its
  exact normalized Hodge-splitting scalar.
\end{itemize}

Walker, M. E. (2007). The morphic Abel--Jacobi map. \emph{Compositio
Mathematica, 143}, 909--944.
\url{https://doi.org/10.1112/S0010437X07002278}

\begin{itemize}
\tightlist
\item
  Primary PDF:
  \url{https://www.cambridge.org/core/services/aop-cambridge-core/content/view/4A44146CE680F4F3C857048FDEDD9F48/S0010437X07002278a.pdf/the-morphic-abel-jacobi-map.pdf}
\item
  Read: Example 5.5, p. 928 (the zero-cycle morphic Jacobian/Albanese);
  Theorem 5.7(2),(4), p. 930 (cycle-map and every-Chow-correspondence
  compatibility); Theorem 6.1, p. 933, and the complete proof of Theorem
  6.2, pp. 933--934 (homologically trivial cycles with nonzero Lawson
  cylinder after an elliptic product); §7 (the separate ordinary
  regular-universality question).
\item
  No unchecked effective hypersurface degree is used. ``Sufficiently
  general of sufficiently large degree'' is the precise retained theorem
  hypothesis.
\end{itemize}

\subsection{Additional primary sources for the exact capping and
compactification
results}\label{additional-primary-sources-for-the-exact-capping-and-compactification-results}

Brock, J. F., \& Margalit, D. (2007). Weil--Petersson isometries via the
pants complex. \emph{Proceedings of the American Mathematical Society,
135}(3), 795--803. \url{https://arxiv.org/abs/math/0412499}

\begin{itemize}
\tightlist
\item
  Operative author version: December 21, 2004, arXiv:math/0412499v1;
  \url{https://jeffbrock.net/home/text/papers/dan/www/wp.pdf} . Main
  Theorem, p. 2, explicitly includes the four-punctured sphere and
  states that every Weil--Petersson isometry is induced by the extended
  mapping class group. The finite kernel in that case does not affect
  the countability argument.
\end{itemize}

Fırat, A. H. (2021). Hyperbolic three-string vertex. \emph{Journal of
High Energy Physics, 2021}(8), Article 35.
\url{https://doi.org/10.1007/JHEP08(2021)035}

\begin{itemize}
\tightlist
\item
  Operative primary version: arXiv:2102.03936v2, August 18, 2021,
  \url{https://arxiv.org/pdf/2102.03936} . Equations (3.11), (3.13),
  (3.14a), (3.23), and (3.30)--(3.34), pp. 14--17 and 24--25, provide
  the hypergeometric projective connection and the arclength
  flat-cylinder cap coordinate. AQ-67-EXACT supplies its own cusp-limit
  and continuous phase proof; no unproved numerical integer-branch
  prescription is used.
\end{itemize}

Kahn, B., \& Sujatha, R. (2008). \emph{Birational geometry and
localisation of categories} (author manuscript, May 22).
\url{https://webusers.imj-prg.fr/~bruno.kahn/preprints/Birat.pdf}

\begin{itemize}
\tightlist
\item
  Theorem 1.6.3(a), p. 11, states the characteristic-zero resolution and
  smooth compactification consequence in the exact form used: a smooth
  quasiprojective variety admits a smooth projective compactification.
  AQ-14-QUASIPROJECTIVE proves the correspondence/Albanese comparison
  from that cited input. This is an inspected primary theorem statement;
  the Hironaka proof is not independently reproduced.
\end{itemize}

Mondello, G. (2011). Riemann surfaces with boundary and natural
triangulations of the Teichmüller space. \emph{Journal of the European
Mathematical Society, 13}(3), 635--684.
\url{https://doi.org/10.4171/JEMS/263}

\begin{itemize}
\tightlist
\item
  Primary PDF:
  \url{https://ems.press/content/serial-article-files/31778} . Section
  2.2, p. 640, recalls Fenchel--Nielsen coordinates with zero boundary
  lengths interpreted as cusps. Section 5.4, Theorem 5.4, p. 660, treats
  the real-analytic infinite-grafting parametrization; the canonical cap
  is its punctured flat cylinder in a disk coordinate.
\end{itemize}

National Institute of Standards and Technology. (n.d.). \emph{Digital
Library of Mathematical Functions}. \url{https://dlmf.nist.gov/}

\begin{itemize}
\tightlist
\item
  Exact inspected identities: (15.4.21),
  \url{https://dlmf.nist.gov/15.4.E21} (zero-balanced hypergeometric
  logarithmic asymptotic); (5.5.5), \url{https://dlmf.nist.gov/5.5.E5}
  (gamma duplication); (5.9.16), \url{https://dlmf.nist.gov/5.9.E16}
  (digamma integral); (5.11.1), \url{https://dlmf.nist.gov/5.11.E1}
  (analytic-log Stirling expansion), and (5.11.13),
  \url{https://dlmf.nist.gov/5.11.E13} (gamma-ratio asymptotic).
  Accessed November 21, 2026. The positivity and range arguments are
  proved in full in AQ-67-EXACT.
\end{itemize}

Wolpert, S. A. (1983). On the symplectic geometry of deformations of a
hyperbolic surface. \emph{Annals of Mathematics, 117}(2), 207--234.
\url{https://doi.org/10.2307/2007075}

\begin{itemize}
\tightlist
\item
  Original scan:
  \url{https://math.umd.edu/~swolpert/preprints/On_the_Symplectic.pdf} .
  Lemma 4.1, printed p. 229, gives the exact twist Hamiltonian identity
  Ω(t\_α)=−dℓ\_α. The finite-area surface scope is stated in §1, pp.
  208--209. AQ-67-NONHOLO uses Brock--Margalit for the exceptional
  four-punctured-sphere isometry assertion, rather than an unqualified
  invocation of Royden.
\end{itemize}

\subsection{Additional primary sources for the all-characteristic
nonproper
implication}\label{additional-primary-sources-for-the-all-characteristic-nonproper-implication}

Deligne, P. (1974). La conjecture de Weil: I. \emph{Publications
Mathématiques de l'IHÉS, 43}, 273--307.
\url{https://doi.org/10.1007/BF02684373}

\begin{itemize}
\tightlist
\item
  Original primary scan:
  \url{https://www.numdam.org/article/PMIHES_1974__43__273_0.pdf} .
  Theorem (1.6), printed p. 276, is the smooth-projective Frobenius
  purity theorem. Twisting H\^{}\{2i−1\} by i gives geometric weight −1;
  a torus Tate module has geometric weight −2. The candidate proves the
  resulting vanishing without requiring semisimplicity.
\end{itemize}

Geisser, T. (2015). Rojtman's theorem for normal schemes.
\emph{Mathematical Research Letters, 22}(4), 1129--1144.
\url{https://arxiv.org/abs/1402.1831}

\begin{itemize}
\tightlist
\item
  Operative directly inspected primary version: arXiv:1402.1831v2,
  February 25, 2015, \url{https://arxiv.org/pdf/1402.1831} . Theorem
  1.1, preprint p. 2, gives prime-to-characteristic torsion
  Suslin--Albanese isomorphism and H₁\^{}S(T,Z)⊗Q\_l/Z\_l=0 for normal
  separated finite-type schemes, without assuming smooth
  compactification. Theorem 1.3, preprint p. 3, supplies the
  smooth-curve special case H₀\^{}S(C,Z)\^{}0≅Alb\_sa(C)(k). Lemma 2.2,
  p. 5, records the Albanese map on Suslin H₀. The exact theorem
  statements were read; the entire Geisser proof is a cited input, not a
  new whole-proof certification.
\end{itemize}

Geisser, T., \& Schmidt, A. (2016). Tame class field theory for singular
varieties over algebraically closed fields. \emph{Documenta Mathematica,
21}, 91--123. \url{https://doi.org/10.4171/DM/528}

\begin{itemize}
\tightlist
\item
  Original journal PDF:
  \url{https://ems.press/content/serial-article-files/26306?nt=1} .
  Theorem 1.1, printed p. 92, gives the finite perfect
  pairing/reciprocity isomorphism for coefficients prime to the
  characteristic. Theorem 7.1, printed p. 118, identifies the pairing
  with the Suslin--Voevodsky comparison. The candidate passes through
  direct and inverse coefficient limits explicitly; it does not claim
  that every finite-coefficient H¹ of a variety equals that of its
  Albanese.
\item
  The published PDF lists authors Thomas Geisser and Alexander Schmidt
  in that order; the publisher metadata happens to reverse them. The
  citation follows the paper itself.
\end{itemize}

Kahn, B. (2014). Algebraic tori as Nisnevich sheaves with transfers.
\emph{Annales de la Faculté des Sciences de Toulouse: Mathématiques,
23}(3), 699--715. \url{https://doi.org/10.5802/afst.1421}

\begin{itemize}
\tightlist
\item
  Directly inspected final journal PDF:
  \url{https://webusers.imj-prg.fr/~bruno.kahn/preprints/tores-nisnevich.pdf}
  . Lemma 4.1(a),(b), printed p. 707, proves that a pointed morphism of
  semiabelian varieties is a homomorphism and that a semiabelian variety
  is its own Serre Albanese. Equivalent earlier locator: July 4, 2013
  author version,
  \url{https://webusers.imj-prg.fr/~bruno.kahn/preprints/tores-nisnevich8.pdf}
  , Lemma 3, p. 8. Both exact statements and the short proof were read.
\end{itemize}

Sun, S. (2017). Generic base change, Artin's comparison theorem, and the
decomposition theorem for complex Artin stacks. \emph{Journal of
Algebraic Geometry, 26}(3), 513--555.
\url{https://doi.org/10.1090/JAG/683}

\begin{itemize}
\tightlist
\item
  Operative inspected primary version: arXiv:1511.01393v1, November 4,
  2015, \url{https://arxiv.org/pdf/1511.01393} . Section 1.3.1, p. 2,
  permits Λ=Z\_l. Section 2.1, p. 3, gives condition (LO) and the
  integral stratifiable adic category. Definition 2.2, p. 3, specifies
  one dense open valid for all subsequent base changes. Lemma 2.3(iii),
  pp. 4--5, and its proof explicitly obtain one dense open for all
  finite coefficient levels by finite Jordan--Hölder dévissage. Theorem
  2.5(i), p. 7, states generic base change; Remark 2.5.1(ii), p. 8,
  recalls preservation of stratifiability. The candidate uses only
  finite-type F\_p bases with l≠p, so (LO) applies. A family of separate
  opens, one for each l\^{}n, would not suffice and is not the argument
  used.
\end{itemize}

Bloch correspondence and specialization inputs are those already
precisely cited in AQ-14-FORWARD: Bloch (1979), Propositions 3.5 and
3.8, and Achter--Casalaina-Martin--Vial, arXiv:2007.07180v2,
Propositions A.23, A.24 and A.26. The latter paper\textquotesingle s
(0.2)--(0.3) and §A.3 also identify the rational Tate module of the
torsion-coefficient cohomology with ordinary Q\_l-cohomology.

Kahn, B., \& Otsubo, N. (2026). \emph{Periods and CM motives of abelian
type} (arXiv:2611.13723v1; manuscript dated November 11).
\url{https://arxiv.org/pdf/2611.13723}

\begin{itemize}
\tightlist
\item
  Exact targeted locators: §2.4, pp. 6--7; Definition 2.5, p. 8;
  Corollaries 3.21, 3.25 and 3.35, pp. 24, 26 and 30--31; Remark 3.22,
  p. 25; Theorems 4.36, 4.39 and 4.44, pp. 46--47 and 49. The new
  comparison statements concern rank-one coefficient motives and period
  classes. AQ-KO-RANK corrects the ancillary rank sentence in Lemma 4.6,
  p. 34, by a complete elementary proof. This targeted reading is not
  certification of the full 64-page preprint.
\end{itemize}

Illusie, L., \& Temkin, M. (2014). Exposé X: Gabber's modification
theorem (log smooth case). In L. Illusie, Y. Laszlo, \& F. Orgogozo
(Eds.), \emph{Travaux de Gabber sur l'uniformisation locale et la
cohomologie étale des schémas quasi-excellents} (\emph{Astérisque,
363--364}, pp. 167--212).
\url{https://www.numdam.org/item/AST_2014__363-364__167_0.pdf}

\begin{itemize}
\tightlist
\item
  Theorem 3.5(ii)--(iii), printed p. 202, gives the separable projective
  prime-to-ℓ alteration over a perfect field used in AQ-14-ALTERATION.
  Definition 3.3.1, p. 197, specifies prime-to-ℓ degree on every generic
  component. The introduction, p. 167, and §3.5.4, p. 207, distinguish
  this separable strengthening from Theorem 2.1. Regular finite-type
  sources over the perfect base are smooth, and projectivity over a
  projective compactification supplies projective parameters. The exact
  theorem was directly inspected; its external proof remains a cited
  input.
\end{itemize}

Milne, J. S. (2008, March 16). \emph{Abelian varieties} (Version 2.0).
\url{https://www.jmilne.org/math/CourseNotes/AV.pdf}

\begin{itemize}
\tightlist
\item
  Chapter I, Theorem 3.2, printed p. 17, extends a rational map from a
  nonsingular variety to an abelian variety over the whole source. This
  exact numbering belongs to the specified course-note version. It is
  the extension input for the fibre-sum trace in AQ-14-ALTERATION.
\end{itemize}

\subsection{Check and scope record}\label{check-and-scope-record}

\nolinkurl{archive_check_events.tsv} records the AQ author's complete
readings of the written proofs and the separate symbolic/finite checks.
Independent readings by the main agent and \nolinkurl{source_and_masullo}
are recorded by those reviewers, not manufactured here. The full prior
SA06/07/08/11/22/23 proof returns are absent from the supplied packet.
Their available intake rows are preserved verbatim in
\nolinkurl{archive_intake_inventory.tsv}; indexing does not certify a
proof.

\subsection{Cote 66 source recovery and affine
comparison}\label{cote-66-source-recovery-and-affine-comparison}

Grothendieck, A. (n.d.). \emph{List of subjects} {[}Unpublished
manuscript, cote 66, physical PDF p. 4, archive footer 3, subject 21{]}.
Archives Grothendieck, Université de Montpellier.
\url{https://grothendieck.umontpellier.fr/66.pdf}

The Stacks Project Authors. (2026). \emph{The Stacks Project}.
\url{https://stacks.math.columbia.edu}

\begin{itemize}
\tightlist
\item
  The parent module \nolinkurl{root_henselian.tex}, BR-66-01--03, uses
  Theorem59.82.7/09ZI (Gabber comparison), Lemma15.11.6/09XI (idempotent
  criterion), Theorem59.91.11/095T (proper base change),
  Lemma15.11.4/0ALJ (complete pairs), Lemma15.12.1/0A02 (henselization),
  and Lemmas21.4.2--3/03AI--03AJ (torsor sections and H1). The complete
  deductions received an independent whole-proof reading by this
  module\textquotesingle s author; the exact Stacks statements were
  directly inspected on21November2026. The parent alone recovered and
  visually checked the original66 sheet; this cross-review does not
  relabel that as a separate source inspection here.
\end{itemize}

\subsection{Source-route references}\label{source-route-references}

\begin{itemize}
\tightlist
\item
  Alpöge, L. (2026, November 20). \emph{Remark} {[}Four-page research
  announcement{]}. \url{https://alpo.ge/remark.pdf}. Adopted theorem
  locators are the five named declarations in
  \nolinkurl{ALPOGE_Main.lean}; definitions in
  \nolinkurl{ALPOGE_Statements.lean}.
\item
  Masullo, A. M. (2026). \emph{The music of the primes} {[}Unfinished
  manuscript released November 20{]}.
  \url{https://math216.wordpress.com/wp-content/uploads/2026/11/the-music-of-the-primes.pdf}.
  Exact internal locators listed above.
\item
  Vakil, R. (2026, November 20). \emph{Hurried announcement: Alessandro
  Masullo GRH}. Math 216: Foundations of Algebraic Geometry.
  \url{https://math216.wordpress.com/2026/11/20/hurried-announcement-alessandro-masullo-grh/}.
\item
  Serre, J.-P. (1977). \emph{Linear representations of finite groups}
  (L. L. Scott, Trans.). Springer.
  \url{https://doi.org/10.1007/978-1-4684-9458-7}. Exact input: §10.5,
  Theorem 20, p.78 (Brauer induction). Authoritative publisher record
  retrieved; theorem text was inspected in a digitized book hosted by
  Rochester, search ref \nolinkurl{}.
\item
  Cogdell, J. W. (n.d.). \emph{On Artin L-functions}. Author manuscript,
  \url{https://people.math.osu.edu/cogdell.1/artin-www.pdf}. Exact
  inputs: §2.1, p.10, equation (2); §2.2, p.12, ramified Euler factors
  and inductivity. Retrieved ref \nolinkurl{}; full
  31-page PDF saved as background. Use n.d. because this inspected
  manuscript has no explicit date. The author homepage now links a
  distinct 33-page PDF,
  \url{https://people.math.osu.edu/cogdell.1/Artin-www.pdf}, under the
  published title \emph{L-functions and non-abelian class field theory,
  from Artin to Langlands} (2015), in D. Dumbaugh and J. Schwermer,
  \emph{Emil Artin and Beyond---Class Field Theory and L-Functions}
  (pp.129--161), EMS, DOI 10.4171/146. Author page and publisher
  contents verified (\nolinkurl{},
  \nolinkurl{}); publication date 31 March 2015
  verified on publisher page (\nolinkurl{}). Do not
  transfer the inspected 31-page manuscript page numbers to that
  differently paginated chapter.
\item
  The Stacks Project Authors. (2026). \emph{Plane curves}, §53.9, Lemma
  53.9.3. \url{https://stacks.math.columbia.edu/tag/0BYA}. Retrieved ref
  \nolinkurl{}. This supplies the exact plane-curve genus
  formula; no new proof of the known formula is claimed.
\end{itemize}

The F-SM-02 zero-counting input can be cited to Iwaniec and Kowalski,
\emph{Analytic Number Theory}, Theorem 5.8 (p.104). This locator was
corroborated in search but the actual theorem page was not yet read
during this subrun; it should remain a citation-check item until
inspected or replaced by an already verified equivalent input. F-SM-01
itself assumes only the explicit summability condition and does not
depend on this unresolved bibliographic detail.

\end{document}
